% Generated mechanically from frozen input p09/ko/src/spaces-more-morphisms.tex.
% Do not edit this generated file; edit the bound locale source instead.

% OK, start here.
%

\setcounter{chapter}{75}
\chapter{대수공간의 사상에 관한 보충}



\phantomsection
\label{spaces-more-morphisms-section-phantom}


\section{서론}
\label{spaces-more-morphisms-section-introduction}

\noindent
이 장에서는 대수공간의 사상들의 성질에 대한 연구를 계속한다.
기본 참고문헌은 \cite{Kn}이다.





\section{규약}
\label{spaces-more-morphisms-section-conventions}

\noindent
모든 스킴은 큰 fppf 자리 $\Sch_{fppf}$에 들어 있다고 항상
가정한다. 또한 고려하는 모든 환 $A$는 $\Spec(A)$가 이 큰 자리의
대상과 (동형이라는 의미에서) 같다는 성질을 가진다.

\medskip\noindent
$S$를 스킴이라 하고 $X$를 $S$ 위의 대수공간이라 하자. 이 장과
다음 장에서는 $X \times_S X$를 $X$와 자기 자신의 곱
($S$ 위 대수공간들의 범주에서 취한 곱)에 대한 표기로 쓰며,
$X \times X$는 쓰지 않을 것이다.







\section{근적 사상}
\label{spaces-more-morphisms-section-radicial}

\noindent
스킴의 사상과 달리, 근적 사상은 보편 단사 사상과 같은 개념이
아니다. 실제로 근적이라는 조건이 조금 더 강하다.

\begin{definition}
\label{spaces-more-morphisms-definition-radicial}
$S$를 스킴이라 하자. $f : X \to Y$를 $S$ 위 대수공간의
사상이라 하자. $f$가 {\it 근적}이라는 것은 다음을 뜻한다. 모든
사상 $\Spec(K) \to Y$에 대하여, 여기서
$K$는 체이고, 축약
$(\Spec(K) \times_Y X)_{red}$가 공집합이거나 $K$의 순수 비분리
체 확대의 스펙트럼으로 표현 가능하다.
\end{definition}

\begin{lemma}
\label{spaces-more-morphisms-lemma-radicial-implies-universally-injective}
대수공간의 근적 사상은 보편 단사이다.
\end{lemma}

\begin{proof}
$S$를 스킴이라 하자. $f : X \to Y$를 $S$ 위 대수공간의 근적
사상이라 하자. 정의로부터 사상 $\Spec(K) \to Y$가 주어졌을 때
이를 $X$로 가는 사상으로 올리는 방법이 많아야 하나임은 분명하다.
따라서 Morphisms of Spaces, Lemma
\ref{spaces-morphisms-lemma-universally-injective}에 의해 $f$는
보편 단사이다.
\end{proof}

\begin{example}
\label{spaces-more-morphisms-example-universally-injective-not-radicial}
이제 보편 단사라는 조건과 근적이라는 조건은 동치가 아니다. 예를
들어 Spaces, Example \ref{spaces-example-Qbar}의 사상
$$
X = [\Spec(\overline{\mathbf{Q}})/
\text{Gal}(\overline{\mathbf{Q}}/\mathbf{Q})]
\longrightarrow
S = \Spec(\mathbf{Q})
$$
은 보편 단사이지만 위 의미에서 근적이지 않다.
\end{example}

\noindent
그럼에도 역함의가 성립하는 경우가 많다.

\begin{lemma}
\label{spaces-more-morphisms-lemma-when-universally-injective-radicial}
$S$를 스킴이라 하자. $f : X \to Y$를 $S$ 위 대수공간의 보편
단사 사상이라 하자.
\begin{enumerate}
\item $f$가 decent이면 $f$는 근적이다.
\item $f$가 준분리이면 $f$는 근적이다.
\item $f$가 국소 분리이면 $f$는 근적이다.
\end{enumerate}
\end{lemma}

\begin{proof}
대수공간의 사상들의 성질 $\mathcal{P}$가 밑변환과 합성에 대해
안정적이고 모든 폐몰입에 대하여 성립한다고 하자.
$f : X \to Y$가 $\mathcal{P}$를 가지며 보편 단사라고 가정하자.
그러면 Definition \ref{spaces-more-morphisms-definition-radicial}의 상황에서 사상
$(\Spec(K) \times_Y X)_{red} \to \Spec(K)$는 보편 단사이고
$\mathcal{P}$를 가진다. 따라서
$$
\mathcal{P} + \text{universally injective}
\Rightarrow
\text{radicial}
$$
을 증명하는 문제는 다음 문제로 환원된다. 체 위의 공집합이 아닌
축약 대수공간 $X$의 구조 사상 $X \to \Spec(K)$가 보편 단사이고
$\mathcal{P}$를 가지면, 이 대수공간이 체의 스펙트럼으로 표현
가능함을 보여야 한다. 실제로 이 경우 $X \to \Spec(K)$는 스킴의
사상이 되고, 스킴의 사상에 대하여 근적과 보편 단사가 동치이므로
결론을 얻는다. Morphisms, Lemma
\ref{morphisms-lemma-universally-injective}를 보라.

\medskip\noindent
(1)을 증명하자. $f$가 decent이고 보편 단사라고 가정하자.
Decent Spaces, Lemmas
\ref{decent-spaces-lemma-base-change-relative-conditions},
\ref{decent-spaces-lemma-composition-relative-conditions},
\ref{decent-spaces-lemma-properties-trivial-implications}에 의해
(몰입이 decent임을 확인하기 위해 마지막 보조정리를 쓴다) 첫
문단의 논의가 적용된다. $X$를 체 $K$ 위의 공집합이 아닌 decent
축약 대수공간이며 보편 단사라고 하자. 특히 $|X|$는 한원소
집합이다. Decent Spaces, Lemma
\ref{decent-spaces-lemma-when-field}에 의해
$X \cong \Spec(L)$임을 얻으며, 여기서 $K \subset L$은 어떤
확대이다. 이는 원하는 결론이다.

\medskip\noindent
준분리 사상은 decent이다. Decent Spaces, Lemma
\ref{decent-spaces-lemma-properties-trivial-implications}를 보라.
따라서 (1)은 (2)를 함의한다.

\medskip\noindent
(3)을 증명하자. 분리 공리들은 밑변환과 합성에 대해 안정적이고,
폐몰입은 분리 사상임을 상기하자. Morphisms of Spaces, Lemmas
\ref{spaces-morphisms-lemma-base-change-separated},
\ref{spaces-morphisms-lemma-composition-separated},
\ref{spaces-morphisms-lemma-immersions-monomorphisms}을 보라.
따라서 증명의 첫 문단의 논의가 적용된다. $X$를 체 $K$ 위의
축약 대수공간이며 보편 단사이고 국소 분리라고 하자. 특히
$|X|$는 한원소 집합이므로 $X$는 준콤팩트이다. Properties of
Spaces, Lemma \ref{spaces-properties-lemma-quasi-compact-space}를
보라. 전사 에탈 사상 $U \to X$를 찾을 수 있고 $U$는 아핀이다.
Properties of Spaces, Lemma
\ref{spaces-properties-lemma-quasi-compact-affine-cover}를 보라.
스킴의 사상
$$
j :
U \times_X U
\longrightarrow
U \times_{\Spec(K)} U
$$
을 생각하자. $X \to \Spec(K)$가 보편 단사이므로 $j$는 전사이고,
$X \to \Spec(K)$가 국소 분리이므로 $j$는 몰입이다. 전사 몰입은
폐몰입이다. Schemes, Lemma
\ref{schemes-lemma-immersion-when-closed}를 보라. 따라서
$R = U \times_X U$는 아핀 스킴의 폐부분스킴이므로 아핀이다.
특히 $R$은 준콤팩트이다. 그러므로 $X = U/R$은 준분리이고,
결과는 (2)에서 따른다.
\end{proof}

\begin{remark}
\label{spaces-more-morphisms-remark-weakly-radicial}
$X \to Y$를 대수공간의 사상이라 하자. (근적 사상의) 어떤
응용에서는 모든 $\Spec(K) \to Y$에 대하여, 여기서 $K$는 체이고,
다음을 요구하는 것으로 충분하다.
\begin{enumerate}
\item 공간 $|\Spec(K) \times_Y X|$는 한원소 집합이다.
\item 단사 사상
$\Spec(L) \to \Spec(K) \times_Y X$가 존재한다.
\item $K \subset L$은 순수 비분리이다.
\end{enumerate}
나중에 필요하면 이러한 사상을 {\it 약하게 근적}이라고 부를 수도
있다. 예를 들어 $X \to Y$가 전사이고 약하게 근적이면
$X(k) \to Y(k)$는 모든 대수적으로 닫힌 체 $k$에 대하여 전사이다.
Example \ref{spaces-more-morphisms-example-universally-injective-not-radicial}의 사상을
밑변환한
$X_{\overline{\mathbf{Q}}} \to \Spec(\overline{\mathbf{Q}})$는
약하게 근적이지만 근적이지 않다. Lemma
\ref{spaces-more-morphisms-lemma-when-universally-injective-radicial}의 유사 명제는
다음과 같다. $X \to Y$가 성질 ($\beta$)를 가지며 보편 단사이면
약하게 근적이다 (증명은 생략한다).
\end{remark}

\begin{lemma}
\label{spaces-more-morphisms-lemma-check-universally-injective}
$S$를 스킴이라 하자. $f : X \to Y$를 $S$ 위 대수공간의 사상이라
하자. 다음을 가정하자.
\begin{enumerate}
\item $f$는 국소 유한형이다.
\item 모든 에탈 사상 $V \to Y$에 대하여 사상
$|X \times_Y V| \to |V|$는 단사이다.
\end{enumerate}
그러면 $f$는 보편 단사이다.
\end{lemma}

\begin{proof}
이 문제는 $Y$ 위에서 에탈 국소적이다. 이는 Morphisms of Spaces,
Lemma \ref{spaces-morphisms-lemma-universally-injective-local}에서
따른다. 따라서 $Y$가 스킴이라고 가정할 수 있다. 그러면 특히
$Y$는 decent이고, Decent Spaces, Lemma
\ref{decent-spaces-lemma-conditions-on-point-in-fibre-and-qf}에 의해
$f$는 국소 준유한이다. $y \in Y$를 한 점이라 하고
$X_y$를 스킴론적 올이라 하자. $X_y$가 공집합이 아니라고
가정하자. Spaces over Fields, Lemma
\ref{spaces-over-fields-lemma-locally-quasi-finite-over-field}에 의해
$X_y$는 $\kappa(y)$ 위에서 국소 준유한인 스킴이다.
$|X_y| \subset |X|$는 $|X| \to |Y|$의 $y$ 위 올이므로,
$X_y$는 유일한 점 $x$를 가진다. 같은 사실은
$X_y \times_{\Spec(\kappa(y))} \Spec(k)$에도 성립한다. 여기서
$k/\kappa(y)$는 임의의 유한 분리 확대이다. 실제로 $k$를 $y$ 위에
놓이는, $Y$ 위의 에탈 스킴의 한 점의 잉여체로 실현할 수 있다.
More on Morphisms, Lemma
\ref{more-morphisms-lemma-realize-prescribed-residue-field-extension-etale}를
보라. 따라서 Varieties, Lemma
\ref{varieties-lemma-characterize-geometrically-disconnected}에 의해
$X_y$는 기하적으로 연결이다. 이는 유한 확대
$\kappa(x)/\kappa(y)$가 순수 비분리임을 함의한다.

\medskip\noindent
(이제 $Y$가 스킴인 경우에) 모든 $y \in Y$에 대하여 올
$X_y$가 공집합이거나,
$(X_y)_{red} = \Spec(\kappa(x))$이고
$\kappa(y) \subset \kappa(x)$가 순수 비분리임을 얻었다.
따라서 $f$는 근적이고 (몇 가지 세부사항은 생략한다), Lemma
\ref{spaces-more-morphisms-lemma-radicial-implies-universally-injective}에 의해 보편
단사이다.
\end{proof}




\section{단사 사상}
\label{spaces-more-morphisms-section-monomorphisms}

\noindent
이 절은 Morphisms of Spaces, Section
\ref{spaces-morphisms-section-monomorphisms}의 계속이다. 대수공간의
모든 단사 사상이 표현 가능한지 알고 싶다. 이것이 참임을 증명할
수 있거나 반례를 알고 있다면
\href{mailto:stacks.project@gmail.com}{stacks.project@gmail.com}으로
전자우편을 보내기 바란다. 지금은 다음 경우에 알려져 있다.
\begin{enumerate}
\item 국소 유한형인 단사 사상
(더 일반적으로 모든 분리 국소 준유한 사상은 Morphisms of Spaces,
Lemma
\ref{spaces-morphisms-lemma-locally-quasi-finite-separated-representable}에
의해 표현 가능하고, 국소 유한형인 단사 사상은 Morphisms of Spaces,
Lemma
\ref{spaces-morphisms-lemma-monomorphism-loc-finite-type-loc-quasi-finite}에
의해 국소 준유한이다),
\item 공역이 0차원 국소환들의 스펙트럼들의 분리합인 경우
(Decent Spaces, Lemma
\ref{decent-spaces-lemma-monomorphism-toward-disjoint-union-dim-0-rings}),
\item 평탄 단사 사상 (아래를 보라).
\end{enumerate}

\begin{lemma}[David Rydh]
\label{spaces-more-morphisms-lemma-flat-case}
대수공간의 평탄 단사 사상은 스킴으로 표현 가능하다.
\end{lemma}

\begin{proof}
$f : X \to Y$를 대수공간의 평탄 단사 사상이라 하자. $f$가 표현
가능함을 증명하려면 $X \times_Y V$가 스킴임을 보여야 한다.
여기서 $V$는 $Y$로 사상되는 임의의 스킴이다. 스킴이라는 성질은
국소적이므로 (Properties of Spaces, Lemma
\ref{spaces-properties-lemma-subscheme}), $V$가 아핀이라고
가정할 수 있다. 따라서 $Y = \Spec(B)$가 아핀 스킴이라고 가정할
수 있다. 다음으로 $X$를 준콤팩트 열린 부분공간으로 바꾸어 $X$가
준콤팩트라고 가정할 수 있다. $X$는 분리이다. 실제로
$X \to X \times_{\Spec(B)} X$가 동형이다. Limits of Spaces,
Lemma \ref{spaces-limits-lemma-enough-local}을 적용하여, $B$가
국소환이고 $X \to \Spec(B)$가 평탄 단사 사상이며 점 $x \in X$가
존재하여 $\Spec(B)$의 닫힌점으로 사상되는 경우로
환원한다. 그러면 분리 대수공간의 평탄 사상을 따라 일반화가
올려지므로 $X \to \Spec(B)$는 전사이다. Decent Spaces, Lemma
\ref{decent-spaces-lemma-generalizations-lift-flat}을 보라.
따라서 $\{X \to \Spec(B)\}$는 fpqc 덮개이다.
$X \to \Spec(B)$는 $X \to \Spec(B)$에 의한 밑변환 뒤 동형이 되는
사상이다. 그러므로 fpqc 내림에 의해 동형이다. Descent on Spaces,
Lemma \ref{spaces-descent-lemma-descending-property-isomorphism}을
보라.
\end{proof}

\noindent
다음 명제는 (어떤 의미에서) 위 보조정리의 변형이다.

\begin{lemma}
\label{spaces-more-morphisms-lemma-ui-case}
$S$를 스킴이라 하자. $f : X \to Y$를 대수공간의 준콤팩트 단사
사상이라 하고, 모든 $T \to Y$에 대하여 사상
$$
\mathcal{O}_T \to f_{T,*}\mathcal{O}_{X \times_Y T}
$$
이 단사라고 가정하자. 그러면 $f$는 동형이다 (따라서 스킴으로
표현 가능하다).
\end{lemma}

\begin{proof}
이 문제는 $Y$ 위에서 에탈 국소적이므로
$Y = \Spec(A)$가 아핀이라고 가정할 수 있다. 그러면 $X$는
준콤팩트이고, 아핀 스킴 $U = \Spec(B)$와 전사 에탈 사상
$U \to X$를 택할 수 있다 (Properties of Spaces, Lemma
\ref{spaces-properties-lemma-quasi-compact-affine-cover}).
$U \times_X U = \Spec(B \otimes_A B)$임을 주목하자. 따라서
준연접 $\mathcal{O}_X$-가군들의 범주는 가군 내림 데이터의 범주
$DD_{B/A}$와 동치이다. 이 내림 데이터는 $A \to B$에 대한 것이다.
Properties of Spaces,
Proposition \ref{spaces-properties-proposition-quasi-coherent},
Descent, Definition \ref{descent-definition-descent-datum-modules},
Descent, Subsection
\ref{descent-subsection-descent-modules-morphisms}를 보라.
한편,
$$
A \to B
$$
는 보편 단사 환 준동형이다. 실제로 $A$-가군 $M$이 주어지면,
보조정리의 가정에 의해
$A \oplus M \to B \otimes_A (A \oplus M)$은 단사이다.
따라서 Descent, Theorem \ref{descent-theorem-descent}에 의해
$DD_{B/A}$는 $A$-가군들의 범주와 동치이다. 그러므로
$f : X \to \Spec(A)$를 따른 당김은 준연접 가군들의 범주의
동치를 결정한다. 특히 $f^*$는 준연접 가군들 위에서 완전함자이고,
따라서 $f$는 평탄이다 (작은 세부사항은 생략한다). 더욱이 $f$가
전사임은 분명하다 (예를 들어
$\Spec(B) \to \Spec(A)$가 전사이기 때문이다). 따라서
$\{X \to \Spec(A)\}$는 fpqc 덮개이다.
$X \to \Spec(A)$는 $X \to \Spec(A)$에 의한 밑변환 뒤 동형이 되는
사상이다. 그러므로 fpqc 내림에 의해 동형이다. Descent on Spaces,
Lemma \ref{spaces-descent-lemma-descending-property-isomorphism}을
보라.
\end{proof}

\begin{lemma}
\label{spaces-more-morphisms-lemma-flat-surjective-monomorphism}
대수공간의 준콤팩트 평탄 전사 단사 사상은 동형이다.
\end{lemma}

\begin{proof}
이러한 사상은 Lemma \ref{spaces-more-morphisms-lemma-ui-case}의 가정을 만족한다.
\end{proof}







\section{몰입의 여법선층}
\label{spaces-more-morphisms-section-conormal-sheaf}

\noindent
$S$를 스킴이라 하자. $i : Z \to X$를 $S$ 위 대수공간의
폐몰입이라 하자. $\mathcal{I} \subset \mathcal{O}_X$를 이에
대응하는 준연접 아이디얼층이라 하자. Morphisms of Spaces, Lemma
\ref{spaces-morphisms-lemma-closed-immersion-ideals}를 보라.
짧은 완전열
$$
0 \to \mathcal{I}^2 \to \mathcal{I} \to \mathcal{I}/\mathcal{I}^2 \to 0
$$
을 생각하자. 이는 $X$ 위 준연접 층들의 완전열이다. 층
$\mathcal{I}/\mathcal{I}^2$는 $\mathcal{I}$에
의해 영멸되므로 Morphisms of Spaces, Lemma
\ref{spaces-morphisms-lemma-i-star-equivalence}에 의해 $Z$ 위의
층에 대응한다. 이 준연접 $\mathcal{O}_Z$-가군을 {\it $Z$의
$X$ 안에서의 여법선층}이라 하며, Morphisms of Spaces, Section
\ref{spaces-morphisms-section-closed-immersions-quasi-coherent}에서
언급한 표기의 남용에 따라 흔히 $\mathcal{I}/\mathcal{I}^2$로
나타낸다.

\medskip\noindent
$i : Z \to X$가 (국소 닫힌) 몰입인 경우, $i$의 여법선층을
폐몰입 $i : Z \to X \setminus \partial Z$의 여법선층으로
정의한다. Morphisms of Spaces, Remark
\ref{spaces-morphisms-remark-immersion}을 보라. 이는 흔히
$\mathcal{I}/\mathcal{I}^2$로 나타내며, 여기서 $\mathcal{I}$는
폐몰입 $i : Z \to X \setminus \partial Z$의 아이디얼층이다.

\begin{definition}
\label{spaces-more-morphisms-definition-conormal-sheaf}
$i : Z \to X$를 몰입이라 하자. $\mathcal{C}_{Z/X}$를
{\it $Z$의 $X$ 안에서의 여법선층}, 또는 {\it $i$의
여법선층}이라 한다. 이는 위에서 설명한 준연접
$\mathcal{O}_Z$-가군 $\mathcal{I}/\mathcal{I}^2$이다.
\end{definition}

\noindent
\cite[IV Definition 16.1.2]{EGA}에서는 이 층을
$\mathcal{N}_{Z/X}$로 나타낸다. 우리는
$\mathcal{N}_{Z/X}$라는 표기를 {\it 몰입의 법선층}을 위해
남겨 두려 하므로 이 규약을 따르지 않는다. 여법선층이 유한
표시이면 법선층을
$$
\mathcal{N}_{Z/X} =
\SheafHom_{\mathcal{O}_Z}(\mathcal{C}_{Z/X}, \mathcal{O}_Z) =
\SheafHom_{\mathcal{O}_Z}(\mathcal{I}/\mathcal{I}^2, \mathcal{O}_Z)
$$
로 정의한다 (그렇지 않으면 법선층이 준연접이 아닐 수도 있다).
법선층은 나중에 다시 다룰 것이다 (미래의 참고문헌을 여기에
삽입하라).

\begin{lemma}
\label{spaces-more-morphisms-lemma-etale-conormal}
$S$를 스킴이라 하자. $i : Z \to X$를 몰입이라 하자.
$\varphi : U \to X$를 에탈 사상이라 하고, 여기서 $U$는
스킴이라고 하자. $Z_U = U \times_X Z$로 놓으면 이는 $U$의
국소 닫힌 부분스킴이다. 그러면
$$
\mathcal{C}_{Z/X}|_{Z_U} = \mathcal{C}_{Z_U/U}
$$
가 표준적으로 성립하며 $U$에 대하여 함자적이다.
\end{lemma}

\begin{proof}
$T \subset X$를 닫힌 부분공간이라 하여 $i$가
$X \setminus T$ 안의 폐몰입을 정의하게 하자.
$\mathcal{I}$를 $X \setminus T$ 위의 준연접 아이디얼층이라 하자.
이 층은 $Z$를 정의한다. 그러면 이 보조정리는 단지
$\mathcal{I}|_{U \setminus \varphi^{-1}(T)}$가 몰입
$Z_U \to U \setminus \varphi^{-1}(T)$의 아이디얼층이라는
명제이다. 이는 Morphisms of Spaces, Lemma
\ref{spaces-morphisms-lemma-closed-immersion-ideals}에서
$\mathcal{I}$를 구성한 방법으로부터 분명하다.
\end{proof}

\begin{lemma}
\label{spaces-more-morphisms-lemma-conormal-functorial}
$S$를 스킴이라 하자.
$$
\xymatrix{
Z \ar[r]_i \ar[d]_f & X \ar[d]^g \\
Z' \ar[r]^{i'} & X'
}
$$
를 $S$ 위 대수공간의 가환 그림이라 하자. $i$, $i'$가 몰입이라고
가정하자. 그러면 $\mathcal{O}_Z$-가군의 표준 사상
$$
f^*\mathcal{C}_{Z'/X'}
\longrightarrow
\mathcal{C}_{Z/X}
$$
이 있다.
\end{lemma}

\begin{proof}
먼저 열린 부분공간 $U' \subset X'$과 $U \subset X$를 찾아
$g(U) \subset U'$이고, $i(Z) \subset U$와 $i(Z') \subset U'$가
닫히게 하자 (존재성의 증명은 생략한다). $X$를 $U$로, $X'$을
$U'$으로 바꾸면 $i$와 $i'$가 폐몰입이라고 가정할 수 있다.
$\mathcal{I}' \subset \mathcal{O}_{X'}$과
$\mathcal{I} \subset \mathcal{O}_X$를 각각 $i'$과 $i$에 대응하는
준연접 아이디얼층이라 하자. Morphisms of Spaces, Lemma
\ref{spaces-morphisms-lemma-closed-immersion-ideals}를 보라.
합성
$$
g^{-1}\mathcal{I}' \to g^{-1}\mathcal{O}_{X'}
\xrightarrow{g^\sharp} \mathcal{O}_X \to
\mathcal{O}_X/\mathcal{I} = i_*\mathcal{O}_Z
$$
을 생각하자. $g(i(Z)) \subset Z'$이므로 이 합성은 영이다
(Morphisms of Spaces, Lemma
\ref{spaces-morphisms-lemma-closed-immersion-ideals}의 분해에 관한
명제를 보라). 따라서 가환 그림
$$
\xymatrix{
0 \ar[r] &
\mathcal{I} \ar[r] &
\mathcal{O}_X \ar[r] &
i_*\mathcal{O}_Z \ar[r] &
0 \\
0 \ar[r] &
g^{-1}\mathcal{I}' \ar[r] \ar[u] &
g^{-1}\mathcal{O}_{X'} \ar[r] \ar[u] &
g^{-1}i'_*\mathcal{O}_{Z'} \ar[r] \ar[u] &
0
}
$$
을 얻는다. $g^{-1}$은 완전함자이므로 아래 행은 완전하다.
완전성에 의해
$(g^{-1}\mathcal{I}')^2 = g^{-1}((\mathcal{I}')^2)$임도 알 수 있다.
그러므로 이 그림은 사상
$g^{-1}(\mathcal{I}'/(\mathcal{I}')^2) \to
\mathcal{I}/\mathcal{I}^2$를 유도한다. 이를 (예를 들어
$i^{-1}$을 사용하여) $Z$로 당기면
$i^{-1}g^{-1}(\mathcal{I}'/(\mathcal{I}')^2) \to
\mathcal{C}_{Z/X}$를 얻는다.
$i^{-1}g^{-1} = f^{-1}(i')^{-1}$이므로 이는 사상
$f^{-1}\mathcal{C}_{Z'/X'} \to \mathcal{C}_{Z/X}$를 주고, 이
사상이 원하는 사상을 유도한다.
\end{proof}

\begin{lemma}
\label{spaces-more-morphisms-lemma-conormal-functorial-more}
$S$를 스킴이라 하자. Definition \ref{spaces-more-morphisms-definition-conormal-sheaf}의
여법선층과 Lemma
\ref{spaces-more-morphisms-lemma-conormal-functorial}의 함자성은 다음 성질들을
만족한다.
\begin{enumerate}
\item $Z \to X$가 $S$ 위 스킴의 몰입이면, 여법선층은 Morphisms,
Definition \ref{morphisms-definition-conormal-sheaf}의 것과
일치한다.
\item Lemma \ref{spaces-more-morphisms-lemma-conormal-functorial}에서 모든 공간이
스킴이면, 사상
$f^*\mathcal{C}_{Z'/X'} \to \mathcal{C}_{Z/X}$는 Morphisms,
Lemma \ref{morphisms-lemma-conormal-functorial}에서 구성한 것과
같다.
\item 가환 그림
$$
\xymatrix{
Z \ar[r]_i \ar[d]_f & X \ar[d]^g \\
Z' \ar[r]^{i'} \ar[d]_{f'} & X' \ar[d]^{g'} \\
Z'' \ar[r]^{i''} & X''
}
$$
이 주어지면, 사상
$(f' \circ f)^*\mathcal{C}_{Z''/X''} \to \mathcal{C}_{Z/X}$는
$f^*\mathcal{C}_{Z'/X'} \to \mathcal{C}_{Z/X}$와
$f$에 의한 사상
$(f')^*\mathcal{C}_{Z''/X''} \to \mathcal{C}_{Z'/X'}$의 당김을
합성한 것과 같다.
\end{enumerate}
\end{lemma}

\begin{proof}
생략한다. (1)은 Lemma \ref{spaces-more-morphisms-lemma-etale-conormal}의 특수한
경우임을 주목하자.
\end{proof}

\begin{lemma}
\label{spaces-more-morphisms-lemma-conormal-functorial-flat}
$S$를 스킴이라 하자.
$$
\xymatrix{
Z \ar[r]_i \ar[d]_f & X \ar[d]^g \\
Z' \ar[r]^{i'} & X'
}
$$
를 $S$ 위 대수공간의 올곱 그림이라 하자. $i$, $i'$가 몰입이라고
가정하자. 그러면 Lemma \ref{spaces-more-morphisms-lemma-conormal-functorial}의 표준 사상
$f^*\mathcal{C}_{Z'/X'} \to \mathcal{C}_{Z/X}$는 전사이다.
$g$가 평탄하면 이 사상은 동형이다.
\end{lemma}

\begin{proof}
가환 그림
$$
\xymatrix{
U \ar[r] \ar[d] & X \ar[d] \\
U' \ar[r] & X'
}
$$
을 택하자. 여기서 $U$, $U'$은 스킴이고 수평 화살표들은 전사
에탈 사상이다. Spaces, Lemma
\ref{spaces-lemma-lift-morphism-presentations}를 보라. 이제 Lemmas
\ref{spaces-more-morphisms-lemma-etale-conormal}와
\ref{spaces-more-morphisms-lemma-conormal-functorial-more}를 사용하면 문제는 스킴의
사상인 경우로 환원된다. 스킴의 경우에는 Morphisms, Lemma
\ref{morphisms-lemma-conormal-functorial-flat}가 그 결과이다.
\end{proof}

\begin{lemma}
\label{spaces-more-morphisms-lemma-transitivity-conormal}
$S$를 스킴이라 하자. $Z \to Y \to X$를 대수공간의 몰입들이라
하자. 그러면 표준 완전열
$$
i^*\mathcal{C}_{Y/X} \to
\mathcal{C}_{Z/X} \to
\mathcal{C}_{Z/Y} \to 0
$$
이 있다. 여기서 사상들은 Lemma
\ref{spaces-more-morphisms-lemma-conormal-functorial}에서 오고, $i : Z \to Y$는 첫째
사상이다.
\end{lemma}

\begin{proof}
$U$를 스킴이라 하고 $U \to X$를 전사 에탈 사상이라 하자.
Lemmas \ref{spaces-more-morphisms-lemma-etale-conormal}와
\ref{spaces-more-morphisms-lemma-conormal-functorial-more}를 통하여, 이 열의 완전성은
스킴의 몰입들
$Z \times_X U \to Y \times_X U \to U$에 대한 대응하는 열의
완전성으로 곧바로 옮겨진다. 따라서 이 보조정리는 Morphisms,
Lemma \ref{morphisms-lemma-transitivity-conormal}에서 따른다.
\end{proof}









\section{몰입의 법선뿔}
\label{spaces-more-morphisms-section-normal-cone}

\noindent
$S$를 스킴이라 하자. $i : Z \to X$를 $S$ 위 대수공간의
폐몰입이라 하자. $\mathcal{I} \subset \mathcal{O}_X$를 이에
대응하는 준연접 아이디얼층이라 하자. Morphisms of Spaces, Lemma
\ref{spaces-morphisms-lemma-closed-immersion-ideals}를 보라.
등급 $\mathcal{O}_X$-대수들의 준연접 층
$\bigoplus_{n \geq 0} \mathcal{I}^n/\mathcal{I}^{n + 1}$을
생각하자. 각 층 $\mathcal{I}^n/\mathcal{I}^{n + 1}$은
$\mathcal{I}$에 의해 영멸되므로, Morphisms of Spaces, Lemma
\ref{spaces-morphisms-lemma-i-star-equivalence}에 의해 이 등급
대수는 등급 $\mathcal{O}_Z$-대수들의 준연접 층에 대응한다.
이 준연접 등급 $\mathcal{O}_Z$-대수를 {\it $Z$의 $X$ 안에서의
여법선 대수}라 하며, Morphisms of Spaces, Section
\ref{spaces-morphisms-section-closed-immersions-quasi-coherent}에서
언급한 표기의 남용에 따라 흔히 단순히
$\bigoplus_{n \geq 0} \mathcal{I}^n/\mathcal{I}^{n + 1}$로
나타낸다.

\medskip\noindent
$i : Z \to X$가 (국소 닫힌) 몰입인 경우, $i$의 여법선 대수를
폐몰입 $i : Z \to X \setminus \partial Z$의 여법선 대수로
정의한다. Morphisms of Spaces, Remark
\ref{spaces-morphisms-remark-immersion}을 보라. 이는 흔히
$\bigoplus_{n \geq 0} \mathcal{I}^n/\mathcal{I}^{n + 1}$로
나타내며, 여기서 $\mathcal{I}$는 폐몰입
$i : Z \to X \setminus \partial Z$의 아이디얼층이다.

\begin{definition}
\label{spaces-more-morphisms-definition-conormal-algebra}
$i : Z \to X$를 몰입이라 하자. $\mathcal{C}_{Z/X, *}$를
{\it $Z$의 $X$ 안에서의 여법선 대수}, 또는 {\it $i$의 여법선
대수}라 한다. 이는 위에서 설명한 등급 $\mathcal{O}_Z$-대수들의
준연접 층
$\bigoplus_{n \geq 0} \mathcal{I}^n/\mathcal{I}^{n + 1}$이다.
\end{definition}

\noindent
따라서 $\mathcal{C}_{Z/X, 1} = \mathcal{C}_{Z/X}$은 몰입의
여법선층이다. 또한
$\mathcal{C}_{Z/X, 0} = \mathcal{O}_Z$이고,
$\mathcal{C}_{Z/X, n}$은 다음 성질로 특징지어지는 준연접
$\mathcal{O}_Z$-가군이다.
\begin{equation}
\label{spaces-more-morphisms-equation-conormal-in-degree-n}
i_*\mathcal{C}_{Z/X, n} = \mathcal{I}^n/\mathcal{I}^{n + 1}
\end{equation}
여기서 $i : Z \to X \setminus \partial Z$이고 $\mathcal{I}$는
위와 같이 $i$의 아이디얼층이다. 끝으로 준연접 등급
$\mathcal{O}_Z$-대수의 표준 전사 사상
\begin{equation}
\label{spaces-more-morphisms-equation-conormal-algebra-quotient}
\text{Sym}^*(\mathcal{C}_{Z/X}) \longrightarrow \mathcal{C}_{Z/X, *}
\end{equation}
이 있으며, 차수 $0$과 $1$에서는 동형임을 주목하자.

\begin{lemma}
\label{spaces-more-morphisms-lemma-etale-conormal-algebra}
$S$를 스킴이라 하자. $i : Z \to X$를 $S$ 위 대수공간의 몰입이라
하자. $\varphi : U \to X$를 에탈 사상이라 하고, 여기서 $U$는
스킴이라고 하자. $Z_U = U \times_X Z$로 놓으면 이는 $U$의 국소
닫힌 부분스킴이다. 그러면
$$
\mathcal{C}_{Z/X, *}|_{Z_U} = \mathcal{C}_{Z_U/U, *}
$$
가 표준적으로 성립하며 $U$에 대하여 함자적이다.
\end{lemma}

\begin{proof}
$T \subset X$를 닫힌 부분공간이라 하여 $i$가
$X \setminus T$ 안의 폐몰입을 정의하게 하자. $\mathcal{I}$를
$X \setminus T$ 위의 준연접 아이디얼층이라 하자. 이 층은 $Z$를
정의한다. 그러면 이 보조정리는
$\mathcal{I}|_{U \setminus \varphi^{-1}(T)}$가 몰입
$Z_U \to U \setminus \varphi^{-1}(T)$의 아이디얼층이라는
사실에서 따른다. 이는 Morphisms of Spaces, Lemma
\ref{spaces-morphisms-lemma-closed-immersion-ideals}에서
$\mathcal{I}$를 구성한 방법으로부터 분명하다.
\end{proof}

\begin{lemma}
\label{spaces-more-morphisms-lemma-conormal-algebra-functorial}
$S$를 스킴이라 하자.
$$
\xymatrix{
Z \ar[r]_i \ar[d]_f & X \ar[d]^g \\
Z' \ar[r]^{i'} & X'
}
$$
를 $S$ 위 대수공간의 가환 그림이라 하자. $i$, $i'$가 몰입이라고
가정하자. 그러면 등급 $\mathcal{O}_Z$-대수의 표준 사상
$$
f^*\mathcal{C}_{Z'/X', *}
\longrightarrow
\mathcal{C}_{Z/X, *}
$$
이 있다.
\end{lemma}

\begin{proof}
먼저 열린 부분공간 $U' \subset X'$과 $U \subset X$를 찾아
$g(U) \subset U'$이고, $i(Z) \subset U$와 $i(Z') \subset U'$가
닫히게 하자 (존재성의 증명은 생략한다). $X$를 $U$로, $X'$을
$U'$으로 바꾸면 $i$와 $i'$가 폐몰입이라고 가정할 수 있다.
$\mathcal{I}' \subset \mathcal{O}_{X'}$과
$\mathcal{I} \subset \mathcal{O}_X$를 각각 $i'$과 $i$에 대응하는
준연접 아이디얼층이라 하자. Morphisms of Spaces, Lemma
\ref{spaces-morphisms-lemma-closed-immersion-ideals}를 보라.
합성
$$
g^{-1}\mathcal{I}' \to g^{-1}\mathcal{O}_{X'}
\xrightarrow{g^\sharp} \mathcal{O}_X \to
\mathcal{O}_X/\mathcal{I} = i_*\mathcal{O}_Z
$$
을 생각하자. $g(i(Z)) \subset Z'$이므로 이 합성은 영이다
(Morphisms of Spaces, Lemma
\ref{spaces-morphisms-lemma-closed-immersion-ideals}의 분해에 관한
명제를 보라). 따라서 가환 그림
$$
\xymatrix{
0 \ar[r] &
\mathcal{I} \ar[r] &
\mathcal{O}_X \ar[r] &
i_*\mathcal{O}_Z \ar[r] &
0 \\
0 \ar[r] &
g^{-1}\mathcal{I}' \ar[r] \ar[u] &
g^{-1}\mathcal{O}_{X'} \ar[r] \ar[u] &
g^{-1}i'_*\mathcal{O}_{Z'} \ar[r] \ar[u] &
0
}
$$
을 얻는다. $g^{-1}$은 완전함자이므로 아래 행은 완전하다.
완전성에 의해
$(g^{-1}\mathcal{I}')^n = g^{-1}((\mathcal{I}')^n)$임도 알 수 있다.
여기서 $n \geq 1$이다.
그러므로 이 그림은 사상
$g^{-1}((\mathcal{I}')^n/(\mathcal{I}')^{n + 1}) \to
\mathcal{I}^n/\mathcal{I}^{n + 1}$을 유도한다. 이를 (예를 들어
$i^{-1}$을 사용하여) $Z$로 당기면
$i^{-1}g^{-1}((\mathcal{I}')^n/(\mathcal{I}')^{n + 1}) \to
\mathcal{C}_{Z/X, n}$을 얻는다.
$i^{-1}g^{-1} = f^{-1}(i')^{-1}$이므로 이는 사상들
$f^{-1}\mathcal{C}_{Z'/X', n} \to \mathcal{C}_{Z/X, n}$을 주고,
이들이 원하는 사상을 유도한다.
\end{proof}

\begin{lemma}
\label{spaces-more-morphisms-lemma-conormal-algebra-functorial-flat}
$S$를 스킴이라 하자.
$$
\xymatrix{
Z \ar[r]_i \ar[d]_f & X \ar[d]^g \\
Z' \ar[r]^{i'} & X'
}
$$
를 $S$ 위 대수공간의 데카르트 사각형이라 하고 $i$, $i'$가
몰입이라고 하자. 그러면 Lemma
\ref{spaces-more-morphisms-lemma-conormal-algebra-functorial}의 표준 사상
$f^*\mathcal{C}_{Z'/X', *} \to \mathcal{C}_{Z/X, *}$는 전사이다.
$g$가 평탄하면 이 사상은 동형이다.
\end{lemma}

\begin{proof}
$X'$ 위에서 에탈 국소화한 뒤 명제를 확인해도 된다. 이 경우
$X' \to X$가 스킴의 사상이라고 가정할 수 있고, 따라서 $Z$와
$Z'$은 스킴이다. 결과는 스킴의 경우인 Divisors, Lemma
\ref{divisors-lemma-conormal-algebra-functorial-flat}에서 따른다.
\end{proof}

\noindent
대수공간 위의 뿔과 벡터다발에 대해서도 스킴의 경우와 같은
규약을 쓴다 (스킴의 경우에는 EGA의 규약을 쓴다). Constructions,
Sections \ref{constructions-section-cone}와
\ref{constructions-section-vector-bundle}을 보라. 특히
벡터다발은 매우 일반적인 대상이다 (그리고 아핀공간 다발과
국소적으로 동형일 필요가 없다).

\begin{definition}
\label{spaces-more-morphisms-definition-normal-cone}
$S$를 스킴이라 하자. $i : Z \to X$를 $S$ 위 대수공간의
몰입이라 하자. $C_ZX$를 $Z$의 $X$ 안에서의 {\it 법선뿔}이라
하며, 이는
$$
C_ZX = \underline{\Spec}_Z(\mathcal{C}_{Z/X, *})
$$
이다. Morphisms of Spaces, Definition
\ref{spaces-morphisms-definition-relative-spec}을 보라. $Z$의
$X$ 안에서의 {\it 법선다발}은 벡터다발
$$
N_ZX = \underline{\Spec}_Z(\text{Sym}(\mathcal{C}_{Z/X}))
$$
이다.
\end{definition}

\noindent
따라서 $C_ZX \to Z$는 $Z$ 위의 뿔이고, $N_ZX \to Z$는 $Z$ 위의
벡터다발이다. 더욱이 등급 대수의 표준 전사
(\ref{spaces-more-morphisms-equation-conormal-algebra-quotient})는 $Z$ 위 뿔의 표준
폐몰입
\begin{equation}
\label{spaces-more-morphisms-equation-normal-cone-in-normal-bundle}
C_ZX \longrightarrow N_ZX
\end{equation}
을 정의한다.










\section{사상의 미분층}
\label{spaces-more-morphisms-section-sheaf-differentials}

\noindent
독자는 가환대수 장의 대응하는 절
(Algebra, Section \ref{algebra-section-differentials}), 스킴의
사상 장의 대응하는 절
(Morphisms, Section \ref{morphisms-section-sheaf-differentials}),
그리고 Modules on Sites, Section
\ref{sites-modules-section-differentials}을 살펴보기를 권한다.
먼저 스킴의 사상에 대한 미분층이라는 개념이 대응하는 작은 에탈
(환 달린) 자리의 사상에 대한 개념과 일치함을 보인다.

\medskip\noindent
다음 보조정리를 명확하게 진술하기 위하여 잠시 다음 표기로
돌아간다. $\mathcal{F}^a$로 $\mathcal{O}_{X_\etale}$-가군의
층을 나타낸다. 이 층은 준연접 $\mathcal{O}_X$-가군
$\mathcal{F}$에 대응하며, 그 가군은 스킴 $X$ 위에 있다. Descent,
Definition \ref{descent-definition-structure-sheaf}를 보라.

\begin{lemma}
\label{spaces-more-morphisms-lemma-match-modules-differentials}
$f : X \to Y$를 스킴의 사상이라 하자.
$f_{small} : X_\etale \to Y_\etale$을 이에 대응하는 작은 에탈
자리의 사상이라 하자. Descent, Remark
\ref{descent-remark-change-topologies-ringed}를 보라. 그러면 보편
미분과 양립하는 표준 동형
$$
(\Omega_{X/Y})^a = \Omega_{X_\etale/Y_\etale}
$$
이 있다. 여기서 첫째 가군은 $X_\etale$ 위의 층으로, 준연접
$\mathcal{O}_X$-가군 $\Omega_{X/Y}$에 대응한다. Morphisms,
Definition \ref{morphisms-definition-sheaf-differentials}를 보라.
둘째 가군은 Modules on Sites, Definition
\ref{sites-modules-definition-module-differentials}의 가군이다.
\end{lemma}

\begin{proof}
$h : U \to X$를 에탈 사상이라 하자. 이 경우 자연 사상
$h^*\Omega_{X/Y} \to \Omega_{U/Y}$는 동형이다. More on Morphisms,
Lemma \ref{more-morphisms-lemma-sheaf-differentials-etale-localization}을
보라. 따라서 자연스러운 $\mathcal{O}_{Y_\etale}$-미분
$$
\text{d}^a : \mathcal{O}_{X_\etale} \longrightarrow (\Omega_{X/Y})^a
$$
이 있다. 실제로 방금 $(\Omega_{X/Y})^a$의 임의의 대상 $U$에서의
값을 생각했는데, 대상의 범주는 $X_\etale$이며 이 값을
$\Gamma(U, \Omega_{U/Y})$와 표준적으로 동일시하였다. 보편 미분
$\text{d}_{X/Y} :
\mathcal{O}_{X_\etale}
\to
\Omega_{X_\etale/Y_\etale}$의 보편 성질에 의해
$\mathcal{O}_{X_\etale}$-선형 사상
$c : \Omega_{X_\etale/Y_\etale} \to (\Omega_{X/Y})^a$가 유일하게
존재한다. 이 사상은 $\text{d}^a = c \circ \text{d}_{X/Y}$를
만족한다.

\medskip\noindent
역으로 $\mathcal{F}$가 $\mathcal{O}_{X_\etale}$-가군이고
$D : \mathcal{O}_{X_\etale} \to \mathcal{F}$가
$\mathcal{O}_{Y_\etale}$-미분이라고 하자. 그러면 $D$를 단순히
작은 자리스키 자리 $X_{Zar}$로 제한할 수 있다. 이 자리는 $X$의
작은 자리스키 자리이다.
$X_{Zar}$ 위의 층은 $X$ 위의 층과 일치하므로 (Descent, Remark
\ref{descent-remark-Zariski-site-space}),
$D|_{X_{Zar}} : \mathcal{O}_X \to \mathcal{F}|_{X_{Zar}}$는
통상적인 $Y$-미분이다. 따라서 사상
$\psi : \Omega_{X/Y} \longrightarrow \mathcal{F}|_{X_{Zar}}$를
얻으며, 이는 $D|_{X_{Zar}} = \psi \circ \text{d}$를 만족한다.
특히 이를 $\mathcal{F} = \Omega_{X_\etale/Y_\etale}$에 적용하면
사상
$$
c' :
\Omega_{X/Y}
\longrightarrow
\Omega_{X_\etale/Y_\etale}|_{X_{Zar}}
$$
을 얻는다. Descent, Remark
\ref{descent-remark-change-topologies-ringed}와 Lemma
\ref{descent-lemma-compare-sites}에서 논의한 환 달린 자리의 사상
$\text{id}_{small, \etale, Zar} : X_\etale \to X_{Zar}$을
생각하자. 제한 함자
$\mathcal{F} \mapsto \mathcal{F}|_{X_{Zar}}$는
$\text{id}_{small, \etale, Zar, *}$와 같고,
$\text{id}_{small, \etale, Zar}^*$는
$\text{id}_{small, \etale, Zar, *}$의 왼쪽 수반이며,
$(\Omega_{X/Y})^a =
\text{id}_{small, \etale, Zar}^*\Omega_{X/Y}$이므로, $c'$은 사상
$$
c'' :
(\Omega_{X/Y})^a
\longrightarrow
\Omega_{X_\etale/Y_\etale}.
$$
에 수반된다. $c''$과 $c'$은 서로 역이라고 주장한다. 이 주장은
보조정리의 증명을 끝낸다. 이를 보이려면
$c''(\text{d}(f)) = \text{d}_{X/Y}(f)$이고
$c(\text{d}_{X/Y}(f)) = \text{d}(f)$임을 확인하면 충분하다.
여기서 $f$는 $\mathcal{O}_X$의 국소 절단이며 $X$의 열린 부분
위에서 주어져 있다. 확인은 생략한다.
\end{proof}

\noindent
이제 다음 정의를 내릴 수 있다. 다른 정의는 Remark
\ref{spaces-more-morphisms-remark-alternative}를 보라.

\begin{definition}
\label{spaces-more-morphisms-definition-sheaf-differentials}
$S$를 스킴이라 하자. $f : X \to Y$를 $S$ 위 대수공간의
사상이라 하자. $\Omega_{X/Y}$를 {\it $X$의 $Y$ 위의 미분층}이라
하며, 이는 미분층이다 (Modules on Sites, Definition
\ref{sites-modules-definition-sheaf-differentials}). 구체적으로
환 달린 토포스 사상
$$
(f_{small}, f^\sharp) :
(X_\etale, \mathcal{O}_X)
\to
(Y_\etale, \mathcal{O}_Y)
$$
에 대한 것이며, 이 사상은 Properties of Spaces, Lemma
\ref{spaces-properties-lemma-morphism-ringed-topoi}의 사상이다.
{\it 보편 $Y$-미분}은
$\text{d}_{X/Y} : \mathcal{O}_X \to \Omega_{X/Y}$로 나타낸다.
\end{definition}

\noindent
Lemma \ref{spaces-more-morphisms-lemma-match-modules-differentials}에 의해 $X$와 $Y$가
표현 가능할 때 이 정의는 이미 있던 개념과 충돌하지 않는다.
이제부터 $X$와 $Y$가 표현 가능하면 위에서 정의한 미분층과
Morphisms, Definition
\ref{morphisms-definition-sheaf-differentials}에서 정의한 미분층을
구별하지 않는다. 이를 스킴의 사상에 대한 통상적인 미분 가군과
연결하고자 한다. 핵심 보조정리는 다음과 같다.

\begin{lemma}
\label{spaces-more-morphisms-lemma-localize-differentials}
$S$를 스킴이라 하자. $f : X \to Y$를 $S$ 위 대수공간의
사상이라 하자. 임의의 가환 그림
$$
\xymatrix{
U \ar[d]_a \ar[r]_\psi & V \ar[d]^b \\
X \ar[r]^f & Y
}
$$
을 생각하자. 수직 화살표들은 대수공간의 에탈 사상이라고 하자.
그러면
$$
\Omega_{X/Y}|_{U_\etale} = \Omega_{U/V}
$$
이다. 특히 $U$, $V$가 스킴이면 이는 스킴의 사상 $U \to V$의
통상적인 미분층과 같다.
\end{lemma}

\begin{proof}
Properties of Spaces, Lemma
\ref{spaces-properties-lemma-etale-morphism-topoi}와 Equation
(\ref{spaces-properties-equation-restrict})에 의해
$X_\etale$ 위 층의 $U_\etale$로의 제한을 $a_{small}$에 의한
당김으로 생각할 수 있다. $b$에 대해서도 마찬가지이다.
Modules on Sites, Lemma
\ref{sites-modules-lemma-localize-differentials}에 의해
$$
\Omega_{X/Y}|_{U_\etale} =
\Omega_{\mathcal{O}_{U_\etale}/
a_{small}^{-1}f_{small}^{-1}\mathcal{O}_{Y_\etale}}
$$
이다. 그리고
$a_{small}^{-1}f_{small}^{-1}\mathcal{O}_{Y_\etale}
= \psi_{small}^{-1}b_{small}^{-1}\mathcal{O}_{Y_\etale}
= \psi_{small}^{-1}\mathcal{O}_{V_\etale}$이므로 보조정리가
성립한다.
\end{proof}

\begin{lemma}
\label{spaces-more-morphisms-lemma-module-differentials-quasi-coherent}
$S$를 스킴이라 하자. $f : X \to Y$를 $S$ 위 대수공간의
사상이라 하자. 그러면 $\Omega_{X/Y}$는 준연접
$\mathcal{O}_X$-가군이다.
\end{lemma}

\begin{proof}
Lemma \ref{spaces-more-morphisms-lemma-localize-differentials}와 같은 그림을 택하되,
$a$와 $b$가 전사이고 $U$와 $V$가 스킴이게 하자. 그러면
$\Omega_{X/Y}|_U = \Omega_{U/V}$이며, 이는 준연접이다 (예를 들어
Morphisms, Lemma \ref{morphisms-lemma-differentials-diagonal}에
의해). 따라서 Properties of Spaces, Lemma
\ref{spaces-properties-lemma-characterize-quasi-coherent}에 의해
$\Omega_{X/Y}$가 준연접임을 얻는다.
\end{proof}

\begin{remark}
\label{spaces-more-morphisms-remark-alternative}
이제 $\Omega_{X/Y}$가 준연접임을 아니, 이를 다른 방법으로
구성해 볼 수 있다. 예를 들어 Properties of Spaces, Section
\ref{spaces-properties-section-quasi-coherent-presentation}의
결과를 사용하여 미분층을 붙임으로 구성할 수 있다. 예컨대 $Y$가
스킴이고 $U \to X$가 스킴에서 $X$로 가는 전사 에탈 사상이면,
$\Omega_{U/Y}$는 준연접 $\mathcal{O}_U$-가군이다. 또한
$s, t : R \to U$가 에탈이므로 Morphisms, Lemma
\ref{morphisms-lemma-triangle-differentials-smooth}를 사용하여 동형
$$
\alpha : s^*\Omega_{U/Y} \to \Omega_{R/Y} \to t^*\Omega_{U/Y}
$$
을 얻는다. 이것이 코사이클 조건을 만족함을 확인하면 끝난다.
$Y$가 스킴이 아니면 $\Omega_{U/Y}$를 사상
$(U \to Y)^*\Omega_{Y/S} \to \Omega_{U/S}$의 여핵으로 정의하고
앞과 같이 진행한다. 이 두 단계 과정은 조금 볼품없다. 또 다른
방법은 Lemma \ref{spaces-more-morphisms-lemma-localize-differentials}와 같은 임의의
그림에 대하여 층들 $\Omega_{U/V}$를 붙이는 것이지만, 이것도
그다지 우아하지 않다. 그래도 두 방법 모두 작동하며 미분층의
조금 더 초등적인 구성을 준다.
\end{remark}

\begin{lemma}
\label{spaces-more-morphisms-lemma-functoriality-differentials}
$S$를 스킴이라 하자.
$$
\xymatrix{
X' \ar[d] \ar[r]_f & X \ar[d] \\
Y' \ar[r] & Y
}
$$
를 대수공간의 가환 그림이라 하자. 사상
$f^\sharp : \mathcal{O}_X \to f_*\mathcal{O}_{X'}$와 사상
$f_*\text{d}_{X'/Y'} : f_*\mathcal{O}_{X'} \to
f_*\Omega_{X'/Y'}$의 합성은 $Y$-미분이다. 따라서 표준
$\mathcal{O}_X$-가군 사상
$\Omega_{X/Y} \to f_*\Omega_{X'/Y'}$를 얻는다. 그리고 $f_*$와
$f^*$의 수반성에 의해 표준 $\mathcal{O}_{X'}$-가군 준동형
$$
c_f : f^*\Omega_{X/Y} \longrightarrow \Omega_{X'/Y'}.
$$
을 얻는다. 이는 $f^*\text{d}_{X/Y}(t)$가
$\text{d}_{X'/Y'}(f^* t)$로 사상된다는 성질로 유일하게
특징지어진다. 여기서 $t$는 $\mathcal{O}_X$의 임의의 국소
절단이다.
\end{lemma}

\begin{proof}
이는 Modules on Sites, Lemma
\ref{sites-modules-lemma-functoriality-differentials-ringed-topoi}의
특수한 경우이다.
\end{proof}

\begin{lemma}
\label{spaces-more-morphisms-lemma-check-functoriality-differentials}
$S$를 스킴이라 하자.
$$
\xymatrix{
X'' \ar[d] \ar[r]_g & X' \ar[d] \ar[r]_f & X \ar[d] \\
Y'' \ar[r] & Y' \ar[r] & Y
}
$$
를 $S$ 위 대수공간의 가환 그림이라 하자. 그러면
$$
c_{f \circ g} = c_g \circ g^* c_f
$$
가 $(f \circ g)^*\Omega_{X/Y} \to \Omega_{X''/Y''}$인 사상으로서
성립한다.
\end{lemma}

\begin{proof}
생략한다. 힌트: 국소 절단에 미치는 작용으로
$c_f, c_g, c_{f \circ g}$를 특징짓는 것을 사용하라.
\end{proof}

\begin{lemma}
\label{spaces-more-morphisms-lemma-triangle-differentials}
$S$를 스킴이라 하자. $f : X \to Y$, $g : Y \to B$를 $S$ 위
대수공간의 사상들이라 하자. 그러면 표준 완전열
$$
f^*\Omega_{Y/B} \to \Omega_{X/B} \to \Omega_{X/Y} \to 0
$$
이 있으며, 사상들은 Lemma
\ref{spaces-more-morphisms-lemma-functoriality-differentials}의 적용에서 온다.
\end{lemma}

\begin{proof}
에탈 국소화를 통해 이 결과의 스킴 판인 Morphisms, Lemma
\ref{morphisms-lemma-triangle-differentials}에서 따른다. Lemma
\ref{spaces-more-morphisms-lemma-localize-differentials}를 보라.
\end{proof}

\begin{lemma}
\label{spaces-more-morphisms-lemma-immersion-differentials}
$S$를 스킴이라 하자. $X \to Y$가 $S$ 위 대수공간의 몰입이면
$\Omega_{X/S}$는 영이다. % P09-ERR-0519
\end{lemma}

\begin{proof}
에탈 국소화를 통해 이 결과의 스킴 판인 Morphisms, Lemma
\ref{morphisms-lemma-immersion-differentials}에서 따른다. Lemma
\ref{spaces-more-morphisms-lemma-localize-differentials}를 보라.
\end{proof}

\begin{lemma}
\label{spaces-more-morphisms-lemma-differentials-relative-immersion}
$S$를 스킴이라 하자. $B$를 $S$ 위의 대수공간이라 하자.
$i : Z \to X$를 $B$ 위 대수공간의 몰입이라 하자. 표준 완전열
$$
\mathcal{C}_{Z/X} \to i^*\Omega_{X/B} \to \Omega_{Z/B} \to 0
$$
이 있다. 첫째 화살표는 $\text{d}_{X/B}$가 유도하고, 둘째
화살표는 Lemma \ref{spaces-more-morphisms-lemma-functoriality-differentials}에서 온다.
\end{lemma}

\begin{proof}
이는 Morphisms, Lemma
\ref{morphisms-lemma-differentials-relative-immersion}의 대수공간
판이며, Lemmas \ref{spaces-more-morphisms-lemma-localize-differentials}와
\ref{spaces-more-morphisms-lemma-etale-conormal}에서 보는 것처럼 에탈 국소화를 통해
그 보조정리에서 따를 것이다. 다만 첫째 화살표를 전역적으로
정의할 수 있음을 확인해야 한다. 따라서 여기서
‘$\text{d}_{X/B}$가 유도한다’의 의미를 설명한다. 즉, $i$가
폐몰입이라고 가정할 수 있다. 이를 위해 $X$를 열린 부분공간으로
바꾼다.
$\mathcal{I} \subset \mathcal{O}_X$를 $Z \subset X$에 대응하는
준연접 아이디얼층이라 하자. 그러면 % P09-ERR-0520
$\text{d}_{X/S} : \mathcal{I} \to \Omega_{X/S}$는 부분층
$\mathcal{I}^2 \subset \mathcal{I}$를
$\mathcal{I}\Omega_{X/S}$로 사상한다. 따라서 사상
$\mathcal{I}/\mathcal{I}^2 \to
\Omega_{X/S}/\mathcal{I}\Omega_{X/S}$를 유도하며, 이 사상은
$\mathcal{O}_X/\mathcal{I}$-선형이다. Morphisms of Spaces,
Lemma \ref{spaces-morphisms-lemma-i-star-equivalence}에 의해 이는
원하는 사상
$\mathcal{C}_{Z/X} \to i^*\Omega_{X/S}$에 대응한다.
\end{proof}

\begin{lemma}
\label{spaces-more-morphisms-lemma-differentials-relative-immersion-section}
$S$를 스킴이라 하자. $B$를 $S$ 위의 대수공간이라 하자.
$i : Z \to X$를 $B$ 위 대수공간의 몰입이라 하고, $i$가 (에탈
국소적으로) 왼쪽 역을 가진다고 가정하자. 그러면 Lemma
\ref{spaces-more-morphisms-lemma-differentials-relative-immersion}의 표준 열
$$
0 \to \mathcal{C}_{Z/X} \to i^*\Omega_{X/B} \to \Omega_{Z/B} \to 0
$$
은 (에탈 국소적으로) 분할 완전하다.
\end{lemma}

\begin{proof}
명확히 말하면 다음을 주장한다. $g : X \to Z$가 $i$의 왼쪽
역이고 이것이 $B$ 위에서 성립하면, $i^*c_g$는 사상
$i^*\Omega_{X/B} \to \Omega_{Z/B}$의 오른쪽 역이다. 이제
Lemmas \ref{spaces-more-morphisms-lemma-localize-differentials}와
\ref{spaces-more-morphisms-lemma-etale-conormal}에서 보는 것처럼 에탈 국소화를 통하여
스킴의 사상에 대한 대응하는 결과에서 결론이 따른다.
\end{proof}

\begin{lemma}
\label{spaces-more-morphisms-lemma-base-change-differentials}
$S$를 스킴이라 하자. $X \to Y$를 $S$ 위 대수공간의 사상이라
하자. $g : Y' \to Y$를 $S$ 위 대수공간의 사상이라 하자.
$X' = X_{Y'}$를 $X$의 밑변환이라 하고, 사영을
$g' : X' \to X$로 나타내자. 그러면 Lemma
\ref{spaces-more-morphisms-lemma-functoriality-differentials}의 사상
$$
(g')^*\Omega_{X/Y} \to \Omega_{X'/Y'}
$$
은 동형이다.
\end{lemma}

\begin{proof}
스킴 판인 Morphisms, Lemma
\ref{morphisms-lemma-base-change-differentials}와 에탈 국소화인
Lemma \ref{spaces-more-morphisms-lemma-localize-differentials}에서 따른다.
\end{proof}

\begin{lemma}
\label{spaces-more-morphisms-lemma-differential-product}
$S$를 스킴이라 하자. $f : X \to B$와 $g : Y \to B$를 공역이
같은 $S$ 위 대수공간의 사상들이라 하자. 사영 사상들을
$p : X \times_B Y \to X$와 $q : X \times_B Y \to Y$라 하자.
Lemma \ref{spaces-more-morphisms-lemma-functoriality-differentials}의 사상들은 동형
$$
p^*\Omega_{X/B} \oplus q^*\Omega_{Y/B}
\longrightarrow
\Omega_{X \times_B Y/B}
$$
을 준다.
\end{lemma}

\begin{proof}
스킴 판인 Morphisms, Lemma
\ref{morphisms-lemma-differential-product}와 에탈 국소화인 Lemma
\ref{spaces-more-morphisms-lemma-localize-differentials}에서 따른다.
\end{proof}

\begin{lemma}
\label{spaces-more-morphisms-lemma-finite-type-differentials}
$S$를 스킴이라 하자. $f : X \to Y$를 $S$ 위 대수공간의
사상이라 하자. $f$가 국소 유한형이면 $\Omega_{X/Y}$는 유한형
$\mathcal{O}_X$-가군이다.
\end{lemma}

\begin{proof}
스킴 판인 Morphisms, Lemma
\ref{morphisms-lemma-finite-type-differentials}와 에탈 국소화인
Lemma \ref{spaces-more-morphisms-lemma-localize-differentials}에서 따른다.
\end{proof}

\begin{lemma}
\label{spaces-more-morphisms-lemma-finite-presentation-differentials}
$S$를 스킴이라 하자. $f : X \to Y$를 $S$ 위 대수공간의
사상이라 하자. $f$가 국소 유한 표시이면 $\Omega_{X/Y}$는 유한
표시 $\mathcal{O}_X$-가군이다.
\end{lemma}

\begin{proof}
스킴 판인 Morphisms, Lemma
\ref{morphisms-lemma-finite-presentation-differentials}와 에탈
국소화인 Lemma \ref{spaces-more-morphisms-lemma-localize-differentials}에서 따른다.
\end{proof}

\begin{lemma}
\label{spaces-more-morphisms-lemma-smooth-omega-finite-locally-free}
$S$를 스킴이라 하자. $f : X \to Y$를 $S$ 위 대수공간의 매끄러운
사상이라 하자. 그러면 미분 가군 $\Omega_{X/Y}$는 유한 국소
자유이다.
\end{lemma}

\begin{proof}
Lemma \ref{spaces-more-morphisms-lemma-localize-differentials}에 의해 이 명제는 $X$와
$Y$ 위에서 에탈 국소적이다. 따라서 스킴의 경우인 Morphisms,
Lemma \ref{morphisms-lemma-smooth-omega-finite-locally-free}에서
따른다.
\end{proof}























\section{에탈 자리의 위상적 불변성}
\label{spaces-more-morphisms-section-topological-invariance}

\noindent
자리 $X_{spaces, \etale}$가 ‘위상적 불변량’임을 보인다. 그러면
$X_\etale$도 위상적 불변량이다. 이는 $X_{spaces, \etale}$의
표현 가능한 대상들로 이루어진다. Lemma
\ref{spaces-more-morphisms-lemma-topological-invariance}를 보라.

\begin{theorem}
\label{spaces-more-morphisms-theorem-topological-invariance}
$S$를 스킴이라 하자. $f : X \to Y$를 $S$ 위 대수공간의
사상이라 하자. $f$가 정수적이고 보편 단사이며 전사라고 가정하자.
함자
$$
V \longmapsto V_X = X \times_Y V
$$
는 범주의 동치
$Y_{spaces, \etale} \to X_{spaces, \etale}$를 정의한다.
\end{theorem}

\begin{proof}
$f$는 표현 가능하고 보편 위상동형이다. Morphisms of Spaces,
Section \ref{spaces-morphisms-section-universal-homeomorphisms}를
보라.

\medskip\noindent
먼저 이 함자가 충실함을 증명하자. $V', V$가
$Y_{spaces, \etale}$의 대상이고
$a, b : V' \to V$가 $Y$ 위에서 서로 다른 사상이라고 하자.
$V', V$는 $Y$ 위에서 에탈이므로 등화자
$$
E =  V' \times_{(a, b), V \times_Y V, \Delta_{V/Y}} V
$$
를 생각할 수 있다. 이는 $a, b$의 등화자이고 $Y$ 위에서
에탈이다. 따라서 $E \to V'$는 에탈 단사 사상
(즉 열린 몰입)이며, 전사일 때 그리고 그때에만 동형이다.
$X \to Y$가 보편 위상동형이므로 이것은
$E_X = V'_X$일 때 그리고 그때에만 성립한다. 즉,
$a_X = b_X$일 때 그리고 그때에만 성립한다.

\medskip\noindent
다음으로 함자가 충만충실함을 증명하자. $V', V$가
$Y_{spaces, \etale}$의 대상이고
$c : V'_X \to V_X$가 $X$ 위의 사상이라고 하자.
사상 $a : V' \to V$를 구성하고자 한다. 이 사상은 $Y$ 위이고
$a_X = c$를 만족해야 한다. 전사 에탈 사상
$a' : V'' \to V'$을 택하되 $V''$이 분리 대수공간이 되게 하자.
사상 $a'' : V'' \to V$를 구성할 수 있고
$a''_X = c \circ a'_X$가 성립한다면, 두 합성
$$
V'' \times_{V'} V'' \xrightarrow{\text{pr}_i} V'' \xrightarrow{a''} V
$$
은 첫 문단에서 증명한 함자의 충실성에 의해 같다. 따라서 $a''$은
유일한 사상 $a : V' \to V$를 통하여 분해된다. 실제로 $V'$은
층으로서 $V''$을 동치관계 $V'' \times_{V'} V''$로 나눈
몫이다. 그러므로 $V'$이 분리라고
가정할 수 있다. 이 경우 그래프
$$
\Gamma_c \subset (V' \times_Y V)_X
$$
는 열리고 닫혀 있다 (세부사항은 생략한다). $X \to Y$가 보편
위상동형이므로 열린닫힌 부분공간
$\Gamma \subset V' \times_Y V$가 존재하여
$\Gamma_X = \Gamma_c$가 된다. 사영 $\Gamma \to V'$은 에탈
사상이고, $X$로 밑변환하면 동형이다. 따라서
$\Gamma \to V'$은 에탈이고 보편 단사이며 전사이므로 Morphisms
of Spaces, Lemma
\ref{spaces-morphisms-lemma-etale-universally-injective-open}에 의해
동형이다. 그러므로 $\Gamma$는 원하는 사상 $a : V' \to V$의
그래프이다.

\medskip\noindent
끝으로 함자가 본질적으로 전사임을 증명하자. $U$를
$X_{spaces, \etale}$의 대상이라 하자. 대상 $V$를 찾아야 한다.
이 대상은 $Y_{spaces, \etale}$에 속하고 $V_X \cong U$를
만족해야 한다. 전사 에탈 사상 $U' \to U$를 택하여
$U' \cong V'_X$이고 $U' \times_U U' \cong V''_X$가 되게 하자.
여기서 $V'', V'$은
$Y_{spaces, \etale}$의 대상이다. 함자의 충만충실성에 의해 사상
$s, t : V'' \to V'$을 얻으며, 이들은
$t_X = \text{pr}_0$와 $s_X = \text{pr}_1$을 만족한다. 이 등식은
$U' \times_U U' \to U'$인 사상으로서의 등식이다.
$(\text{pr}_0, \text{pr}_1) : U' \times_U U' \to U' \times_S U'$
이 에탈 동치관계라는 사실과
$U' \to V'$ 및 $U' \times_U U' \to V''$이 보편 단사이고
전사라는 사실로부터
$(t, s) : V'' \to V' \times_S V'$이 에탈 동치관계임을 얻는다.
그러면 몫 $V = V'/V''$은 (Spaces, Theorem
\ref{spaces-theorem-presentation}을 보라) 대수공간 $V$이며 $Y$ 위에
놓인다. 사상 $V' \to V$가 존재하여
$V'' = V' \times_V V'$가 된다. 따라서 사상 $V \to Y$를 얻는다
(Descent on Spaces, Lemma
\ref{spaces-descent-lemma-fpqc-universal-effective-epimorphisms}을
보라). $X$로 밑변환하면 사상 $U' \to V_X$와 양립하는 동형
$U' \times_{V_X} U' = U' \times_U U'$을 얻고, 이는
$V_X \cong U$를 함의한다 (방금 인용한 보조정리를 한 번 더
사용한다).

\medskip\noindent
스킴 $W$와 전사 에탈 사상 $W \to Y$를 택하자. 스킴 $U'$과
전사 에탈 사상 $U' \to U \times_X W_X$를 택하자.
$U'$과 $U' \times_U U'$은 $X$ 위에서 에탈인 스킴이고, 이들의
$X$로 가는 구조 사상은 스킴 $W_X$를 통하여 분해됨을 주목하자.
따라서 \'Etale Cohomology, Theorem
\ref{etale-cohomology-theorem-topological-invariance}에 의해
스킴들 $V', V''$이 존재하며 $W$ 위에서 에탈이다. 이들을 $W_X$로
밑변환하면 각각 $U'$ 및 $U' \times_U U'$과 동형이다. 이로써
증명이 끝난다.
\end{proof}

\begin{lemma}
\label{spaces-more-morphisms-lemma-topological-invariance}
Theorem \ref{spaces-more-morphisms-theorem-topological-invariance}과 같은 가정과 표기를
사용하자. 범주의 동치
$Y_{spaces, \etale} \to X_{spaces, \etale}$는 범주의 동치들
$Y_\etale \to X_\etale$ 및
$Y_{affine, \etale} \to X_{affine, \etale}$로 제한된다.
\end{lemma}

\begin{proof}
이는 단지 다음 명제이다. 대상 $V \in Y_{spaces, \etale}$가
주어졌을 때, $V$가 (아핀) 스킴인 것은
$V \times_Y X$가 (아핀) 스킴인 것과 동치이다.
$V \times_Y X \to V$는 $X \to Y$의 밑변환이므로 정수적이고
보편 단사이며 전사이다. 따라서 이는 Limits of Spaces, Lemma
\ref{spaces-limits-lemma-integral-universally-bijective-scheme}와
Proposition \ref{spaces-limits-proposition-affine}에서 따른다.
\end{proof}

\begin{remark}
\label{spaces-more-morphisms-remark-topological-invariance-etale-site}
\begin{reference}
Lenny Taelman이 2016년 5월 1일에 보낸 전자우편.
\end{reference}
대수공간의 보편 위상동형은 표현 가능할 필요가 없다. Morphisms
of Spaces, Example
\ref{spaces-morphisms-example-universal-homeomorphism}을 보라.
실제로 Theorem \ref{spaces-more-morphisms-theorem-topological-invariance}은 보편
위상동형에 대해서는 성립하지 않는다. 이를 보기 위해 대수적으로
닫힌 체 $k$를 택하되 표수는 $0$이라 하고,
$$
\mathbf{A}^1 \to X \to \mathbf{A}^1
$$
을 Morphisms of Spaces, Example
\ref{spaces-morphisms-example-universal-homeomorphism}과 같이
두자. 첫째 사상은 에탈이고 $t$와 $-t$를 동일시한다. 여기서
$t \in \mathbf{A}^1_k \setminus \{0\}$이다. 둘째 사상은 우리의
보편 위상동형임을 상기하자.
$\mathbf{A}^1_k$에는 자명하지 않은 연결 유한 에탈 덮개가 없으므로
($k$가 표수 0인 대수적으로 닫힌 체이기 때문이다. 세부사항은
생략한다), 자명하지 않은 연결 유한 에탈 덮개 $Y \to X$를
구성하면 충분하다. 이를 위해 $Y$를 원점을 두 겹으로 만든
아핀직선이라 하자 (Schemes, Example
\ref{schemes-example-affine-space-zero-doubled}). 그러면
$Y = Y_1 \cup Y_2$이고 $Y_i = \mathbf{A}^1_k$이며,
$\mathbf{A}^1_k \setminus \{0\}$을 따라 붙인다. 사상 $Y \to X$를
정의하기 위해 사상들
$$
Y_1 \xrightarrow{1} \mathbf{A}^1_k \to X
\quad\text{and}\quad
Y_2 \xrightarrow{-1} \mathbf{A}^1_k \to X.
$$
을 사용한다. 이들은 $Y_1 \cap Y_2$ 위에서 붙는다. 이는 $X$의
구성에서 따르므로 사상 $Y \to X$를 정의한다. 실제로
$$
\xymatrix{
Y \ar[d] & Y_1 \amalg Y_2 \ar[l] \ar[d] \\
X & \mathbf{A}^1_k \ar[l]
}
$$
가 데카르트 사각형이라고 주장한다. 세부사항은 생략한다. 예를
들어 Groupoids, Lemma
\ref{groupoids-lemma-criterion-fibre-product}를 사용할 수 있다.
$\mathbf{A}^1_k \to X$가 에탈이고 전사이므로, 이로부터
$Y \to X$는 차수 $2$의 유한 에탈 사상이고 원하는 예를 얻는다.

\medskip\noindent
더 간단히 다음과 같이 논할 수도 있다. 스킴 $Y$에는 군
$G = \{+1, -1\}$의 자유 작용이 있다. 여기서 $-1$은 $Y_1$과
$Y_2$를 서로 바꾸고 좌표의 부호를 바꾼다. 그러면
$X = Y/G$이고 (Spaces, Definition
\ref{spaces-definition-quotient}을 보라), 따라서 $Y \to X$는
유한 에탈이다. 또한 보편 위상동형 $X \to \mathbf{A}^1_k$가
존재함을 직접 보일 수 있다. 이를 위해 아핀직선 위의
$t \mapsto t^2$를 사용한다. 실제로 이 $X$는 위의 $X$와 같다.
\end{remark}








\section{두껍게 하기}
\label{spaces-more-morphisms-section-thickenings}

\noindent
다음 용어는 완전히 표준적이지 않을 수도 있지만 편리하다.

\begin{definition}
\label{spaces-more-morphisms-definition-thickening}
두껍게 하기. $S$를 스킴이라 하자.
\begin{enumerate}
\item 대수공간 $X'$을 대수공간 $X$의 {\it 두껍게 하기}라 함은
$X$가 $X'$의 닫힌 부분공간이고 대응하는 위상공간들이 같다는 뜻이다.
\item $X'$을 $X$의 {\it 일차 두껍게 하기}라 함은
$X$가 $X'$의 닫힌 부분공간이고 준연접 아이디얼 층
$\mathcal{I} \subset \mathcal{O}_{X'}$, 즉 $X$를 정의하는 아이디얼 층의
제곱이 영이라는 뜻이다.
% P09-ERR-1055: frozen source's self-containment X subset X is retained.
\item $X'$을 $X$의 {\it 유한 차수 두껍게 하기}라 함은
$X$가 $X$의 닫힌 부분공간이고 준연접 아이디얼 층
$\mathcal{I} \subset \mathcal{O}_{X'}$, 즉 $X$를 정의하는 아이디얼 층이
멱영이라는 뜻이다. 즉 어떤 정수
$n \geq 0$가 존재하여 $\mathcal{I}^{n + 1} = 0$이다.
% P09-ERR-1056: frozen source's self-containment X subset X is retained.
\item $X'$을 {\it $n$차 두껍게 하기}라 하되 그 대상이 $X$라는 것은,
$X$가 $X$의 닫힌 부분공간이고 $\mathcal{I}^{n + 1} = 0$이라는 뜻이다.
여기서 $\mathcal{I} \subset \mathcal{O}_{X'}$는 $X$를 정의하는
준연접 아이디얼 층이다.
% P09-ERR-1057: frozen source's f(X) before f is defined is retained.
\item 두 두껍게 하기 $X \subset X'$과 $Y \subset Y'$가 주어졌을 때,
{\it 두껍게 하기의 사상}은 사상 $f' : X' \to Y'$로서
$f(X) \subset Y$를 만족하는 것, 즉 $f'|_X$가 닫힌 부분공간 $Y$를
통해 인수분해되는 것이다. 이 상황에서 $f = f'|_X : X \to Y$로 놓고,
$(f, f') : (X \subset X') \to (Y \subset Y')$를 두껍게 하기의
사상이라 한다.
\item $B$를 대수공간이라 하자. 같은 방식으로
{$B$ 위의 \it 두껍게 하기}와 {$B$ 위의 \it 두껍게 하기의 사상}을
정의한다. 이는 위의 공간 $X, X', Y, Y'$가 $B$로 가는 구조 사상을
갖춘 대수공간이고, 사상
$X \to X'$, $Y \to Y'$ 및 $f' : X' \to Y'$가 $B$ 위의 사상이라는
뜻이다.
\end{enumerate}
\end{definition}

\noindent
기본 동치. $X \subset X'$가 두껍게 하기이면 $X \to X'$는 정수적이고
보편 전단사임을 주목하자. 따라서 당김 함자를 통해
\begin{equation}
\label{spaces-more-morphisms-equation-equivalence-etale-spaces}
X_{spaces, \etale} = X'_{spaces, \etale}
\end{equation}
를 얻는다. Theorem \ref{spaces-more-morphisms-theorem-topological-invariance}을 보라.
그러므로 $\mathcal{O}_{X'}$를 $X_{spaces, \etale}$ 위의 층으로
생각할 수 있다. 따라서 국소환 달린 토포스의 표준 동치
\begin{equation}
\label{spaces-more-morphisms-equation-fundamental-equivalence}
(\Sh(X'_{spaces, \etale}), \mathcal{O}_{X'})
\cong
(\Sh(X_{spaces, \etale}), \mathcal{O}_{X'})
\end{equation}
가 있다. 아래에서는 이를 Properties of Spaces, Theorem
\ref{spaces-properties-theorem-fully-faithful}의 완전충실성 결과와
자주 함께 사용할 것이다. 예를 들어 닫힌 몰입 $i_X : X \to X'$는
전사 사상 $i_X^\sharp : \mathcal{O}_{X'} \to \mathcal{O}_X$에 대응한다.

\medskip\noindent
$S$를 스킴이라 하고 $B$를 $S$ 위의 대수공간이라 하자.
$(f, f') : (X \subset X') \to (Y \subset Y')$를 $B$ 위의
두껍게 하기의 사상이라 하자. 연속 함자들의 도식
$$
\xymatrix{
X_{spaces, \etale} &
Y_{spaces, \etale} \ar[l] \\
X'_{spaces, \etale} \ar[u] &
Y'_{spaces, \etale} \ar[u] \ar[l]
}
$$
은 가환하고 수직 화살표들은 동치이다. 따라서
$f_{spaces, \etale}$, $f_{small}$,
$f'_{spaces, \etale}$ 및 $f'_{small}$은 모두 같은 토포스의 사상을
정의한다. 그러므로
$$
(f')^\sharp :
f_{spaces, \etale}^{-1}\mathcal{O}_{Y'}
\longrightarrow
\mathcal{O}_{X'}
$$
를 다음 가환 도식에 들어맞는 $\mathcal{O}_B$-대수 층의 사상으로
생각할 수 있다.
$$
\xymatrix{
f_{spaces, \etale}^{-1}\mathcal{O}_Y \ar[r]_-{f^\sharp} \ar[r] &
\mathcal{O}_X \\
f_{spaces, \etale}^{-1}\mathcal{O}_{Y'} \ar[r]^-{(f')^\sharp}
\ar[u]^{i_Y^\sharp} &
\mathcal{O}_{X'} \ar[u]_{i_X^\sharp}
}
$$
여기서 $i_X : X \to X'$와 $i_Y : Y \to Y'$는 주어진 닫힌 몰입들의
이름이다.

\begin{lemma}
\label{spaces-more-morphisms-lemma-first-order-thickening-maps}
$S$를 스킴이라 하자. $B$를 $S$ 위의 대수공간이라 하자.
$X \subset X'$과 $Y \subset Y'$를 $B$ 위의 대수공간들의
두껍게 하기라 하자. $f : X \to Y$를 $B$ 위의 대수공간의 사상이라
하자. 다음을 만족하는 임의의 $\mathcal{O}_B$-대수 사상
$$
\alpha : f_{spaces, \etale}^{-1}\mathcal{O}_{Y'} \to \mathcal{O}_{X'}
$$
가 주어졌다고 하자. 즉 도식
$$
\xymatrix{
f_{spaces, \etale}^{-1}\mathcal{O}_Y \ar[r]_-{f^\sharp} \ar[r] &
\mathcal{O}_X \\
f_{spaces, \etale}^{-1}\mathcal{O}_{Y'} \ar[r]^-\alpha
\ar[u]^{i_Y^\sharp} &
\mathcal{O}_{X'} \ar[u]_{i_X^\sharp}
}
$$
이 가환한다고 하자.
% P09-ERR-0524 / P09-ERR-1058: malformed frozen uniqueness phrase is recorded.
그러면 두껍게 하기의 사상 $(f, f')$가 $B$ 위에서 유일하게 존재하며
$\alpha = (f')^\sharp$를 만족한다.
\end{lemma}

\begin{proof}
$f'$을 찾기 위해 Properties of Spaces, Theorem
\ref{spaces-properties-theorem-fully-faithful}에 의하면 환 달린 토포스의 사상
$$
(f_{spaces, \etale}, \alpha) :
(\Sh(X_{spaces, \etale}), \mathcal{O}_{X'})
\longrightarrow
(\Sh(Y_{spaces, \etale}), \mathcal{O}_{Y'})
$$
이 국소환 달린 토포스의 사상임을 보이기만 하면 된다. 이는 국소환 달린
토포스의 사상의 정의(Modules on Sites, Definition
\ref{sites-modules-definition-morphism-locally-ringed-topoi}),
$(f, f^\sharp)$가 국소환 달린 토포스의 사상이라는 사실(Properties of
Spaces, Lemma \ref{spaces-properties-lemma-morphism-locally-ringed}),
$\alpha$가 주어진 가환 도식에 들어맞는다는 사실, 그리고
$i_X^\sharp$와 $i_Y^\sharp$의 핵이 국소적으로 멱영이라는 사실에서
곧바로 따른다. 마지막으로 $f' \circ i_X = i_Y \circ f$라는 사실은
도식의 가환성과 Properties of Spaces, Theorem
\ref{spaces-properties-theorem-fully-faithful}을 한 번 더 적용하여 얻는다.
$f'$이 $B$ 위의 사상이라는 확인은 생략한다.
\end{proof}

\begin{lemma}
\label{spaces-more-morphisms-lemma-open-subspace-thickening}
$S$를 스킴이라 하자. $X \subset X'$을 $S$ 위의 대수공간들의
두껍게 하기라 하자. 임의의 열린 부분공간 $U \subset X$에 대해
열린 부분공간 $U' \subset X'$이 유일하게 존재하여
$U = X \times_{X'} U'$을 만족한다.
\end{lemma}

\begin{proof}
$U' \to X'$를 $X'_{spaces, \etale}$의 대상으로 하되,
$U \to X$라는 $X_{spaces, \etale}$의 대상에
(\ref{spaces-more-morphisms-equation-equivalence-etale-spaces})를 통해 대응하도록 하자.
사상 $U' \to X'$는 에탈이고 보편 단사이므로 열린 몰입이다.
Morphisms of Spaces, Lemma
\ref{spaces-morphisms-lemma-etale-universally-injective-open}을 보라.
\end{proof}

\noindent
$i_X : X \to X'$를 대수공간들의 두껍게 하기라 하자. 핵
$\mathcal{I} = \Ker(i_X^\sharp) \subset \mathcal{O}_{X'}$의 모든
국소 단면은 국소적으로 멱영이다. 이는 $\mathcal{O}_{X'}$의 구조층에
대한 어느 정도의 제어를 주지만, 기술적으로는 유한 차수 두껍게 하기
$X \subset X'$의 부류가 훨씬 다루기 쉽다. 실제로 $X'$이
$n$차 두껍게 하기이면 여과
$$
0 \subset \mathcal{I}^n \subset \mathcal{I}^{n - 1} \subset
\ldots \subset \mathcal{I} \subset \mathcal{O}_{X'}
$$
를 가지며, $X'$이 닫힌 부분공간들로 여과됨을 알 수 있다.
% P09-ERR-1059: frozen filtration indices are retained.
$$
X = X_0 \subset X_1 \subset \ldots \subset X_{n - 1} \subset X_{n + 1} = X'
$$
여기서 각 쌍 $X_i \subset X_{i + 1}$는 $B$ 위의 일차 두껍게 하기이다.
간단한 귀납 논증을 사용하면 일차 두껍게 하기에 대해 증명한 많은 결과를
유한 차수 두껍게 하기에 대한 결과로 다시 진술할 수 있다.

\begin{lemma}
\label{spaces-more-morphisms-lemma-first-order-thickening-surjective}
$S$를 스킴이라 하자. $X \subset X'$을 $S$ 위의 대수공간들의
두껍게 하기라 하자. $U$를 $X_{spaces, \etale}$의 아핀 대상이라 하자.
그러면
$$
\Gamma(U, \mathcal{O}_{X'}) \to \Gamma(U, \mathcal{O}_X)
$$
는 전사이다. 여기서는 (\ref{spaces-more-morphisms-equation-fundamental-equivalence})를 통해
$\mathcal{O}_{X'}$를 $X_{spaces, \etale}$ 위의 층으로 생각한다.
\end{lemma}

\begin{proof}
대수공간의 에탈 사상 $U' \to X'$를 택하되
$U = X \times_{X'} U'$을 만족하도록 하자.
Theorem \ref{spaces-more-morphisms-theorem-topological-invariance}을 보라.
Limits of Spaces, Lemma \ref{spaces-limits-lemma-affine}에 의해
$U'$은 아핀 스킴이다. 따라서
$\Gamma(U, \mathcal{O}_{X'}) = \Gamma(U', \mathcal{O}_{U'}) \to
\Gamma(U, \mathcal{O}_U)$
는 전사이다. 실제로 $U \to U'$는 아핀 스킴들의 닫힌 몰입이다.
아래에서는 실제로 가장 자주 쓰이는 경우인 유한 차수 두껍게 하기에 대한
직접 증명을 준다.
\end{proof}

\begin{proof}[유한 차수 두껍게 하기에 대한 증명]
위에서 설명한 원리에 의해 $X \subset X'$이 일차 두껍게 하기라고
가정해도 된다. $\mathcal{I}$로 전사
$\mathcal{O}_{X'} \to \mathcal{O}_X$의 핵을 나타내자.
$\mathcal{I}$는 준연접 $\mathcal{O}_{X'}$-가군이고 일차 두껍게 하기의
정의에 의해 $\mathcal{I}^2 = 0$이므로, Morphisms of Spaces, Lemma
\ref{spaces-morphisms-lemma-i-star-equivalence}을 적용하여
$\mathcal{I}$가 준연접 $\mathcal{O}_X$-가군임을 알 수 있다.
따라서 보조정리는 짧은 완전열
$$
0 \to \mathcal{I} \to \mathcal{O}_{X'} \to \mathcal{O}_X \to 0
$$
에 딸린 긴 코호몰로지 완전열과
$H^1_\etale(U, \mathcal{I}) = 0$이라는 사실에서 따른다. 여기서
$\mathcal{I}$가 준연접이기 때문이다. Descent, Proposition
\ref{descent-proposition-same-cohomology-quasi-coherent}과
Cohomology of Schemes, Lemma
\ref{coherent-lemma-quasi-coherent-affine-cohomology-zero}를 보라.
\end{proof}

\begin{lemma}
\label{spaces-more-morphisms-lemma-thickening-scheme}
$S$를 스킴이라 하자. $X \subset X'$을 $S$ 위의 대수공간들의
두껍게 하기라 하자. $X$가 스킴으로 표현 가능하면 $X'$도 그러하다.
\end{lemma}

\begin{proof}
$X'_{red} = X_{red}$임을 주목하자. 따라서 $X$가 스킴이면
$X'_{red}$도 스킴이다. 그러므로 결과는 Limits of Spaces, Lemma
\ref{spaces-limits-lemma-reduction-scheme}에서 따른다.
아래에서는 실제로 가장 자주 쓰이는 경우인 유한 차수 두껍게 하기에 대한
직접 증명을 준다.
\end{proof}

\begin{proof}[유한 차수 두껍게 하기에 대한 증명]
$X'$이 $X$의 일차 두껍게 하기인 경우를 증명하면 충분하다.
Properties of Spaces, Lemma \ref{spaces-properties-lemma-subscheme}에 의해
$X'$에는 스킴인 가장 큰 열린 부분공간이 존재한다. 따라서 모든 점
$x$가 $|X'| = |X|$ 안에서 스킴인 $X'$의 열린 부분공간에 들어감을
보이면 된다. Lemma \ref{spaces-more-morphisms-lemma-open-subspace-thickening}을 사용하여
$X \subset X'$을 $U \subset U'$로 바꿀 수 있다. 여기서
$x \in U$이고 $U$는 아핀 스킴이다. 따라서 $X$가 아핀이라고
가정해도 된다. 이로써 다음 문단에서 다루는 경우로 환원된다.

\medskip\noindent
$X \subset X'$이 일차 두껍게 하기이고 $X$가 아핀 스킴이라고 하자.
$A = \Gamma(X, \mathcal{O}_X)$ 및
$A' = \Gamma(X', \mathcal{O}_{X'})$로 놓자. Lemma
\ref{spaces-more-morphisms-lemma-first-order-thickening-surjective}에 의해
사상 $A' \to A$는 전사이다. 핵 $I$는 제곱이 영인 아이디얼이다.
Properties of Spaces, Lemma
\ref{spaces-properties-lemma-morphism-to-affine-scheme}에 의해 다음 가환
도식에 들어맞는 표준 사상 $f : X' \to \Spec(A')$를 얻는다.
$$
\xymatrix{
X \ar@{=}[d] \ar[r] &  X' \ar[d]^f \\
\Spec(A) \ar[r] & \Spec(A')
}
$$
수평 화살표들이 두껍게 하기이므로 $f$가 보편 단사이고 전사임은
명백하다. 따라서 $f$가 에탈임을 보이면 충분하다. 그러면
Morphisms of Spaces, Lemma
\ref{spaces-morphisms-lemma-etale-universally-injective-open}에 의해
$f$는 동형사상이기 때문이다.

\medskip\noindent
$f$가 에탈임을 보이기 위해 아핀 스킴 $U'$과 에탈 사상
$U' \to X'$를 택하자. $U' \to X' \to \Spec(A')$가 에탈임을
보이면 충분하다. Properties of Spaces, Definition
\ref{spaces-properties-definition-etale}을 보라.
$U' = \Spec(B')$라 쓰자. $U = X \times_{X'} U'$로 놓자.
$U$는 $U'$의 닫힌 부분공간이므로 닫힌 부분스킴이고, 따라서
$U = \Spec(B)$이며 $B' \to B$는 전사이다.
$J = \Ker(B' \to B)$라 하고 다음을 주목하자.
$J = \Gamma(U, \mathcal{I})$이며, 여기서
$\mathcal{I} = \Ker(\mathcal{O}_{X'} \to \mathcal{O}_X)$는
Lemma \ref{spaces-more-morphisms-lemma-first-order-thickening-surjective}의 증명에서와 같이
$X_{spaces, \etale}$ 위에 있다.
사상 $U' \to X' \to \Spec(A')$는 가환 도식
$$
\xymatrix{
0 \ar[r] &
J \ar[r] &
B' \ar[r] &
B \ar[r] & 0 \\
0 \ar[r] &
I \ar[r] \ar[u] &
A' \ar[r] \ar[u] &
A \ar[r] \ar[u] & 0
}
$$
을 유도한다. 이제 $\mathcal{I}$는 준연접 $\mathcal{O}_X$-가군이므로
$\mathcal{I} = (\widetilde I)^a$이다. 표기에 대해서는 Descent,
Definition \ref{descent-definition-structure-sheaf}을, 이것이 참인 이유에
대해서는 Descent, Proposition
\ref{descent-proposition-equivalence-quasi-coherent}을 보라.
따라서 $J = I \otimes_A B$임을 알 수 있다. 마지막으로
$A \to B$가 에탈임을 주목하자. 실제로 $U \to X$는 에탈 사상
$U' \to X'$의 밑변환이므로 에탈이다. Algebra, Lemma
\ref{algebra-lemma-lift-etale-infinitesimal}에 의해
$A' \to B'$가 에탈이라고 결론 내린다.
\end{proof}

\begin{lemma}
\label{spaces-more-morphisms-lemma-thickening-equivalence}
$S$를 스킴이라 하자. $X \subset X'$을 $S$ 위의 대수공간들의
두껍게 하기라 하자. 함자
$$
V' \longmapsto V = X \times_{X'} V'
$$
는 범주의 동치 $X'_\etale \to X_\etale$를 정의한다.
\end{lemma}

\begin{proof}
함자 $V' \mapsto V$는 범주의 동치
$X'_{spaces, \etale} \to X_{spaces, \etale}$를 정의한다.
Theorem \ref{spaces-more-morphisms-theorem-topological-invariance}을 보라.
따라서 $V$가 스킴인 것과 $V'$이 스킴인 것이 동치임을 보이면 충분하다.
이는 Lemma \ref{spaces-more-morphisms-lemma-thickening-scheme}의 내용이다.
\end{proof}

\noindent
% P09-ERR-0528 / P09-ERR-1061: frozen singular/plural mismatch is recorded.
일차 두껍게 하기는 다음과 같이 기술된다.

\begin{lemma}
\label{spaces-more-morphisms-lemma-first-order-thickening}
$S$를 스킴이라 하자.
$f : X \to B$를 $S$ 위의 대수공간의 사상이라 하자.
짧은 완전열
$$
0 \to \mathcal{I} \to \mathcal{A} \to \mathcal{O}_X \to 0
$$
을 생각하자. 이는 $X_\etale$ 위의 층들의 열이다. 여기서
$\mathcal{A}$는 $f^{-1}\mathcal{O}_B$-대수의 층이고,
$\mathcal{A} \to \mathcal{O}_X$는 $f^{-1}\mathcal{O}_B$-대수 층의
전사이며, $\mathcal{I}$는 그 핵이다. 다음을 가정하자.
\begin{enumerate}
\item $\mathcal{I}$는 $\mathcal{A}$ 안에서 제곱이 영인 아이디얼이다.
\item $\mathcal{I}$는 $\mathcal{O}_X$-가군으로서 준연접이다.
\end{enumerate}
그러면 일차 두껍게 하기 $X \subset X'$이 $B$ 위에 존재하고,
동형사상 $\mathcal{O}_{X'} \to \mathcal{A}$가
$f^{-1}\mathcal{O}_B$-대수의 동형사상으로 존재하며, 이는
$\mathcal{O}_X$로 가는 전사들과 양립한다.
\end{lemma}

\begin{proof}
이 증명에서는 Lemmas \ref{spaces-more-morphisms-lemma-first-order-thickening-surjective}와
\ref{spaces-more-morphisms-lemma-thickening-scheme}의 증명에 사용한 논증 일부를 다시 전개한다.
먼저 $B = S = \Spec(\mathbf{Z})$인 경우를 다룬다.
$U$를 아핀 스킴이라 하고 $U \to X$를 에탈이라 하자. 그러면
$$
0 \to \mathcal{I}(U) \to \mathcal{A}(U) \to \mathcal{O}_X(U) \to 0
$$
은 완전하다. 실제로 $H^1(U_\etale, \mathcal{I}) = 0$인데,
$\mathcal{I}$가 준연접이기 때문이다. Descent, Proposition
\ref{descent-proposition-same-cohomology-quasi-coherent}과
Cohomology of Schemes, Lemma
\ref{coherent-lemma-quasi-coherent-affine-cohomology-zero}를 보라.
$V \to U$가 $X_{spaces, \etale}$의 아핀 대상들 사이의 사상이면
$$
\mathcal{I}(V) = \mathcal{I}(U) \otimes_{\mathcal{O}_X(U)} \mathcal{O}_X(V)
$$
이다. 이는 $\mathcal{I}$가 준연접 $\mathcal{O}_X$-가군이기 때문이다.
Descent, Proposition
\ref{descent-proposition-equivalence-quasi-coherent}을 보라.
따라서 $\mathcal{A}(U) \to \mathcal{A}(V)$는 에탈 환 사상이다.
Algebra, Lemma \ref{algebra-lemma-lift-etale-infinitesimal}을 보라.
그러므로
$$
U \longmapsto U' = \Spec(\mathcal{A}(U))
$$
는 $X_{affine, \etale}$에서 아핀 스킴과 에탈 사상의 범주로 가는
함자임을 알 수 있다. 실제로 이 함자는 함자 $U \mapsto U'$으로
$X_\etale$ 전체에 확장된다고 주장한다. 이를 보이기 위해
$U$가 $X_\etale$의 대상이면
$$
0 \to \mathcal{I}|_{U_{Zar}} \to \mathcal{A}|_{U_{Zar}} \to
\mathcal{O}_X|_{U_{Zar}} \to 0
$$
이고 $\mathcal{I}|_{U_{Zar}}$가 $U$ 위의 준연접 층임을 주목하자.
Descent, Proposition
\ref{descent-proposition-equivalence-quasi-coherent-functorial}을 보라.
따라서 More on Morphisms, Lemma
\ref{more-morphisms-lemma-first-order-thickening}에 의해 스킴들의 일차
두껍게 하기 $U \subset U'$을 얻으며, $\mathcal{O}_{U'}$는
$\mathcal{A}|_{U_{Zar}}$와 동형이다. 이 구성이 위의 아핀인 경우의
구성과 양립함은 명백하다.

\medskip\noindent
표시 $X = U/R$을 택하자. Spaces, Definition
\ref{spaces-definition-presentation}을 보라. 이때
$s, t : R \to U$는 에탈 동치관계를 정의한다. 위의 함자를 적용하면
스킴들에서의 에탈 동치관계 $s', t' : R' \to U'$를 얻는다.
대수공간 $X' = U'/R'$을 생각하자(Spaces, Theorem
\ref{spaces-theorem-presentation}을 보라).
사상 $X = U/R \to U'/R' = X'$은 일차 두껍게 하기이다.
$\mathcal{O}_{X'}$를 $X_\etale$ 위의 층으로 보자.
구성에 의해 동형사상
$$
\gamma :
\mathcal{O}_{X'}|_{U_\etale}
\longrightarrow
\mathcal{A}|_{U_\etale}
$$
을 얻으며, $s^{-1}\gamma$는 $t^{-1}\gamma$와 일치한다. 이 일치는
$R_\etale$ 위에서 성립한다. 따라서 Properties of Spaces, Lemma
\ref{spaces-properties-lemma-descent-sheaf}에 의해 $\gamma$는 원하는
유일한 동형사상 $\mathcal{O}_{X'} \to \mathcal{A}$에서 온다.

\medskip\noindent
일반적인 밑 대수공간 $B$의 경우를 다루기 위해, 먼저 위와 같이
$X'$을 $\mathbf{Z}$ 위의 대수공간으로 구성한다. 그런 다음 동형사상
$\mathcal{O}_{X'} \to \mathcal{A}$를 사용하여
$f^{-1}\mathcal{O}_B \to \mathcal{O}_{X'}$를 정의한다. Lemma
\ref{spaces-more-morphisms-lemma-first-order-thickening-maps}에 따르면 이는 사상 $X' \to B$를
정의하고, 이 사상은 주어진 사상 $X \to B$와 양립한다. 증명이 끝난다.
\end{proof}

\begin{lemma}
\label{spaces-more-morphisms-lemma-base-change-thickening}
$S$를 스킴이라 하자. $Y \subset Y'$을 $S$ 위의 대수공간들의
두껍게 하기라 하자. $X' \to Y'$를 사상이라 하고
$X = Y \times_{Y'} X'$로 놓자. 그러면
$(X \subset X') \to (Y \subset Y')$는 두껍게 하기의 사상이다.
$Y \subset Y'$가 일차 두껍게 하기(각각 유한 차수 두껍게 하기)이면
$X \subset X'$도 일차 두껍게 하기(각각 유한 차수 두껍게 하기)이다.
\end{lemma}

\begin{proof}
생략한다.
\end{proof}

\begin{lemma}
\label{spaces-more-morphisms-lemma-composition-thickening}
$S$를 스킴이라 하자. $X \subset X'$과 $X' \subset X''$이
$S$ 위의 대수공간들의 두껍게 하기이면 $X \subset X''$도 그러하다.
\end{lemma}

\begin{proof}
생략한다.
\end{proof}

\begin{lemma}
\label{spaces-more-morphisms-lemma-descending-property-thickening}
두껍게 하기라는 성질은 fpqc 국소적이다.
일차 두껍게 하기에 대해서도 마찬가지이다.
\end{lemma}

\begin{proof}
명제의 뜻은 다음과 같다. $S$를 스킴이라 하고
$X \to X'$를 $S$ 위의 대수공간의 사상이라 하자.
$\{g_i : X'_i \to X'\}$를 대수공간의 fpqc 덮개라 하고,
밑변환 $X_i \to X'_i$가 모든 $i$에 대해 두껍게 하기라고 하자.
그러면 $X \to X'$는 두껍게 하기이다. 사상들 $g_i$가 함께 전사이므로
$X \to X'$는 전사이다. Descent on Spaces, Lemma
\ref{spaces-descent-lemma-descending-property-closed-immersion}에 의해
$X \to X'$는 닫힌 몰입이다. 따라서 $X \to X'$는 두껍게 하기이다.
일차 두껍게 하기인 경우의 증명은 생략한다.
\end{proof}








\section{두껍게 하기의 사상}
\label{spaces-more-morphisms-section-morphisms-thickenings}

\noindent
$(f, f') : (X \subset X') \to (Y \subset Y')$가 대수공간들의 두껍게
하기 사이의 사상이면, 사상 $f$의 성질을 $f'$이 물려받는 경우가 많다.
몇 가지 변형이 있다.

\begin{lemma}
\label{spaces-more-morphisms-lemma-thicken-property-morphisms}
$S$를 스킴이라 하자. $(f, f') : (X \subset X') \to (Y \subset Y')$를
$S$ 위의 대수공간들의 두껍게 하기 사이의 사상이라 하자. 그러면
\begin{enumerate}
\item $f$가 아핀 사상인 것과 $f'$이 아핀 사상인 것은 동치이다.
\item $f$가 전사 사상인 것과 $f'$이 전사 사상인 것은 동치이다.
% P09-ERR-0529 / P09-ERR-1062: frozen missing copula is recorded.
\item $f$가 준콤팩트인 것과 $f'$이 준콤팩트인 것은 동치이다.
\item $f$가 보편 닫힌인 것과 $f'$이 보편 닫힌인 것은 동치이다.
\item $f$가 정수적인 것과 $f'$이 정수적인 것은 동치이다.
\item $f$가 (준)분리인 것과 $f'$이 (준)분리인 것은 동치이다.
\item $f$가 보편 단사인 것과 $f'$이 보편 단사인 것은 동치이다.
\item $f$가 보편 열린인 것과 $f'$이 보편 열린인 것은 동치이다.
\item $f$가 표현 가능한 것과 $f'$이 표현 가능한 것은 동치이다.
% P09-ERR-0530: frozen editorial placeholder is retained.
\item 여기에 더 추가할 것.
\end{enumerate}
\end{lemma}

\begin{proof}
$Y \to Y'$와 $X \to X'$는 정수적 보편 위상동형임을 관찰하자.
이로부터 (2), (3), (4), (7), (8)이 즉시 따른다.
(1)은 Limits of Spaces, Proposition
\ref{spaces-limits-proposition-affine}에서 따른다. 이 명제는
$X$와 $X'$ 위에서 에탈인 아핀 스킴들 사이, 그리고 $Y$와 $Y'$ 위에서
에탈인 아핀 스킴들 사이에 일대일 대응이 있음을 말한다.
(5)는 (1), (4)와 Morphisms of Spaces, Lemma
\ref{spaces-morphisms-lemma-integral-universally-closed}에서 따른다.
마지막으로
$$
X \times_Y X = X \times_{Y'} X \to X \times_{Y'} X' \to X' \times_{Y'} X'
$$
가 두껍게 하기임을 주목하자. 두 화살표는 Lemma
\ref{spaces-more-morphisms-lemma-base-change-thickening}에 의해 두껍게 하기이다.
따라서 (3), (4)를 사상
% P09-ERR-0531 / P09-ERR-1064: frozen inner arrow in the target pair is retained.
$(X \subset X') \to (X \times_Y X \to X' \times_{Y'} X')$
에 적용하면 (6)을 얻는다. 마지막으로 (9)는 스킴의 대수공간 두껍게
하기가 다시 스킴이라는 사실에서 따른다. Lemma
\ref{spaces-more-morphisms-lemma-thickening-scheme}을 보라.
\end{proof}

\begin{lemma}
\label{spaces-more-morphisms-lemma-thicken-property-morphisms-cartesian}
$S$를 스킴이라 하자. $(f, f') : (X \subset X') \to (Y \subset Y')$를
$S$ 위의 대수공간들의 두껍게 하기 사이의 사상이라 하고
$X = Y \times_{Y'} X'$이라고 하자. $X \subset X'$이 유한 차수
두껍게 하기이면 다음이 성립한다.
\begin{enumerate}
\item $f$가 닫힌 몰입인 것과 $f'$이 닫힌 몰입인 것은 동치이다.
\item $f$가 국소 유한형인 것과 $f'$이 국소 유한형인 것은 동치이다.
\item $f$가 국소 준유한인 것과 $f'$이 국소 준유한인 것은 동치이다.
\item $f$가 상대 차원 $d$인 국소 유한형인 것과 $f'$이 상대 차원
$d$인 국소 유한형인 것은 동치이다.
\item $\Omega_{X/Y} = 0$인 것과 $\Omega_{X'/Y'} = 0$인 것은 동치이다.
\item $f$가 비분기인 것과 $f'$이 비분기인 것은 동치이다.
\item $f$가 고유인 것과 $f'$이 고유인 것은 동치이다.
% P09-ERR-0532 / P09-ERR-1063: frozen article error is recorded.
\item $f$가 유한 사상인 것과 $f'$이 유한 사상인 것은 동치이다.
\item $f$가 단사상인 것과 $f'$이 단사상인 것은 동치이다.
\item $f$가 몰입인 것과 $f'$이 몰입인 것은 동치이다.
% P09-ERR-0533: frozen editorial placeholder is retained.
\item 여기에 더 추가할 것.
\end{enumerate}
\end{lemma}

\begin{proof}
스킴 $V'$과 전사 에탈 사상 $V' \to Y'$를 택하자.
스킴 $U'$과 전사 에탈 사상
$U' \to X' \times_{Y'} V'$를 택하자.
$V = Y \times_{Y'} V'$ 및 $U = X \times_{X'} U'$로 놓자.
그러면 사상의 에탈 국소적 성질에 대해서는 스킴들의 두껍게 하기 사이의
사상 $(U \subset U') \to (V \subset V')$로 환원하고 More on
Morphisms, Lemma
\ref{more-morphisms-lemma-thicken-property-morphisms-cartesian}을 적용할 수
있다. 이로써 (2), (3), (4), (5), (6)이 증명된다.

\medskip\noindent
(1), (7), (8), (9), (10)의 사상 성질들은 밑변환 아래 안정하다.
따라서 $f'$이 성질 $\mathcal{P}$를 가지면 $f$도 그러하다.
Spaces, Lemma \ref{spaces-lemma-base-change-immersions} 및
Morphisms of Spaces, Lemmas
\ref{spaces-morphisms-lemma-base-change-proper},
\ref{spaces-morphisms-lemma-base-change-integral},
\ref{spaces-morphisms-lemma-base-change-monomorphism}을 보라.

\medskip\noindent
(1), (7), (8), (9), (10)에서 흥미로운 방향은 $f$가 해당 성질을
가진다고 가정하여 $f'$도 그 성질을 가짐을 보이는 것이다.
두껍게 하기의 차수에 대한 귀납법으로 $Y \subset Y'$이 일차 두껍게
하기라고 가정해도 된다. 위의 유한 차수 두껍게 하기에 관한 논의를 보라.

\medskip\noindent
(1)의 증명. 스킴 $V'$과 전사 에탈 사상 $V' \to Y'$를 택하자.
$V = Y \times_{Y'} V'$, $U' = X' \times_{Y'} V'$ 및
$U = X \times_Y V$로 놓자. 그러면 $U \to V$는 닫힌 몰입이므로
$U$는 스킴이고, 다시 이로부터 $U'$이 스킴임을 알 수 있다
(Lemma \ref{spaces-more-morphisms-lemma-thickening-scheme}). 따라서 스킴인 경우의 보조정리
(More on Morphisms, Lemma
\ref{more-morphisms-lemma-thicken-property-morphisms-cartesian})를
$(U \subset U') \to (V \subset V')$에 적용하여 결론을 얻는다.

\medskip\noindent
(7)의 증명. (2)를 Lemma \ref{spaces-more-morphisms-lemma-thicken-property-morphisms}의 결과와
결합하고, 고유가 준콤팩트 $+$ 분리 $+$ 국소 유한형 $+$ 보편 닫힌과
같다는 사실을 사용하면 된다.

\medskip\noindent
(8)의 증명. (2)를 Lemma \ref{spaces-more-morphisms-lemma-thicken-property-morphisms}의 결과와
결합하고, 유한이 정수적 $+$ 국소 유한형과 같다는 사실을 사용하면 된다
(Morphisms, Lemma \ref{morphisms-lemma-finite-integral}).

\medskip\noindent
(9)의 증명. $f$가 단사상이므로 $X = X \times_Y X$이다.
지금까지 증명한 결과를 두껍게 하기의 사상
$(X \subset X') \to (X \times_Y X \subset X' \times_{Y'} X')$에
적용할 수 있다. (1)에 의해 $X' \to X' \times_{Y'} X'$가 닫힌
몰입이라고 결론 내린다. 실제로 닫힌 몰입
$X' \to X' \times_{Y'} X'$을 정의하는 아이디얼은 아이디얼
$\mathcal{I} \subset \mathcal{O}_{Y'}$, 즉 $Y$를 $Y'$ 안에서
잘라내는 아이디얼의 당김에 포함되므로,
이는 일차 두껍게 하기이다. 실제로
$X = X \times_Y X = (X' \times_{Y'} X') \times_{Y'} Y$는
$X'$에 포함된다. 닫힌 몰입
$\Delta : X' \to X' \times_{Y'} X'$의 여법층은
$\Omega_{X'/Y'}$와 같다(이는 대수공간에 대한 Morphisms, Lemma
\ref{morphisms-lemma-differentials-diagonal}의 유사 명제이며,
에탈 국소화를 사용하거나 Lemmas
\ref{spaces-more-morphisms-lemma-differentials-relative-immersion-section}과
\ref{spaces-more-morphisms-lemma-differential-product}를 결합하여 얻을 수 있다. 몇 가지
세부사항은 생략한다). 따라서 $\Omega_{X'/Y'} = 0$임을 보이면 충분하며,
이는 (5)와 $X/Y$에 대한 대응 명제에서 따른다.

\medskip\noindent
(10)의 증명. $f : X \to Y$가 몰입이면
$X \to V \to Y$로 인수분해된다. 여기서 $V \to Y$는 열린 부분공간이고
$X \to V$는 닫힌 몰입이다. Morphisms of Spaces, Remark
\ref{spaces-morphisms-remark-immersion}을 보라.
$V' \subset Y'$를 기저 위상공간 $|V'|$이
$|V| \subset |Y| = |Y'|$와 같은 열린 부분공간이라 하자.
그러면 $X' \to Y'$는 $V'$를 통해 인수분해되고, (1)에 의해
$X' \to V'$가 닫힌 몰입이라고 결론 내린다. 이로써 증명이 끝난다.
\end{proof}

\noindent
다음 보조정리는 앞의 보조정리의 한 변형이다. 여기서는 관련된 두껍게
하기들이 유한 차수라고 가정하여 국소 유한형이라는 성질을 $f$에서
$f'$으로 옮기는 대신, $f$와 $f'$ 각각이 국소 유한형이라는 것을
가정으로 둔다.

\begin{lemma}
\label{spaces-more-morphisms-lemma-properties-that-extend-over-thickenings}
$S$를 스킴이라 하자.
% P09-ERR-0534 / P09-ERR-1065: frozen target pair Y -> Y' is retained.
$(f, f') : (X \subset X') \to (Y \to Y')$를 $S$ 위의 대수공간들의
두껍게 하기 사이의 사상이라 하자. $f$와 $f'$이 국소 유한형이고
$X = Y \times_{Y'} X'$이라고 가정하자. 그러면
\begin{enumerate}
\item $f$가 국소 준유한인 것과 $f'$이 국소 준유한인 것은 동치이다.
\item $f$가 유한인 것과 $f'$이 유한인 것은 동치이다.
\item $f$가 닫힌 몰입인 것과 $f'$이 닫힌 몰입인 것은 동치이다.
\item $\Omega_{X/Y} = 0$인 것과 $\Omega_{X'/Y'} = 0$인 것은 동치이다.
\item $f$가 비분기인 것과 $f'$이 비분기인 것은 동치이다.
\item $f$가 단사상인 것과 $f'$이 단사상인 것은 동치이다.
\item $f$가 몰입인 것과 $f'$이 몰입인 것은 동치이다.
\item $f$가 고유인 것과 $f'$이 고유인 것은 동치이다.
% P09-ERR-0535: frozen editorial placeholder is retained.
\item 여기에 더 추가할 것.
\end{enumerate}
\end{lemma}

\begin{proof}
스킴 $V'$과 전사 에탈 사상 $V' \to Y'$를 택하자.
스킴 $U'$과 전사 에탈 사상
$U' \to X' \times_{Y'} V'$를 택하자.
$V = Y \times_{Y'} V'$ 및 $U = X \times_{X'} U'$로 놓자.
그러면 사상의 에탈 국소적 성질에 대해서는 스킴들의 두껍게 하기 사이의
사상 $(U \subset U') \to (V \subset V')$로 환원하고 More on
Morphisms, Lemma
\ref{more-morphisms-lemma-properties-that-extend-over-thickenings}을 적용할
수 있다. 이로써 (1), (4), (5)가 증명된다.

\medskip\noindent
(2), (3), (6), (7), (8)의 성질들은 밑변환 아래 안정하다. 따라서
$f'$이 성질 $\mathcal{P}$를 가지면 $f$도 그러하다. Spaces, Lemma
\ref{spaces-lemma-base-change-immersions} 및 Morphisms of Spaces, Lemmas
\ref{spaces-morphisms-lemma-base-change-proper},
\ref{spaces-morphisms-lemma-base-change-integral},
\ref{spaces-morphisms-lemma-base-change-monomorphism}을 보라.
따라서 각 경우에 $f$가 원하는 성질을 가지면 $f'$도 그 성질을 가짐만
증명하면 된다.

\medskip\noindent
(2)의 경우는 Lemma \ref{spaces-more-morphisms-lemma-thicken-property-morphisms}의 (5)와,
유한 사상이 정확히 국소 유한형인 정수적 사상이라는 사실에서 따른다
(Morphisms of Spaces, Lemma
\ref{spaces-morphisms-lemma-finite-integral}).

\medskip\noindent
(3)의 경우. 이는 Limits of Spaces, Lemma
\ref{spaces-limits-lemma-check-closed-infinitesimally}에서 즉시 따른다.

\medskip\noindent
(6)의 증명. $f$가 단사상이므로 $X = X \times_Y X$이다.
지금까지 증명한 결과를 두껍게 하기의 사상
$(X \subset X') \to (X \times_Y X \subset X' \times_{Y'} X')$에
적용할 수 있다. (3)에 의해
$\Delta_{X'/Y'} : X' \to X' \times_{Y'} X'$가 닫힌 몰입이라고
결론 내린다. 실제로 $\Delta_{X'/Y'}$는 전단사
$|X'| \to |X' \times_{Y'} X'|$를 유도하므로
$\Delta_{X'/Y'}$는 두껍게 하기이다. 한편 $\Delta_{X'/Y'}$는
Morphisms of Spaces, Lemma
\ref{spaces-morphisms-lemma-diagonal-morphism-finite-type}에 의해
국소 유한 표시이다. 다시 말해 $\Delta_{X'/Y'}(X')$는 유한형
준연접 아이디얼 층
$\mathcal{J} \subset \mathcal{O}_{X' \times_{Y'} X'}$로 잘려 나온다.
(5)에 의해 $\Omega_{X'/Y'} = 0$이므로
$X' \to X' \times_{Y'} X'$의 여법층이 영임을 알 수 있다.
(닫힌 몰입 $\Delta_{X'/Y'}$의 여법층은 $\Omega_{X'/Y'}$와 같다.
이는 대수공간에 대한 Morphisms, Lemma
\ref{morphisms-lemma-differentials-diagonal}의 유사 명제이며,
에탈 국소화를 사용하거나 Lemmas
\ref{spaces-more-morphisms-lemma-differentials-relative-immersion-section}과
\ref{spaces-more-morphisms-lemma-differential-product}를 결합하여 얻을 수 있다. 몇 가지
세부사항은 생략한다.) 다시 말해 $\mathcal{J}/\mathcal{J}^2 = 0$이다.
이로부터 $\Delta_{X'/Y'}$가 동형사상임을 알 수 있다. 예를 들어
Algebra, Lemma \ref{algebra-lemma-ideal-is-squared-union-connected}를 보라.

\medskip\noindent
(7)의 증명. $f : X \to Y$가 몰입이면
$X \to V \to Y$로 인수분해된다. 여기서 $V \to Y$는 열린 부분공간이고
$X \to V$는 닫힌 몰입이다. Morphisms of Spaces, Remark
\ref{spaces-morphisms-remark-immersion}을 보라.
$V' \subset Y'$를 기저 위상공간 $|V'|$이
$|V| \subset |Y| = |Y'|$와 같은 열린 부분공간이라 하자.
그러면 $X' \to Y'$는 $V'$를 통해 인수분해되고, (3)에 의해
$X' \to V'$가 닫힌 몰입이라고 결론 내린다.

\medskip\noindent
(8)의 경우는 Lemma \ref{spaces-more-morphisms-lemma-thicken-property-morphisms}와, 고유 사상이
국소 유한형이면서 준콤팩트, 보편 닫힌, 분리인 사상이라는 정의에서
따른다.
\end{proof}








\section{두껍게 하기의 Picard 군}
\label{spaces-more-morphisms-section-picard-group-thickening}

\noindent
두껍게 하기의 Picard 군에 관한 몇 가지 내용이다.

\begin{lemma}
\label{spaces-more-morphisms-lemma-picard-group-first-order-thickening}
$S$를 스킴이라 하자. $X \subset X'$을 $S$ 위의 대수공간들의
일차 두껍게 하기라 하고, 그 아이디얼 층을 $\mathcal{I}$라 하자.
그러면 아벨 군들의 표준 완전열
$$
\xymatrix{
0 \ar[r] &
H^0(X, \mathcal{I}) \ar[r] &
H^0(X', \mathcal{O}_{X'}^*) \ar[r] &
H^0(X, \mathcal{O}^*_X) \ar `r[d] `d[l] `l[llld] `d[dll] [dll] \\
& H^1(X, \mathcal{I}) \ar[r] &
\Pic(X') \ar[r] &
\Pic(X) \ar `r[d] `d[l] `l[llld] `d[dll] [dll] \\
& H^2(X, \mathcal{I}) \ar[r] & \ldots \ar[r] & \ldots
}
$$
이 있다.
\end{lemma}

\begin{proof}
$X_\etale = X'_\etale$임을 상기하자. Lemma
\ref{spaces-more-morphisms-lemma-thickening-equivalence} 및 더 일반적으로 Section
\ref{spaces-more-morphisms-section-thickenings}의 논의를 보라. 보조정리의 열은 아벨 군의 층들의
짧은 완전열
$$
0 \to \mathcal{I} \to \mathcal{O}_{X'}^* \to \mathcal{O}_X^* \to 0
$$
에 딸린 긴 코호몰로지 완전열이다. 이 열은 $X_\etale$ 위에 있으며,
첫째 사상은 국소 단면 $f$, 즉 $\mathcal{I}$의 단면을 가역 단면
$1 + f$, 즉 $\mathcal{O}_{X'}$의 단면으로 보낸다. 또한 환 달린
사이트의 Picard 군과 가역함수 층의 일차 코호몰로지 군을 동일시한다.
Cohomology on Sites, Lemma
\ref{sites-cohomology-lemma-h1-invertible}을 보라.
\end{proof}








\section{무한소 근방}
\label{spaces-more-morphisms-section-first-order-infinitesimal-neighbourhood}

\noindent
유한 차수 두껍게 하기의 자연스러운 구성은 다음과 같다.
% P09-ERR-0536 / P09-ERR-1066: frozen malformed supposition is recorded.
$i : Z \to X$가 대수공간의 몰입이라고 하자. 열린 부분공간
$U \subset X$를 택하되, $i$가 $Z$를 닫힌 부분공간
$Z \subset U$와 동일시하도록 하자(Morphisms of Spaces, Remark
\ref{spaces-morphisms-remark-immersion}을 보라).
$\mathcal{I} \subset \mathcal{O}_U$를 $Z$를 $U$ 안에서 정의하는
준연접 아이디얼 층이라 하자. Morphisms of Spaces, Lemma
\ref{spaces-morphisms-lemma-closed-immersion-ideals}을 보라.
$n \geq 1$에 대해 닫힌 부분공간 $Z_n \subset U$를 생각할 수 있으며,
이는 준연접 아이디얼 층 $\mathcal{I}^{n + 1}$로 정의된다.

\begin{definition}
\label{spaces-more-morphisms-definition-first-order-infinitesimal-neighbourhood}
$i : Z \to X$를 대수공간의 몰입이라 하자.
\begin{enumerate}
\item $Z$의 $X$ 안에서의 {\it 일차 무한소 근방}은 위에서 기술한
일차 두껍게 하기 $Z \subset Z_1$이며, 이는 $X$ 위에 있다.
\item {\it $n$차 무한소 근방}은 $Z$의 $X$ 안에서 위와 같이 기술한
$n$차 두껍게 하기 $Z \subset Z_n$이며, 이는 $X$ 위에 있다.
\end{enumerate}
\end{definition}

\noindent
이 두껍게 하기는 다음 보편 성질을 가진다. 이 성질은 위의 구성이 열린집합
$U$의 선택에 의존할지 모른다는 모든 우려를 없애 준다.

\begin{lemma}
\label{spaces-more-morphisms-lemma-first-order-infinitesimal-neighbourhood}
$i : Z \to X$를 대수공간의 몰입이라 하자.
\begin{enumerate}
\item $Z'$을 $Z$의 $X$ 안에서의 일차 무한소 근방이라 하자. 이는
다음 보편 성질을 가진다. 임의의 가환 도식
$$
\xymatrix{
Z \ar[d]_i & T \ar[l]^a \ar[d] \\
X & T' \ar[l]_b
}
$$
이 주어지고 $T \subset T'$이 $X$ 위의 일차 두껍게 하기이면,
% P09-ERR-0537: frozen reversed pair (a',a) is retained.
두껍게 하기 사이의 사상
$(a', a) : (T \subset T') \to (Z \subset Z')$가 $X$ 위에서
유일하게 존재한다.
\item $n \geq 1$에 대해 $n$차 무한소 근방 $Z_n$을 $Z$의 $X$ 안에서
취하자. 이는 다음 보편 성질을 가진다. 임의의 가환 도식
$$
\xymatrix{
Z \ar[d]_i & T \ar[l]^a \ar[d] \\
X & T' \ar[l]_b
}
$$
이 주어지고 $T \subset T'$이 $n$차 두껍게 하기이고 $X$ 위에 있으면,
% P09-ERR-0537: frozen reversed pair (a',a) is retained.
두껍게 하기 사이의 사상
$(a', a) : (T \subset T') \to (Z \subset Z_n)$가 $X$ 위에서
유일하게 존재한다.
\end{enumerate}
\end{lemma}

\begin{proof}
(1)만 증명한다. $U \subset X$를 $Z'$의 구성에 사용한 열린 부분공간,
즉 $Z$가 $U$의 닫힌 부분공간과 동일시되도록 하는 열린집합이라 하자.
이 닫힌 부분공간은 준연접 아이디얼 층 $\mathcal{I}$로 잘려 나온다.
$|T| = |T'|$이므로 $|b|(|T'|) \subset |U|$임을 알 수 있다.
따라서 $b$를 $U$로 가는 사상으로 생각할 수 있다. Properties of
Spaces, Lemma
\ref{spaces-properties-lemma-factor-through-open-subspace}를 보라.
$\mathcal{J} \subset \mathcal{O}_{T'}$를 $T$를 잘라내는, 제곱이 영인
준연접 아이디얼 층이라 하자. 도식의 가환성에 의해
$b|_T = i \circ a$이며, 여기서 $i : Z \to U$는 닫힌 몰입이다.
그러므로 Morphisms of Spaces, Lemma
\ref{spaces-morphisms-lemma-closed-immersion-ideals}에 의해
$b^\sharp(b^{-1}\mathcal{I}) \subset \mathcal{J}$이다.
$T'$은 $T$의 일차 두껍게 하기이므로 $\mathcal{J}^2 = 0$이고, 따라서
$b^\sharp(b^{-1}(\mathcal{I}^2)) = 0$이다. Morphisms of Spaces, Lemma
\ref{spaces-morphisms-lemma-closed-immersion-ideals}에 의해 이는 $b$가
$Z'$를 통해 인수분해됨을 뜻한다. 이 인수분해를
$a' : T' \to Z'$라 두면 원하는 결론을 얻는다.
\end{proof}

\begin{lemma}
\label{spaces-more-morphisms-lemma-infinitesimal-neighbourhood-conormal}
$i : Z \to X$를 대수공간의 몰입이라 하자.
$Z \subset Z'$을 $Z$의 $X$ 안에서의 일차 무한소 근방이라 하자.
그러면 도식
$$
\xymatrix{
Z \ar[r] \ar[d] & Z' \ar[d] \\
Z \ar[r] & X
}
$$
은 Lemma \ref{spaces-more-morphisms-lemma-conormal-functorial}에 의해 여법층의 사상
$\mathcal{C}_{Z/X} \to \mathcal{C}_{Z/Z'}$을 유도한다.
이 사상은 동형사상이다.
\end{lemma}

\begin{proof}
이는 위의 $Z'$의 구성에서 명백하다.
\end{proof}





\section{형식적으로 매끄러운, 에탈인, 비분기인 변환}
\label{spaces-more-morphisms-section-formally-smooth-etale-unramified}

\noindent
환 사상 $R \to A$가 {\it 형식적으로 매끄럽다}, 각각
{\it 형식적으로 에탈이다}, 각각 {\it 형식적으로 비분기이다}라는
정의를 상기하자(Algebra, Definition
\ref{algebra-definition-formally-smooth}, 각각 Definition
\ref{algebra-definition-formally-etale}, 각각 Definition
\ref{algebra-definition-formally-unramified}을 보라). 이는 모든 가환 실선
도식
$$
\xymatrix{
A \ar[r] \ar@{-->}[rd] & B/I \\
R \ar[r] \ar[u] & B \ar[u]
}
$$
에서 $I \subset B$가 제곱이 영인 아이디얼일 때, 도식을 가환으로 만드는
점선 화살표가 존재하는 것, 각각 유일하게 존재하는 것, 각각 많아야 하나
존재하는 것을 뜻한다. 이는 대수공간의 사상, 더 일반적으로 함자에 대한
다음 유사 개념의 동기가 된다.

\begin{definition}
\label{spaces-more-morphisms-definition-formally-smooth-etale-unramified}
$S$를 스킴이라 하자.
$a : F \to G$를 함자들의 변환이라 하고
$F, G : (\Sch/S)_{fppf}^{opp} \to \textit{Sets}$라 하자.
다음 꼴의 가환 실선 도식을 생각하자.
$$
\xymatrix{
F \ar[d]_a & T \ar[d]^i \ar[l] \\
G & T' \ar[l] \ar@{-->}[lu]
}
$$
여기서 $T$와 $T'$은 아핀 스킴이고 $i$는 제곱이 영인 아이디얼로
정의된 닫힌 몰입이다.
\begin{enumerate}
\item 위와 같은 임의의 실선 도식이 주어졌을 때 도식을 가환으로 만드는
점선 화살표가 존재하면 $a$가 {\it 형식적으로 매끄럽다}고 한다\footnote{이는
여기서 둘 수 있는 한 가지 가능한 정의일 뿐이다. 조금 더 약한 다른
조건은 점선 화살표가 $T'$ 위에서 fppf 국소적으로 존재한다고 요구하는
것이다. 이 약한 개념은 어떤 의미에서 더 좋은 형식적 성질을 가진다.}.
\item 위와 같은 임의의 실선 도식이 주어졌을 때 도식을 가환으로 만드는
점선 화살표가 정확히 하나 존재하면 $a$가 {\it 형식적으로 에탈이다}고
한다.
\item 위와 같은 임의의 실선 도식이 주어졌을 때 도식을 가환으로 만드는
점선 화살표가 많아야 하나 존재하면 $a$가 {\it 형식적으로 비분기이다}고
한다.
\end{enumerate}
\end{definition}

\begin{lemma}
\label{spaces-more-morphisms-lemma-formally-etale-is-combination}
$S$를 스킴이라 하자.
$a : F \to G$를 함자들의 변환이라 하고
$F, G : (\Sch/S)_{fppf}^{opp} \to \textit{Sets}$라 하자.
그러면 $a$가 형식적으로 에탈인 것과 $a$가 형식적으로 매끄럽고 동시에
형식적으로 비분기인 것은 동치이다.
\end{lemma}

\begin{proof}
정의에서 형식적으로 따른다.
\end{proof}

\begin{lemma}
\label{spaces-more-morphisms-lemma-composition-formally-smooth-etale-unramified}
합성.
\begin{enumerate}
\item 형식적으로 매끄러운 함자 변환들의 합성은 형식적으로 매끄럽다.
\item 형식적으로 에탈인 함자 변환들의 합성은 형식적으로 에탈이다.
\item 형식적으로 비분기인 함자 변환들의 합성은 형식적으로 비분기이다.
\end{enumerate}
\end{lemma}

\begin{proof}
이는 형식적이다.
\end{proof}

\begin{lemma}
\label{spaces-more-morphisms-lemma-base-change-formally-smooth-etale-unramified}
$S$를 $\Sch_{fppf}$에 포함되는 스킴이라 하자.
$F, G, H : (\Sch/S)_{fppf}^{opp} \to \textit{Sets}$라 하자.
$a : F \to G$, $b : H \to G$를 함자들의 변환이라 하자.
섬유곱 도식
$$
\xymatrix{
H \times_{b, G, a} F \ar[r]_-{b'} \ar[d]_{a'} & F \ar[d]^a \\
H \ar[r]^b & G
}
$$
을 생각하자.
\begin{enumerate}
\item $a$가 형식적으로 매끄러우면 밑변환 $a'$도 형식적으로 매끄럽다.
\item $a$가 형식적으로 에탈이면 밑변환 $a'$도 형식적으로 에탈이다.
\item $a$가 형식적으로 비분기이면 밑변환 $a'$도 형식적으로 비분기이다.
\end{enumerate}
\end{lemma}

\begin{proof}
이는 형식적이다.
\end{proof}

\begin{lemma}
\label{spaces-more-morphisms-lemma-representable-property-formally-property}
$S$를 스킴이라 하자.
$F, G : (\Sch/S)_{fppf}^{opp} \to \textit{Sets}$라 하자.
$a : F \to G$를 표현 가능한 함자 변환이라 하자.
\begin{enumerate}
\item $a$가 매끄러우면 $a$는 형식적으로 매끄럽다.
\item $a$가 에탈이면 $a$는 형식적으로 에탈이다.
\item $a$가 비분기이면 $a$는 형식적으로 비분기이다.
\end{enumerate}
\end{lemma}

\begin{proof}
Definition \ref{spaces-more-morphisms-definition-formally-smooth-etale-unramified}에서와 같은
가환 실선 도식
$$
\xymatrix{
F \ar[d]_a & T \ar[d]^i \ar[l] \\
G & T' \ar[l] \ar@{-->}[lu]
}
$$
을 생각하자. 그러면 $F \times_G T'$은 $T'$ 위에서 매끄러운
(각각 에탈인, 각각 비분기인) 스킴이다. 따라서 More on Morphisms,
Lemma \ref{more-morphisms-lemma-smooth-formally-smooth}
(각각 Lemma \ref{more-morphisms-lemma-etale-formally-etale},
각각 Lemma \ref{more-morphisms-lemma-unramified-formally-unramified})에 의해
도식
$$
\xymatrix{
F \times_G T' \ar[d] & T \ar[d]^i \ar[l] \\
T' & T' \ar[l] \ar@{-->}[lu]
}
$$
의 점선 화살표를 채울 수 있다(각각 유일하게 채울 수 있다, 각각 많아야
한 가지 방식으로 채울 수 있다).
% P09-ERR-0538 / P09-ERR-1067: frozen "an hence" typo is recorded.
따라서 첫째 도식에서도 대응하는 주장을 얻는다.
\end{proof}

\begin{lemma}
\label{spaces-more-morphisms-lemma-etale-on-top}
$S$를 $\Sch_{fppf}$에 포함되는 스킴이라 하자.
$F, G, H : (\Sch/S)_{fppf}^{opp} \to \textit{Sets}$라 하자.
$a : F \to G$, $b : G \to H$를 함자들의 변환이라 하자.
$a$가 표현 가능하고 전사이며 에탈이라고 가정하자.
\begin{enumerate}
\item $b$가 형식적으로 매끄러우면 $b \circ a$는 형식적으로 매끄럽다.
\item $b$가 형식적으로 에탈이면 $b \circ a$는 형식적으로 에탈이다.
\item $b$가 형식적으로 비분기이면 $b \circ a$는 형식적으로 비분기이다.
\end{enumerate}
역으로, 가환 실선 도식
$$
\xymatrix{
G \ar[d]_b & T \ar[d]^i \ar[l] \\
H & T' \ar[l] \ar@{-->}[lu]
}
$$
을 생각하자. 여기서 $T'$은 $S$ 위의 아핀 스킴이고
$i : T \to T'$는 제곱이 영인 아이디얼로 정의되는 닫힌 몰입이다.
\begin{enumerate}
\item[(4)] $b \circ a$가 형식적으로 매끄러우면, 모든 $t \in T$에 대해
아핀들의 에탈 사상 $U' \to T'$와 사상 $U' \to G$가 존재하여
$$
\xymatrix{
G \ar[d]_b & T \ar[l] & T \times_{T'} U' \ar[d] \ar[l]\\
H & T' \ar[l] & U' \ar[llu] \ar[l]
}
$$
이 가환하고 $t$는 $U' \to T'$의 상에 속한다.
\item[(5)] $b \circ a$가 형식적으로 비분기이면 위 도식의 점선 화살표는
많아야 하나 존재한다. 즉 $b$는 형식적으로 비분기이다.
\item[(6)] $b \circ a$가 형식적으로 에탈이면 위 도식의 점선 화살표는
정확히 하나 존재한다. 즉 $b$는 형식적으로 에탈이다.
\end{enumerate}
\end{lemma}

\begin{proof}
$b$가 형식적으로 매끄럽다(각각 형식적으로 에탈이다, 각각 형식적으로
비분기이다)고 하자. 에탈 사상은 매끄럽고 비분기이므로, $a$는 표현
가능하고 매끄럽다(각각 에탈이다, 각각 비분기이다).
따라서 (1), (2), (3)은 Lemma
\ref{spaces-more-morphisms-lemma-representable-property-formally-property}와 Lemma
\ref{spaces-more-morphisms-lemma-composition-formally-smooth-etale-unramified}를 결합하여 얻는다.

\medskip\noindent
$b \circ a$가 형식적으로 매끄럽다고 하자. 보조정리의 명제에 나오는
도식을 생각하자. $W = F \times_G T$로 놓자. 가정에 의해 $W$는
$T$ 위의 전사 에탈 스킴이다. 따라서 \'Etale Morphisms, Theorem
\ref{etale-theorem-remarkable-equivalence}에 의해 스킴 $W'$이 존재하고
$T'$ 위에서 에탈이며 $W = T \times_{T'} W'$이 된다. 아핀 열린 부분스킴
$U' \subset W'$를 택하되 $t$가 $U' \to T'$의 상에 속하도록 하자.
$b \circ a$가 형식적으로 매끄러우므로,
% P09-ERR-0539 / P09-ERR-1068: frozen existence typo is recorded.
사상 $U' \to F$가 존재하여
도식
$$
\xymatrix{
F \ar[d]_{b \circ a} & W \ar[l] & T \times_{T'} U' \ar[d] \ar[l]\\
H & T' \ar[l] & U' \ar[llu] \ar[l]
}
$$
이 가환한다. 합성
$U' \to F \to G$를 취하면 보조정리의 (5)에 나오는 사상을 얻는다.
% P09-ERR-0540 / P09-ERR-1069: frozen wrong item number (5) is retained.

\medskip\noindent
$f, g : T' \to G$가 보조정리의 도식에 들어맞는 두 점선 화살표라고 하자.
$W = F \times_G T$로 놓자. 가정에 의해 $W$는 $T$ 위의 전사 에탈
스킴이다. \'Etale Morphisms, Theorem
\ref{etale-theorem-remarkable-equivalence}에 의해 스킴 $W'$이 존재하고
$T'$ 위에서 에탈이며 $W = T \times_{T'} W'$이 된다. $a$가 형식적으로
에탈이므로 합성
$$
W' \to T' \xrightarrow{f} G
\quad\text{and}\quad
W' \to T' \xrightarrow{g} G
$$
은 사상 $f', g' : W' \to F$로 올라간다. 아핀 열린집합들 위에서
올린 뒤 유일성으로 붙이면 된다. 이제
$b \circ a : F \to H$가 형식적으로 비분기이면 $f' = g'$이고,
따라서 $f = g$이다. 실제로 $W' \to T'$은 에탈 덮개이다.
이는 보조정리의 (6)을
증명한다.
% P09-ERR-0541 / P09-ERR-1070: frozen wrong item number (6) is retained.

\medskip\noindent
$b \circ a$가 형식적으로 에탈이라고 하자. 그러면 (4)에 의해
$T'$ 위에서 에탈 국소적으로 도식에 들어맞는 점선 화살표를 찾을 수 있고,
(5)에 의해 이 점선 화살표는 유일하다. 따라서 국소 해들을 붙여 (6)을
얻을 수 있다. 몇 가지 세부사항은 생략한다.
\end{proof}

\begin{remark}
\label{spaces-more-morphisms-remark-tempting}
Lemma \ref{spaces-more-morphisms-lemma-etale-on-top}의 상황에서
``$b$가 형식적으로 매끄럽다'' $\Leftrightarrow$
``$b \circ a$가 형식적으로 매끄럽다''라고 생각하고 싶어진다.
그러나 이는 일반적으로 참이 아닐 가능성이 크다.
\end{remark}

\begin{lemma}
\label{spaces-more-morphisms-lemma-formally-permanence}
$S$를 스킴이라 하자.
$F, G, H : (\Sch/S)_{fppf}^{opp} \to \textit{Sets}$라 하자.
$a : F \to G$, $b : G \to H$를 함자들의 변환이라 하자.
$b$가 형식적으로 비분기라고 가정하자.
\begin{enumerate}
\item $b \circ a$가 형식적으로 비분기이면 $a$는 형식적으로 비분기이다.
\item $b \circ a$가 형식적으로 에탈이면 $a$는 형식적으로 에탈이다.
\item $b \circ a$가 형식적으로 매끄러우면 $a$는 형식적으로 매끄럽다.
\end{enumerate}
\end{lemma}

\begin{proof}
$T \subset T'$을 제곱이 영인 아이디얼로 정의되는 아핀 스킴들의 닫힌
몰입이라 하자. $g' : T' \to G$와 $f : T \to F$가 주어지고
$g'|_T = a \circ f$를 만족한다고 하자. $b$가 형식적으로 비분기이므로
$$
\{f' : T' \to F \mid f = f'|_T\text{ and }a \circ f' = g'\}
$$
와
$$
\{f' : T' \to F \mid f = f'|_T\text{ and }b \circ a \circ f' = b \circ g'\}.
$$
사이에 일대일 대응이 있다. 이로부터 보조정리가 형식적으로 따른다.
\end{proof}








\section{형식적으로 비분기인 사상}
\label{spaces-more-morphisms-section-formally-unramified}

\noindent
이 절에서는 대수공간의 사상이 형식적으로 비분기라는 말의 뜻을 밝힌다.

\begin{definition}
\label{spaces-more-morphisms-definition-formally-unramified}
$S$를 스킴이라 하자. 대수공간의 사상 $f : X \to Y$가 $S$ 위에 있고
Definition \ref{spaces-more-morphisms-definition-formally-smooth-etale-unramified}에서와 같이
함자 변환으로서 형식적으로 비분기이면, 이 사상이
{\it 형식적으로 비분기이다}라고 한다.
\end{definition}

\noindent
Section \ref{spaces-more-morphisms-section-formally-smooth-etale-unramified}에서 형식적으로
비분기인 함자 변환이라는 더 일반적인 설정에서 증명한 결과들을 다시
진술하지는 않겠다. 이 성질은 상대 미분 가군의 소멸로 특징지을 수 있음이
드러난다. Lemma \ref{spaces-more-morphisms-lemma-characterize-formally-unramified}를 보라.

\begin{lemma}
\label{spaces-more-morphisms-lemma-formally-unramified}
$S$를 스킴이라 하자. $f : X \to Y$를 대수공간의 사상이라 하되
$S$ 위에 있다고 하자.
다음은 동치이다.
\begin{enumerate}
\item $f$는 형식적으로 비분기이다.
\item 모든 도식
$$
\xymatrix{
U \ar[d] \ar[r]_\psi & V \ar[d] \\
X \ar[r]^f & Y
}
$$
에서 $U$와 $V$가 스킴이고 수직 화살표들이 에탈이면, 스킴의 사상
$\psi$는 형식적으로 비분기이다(More on Morphisms, Definition
\ref{more-morphisms-definition-formally-unramified}에서와 같은 의미이다).
\item 수직 화살표들이 전사인 그러한 도식 하나에 대해 사상 $\psi$는
형식적으로 비분기이다.
\end{enumerate}
\end{lemma}

\begin{proof}
$f$가 형식적으로 비분기라고 하자. Lemma
\ref{spaces-more-morphisms-lemma-representable-property-formally-property}에 의해 사상
$U \to X$와 $V \to Y$는 형식적으로 비분기이다. 따라서 Lemma
\ref{spaces-more-morphisms-lemma-composition-formally-smooth-etale-unramified}에 의해 합성
$U \to Y$는 형식적으로 비분기이다. 그러면 Lemma
\ref{spaces-more-morphisms-lemma-formally-permanence}에 의해 $U \to V$가 형식적으로
비분기이다. 따라서 (1)은 (2)를 함의한다. 또한 (2)는 (3)을 자명하게
함의한다.
% P09-ERR-0542: frozen missing terminal punctuation is recorded.

\medskip\noindent
(3)에서와 같은 도식이 주어졌다고 하자. Lemma
\ref{spaces-more-morphisms-lemma-representable-property-formally-property}에 의해 사상
$V \to Y$는 형식적으로 비분기이다. 따라서 Lemma
\ref{spaces-more-morphisms-lemma-composition-formally-smooth-etale-unramified}에 의해 합성
$U \to Y$는 형식적으로 비분기이다. 그러면 Lemma
\ref{spaces-more-morphisms-lemma-etale-on-top}에 의해 $X \to Y$는 형식적으로 비분기이다.
즉 (1)이 성립한다.
\end{proof}

\begin{lemma}
\label{spaces-more-morphisms-lemma-formally-unramified-not-affine}
$S$를 스킴이라 하자.
$f : X \to Y$가 $S$ 위의 대수공간의 형식적으로 비분기인 사상이면,
임의의 가환 실선 도식
$$
\xymatrix{
X \ar[d]_f & T \ar[d]^i \ar[l] \\
Y & T' \ar[l] \ar@{-->}[lu]
}
$$
에서 $T \subset T'$이 $S$ 위의 대수공간들의 일차 두껍게 하기일 때,
도식을 가환으로 만드는 점선 화살표는 많아야 하나 존재한다. 다시 말해
Definition \ref{spaces-more-morphisms-definition-formally-unramified}에서 $T$가 아핀 스킴이라는
조건을 버려도 된다.
\end{lemma}

\begin{proof}
이는 전사 에탈 사상 $U' \to T'$가 존재하여 $U'$이 아핀 스킴들의
분리합이 되기 때문이다(Properties of Spaces, Lemma
\ref{spaces-properties-lemma-cover-by-union-affines}을 보라). 또한 사상
$T' \to X$는 $U'$에 대한 제한으로 결정된다.
\end{proof}

\begin{lemma}
\label{spaces-more-morphisms-lemma-composition-formally-unramified}
형식적으로 비분기인 사상들의 합성은 형식적으로 비분기이다.
\end{lemma}

\begin{proof}
이는 형식적이다.
\end{proof}

\begin{lemma}
\label{spaces-more-morphisms-lemma-base-change-formally-unramified}
형식적으로 비분기인 사상의 밑변환은 형식적으로 비분기이다.
\end{lemma}

\begin{proof}
이는 형식적이다.
\end{proof}

\begin{lemma}
\label{spaces-more-morphisms-lemma-characterize-formally-unramified}
$S$를 스킴이라 하자. $f : X \to Y$를 $S$ 위의 대수공간의 사상이라
하자. 다음은 동치이다.
\begin{enumerate}
\item $f$는 형식적으로 비분기이다.
\item $\Omega_{X/Y} = 0$이다.
\end{enumerate}
\end{lemma}

\begin{proof}
이는 Lemma \ref{spaces-more-morphisms-lemma-formally-unramified}, More on Morphisms, Lemma
\ref{more-morphisms-lemma-formally-unramified-differentials} 및 Lemma
\ref{spaces-more-morphisms-lemma-localize-differentials}를 결합한 것이다.
\end{proof}

\begin{lemma}
\label{spaces-more-morphisms-lemma-unramified-formally-unramified}
$S$를 스킴이라 하자.
$f : X \to Y$를 $S$ 위의 대수공간의 사상이라 하자.
다음은 동치이다.
\begin{enumerate}
\item 사상 $f$는 비분기이다.
\item 사상 $f$는 국소 유한형이고 $\Omega_{X/Y} = 0$이다.
\item 사상 $f$는 국소 유한형이고 형식적으로 비분기이다.
\end{enumerate}
\end{lemma}

\begin{proof}
도식
$$
\xymatrix{
U \ar[d] \ar[r]_\psi & V \ar[d] \\
X \ar[r]^f & Y
}
$$
을 택하되, $U$와 $V$는 스킴이고 수직 화살표들은 에탈이며 전사라고
하자. 그러면 다음을 얻는다.
% P09-ERR-0543 / P09-ERR-1072: frozen "locally finite type" wording is recorded.
\begin{align*}
f\text{는 비분기}
& \Leftrightarrow
\psi\text{는 비분기} \\
& \Leftrightarrow
\psi\text{는 국소 유한형이고 }\Omega_{U/V} = 0 \\
& \Leftrightarrow
f\text{는 국소 유한형이고 }\Omega_{X/Y} = 0 \\
& \Leftrightarrow
f\text{는 국소 유한형이고 형식적으로 비분기}
\end{align*}
여기서 Morphisms, Lemma \ref{morphisms-lemma-unramified-omega-zero}와
Lemma \ref{spaces-more-morphisms-lemma-characterize-formally-unramified}를 사용했다.
\end{proof}

\begin{lemma}
\label{spaces-more-morphisms-lemma-universally-injective-unramified}
$S$를 스킴이라 하자.
$f : X \to Y$를 $S$ 위의 대수공간의 사상이라 하자.
다음은 동치이다.
\begin{enumerate}
\item $f$는 비분기이고 단사상이다.
\item $f$는 비분기이고 보편 단사이다.
\item $f$는 국소 유한형이고 단사상이다.
\item $f$는 보편 단사이고 국소 유한형이며 형식적으로 비분기이다.
\end{enumerate}
더욱이 이 경우 $f$는 표현 가능하고 분리이며 국소 준유한이다.
\end{lemma}

\begin{proof}
Lemma \ref{spaces-more-morphisms-lemma-unramified-formally-unramified}에서 형식적으로 비분기이고
국소 유한형이라는 것이 비분기라는 것과 같음을 보았다. 따라서 (4)는
(2)와 동치이다. 단사상은 분명히 형식적으로 비분기이므로 (3)은 (4)를
함의한다.
% P09-ERR-0544: frozen missing clause punctuation is recorded.
(1)이 (3)을 함의함은 명백하다. 마지막으로 (2)가 성립하면
$\Delta : X \to X \times_Y X$는 열린 몰입이고(Morphisms of Spaces,
Lemma \ref{spaces-morphisms-lemma-diagonal-unramified-morphism}), 동시에
전사이다(Morphisms of Spaces, Lemma
\ref{spaces-morphisms-lemma-universally-injective}). 따라서 동형사상이다.
즉 $f$는 단사상이다. 이로써 (2)가 (1)을 함의함을 알 수 있다.
마지막으로 Morphisms of Spaces, Lemmas
\ref{spaces-morphisms-lemma-monomorphism-loc-finite-type-loc-quasi-finite}와
\ref{spaces-morphisms-lemma-locally-quasi-finite-separated-representable}에
의해 $f$가 표현 가능하고 분리이며 국소 준유한임을 알 수 있다.
\end{proof}

\begin{lemma}
\label{spaces-more-morphisms-lemma-characterize-closed-immersion}
$S$를 스킴이라 하자.
$f : X \to Y$를 $S$ 위의 대수공간의 사상이라 하자.
다음은 동치이다.
\begin{enumerate}
\item $f$는 닫힌 몰입이다.
\item $f$는 보편 닫힌이고 비분기이며 단사상이다.
\item $f$는 보편 닫힌이고 비분기이며 보편 단사이다.
\item $f$는 보편 닫힌이고 국소 유한형이며 단사상이다.
\item $f$는 보편 닫힌이고 보편 단사이며 국소 유한형이고 형식적으로
비분기이다.
\end{enumerate}
\end{lemma}

\begin{proof}
(2)--(5)의 동치는 Lemma
\ref{spaces-more-morphisms-lemma-universally-injective-unramified}에서 즉시 따른다.
더욱이 (2)--(5)가 성립하면 $f$는 표현 가능하다. 마찬가지로 (1)이
성립하면 $f$는 표현 가능하다. 따라서 결과는 스킴인 경우에서 따른다.
\'Etale Morphisms, Lemma
\ref{etale-lemma-characterize-closed-immersion}을 보라.
\end{proof}





\section{보편 일차 두껍게 하기}
\label{spaces-more-morphisms-section-universal-thickening}

\noindent
$S$를 스킴이라 하자.
$h : Z \to X$를 $S$ 위의 대수공간의 사상이라 하자.
$Z$의 $X$ 위의 {\it 보편 일차 두껍게 하기}는 일차 두껍게 하기
$Z \subset Z'$이며, 이는 $X$ 위에 있다. 임의의 일차 두껍게 하기
$T \subset T'$이 $X$ 위에 있고 가환 실선 도식
\begin{equation}
\label{spaces-more-morphisms-equation-universal-first-order-thickening}
\vcenter{
\xymatrix{
& Z \ar[ld] & & T \ar[rd] \ar[ll]^a \\
Z' \ar[rrd] & & & & T' \ar@{..>}[llll]_{a'} \ar[lld]^b \\
 & & X
}
}
\end{equation}
이 주어졌을 때 도식을 가환으로 만드는 점선 화살표가 유일하게 존재하는
것을 말한다. 이 상황에서
$(a, a') : (T \subset T') \to (Z \subset Z')$는 $X$ 위의 두껍게
하기 사이의 사상임을 주목하자. 따라서 보편 일차 두껍게 하기가 존재하면
유일한 동형사상까지 유일하다. 일반적으로 보편 일차 두껍게 하기는
존재하지 않지만, $h$가 형식적으로 비분기이면 존재한다. 이를 증명하기
전에, 스킴의 범주에서의 보편 일차 두껍게 하기가 대수공간의 범주에서도
보편 일차 두껍게 하기임을 보이자.

\begin{lemma}
\label{spaces-more-morphisms-lemma-check-universal-first-order-thickening}
$S$를 스킴이라 하자.
$h : Z \to X$를 $S$ 위의 대수공간의 사상이라 하자.
$Z \subset Z'$을 $X$ 위의 일차 두껍게 하기라 하자.
다음은 동치이다.
% P09-ERR-0545: frozen missing list-introducing colon is recorded.
\begin{enumerate}
\item $Z \subset Z'$은 보편 일차 두껍게 하기이다.
\item $T'$이 스킴인 임의의 도식
(\ref{spaces-more-morphisms-equation-universal-first-order-thickening})에 대해 도식을 가환으로
만드는 점선 화살표가 유일하게 존재한다.
\item $T'$이 아핀 스킴인 임의의 도식
(\ref{spaces-more-morphisms-equation-universal-first-order-thickening})에 대해 도식을 가환으로
만드는 점선 화살표가 유일하게 존재한다.
\end{enumerate}
\end{lemma}

\begin{proof}
함의 (1) $\Rightarrow$ (2) $\Rightarrow$ (3)은 형식적이다.
% P09-ERR-0546 / P09-ERR-1073: frozen malformed conjunction is recorded.
(3)을 가정하고 임의의 도식
(\ref{spaces-more-morphisms-equation-universal-first-order-thickening})이 주어졌다고 하자.
표시 $T' = U'/R'$을 택하자. Spaces, Definition
\ref{spaces-definition-presentation}을 보라.
$U' = \coprod U'_i$가 아핀들의 분리합이라고 가정해도 된다. 따라서
% P09-ERR-0547 / P09-ERR-1074: frozen malformed \times_T' products are retained.
$R' = U' \times_{T'} U' = \coprod_{i, j} U'_i \times_T' U'_j$이다.
각 쌍 $(i, j)$에 대해 아핀 열린 덮개
$U'_i \times_T' U'_j = \bigcup_k R'_{ijk}$를 택하자.
$U_i, R_{ijk}$로 $T$와 취한 섬유곱들을 나타내되, 그 밑은 $T'$이라 하자.
그러면 각각의 $U_i \subset U'_i$와 $R_{ijk} \subset R'_{ijk}$는
아핀 스킴들의 일차 두껍게 하기이다.
$a_i : U_i \to Z$, 각각 $a_{ijk} : R_{ijk} \to Z$로
$a : T \to Z$와 사상 $U_i \to T$, 각각 $R_{ijk} \to T$의 합성을
나타내자. (3)을 $a_i : U_i \to Z$에 적용하여 유일한 사상
$a'_i : U'_i \to Z'$을 얻는다. (3)을 $a_{ijk}$에 적용하면 두 합성
% P09-ERR-0548 / P09-ERR-1075: frozen undefined R'_i and R'_j are retained.
$R'_{ijk} \to R'_i \to Z'$와 $R'_{ijk} \to R'_j \to Z'$가 같음을
알 수 있다. 따라서
$a' = \coprod a'_i : U' = \coprod U'_i \to Z'$은 몫 층
$T' = U'/R'$로 내려가며, 원하는 결론을 얻는다.
\end{proof}

\begin{lemma}
\label{spaces-more-morphisms-lemma-universal-thickening-over-formally-etale}
$S$를 스킴이라 하자.
$Z \to Y \to X$를 $S$ 위의 대수공간의 사상이라 하자.
$Z \subset Z'$이 $Z$의 $Y$ 위의 보편 일차 두껍게 하기이고
$Y \to X$가 형식적으로 에탈이면, $Z \subset Z'$은
$Z$의 $X$ 위의 보편 일차 두껍게 하기이다.
\end{lemma}

\begin{proof}
이는 형식적이다. 실제로 Lemma
\ref{spaces-more-morphisms-lemma-check-universal-first-order-thickening}에 의해 $T'$이
아핀 스킴인 가환 실선 도식
(\ref{spaces-more-morphisms-equation-universal-first-order-thickening})만 고려하면 충분하다.
합성 $T \to Z \to Y$는 $T' \to Y$로 유일하게 올라간다. 이는
$Y \to X$가 형식적으로 에탈이라고 가정했기 때문이다. 따라서
$Z \subset Z'$이 $Y$ 위의 보편 일차 두껍게 하기라는 사실로부터
원하는 사상 $a' : T' \to Z'$을 얻는다.
\end{proof}

\begin{lemma}
\label{spaces-more-morphisms-lemma-etale-morphism-of-universal-thickenings}
$S$를 스킴이라 하자.
$Z \to Y \to X$를 $S$ 위의 대수공간의 사상이라 하자.
$Z \to Y$가 에탈이라고 가정하자.
\begin{enumerate}
\item $Y \subset Y'$이 $Y$의 $X$ 위의 보편 일차 두껍게 하기이면,
유일한 에탈 사상 $Z' \to Y'$ 가운데
$Z = Y \times_{Y'} Z'$을 만족하는 것(Theorem
\ref{spaces-more-morphisms-theorem-topological-invariance}을 보라)은
$Z$의 $X$ 위의 보편 일차 두껍게 하기이다.
\item $Z \to Y$가 전사이고
$(Z \subset Z') \to (Y \subset Y')$이 $X$ 위 일차 두껍게 하기의
에탈 사상이며, $Z'$이 $Z$의 $X$ 위의 보편 일차 두껍게 하기이면,
$Y'$은 $Y$의 $X$ 위의 보편 일차 두껍게 하기이다.
\end{enumerate}
\end{lemma}

\begin{proof}
(1)의 증명. Lemma
\ref{spaces-more-morphisms-lemma-check-universal-first-order-thickening}에 의해 $T'$이
아핀 스킴인 가환 실선 도식
(\ref{spaces-more-morphisms-equation-universal-first-order-thickening})만 고려하면 충분하다.
합성 $T \to Z \to Y$는 $T' \to Y'$로 유일하게 올라간다. 이는
$Y'$이 보편 일차 두껍게 하기이기 때문이다. 그러면 $Z' \to Y'$이
에탈이라는 사실로부터(Lemma
\ref{spaces-more-morphisms-lemma-representable-property-formally-property}을 보라)
$T' \to Y'$는 원하는 사상 $a' : T' \to Z'$로 올라간다.

\medskip\noindent
(2)의 증명. $T \subset T'$을 $X$ 위의 일차 두껍게 하기라 하고,
$a : T \to Y$를 사상이라 하자. $W = T \times_Y Z$라 놓고
% P09-ERR-0549: frozen missing terminal punctuation is recorded.
$c : W \to Z$로 사영을 나타내자. 유일한 에탈 사상 $W' \to T'$ 가운데
$W = T \times_{T'} W'$을 만족하는 것을 취하자. Theorem
\ref{spaces-more-morphisms-theorem-topological-invariance}을 보라.
$W' \to T'$은 전사이다. 이는 $Z \to Y$가 전사이기 때문이다.
가정에 의해 유일한 사상 $c' : W' \to Z'$을 얻으며, 이는
$X$ 위의 사상이고 $c$로 제한되는데 그 제한은 $W$ 위에서 취한다.
유일성에 의해 $c'$의
$W' \times_{T'} W'$로의 두 제한은 같다(이는 $c$의
$W \times_T W$로의 두 제한이 같기 때문이다). 따라서 $c'$은
유일한 사상 $a' : T' \to Y'$로 내려가며, 원하는 결론을 얻는다.
\end{proof}

\begin{lemma}
\label{spaces-more-morphisms-lemma-universal-thickening}
$S$를 스킴이라 하자.
$h : Z \to X$를 $S$ 위의 대수공간의 형식적으로 비분기인 사상이라
하자. 보편 일차 두껍게 하기 $Z \subset Z'$이 존재하며, 이는
$Z$의 $X$ 위에 있다.
\end{lemma}

\begin{proof}
임의의 가환 도식
$$
\xymatrix{
V \ar[d] \ar[r] & U \ar[d] \\
Z \ar[r] & X
}
$$
을 택하되, $V$와 $U$는 스킴이고 수직 화살표들은 에탈이라 하자.
$V \to U$가 스킴의 형식적으로 비분기인 사상임을 주목하자.
Lemma \ref{spaces-more-morphisms-lemma-formally-unramified}을 보라. Lemma
\ref{spaces-more-morphisms-lemma-check-universal-first-order-thickening}과
More on Morphisms, Lemma
\ref{more-morphisms-lemma-universal-thickening}을 결합하면
보편 일차 두껍게 하기 $V \subset V'$이 존재하며, 이는
$V$의 $U$ 위에 있음을 알 수 있다.
% P09-ERR-0550 / P09-ERR-1076: frozen nonexistent part (1) is retained.
Lemma \ref{spaces-more-morphisms-lemma-universal-thickening-over-formally-etale} part (1)에 의해
$V'$은 $V$의 $X$ 위의 보편 일차 두껍게 하기이다.

\medskip\noindent
스킴 $U$와 전사 에탈 사상 $U \to X$를 고정하자. 위 논증은
임의의 에탈 사상 $V \to Z$에 대해 다음을 보인다. $V$가 스킴이고
$V \to Z \to X$가 $U$를 통해 인수분해되면 보편 일차 두껍게 하기
$V \subset V'$이 존재하며, 이는 $V$의 $X$ 위에 있다(그러나 이는
$V \to X$를 $U$를 통해 인수분해하는 방식의 선택에는 의존하지 않는다). 이제
$V$를 택하되 $V \to Z$가 전사 에탈이 되게 할 수 있다(Spaces, Lemma
\ref{spaces-lemma-lift-morphism-presentations}을 보라). 그러면
% P09-ERR-0551 / P09-ERR-1077: frozen omitted copula is recorded.
$R = V \times_Z V$는 $Z$ 위에서 에탈인 스킴이고, $R \to X$도
$U$를 통해 인수분해된다. 따라서 보편 일차 두껍게 하기
$V \subset V'$과 $R \subset R'$을 얻으며, 둘 다 $X$ 위에 있다.
$V \subset V'$이
보편 일차 두껍게 하기이므로 두 사영 $s, t : R \to V$는
사상 $s', t': R' \to V'$로 올라간다. Lemma
\ref{spaces-more-morphisms-lemma-etale-morphism-of-universal-thickenings}에 의해 $R'$이
$R$의 $X$ 위의 보편 일차 두껍게 하기이므로 이 사상들은 에탈이다.
% P09-ERR-0552 / P09-ERR-1078: frozen codomain V' of the relation pair is retained.
그러면 $(t', s') : R' \to V'$는 에탈 동치관계이고
$Z' = V'/R'$라 놓을 수 있다. $V' \to Z'$이 전사 에탈이고
% P09-ERR-0553 / P09-ERR-1079: frozen undefined lowercase v' is retained.
$v'$이 $V$의 $X$ 위의 보편 일차 두껍게 하기이므로
% P09-ERR-0554 / P09-ERR-1080: frozen wrong lemma reference and part (2) are retained.
Lemma \ref{spaces-more-morphisms-lemma-universal-thickening-over-formally-etale} part (2)에 의해
$Z'$은 $Z$의 $X$ 위의 보편 일차 두껍게 하기이다.
\end{proof}

\begin{definition}
\label{spaces-more-morphisms-definition-universal-thickening}
$S$를 스킴이라 하자.
$h : Z \to X$를 $S$ 위의 대수공간의 형식적으로 비분기인 사상이라
하자.
\begin{enumerate}
\item $Z$의 $X$ 위의 {\it 보편 일차 두껍게 하기}는 Lemma
\ref{spaces-more-morphisms-lemma-universal-thickening}에서 구성한 두껍게 하기
$Z \subset Z'$이다.
\item {\it $Z$의 $X$ 위의 여법선층}은 $Z$의 보편 일차 두껍게 하기
$Z'$ 안에서 $X$ 위로 본 여법선층이다.
\end{enumerate}
이 상황에서 여법선층은 흔히 $\mathcal{C}_{Z/X}$로 나타낸다.
\end{definition}

\noindent
따라서 층들의 짧은 완전열
$$
0 \to \mathcal{C}_{Z/X} \to \mathcal{O}_{Z'} \to \mathcal{O}_Z \to 0
$$
이 있으며, 이는 $Z_\etale$ 위에 있다. 또한 $\mathcal{C}_{Z/X}$는
준연접 $\mathcal{O}_Z$-가군이다.
다음 보조정리는 이 정의가 $Z \to X$가 몰입인 경우의 정의와
충돌하지 않음을 증명한다.

\begin{lemma}
\label{spaces-more-morphisms-lemma-immersion-universal-thickening}
$S$를 스킴이라 하자.
$i : Z \to X$를 $S$ 위 대수공간의 몰입이라 하자. 그러면
\begin{enumerate}
\item $i$는 형식적으로 비분기이다.
\item $Z$의 $X$ 위의 보편 일차 두껍게 하기는 Definition
\ref{spaces-more-morphisms-definition-first-order-infinitesimal-neighbourhood}의
$Z$의 $X$ 안에서의 일차 무한소 근방이다.
\item Definition \ref{spaces-more-morphisms-definition-conormal-sheaf}의 의미에서의
$i$의 여법선층은 Definition \ref{spaces-more-morphisms-definition-universal-thickening}의
의미에서의 $i$의 여법선층과 일치한다.
\end{enumerate}
\end{lemma}

\begin{proof}
대수공간의 몰입은 정의에 의해 표현 가능한 사상이다. 따라서
Morphisms, Lemmas \ref{morphisms-lemma-open-immersion-unramified} and
\ref{morphisms-lemma-closed-immersion-unramified}에 의해 몰입은
비분기이다(Spaces, Lemma
\ref{spaces-lemma-representable-transformations-property-implication}의
추상 원리를 사용한다). 그러므로 Lemma
\ref{spaces-more-morphisms-lemma-unramified-formally-unramified}에 의해 형식적으로
비분기이다. 나머지 주장들은 Lemmas
\ref{spaces-more-morphisms-lemma-first-order-infinitesimal-neighbourhood} and
\ref{spaces-more-morphisms-lemma-infinitesimal-neighbourhood-conormal}과 정의들을 결합하면
따른다.
\end{proof}

\begin{lemma}
\label{spaces-more-morphisms-lemma-universal-thickening-unramified}
$S$를 스킴이라 하자.
$Z \to X$를 $S$ 위 대수공간의 형식적으로 비분기인 사상이라 하자.
그러면 보편 일차 두껍게 하기 $Z'$은 $X$ 위에서 형식적으로
비분기이다.
\end{lemma}

\begin{proof}
$T \subset T'$을 $X$ 위 아핀 스킴의 일차 두껍게 하기라 하자.
가환 도식
$$
\xymatrix{
Z' \ar[d] & T \ar[l]^c \ar[d] \\
X & T' \ar[l] \ar[lu]^{a, b}
}
$$
을 취하자. $T_0 = c^{-1}(Z) \subset T$와
$T'_a = a^{-1}(Z)$를 놓자(스킴론적으로 취한다).
$Z'$이 $Z$의 일차 두껍게 하기이므로 $T'$은 $T'_a$의 일차
두껍게 하기이다. 더욱이 $c = a|_T$이므로
$T_0 = T \cap T'_a$이다(스킴론적으로 취한다). $T'$이 $T$의
일차 두껍게 하기이므로 $T'_a$는 $T_0$의 일차 두껍게 하기이다.
이제 $a|_{T'_a}$와 $b|_{T'_a}$는 $T'_a$에서 $Z'$으로 가는
$X$ 위의 사상들이고, $T_0$ 위에서는 $Z$으로 가는 사상으로서
일치한다. 따라서 $Z'$의 보편 성질에 의해
$a|_{T'_a} = b|_{T'_a}$임을 얻는다.
% P09-ERR-0555: frozen singular "morphism" with plural subject is recorded.
따라서 $a$와 $b$는 일차 두껍게 하기 $T'$에서 출발하는 사상들이며,
이는 $T'_a$의 두껍게 하기이다. 이들의 $T'_a$로의 제한은 $Z$으로
가는 사상으로서 일치한다. 그러므로 $Z'$의 보편 성질을 다시 한 번 사용하면
$a = b$를 얻는다. 다시 말해 형식적으로 비분기인 사상의 정의 성질이
원하는 대로 $Z' \to X$에 대해 성립한다.
\end{proof}

\begin{lemma}
\label{spaces-more-morphisms-lemma-universal-thickening-functorial}
% P09-ERR-0556 / P09-ERR-1082: frozen missing terminal punctuation is recorded.
$S$를 스킴이라 하자.
$S$ 위 대수공간의 가환 도식
$$
\xymatrix{
Z \ar[r]_h \ar[d]_f & X \ar[d]^g \\
W \ar[r]^{h'} & Y
}
$$
을 생각하되, $h$와 $h'$은 형식적으로 비분기라 하자.
$Z \subset Z'$을 $Z$의 $X$ 위의 보편 일차 두껍게 하기라 하자.
$W \subset W'$을 $W$의 $Y$ 위의 보편 일차 두껍게 하기라 하자.
두껍게 하기의 표준 사상
$(f, f') : (Z, Z') \to (W, W')$이 존재하며, 이는 $Y$ 위에 있고
다음 가환 도식에 들어맞는다.
$$
\xymatrix{
& & & Z' \ar[ld] \ar[d]^{f'} \\
Z \ar[rr] \ar[d]_f \ar[rrru] & & X \ar[d] & W' \ar[ld] \\
W \ar[rrru]|!{[rr];[rruu]}\hole \ar[rr] & & Y
}
$$
특히 두껍게 하기의 사상 $(f, f')$은 여법선층의 사상
$f^*\mathcal{C}_{W/Y} \to \mathcal{C}_{Z/X}$를 유도한다.
\end{lemma}

\begin{proof}
첫째 주장은 $W'$의 보편 성질에서 명백하다. 여법선층 위의 유도
사상은 $(Z \subset Z') \to (W \subset W')$에 Lemma
\ref{spaces-more-morphisms-lemma-conormal-functorial}을 적용하여 얻는 사상이다.
\end{proof}

\begin{lemma}
\label{spaces-more-morphisms-lemma-universal-thickening-fibre-product}
$S$를 스킴이라 하자. 대수공간의 도식
$$
\xymatrix{
Z \ar[r]_h \ar[d]_f & X \ar[d]^g \\
W \ar[r]^{h'} & Y
}
$$
이 $S$ 위의 섬유곱 도식이고 $h'$이 형식적으로 비분기라고 하자.
그러면 $h$는 형식적으로 비분기이며, $W \subset W'$이
$W$의 $Y$ 위의 보편 일차 두껍게 하기이면
$Z = X \times_Y W \subset X \times_Y W'$은 $Z$의 $X$ 위의
보편 일차 두껍게 하기이다. 특히 Lemma
\ref{spaces-more-morphisms-lemma-universal-thickening-functorial}의 표준 사상
$f^*\mathcal{C}_{W/Y} \to \mathcal{C}_{Z/X}$는 전사이다.
\end{lemma}

\begin{proof}
사상 $h$는 Lemma
\ref{spaces-more-morphisms-lemma-base-change-formally-unramified}에 의해 형식적으로
비분기이다. $X \times_Y W'$이 일차 두껍게 하기임은 명백하다.
$W'$이 보편 성질을 가지므로 이것도 보편 성질을 가짐을
섬유곱의 사상 성질로 곧바로 확인할 수 있다. 이것이 여법선층
사상의 전사성을 함의하는 이유는 Lemma
\ref{spaces-more-morphisms-lemma-conormal-functorial-flat}을 보라.
\end{proof}

\begin{lemma}
\label{spaces-more-morphisms-lemma-universal-thickening-fibre-product-flat}
$S$를 스킴이라 하자. 대수공간의 도식
$$
\xymatrix{
Z \ar[r]_h \ar[d]_f & X \ar[d]^g \\
W \ar[r]^{h'} & Y
}
$$
이 $S$ 위의 섬유곱 도식이고 $h'$이 형식적으로 비분기이며
$g$가 평탄하다고 하자. 이 경우 보편 일차 두껍게 하기들의 대응하는
사상 $Z' \to W'$은 평탄하고,
$f^*\mathcal{C}_{W/Y} \to \mathcal{C}_{Z/X}$는 동형사상이다.
\end{lemma}

\begin{proof}
평탄성은 밑변환 아래 보존된다. Morphisms of Spaces, Lemma
\ref{spaces-morphisms-lemma-base-change-flat}을 보라.
따라서 첫째 주장은 Lemma
\ref{spaces-more-morphisms-lemma-universal-thickening-fibre-product}에 나오는 $W'$의
서술에서 따른다. $X \times_Y W'$이 일차 두껍게 하기임은 명백하다.
$W'$이 보편 성질을 가지므로 이것도 보편 성질을 가짐을
섬유곱의 사상 성질로 곧바로 확인할 수 있다. 이것이 여법선층
사상의 동형성을 함의하는 이유는 Lemma
\ref{spaces-more-morphisms-lemma-conormal-functorial-flat}을 보라.
\end{proof}

\begin{lemma}
\label{spaces-more-morphisms-lemma-universal-thickening-localize}
보편 일차 두껍게 하기를 취하는 것은 에탈 국소화와 가환한다.
더 정확히, $h : Z \to X$를 밑 스킴 $S$ 위 대수공간의
형식적으로 비분기인 사상이라 하자. 수직 화살표들이 에탈인
가환 도식
$$
\xymatrix{
V \ar[d] \ar[r] & U \ar[d] \\
Z \ar[r] & X
}
$$
을 취하자. $Z'$을 $Z$의 $X$ 위의 보편 일차 두껍게 하기라 하자.
그러면 $V \to U$는 형식적으로 비분기이고, 보편 일차 두껍게 하기
$V'$은 $V$의 $U$ 위에 있으며 $Z'$ 위에서 에탈이다. 특히
$\mathcal{C}_{Z/X}|_V = \mathcal{C}_{V/U}$이다.
\end{lemma}

\begin{proof}
첫째 주장은 Lemma \ref{spaces-more-morphisms-lemma-formally-unramified}이다. 보편 일차
두껍게 하기들의 양립성은 Lemmas
\ref{spaces-more-morphisms-lemma-universal-thickening-over-formally-etale} and
\ref{spaces-more-morphisms-lemma-etale-morphism-of-universal-thickenings}의 결과이다.
\end{proof}

\begin{lemma}
\label{spaces-more-morphisms-lemma-differentials-universally-unramified}
$S$를 스킴이라 하자. $B$를 $S$ 위의 대수공간이라 하자.
$h : Z \to X$를 $B$ 위 대수공간의 형식적으로 비분기인 사상이라
하자. $Z \subset Z'$을 $Z$의 $X$ 위의 보편 일차 두껍게 하기라
하고, 그 구조 사상을 $h' : Z' \to X$라 하자. 표준 사상
% P09-ERR-0557: frozen d h' notation, inconsistent with c_{h'} in the proof, is retained.
$$
\text{d}h' : (h')^*\Omega_{X/B} \to \Omega_{Z'/B}
$$
은 동형사상
$h^*\Omega_{X/B} \to \Omega_{Z'/B} \otimes \mathcal{O}_Z$을
유도한다.
\end{lemma}

\begin{proof}
사상 $c_{h'}$는 Lemma
\ref{spaces-more-morphisms-lemma-functoriality-differentials}에서 정의한 사상이다.
$i : Z \to Z'$을 주어진 폐몰입이라 하면 $i^*c_{h'}$는 사상
% P09-ERR-0558 / P09-ERR-1084: frozen change from base B to S is retained.
$h^*\Omega_{X/S} \to \Omega_{Z'/S} \otimes \mathcal{O}_Z$이다.
이것이 동형사상임을 확인하는 일은 에탈 국소화에 의해 스킴의 경우로
환원된다. Lemma \ref{spaces-more-morphisms-lemma-universal-thickening-localize}와
Lemma \ref{spaces-more-morphisms-lemma-localize-differentials}를 보라.
이 경우 결과는 More on Morphisms, Lemma
\ref{more-morphisms-lemma-differentials-universally-unramified}이다.
\end{proof}

\begin{lemma}
\label{spaces-more-morphisms-lemma-universally-unramified-differentials-sequence}
$S$를 스킴이라 하자. $B$를 $S$ 위의 대수공간이라 하자.
$h : Z \to X$를 $B$ 위 대수공간의 형식적으로 비분기인 사상이라
하자. 표준 완전열
$$
\mathcal{C}_{Z/X} \to h^*\Omega_{X/B} \to \Omega_{Z/B} \to 0.
$$
이 존재한다. 첫째 화살표는 $\text{d}_{Z'/B}$에 의해 유도된다.
% P09-ERR-0559: frozen use of "neighbourhood" for the universal thickening is retained.
여기서 $Z'$은 $Z$의 $X$ 위의 보편 일차 근방이다.
\end{lemma}

\begin{proof}
표준 완전열
$$
\mathcal{C}_{Z/Z'} \to
\Omega_{Z'/S} \otimes \mathcal{O}_Z \to
\Omega_{Z/S} \to 0.
$$
이 존재함을 알고 있다.
% P09-ERR-0560 / P09-ERR-1085: frozen S-relative sequence in a B-relative lemma is retained.
Lemma \ref{spaces-more-morphisms-lemma-differentials-relative-immersion}을 보라. 따라서
Lemma \ref{spaces-more-morphisms-lemma-differentials-universally-unramified}를 적용하면
결과가 따른다.
\end{proof}

\begin{lemma}
\label{spaces-more-morphisms-lemma-two-unramified-morphisms}
$S$를 스킴이라 하자. 대수공간의 가환 도식
$$
\xymatrix{
Z \ar[r]_i \ar[rd]_j & X \ar[d] \\
& Y
}
$$
이 $S$ 위에 있고 $i$와 $j$가 형식적으로 비분기라고 하자.
그러면 표준 완전열
$$
\mathcal{C}_{Z/Y} \to
\mathcal{C}_{Z/X} \to
i^*\Omega_{X/Y} \to 0
$$
이 존재한다. 첫째 화살표는 Lemma
\ref{spaces-more-morphisms-lemma-universal-thickening-functorial}에서, 둘째 화살표는
Lemma \ref{spaces-more-morphisms-lemma-universally-unramified-differentials-sequence}에서
온다.
\end{lemma}

\begin{proof}
사상들이 이미 정의되었으므로 열이 완전함을 확인하는 일은 에탈
국소화에 의해 스킴의 경우로 환원된다. Lemma
\ref{spaces-more-morphisms-lemma-universal-thickening-localize}와 Lemma
\ref{spaces-more-morphisms-lemma-localize-differentials}를 보라. 이 경우 결과는
More on Morphisms, Lemma
\ref{more-morphisms-lemma-two-unramified-morphisms}이다.
\end{proof}

\begin{lemma}
\label{spaces-more-morphisms-lemma-transitivity-conormal-unramified}
$S$를 스킴이라 하자.
$Z \to Y \to X$를 $S$ 위 대수공간의 형식적으로 비분기인
사상들이라 하자.
\begin{enumerate}
\item $Z \subset Z'$이 $Z$의 $X$ 위의 보편 일차 두껍게 하기이고
$Y \subset Y'$이 $Y$의 $X$ 위의 보편 일차 두껍게 하기이면,
사상 $Z' \to Y'$이 존재하고 $Y \times_{Y'} Z'$은
$Z$의 $Y$ 위의 보편 일차 두껍게 하기이다.
\item 표준 완전열
$$
i^*\mathcal{C}_{Y/X} \to
\mathcal{C}_{Z/X} \to
\mathcal{C}_{Z/Y} \to 0
$$
이 존재하며, 여기의 사상들은 Lemma
\ref{spaces-more-morphisms-lemma-universal-thickening-functorial}에서 온다. 여기서
$i : Z \to Y$는 첫째 사상이다.
\end{enumerate}
\end{lemma}

\begin{proof}
(1)의 사상 $h : Z' \to Y'$은 Lemma
\ref{spaces-more-morphisms-lemma-universal-thickening-functorial}에서 온다.
$Y \times_{Y'} Z'$이 $Z$의 $Y$ 위의 보편 일차 두껍게 하기라는
주장은 $Z'$과 $Y'$의 보편 성질에서 명백하다. Lemma
\ref{spaces-more-morphisms-lemma-transitivity-conormal}에 의해 완전열
$$
(i')^*\mathcal{C}_{Y \times_{Y'} Z'/Z'} \to
\mathcal{C}_{Z/Z'} \to
\mathcal{C}_{Z/Y \times_{Y'} Z'} \to 0
$$
이 있다. 여기서 $i' : Z \to Y \times_{Y'} Z'$은 주어진 사상이다.
Lemma \ref{spaces-more-morphisms-lemma-conormal-functorial-flat}에 의해 전사
$h^*\mathcal{C}_{Y/Y'} \to \mathcal{C}_{Y \times_{Y'} Z'/Z'}$가
존재한다. 이를 등식
$\mathcal{C}_{Y/Y'} = \mathcal{C}_{Y/X}$,
$\mathcal{C}_{Z/Z'} = \mathcal{C}_{Z/X}$, 그리고
$\mathcal{C}_{Z/Y \times_{Y'} Z'} = \mathcal{C}_{Z/Y}$와
결합하면 보조정리를 얻는다.
\end{proof}













\section{형식적 에탈 사상}
\label{spaces-more-morphisms-section-formally-etale}

\noindent
이 절에서는 대수공간의 사상이 형식적으로 에탈이라는 것이 무엇을
뜻하는지 살펴본다.

\begin{definition}
\label{spaces-more-morphisms-definition-formally-etale}
$S$를 스킴이라 하자. $f : X \to Y$를 $S$ 위 대수공간의 사상이라
하자. 이 사상이 {\it 형식적으로 에탈}이라는 것은 Definition
\ref{spaces-more-morphisms-definition-formally-smooth-etale-unramified}에서처럼 함자의
변환으로서 형식적으로 에탈이라는 뜻이다.
\end{definition}

\noindent
Section \ref{spaces-more-morphisms-section-formally-smooth-etale-unramified}에서 함자의
형식적 에탈 변환이라는 더 일반적인 상황에서 증명한 결과들은
여기서 다시 진술하지 않는다.

\begin{lemma}
\label{spaces-more-morphisms-lemma-formally-etale}
$S$를 스킴이라 하자. $f : X \to Y$를 $S$ 위 대수공간의 사상이라
하자. 다음은 동치이다.
\begin{enumerate}
\item $f$는 형식적으로 에탈이다.
\item 모든 도식
$$
\xymatrix{
U \ar[d] \ar[r]_\psi & V \ar[d] \\
X \ar[r]^f & Y
}
$$
에서 $U$와 $V$가 스킴이고 수직 화살표들이 에탈이면, 스킴의 사상
$\psi$는 형식적으로 에탈이다(More on Morphisms, Definition
\ref{more-morphisms-definition-formally-etale}의 의미에서).
\item 수직 화살표들이 전사인 그러한 도식 하나에 대해 사상
$\psi$가 형식적으로 에탈이다.
\end{enumerate}
\end{lemma}

\begin{proof}
$f$가 형식적으로 에탈이라고 가정하자. Lemma
\ref{spaces-more-morphisms-lemma-representable-property-formally-property}에 의해 사상
$U \to X$와 $V \to Y$는 형식적으로 에탈이다. 따라서 Lemma
\ref{spaces-more-morphisms-lemma-composition-formally-smooth-etale-unramified}에 의해 합성
$U \to Y$는 형식적으로 에탈이다. 그러면 Lemma
\ref{spaces-more-morphisms-lemma-formally-permanence}에서 $U \to V$가 형식적으로
에탈임이 따른다. 따라서 (1)은 (2)를 함의한다. 그리고 (2)는 (3)을
자명하게 함의한다.
% P09-ERR-0561 / P09-ERR-1089: frozen missing terminal punctuation is recorded.

\medskip\noindent
(3)에서와 같은 도식이 주어졌다고 가정하자. Lemma
\ref{spaces-more-morphisms-lemma-representable-property-formally-property}에 의해 사상
$V \to Y$는 형식적으로 에탈이다. 따라서 Lemma
\ref{spaces-more-morphisms-lemma-composition-formally-smooth-etale-unramified}에 의해 합성
$U \to Y$는 형식적으로 에탈이다. 그러면 Lemma
\ref{spaces-more-morphisms-lemma-etale-on-top}에서 $X \to Y$가 형식적으로 에탈임, 즉
(1)이 성립함이 따른다.
\end{proof}

\begin{lemma}
\label{spaces-more-morphisms-lemma-formally-etale-not-affine}
$S$를 스킴이라 하자.
$f : X \to Y$를 $S$ 위 대수공간의 형식적으로 에탈인 사상이라 하자.
그러면 임의의 가환 실선 도식
$$
\xymatrix{
X \ar[d]_f & T \ar[d]^i \ar[l]_a \\
Y & T' \ar[l] \ar@{-->}[lu]
}
$$
이 주어지고 $T \subset T'$이 $Y$ 위 대수공간의 일차 두껍게
하기이면, 도식을 가환으로 만드는 점선 화살표가 정확히 하나
존재한다. 다시 말해 Definition \ref{spaces-more-morphisms-definition-formally-etale}에서
$T$가 아핀이라는 조건을 버릴 수 있다.
\end{lemma}

\begin{proof}
$U' \to T'$을 전사 에탈 사상이라 하되,
$U' = \coprod U'_i$가 아핀 스킴들의 분리합이라 하자.
$U_i = T \times_{T'} U'_i$라 놓자. 그러면
$a'_i : U'_i \to X$라는 사상들을 얻으며, $a'_i|_{U_i}$는 합성
$U_i \to T \to X$와 같다. 유일성에 의해(Lemma
\ref{spaces-more-morphisms-lemma-formally-unramified-not-affine}을 보라)
$a'_i$와 $a'_j$는 섬유곱 $U'_i \times_{T'} U'_j$ 위에서 일치한다.
따라서 $\coprod a'_i : U' \to X$는 내려가서 유일한 사상
$a' : T' \to X$를 준다.
\end{proof}

\begin{lemma}
\label{spaces-more-morphisms-lemma-composition-formally-etale}
형식적으로 에탈인 사상들의 합성은 형식적으로 에탈이다.
\end{lemma}

\begin{proof}
이는 형식적이다.
\end{proof}

\begin{lemma}
\label{spaces-more-morphisms-lemma-base-change-formally-etale}
형식적으로 에탈인 사상의 밑변환은 형식적으로 에탈이다.
\end{lemma}

\begin{proof}
이는 형식적이다.
\end{proof}

\begin{lemma}
\label{spaces-more-morphisms-lemma-characterize-formally-etale}
$S$를 스킴이라 하자.
% P09-ERR-0562 / P09-ERR-1086: frozen missing terminal punctuation is recorded.
$f : X \to Y$를 $S$ 위 대수공간의 사상이라 하자.
다음은 동치이다.
\begin{enumerate}
\item $f$는 형식적으로 에탈이다.
\item $f$는 형식적으로 비분기이고 $X$의 $Y$ 위의 보편 일차
두껍게 하기가 $X$와 같다.
\item $f$는 형식적으로 비분기이고 $\mathcal{C}_{X/Y} = 0$이다.
\item $\Omega_{X/Y} = 0$이고 $\mathcal{C}_{X/Y} = 0$이다.
\end{enumerate}
\end{lemma}

\begin{proof}
% P09-ERR-0563 / P09-ERR-1087: frozen subject-verb disagreement is recorded.
사실 마지막 주장은 $\Omega_{X/Y} = 0$이 Lemma
\ref{spaces-more-morphisms-lemma-characterize-formally-unramified}와 Definition
\ref{spaces-more-morphisms-definition-universal-thickening}을 통해
$\mathcal{C}_{X/Y}$가 정의된다는 것을 함의하기 때문에만 의미가 있다.
이것으로 (3)과 (4)가 동치임도 명백해진다.

\medskip\noindent
% P09-ERR-0564: frozen three-way "Either ... imply" wording is recorded.
가정 (1), (2), (3) 각각은 $f$가 형식적으로 비분기임을 함의한다.
따라서 $f$가 형식적으로 비분기라고 가정해도 된다.
% P09-ERR-0565: frozen singular/plural disagreement is recorded.
(1), (2), (3)의 동치는 보편 일차 두껍게 하기
% P09-ERR-0566 / P09-ERR-1088: frozen wrong base S is retained.
$X'$의 보편 성질에서 따른다. 여기서 이는 $X$의 $S$ 위의 것이며,
또한
$X = X' \Leftrightarrow \mathcal{C}_{X/Y} = 0$이라는 사실을 쓴다.
왜냐하면 정의에 의해
$\mathcal{C}_{X/Y} = \mathcal{C}_{X/X'}$는 결국 $X$의 아이디얼
층이며, 이를 $X'$ 안에서 취하기 때문이다.
\end{proof}

\begin{lemma}
\label{spaces-more-morphisms-lemma-unramified-flat-formally-etale}
비분기 평탄 사상은 형식적으로 에탈이다.
\end{lemma}

\begin{proof}
스킴의 경우에서 따른다. More on Morphisms, Lemma
\ref{more-morphisms-lemma-unramified-flat-formally-etale}을 보라.
또한 에탈 국소화에 대해서는 Lemmas
\ref{spaces-more-morphisms-lemma-formally-unramified} and \ref{spaces-more-morphisms-lemma-formally-etale}와
Morphisms of Spaces, Lemma
\ref{spaces-morphisms-lemma-flat-local}을 보라.
\end{proof}

\begin{lemma}
\label{spaces-more-morphisms-lemma-etale-formally-etale}
$S$를 스킴이라 하자.
$f : X \to Y$를 $S$ 위 대수공간의 사상이라 하자.
다음은 동치이다.
\begin{enumerate}
\item 사상 $f$는 에탈이다.
\item 사상 $f$는 국소 유한 표시이고 형식적으로 에탈이다.
\end{enumerate}
\end{lemma}

\begin{proof}
스킴의 경우에서 따른다. More on Morphisms, Lemma
\ref{more-morphisms-lemma-etale-formally-etale}을 보라. 또한 에탈
국소화에 대해서는 Lemma \ref{spaces-more-morphisms-lemma-formally-etale}와
Morphisms of Spaces, Lemmas
\ref{spaces-morphisms-lemma-finite-presentation-local} and
\ref{spaces-morphisms-lemma-etale-local}을 보라.
\end{proof}















\section{사상의 무한소 변형}
\label{spaces-more-morphisms-section-action-by-derivations}

\noindent
이 절에서는 미분을 사용하여 사상을 무한소적으로 움직이는 방법을
설명한다. 이 절 전체에서 두껍게 하기 $X'$ 위의 층을 생각하되,
이는 $X$의 두껍게 하기이고, 그 층을 $X$ 위의 층으로 볼 수 있다는
사실을 사용한다. 식
(\ref{spaces-more-morphisms-equation-equivalence-etale-spaces})와
(\ref{spaces-more-morphisms-equation-fundamental-equivalence})를 보라.

\begin{lemma}
\label{spaces-more-morphisms-lemma-difference-derivation}
$S$를 스킴이라 하자. $B$를 $S$ 위의 대수공간이라 하자.
$X \subset X'$과 $Y \subset Y'$을 $B$ 위 대수공간의 두 일차
두껍게 하기라 하자.
$(a, a'), (b, b') : (X \subset X') \to (Y \subset Y')$를
$B$ 위 두껍게 하기의 두 사상이라 하자. 다음을 가정하자.
\begin{enumerate}
\item $a = b$이다.
\item 두 사상 $a^*\mathcal{C}_{Y/Y'} \to \mathcal{C}_{X/X'}$
(Lemma \ref{spaces-more-morphisms-lemma-conormal-functorial})가 같다.
\end{enumerate}
그러면 사상 $(a')^\sharp - (b')^\sharp$는 다음과 같이 인수분해된다.
$$
\mathcal{O}_{Y'} \to \mathcal{O}_Y \xrightarrow{D}
a_*\mathcal{C}_{X/X'} \to a_*\mathcal{O}_{X'}
$$
여기서 $D$는 $\mathcal{O}_B$-미분이다.
\end{lemma}

\begin{proof}
$Y$ 위에서 작업하는 대신 $X$ 위에서 작업하자. 장점은 당김 함자
$a^{-1}$가 완전하다는 점이다. (1)과 (2)를 사용하여 행들이 완전한
가환 도식
$$
\xymatrix{
0 \ar[r] &
\mathcal{C}_{X/X'} \ar[r] &
\mathcal{O}_{X'} \ar[r] &
\mathcal{O}_X \ar[r] & 0 \\
0 \ar[r] &
a^{-1}\mathcal{C}_{Y/Y'} \ar[r] \ar[u] &
a^{-1}\mathcal{O}_{Y'}
\ar[r] \ar@<1ex>[u]^{(a')^\sharp} \ar@<-1ex>[u]_{(b')^\sharp} &
a^{-1}\mathcal{O}_Y \ar[r] \ar[u] & 0
}
$$
을 얻는다. 이제 이러한 상황에서
$\mathcal{O}_B$-대수 사상 $(a')^\sharp$와 $(b')^\sharp$의 차가
$\mathcal{O}_B$-미분이며, $a^{-1}\mathcal{O}_Y$에서
$\mathcal{C}_{X/X'}$로 간다는 것은 일반적 사실이다. 함자 $a^{-1}$와
$a_*$의 수반성에 의해 이는 $\mathcal{O}_B$-미분으로서
$\mathcal{O}_Y$에서 $a_*\mathcal{C}_{X/X'}$로 가는 것과 같다.
일부 세부사항은 생략한다.
\end{proof}

\noindent
위 보조정리의 상황에서 $D$를
\begin{equation}
\label{spaces-more-morphisms-equation-D}
% P09-ERR-0567: frozen ill-typed composition order is retained.
D = \text{d}_{Y/B} \circ \theta
\end{equation}
로 쓸 수 있음을 주목하자. 여기서 $\theta$는
$\mathcal{O}_Y$-선형 사상
$\theta : \Omega_{Y/B} \to a_*\mathcal{C}_{X/X'}$이다. 물론 다시
수반성을 사용하면 $\theta$를 $\mathcal{O}_X$-선형 사상
$\theta : a^*\Omega_{Y/B} \to \mathcal{C}_{X/X'}$로 볼 수 있다.

\begin{lemma}
\label{spaces-more-morphisms-lemma-action-by-derivations}
$S$를 스킴이라 하자. $B$를 $S$ 위의 대수공간이라 하자.
$(a, a') : (X \subset X') \to (Y \subset Y')$를 $B$ 위 일차
두껍게 하기의 사상이라 하자. 사상
$$
\theta : a^*\Omega_{Y/B} \to \mathcal{C}_{X/X'}
$$
이 $\mathcal{O}_X$-선형이라고 하자. 그러면 쌍들의 유일한 사상
$(b, b') : (X \subset X') \to (Y \subset Y')$가 존재하여
Lemma \ref{spaces-more-morphisms-lemma-difference-derivation}의 (1)과 (2)를 만족하고,
미분 $D$와 $\theta$는 식 (\ref{spaces-more-morphisms-equation-D})으로 관련된다.
\end{lemma}

\begin{proof}
사상
$$
\alpha = (a')^\sharp + D : a^{-1}\mathcal{O}_{Y'} \to \mathcal{O}_{X'}
$$
을 생각하자. 여기서 $D$는 식 (\ref{spaces-more-morphisms-equation-D})에서와 같다.
$D$가 $\mathcal{O}_B$-미분이므로 $\alpha$는
$\mathcal{O}_B$-대수들의 층의 사상이다. 구성에 의해
$i_X^\sharp \circ \alpha = a^\sharp \circ i_Y^\sharp$이다. 여기서
$i_X : X \to X'$과 $i_Y : Y \to Y'$은 주어진 폐몰입들이다.
Lemma \ref{spaces-more-morphisms-lemma-first-order-thickening-maps}에 의해 유일한 사상
$(a, b') : (X  \subset X') \to (Y \subset Y')$을 얻으며, 이는
$B$ 위 두껍게 하기의 사상이고 $\alpha = (b')^\sharp$이다.
$b = a$라 놓으면 원하는 결론을 얻는다.
\end{proof}

\begin{remark}
\label{spaces-more-morphisms-remark-action-by-derivations}
가정과 표기는 Lemma \ref{spaces-more-morphisms-lemma-action-by-derivations}에서와 같다.
국소 절단 $\theta$가 $a'$에 작용한 결과를 때때로
$\theta \cdot a'$로 나타낸다. 이는 $(a')^\sharp$와
$(\theta \cdot a')^\sharp$의 차가 미분 $D$이며, 이 미분이
$\theta$와 식 (\ref{spaces-more-morphisms-equation-D})으로 관련된다는 것 외에 아무
의미도 없다.
\end{remark}

\begin{lemma}
\label{spaces-more-morphisms-lemma-sheaf}
$S$를 스킴이라 하자. $B$를 $S$ 위의 대수공간이라 하자.
$X \subset X'$과 $Y \subset Y'$을 $B$ 위의 일차 두껍게 하기라
하자. 사상 $a : X \to Y$와 사상
$A : a^*\mathcal{C}_{Y/Y'} \to \mathcal{C}_{X/X'}$가 주어졌다고
가정하자. 후자는 $\mathcal{O}_X$-가군의 사상이다.
대상 $U'$을 취하되 이는 $(X')_{spaces, \etale}$의 대상이라 하자.
$U = X \times_{X'} U'$이라 하고, 다음을 만족하는 사상
$a' : U' \to Y'$들을 생각하자.
\begin{enumerate}
\item $a'$은 $B$ 위의 사상이다.
\item $a'|_U = a|_U$이다.
\item 유도 사상
$a^*\mathcal{C}_{Y/Y'}|_U \to \mathcal{C}_{X/X'}|_U$는
$A$를 $U$에 제한한 것이다.
\end{enumerate}
그러면 규칙
\begin{equation}
\label{spaces-more-morphisms-equation-sheaf}
U' \mapsto
\{a' : U' \to Y'\text{로서 (1), (2), (3)을 만족하는 것}\}
\end{equation}
은 $(X')_{spaces, \etale}$ 위 집합의 층을 정의한다.
\end{lemma}

\begin{proof}
$\mathcal{F}$로 보조정리의 규칙을 나타내자.
제한 사상 $\mathcal{F}(U') \to \mathcal{F}(V')$을 생각하자.
이는 $V' \subset U' \subset X'$에 대한 $\mathcal{F}$의 것이고,
실제로 제한 사상 $a' \mapsto a'|_{V'}$이다. 이 정의 아래에서
$\mathcal{F}$가 층임은
명백하다. 대수공간의 사상들이 에탈 내림을 만족하기 때문이다.
Descent on Spaces, Lemma
\ref{spaces-descent-lemma-fpqc-universal-effective-epimorphisms}를 보라.
\end{proof}

\begin{lemma}
\label{spaces-more-morphisms-lemma-action-sheaf}
표기와 가정은 Lemma \ref{spaces-more-morphisms-lemma-sheaf}에서와 같다.
식 (\ref{spaces-more-morphisms-equation-equivalence-etale-spaces})을 통해 $X$와 $X'$ 위의
층들을 동일시한다. 층
$$
\SheafHom_{\mathcal{O}_X}(a^*\Omega_{Y/B}, \mathcal{C}_{X/X'})
$$
은 식 (\ref{spaces-more-morphisms-equation-sheaf})의 층에 작용한다. 더욱이
대상 $U'$을 취하되 이는 $(X')_{spaces, \etale}$의 대상이라 하자.
이 대상 위에서 식 (\ref{spaces-more-morphisms-equation-sheaf})의 층이 절단을 가질
때마다 이 작용은 단순 추이적이다.
\end{lemma}

\begin{proof}
Lemmas \ref{spaces-more-morphisms-lemma-difference-derivation},
\ref{spaces-more-morphisms-lemma-action-by-derivations}, and \ref{spaces-more-morphisms-lemma-sheaf}를 결합하면 된다.
\end{proof}

\begin{remark}
\label{spaces-more-morphisms-remark-special-case}
Lemmas \ref{spaces-more-morphisms-lemma-difference-derivation},
\ref{spaces-more-morphisms-lemma-action-by-derivations},
\ref{spaces-more-morphisms-lemma-sheaf}, and
\ref{spaces-more-morphisms-lemma-action-sheaf}의 특수한 경우로 $Y = Y'$인 경우가 있다.
이 경우 사상 $A$는 항상 영이다. Lemma \ref{spaces-more-morphisms-lemma-sheaf}의 층은
단지 규칙
$$
U' \mapsto
\{a' : U' \to Y\text{는 }B\text{ 위에 있고 } a'|_U = a|_U\}
$$
으로 주어지며, 여기에 층
$\SheafHom_{\mathcal{O}_X}(a^*\Omega_{Y/B}, \mathcal{C}_{X/X'})$이
작용한다.
\end{remark}

\begin{remark}
\label{spaces-more-morphisms-remark-another-special-case}
Lemmas \ref{spaces-more-morphisms-lemma-difference-derivation},
\ref{spaces-more-morphisms-lemma-action-by-derivations},
\ref{spaces-more-morphisms-lemma-sheaf}, and
\ref{spaces-more-morphisms-lemma-action-sheaf}의 또 다른 특수한 경우로, $B$ 자체가 두껍게
하기 $Z \subset Z' = B$이고 $Y = Z \times_{Z'} Y'$인 경우가 있다.
그림으로는
$$
\xymatrix{
(X \subset X') \ar@{..>}[rr]_{(a, ?)} \ar[rd]_{(g, g')} & &
(Y \subset Y') \ar[ld]^{(h, h')} \\
& (Z \subset Z')
}
$$
이다. 이 경우 사상
$A : a^*\mathcal{C}_{Y/Y'} \to \mathcal{C}_{X/X'}$는 $a$에 의해
결정된다. 사상
$h^*\mathcal{C}_{Z/Z'} \to \mathcal{C}_{Y/Y'}$는 전사이다
(우리가 $Y = Z \times_{Z'} Y'$이라 가정했기 때문이다). 따라서 당김
$g^*\mathcal{C}_{Z/Z'} = a^*h^*\mathcal{C}_{Z/Z'} \to
a^*\mathcal{C}_{Y/Y'}$은 전사이고, 합성
$g^*\mathcal{C}_{Z/Z'} \to a^*\mathcal{C}_{Y/Y'} \to
\mathcal{C}_{X/X'}$은 $g'$이 유도하는 표준 사상이어야 한다. 따라서
Lemma \ref{spaces-more-morphisms-lemma-sheaf}의 층은 단지 규칙
$$
U' \mapsto
\{a' : U' \to Y'\text{는 }Z'\text{ 위에 있고 } a'|_U = a|_U\}
$$
으로 주어지며, 여기에 층
$\SheafHom_{\mathcal{O}_X}(a^*\Omega_{Y/Z}, \mathcal{C}_{X/X'})$이
작용한다.
\end{remark}

\begin{lemma}
\label{spaces-more-morphisms-lemma-action-by-derivations-etale-localization}
$S$를 스킴이라 하자. 일차 두껍게 하기들의 가환 도식
$$
\vcenter{
\xymatrix{
(T_2 \subset T_2') \ar[d]_{(h, h')} \ar[rr]_{(a_2, a_2')} & &
(X_2 \subset X_2') \ar[d]^{(f, f')} \\
(T_1 \subset T_1') \ar[rr]^{(a_1, a_1')} & &
(X_1 \subset X_1')
}
}
\quad
\begin{matrix}
\text{그리고 가환} \\
\text{도식}
\end{matrix}
\quad
\vcenter{
\xymatrix{
X_2' \ar[r] \ar[d] & B_2 \ar[d] \\
X_1' \ar[r] & B_1
}
}
$$
을 생각하자. 이들은 $S$ 위의 대수공간들이며,
$X_2 \to X_1$과 $B_2 \to B_1$은 에탈이라 하자.
임의의 $\mathcal{O}_{T_1}$-선형 사상
$\theta_1 : a_1^*\Omega_{X_1/B_1} \to \mathcal{C}_{T_1/T'_1}$에 대해
$\theta_2$를 합성
$$
\xymatrix{
a_2^*\Omega_{X_2/B_2} \ar@{=}[r] &
h^*a_1^*\Omega_{X_1/B_1} \ar[r]^-{h^*\theta_1} &
h^*\mathcal{C}_{T_1/T'_1} \ar[r] &
\mathcal{C}_{T_2/T'_2}
}
$$
이라 하자(등호는 증명에서 설명한다). 그러면 도식
$$
\xymatrix{
T_2' \ar[rr]_{\theta_2 \cdot a_2'} \ar[d] & & X'_2 \ar[d] \\
T_1' \ar[rr]^{\theta_1 \cdot a_1'} & & X'_1
}
$$
은 가환한다. 여기서 작용 $\theta_2 \cdot a_2'$과
$\theta_1 \cdot a_1'$은 Remark
\ref{spaces-more-morphisms-remark-action-by-derivations}에서와 같다.
\end{lemma}

\begin{proof}
등호는 동일시
% P09-ERR-0568: frozen undefined S_1 and S_2 bases are retained.
$f^*\Omega_{X_1/S_1} = \Omega_{X_2/S_2}$에서 온다. 미분층의 구성이
에탈 국소화와 양립하기 때문에 이 동일시를 얻는다(원천과 타깃 모두에
대해). Lemma \ref{spaces-more-morphisms-lemma-localize-differentials}를 보라.
구체적으로 이를 사용하면
% P09-ERR-0569: frozen undefined S_i bases in the identification chain are retained.
$a_2^*\Omega_{X_2/S_2} = a_2^*f^*\Omega_{X_1/S_1} =
h^*a_1^*\Omega_{X_1/S_1}$을 얻는다. 이는
$f \circ a_2 = a_1 \circ h$이기 때문이다.
% P09-ERR-0570: frozen spurious preposition in "checked on etale locally" is recorded.
이제 도식의 가환성은 에탈 국소적으로 확인할 수 있다. 따라서
$T'_i$, $X'_i$, $B_2$, $B_1$이 스킴이라고 가정해도 되며, 이 경우
보조정리는 More on Morphisms, Lemma
\ref{more-morphisms-lemma-action-by-derivations-etale-localization}에서
따른다. 다른 증명은 다음과 같다. Lemma
\ref{spaces-more-morphisms-lemma-first-order-thickening-maps}를 사용하면 $X_1'$ 위
환의 층들의 어떤 도식이 가환임을 보이면 충분하다. 그런 다음 앞서 말한
More on Morphisms, Lemma
\ref{more-morphisms-lemma-action-by-derivations-etale-localization}의
증명과 정확히 같은 논증으로 실제로 그렇다는 것을 알 수 있다.
\end{proof}










\section{대수공간의 무한소 변형}
\label{spaces-more-morphisms-section-deform}

\noindent
다음의 간단한 보조정리는 사상의 무한소 변형이 평탄한지 확인할 때
흔히 편리한 도구가 된다.

\begin{lemma}
\label{spaces-more-morphisms-lemma-deform}
$S$를 스킴이라 하자.
$(f, f') : (X \subset X') \to (Y \subset Y')$를 $S$ 위 대수공간의
일차 두껍게 하기 사이의 사상이라 하자. $f$가 평탄하다고 가정하자.
다음은 동치이다.
\begin{enumerate}
\item $f'$은 평탄하고 $X = Y \times_{Y'} X'$이다.
\item 표준 사상 $f^*\mathcal{C}_{Y/Y'} \to \mathcal{C}_{X/X'}$는
동형사상이다.
\end{enumerate}
\end{lemma}

\begin{proof}
스킴 $V'$과 전사 에탈 사상 $V' \to Y'$을 택하자. 스킴 $U'$과
전사 에탈 사상 $U' \to X' \times_{Y'} V'$을 택하자.
$U = X \times_{X'} U'$와 $V = Y \times_{Y'} V'$라 놓자.
대수공간의 평탄 사상에 대한 우리의 정의에 따르면 유도 사상
$g : U \to V$가 스킴의 평탄 사상이고, $f'$이 평탄인 것과 대응하는
사상 $g' : U' \to V'$이 평탄인 것은 동치이다. 또한
$X = Y \times_{Y'} X'$인 것과
% P09-ERR-0571: frozen cartesian identity repeats V' instead of U'.
$U = V \times_{V'} V'$인 것은 동치이다. 마지막으로 사상
$f^*\mathcal{C}_{Y/Y'} \to \mathcal{C}_{X/X'}$가 동형사상인 것과
$g^*\mathcal{C}_{V/V'} \to \mathcal{C}_{U/U'}$가 동형사상인 것은
동치이다. 따라서 보조정리는 스킴 사상에 대한 유사 명제에서 따른다.
More on Morphisms, Lemma
\ref{more-morphisms-lemma-deform}을 보라.
\end{proof}

\noindent
다음 보조정리는 “섬유별 평탄성 판정법”의 “멱영” 판이다.
Section \ref{spaces-more-morphisms-section-criterion-flat-fibres}를 보라.

\begin{lemma}
\label{spaces-more-morphisms-lemma-flatness-morphism-thickenings}
$S$를 스킴이라 하자. 가환 도식
$$
\xymatrix{
(X \subset X') \ar[rr]_{(f, f')} \ar[rd] & & (Y \subset Y') \ar[ld] \\
& (B \subset B')
}
$$
을 생각하자. 이는 $S$ 위 대수공간의 두껍게 하기들의 도식이다.
다음을 가정하자.
\begin{enumerate}
\item $X'$은 $B'$ 위에서 평탄하다.
\item $f$는 평탄하다.
\item $B \subset B'$은 유한 차수 두껍게 하기이다.
\item $X = B \times_{B'} X'$이고 $Y = B \times_{B'} Y'$이다.
\end{enumerate}
그러면 $f'$은 평탄하고, $Y'$은 $B'$ 위에서 $f'$의 상에 속하는
모든 점에서 평탄하다.
\end{lemma}

\begin{proof}
스킴 $U'$과 전사 에탈 사상 $U' \to B'$을 택하자.
스킴 $V'$과 전사 에탈 사상
$V' \to U' \times_{B'} Y'$을 택하자.
스킴 $W'$과 전사 에탈 사상
$W' \to V' \times_{Y'} X'$을 택하자. $U, V, W$를
$U', V', W'$의 밑변환이라 하되, 이 밑변환은 $B \to B'$에
의해 취하자. 그러면
$f'$의 평탄성은 $W' \to V'$의 평탄성과 동치이고, $W \to V$가
평탄하다는 것이 주어졌다. 따라서 스킴의 두껍게 하기들의 도식
$$
\xymatrix{
(W \subset W') \ar[rr] \ar[rd] & & (V \subset V') \ar[ld] \\
& (U \subset U')
}
$$
에 스킴의 경우의 보조정리를 적용할 수 있다. More on Morphisms,
Lemma
\ref{more-morphisms-lemma-flatness-morphism-thickenings}를 보라.
$Y'/B'$의 평탄성에 관한 주장도 $f'$의 상에 속하는 점들에서 같은
방식으로 따른다.
\end{proof}

\noindent
스킴 사상의 많은 성질은 평탄 변형 아래 보존된다.

\begin{lemma}
\label{spaces-more-morphisms-lemma-deform-property}
$S$를 스킴이라 하자. 가환 도식
$$
\xymatrix{
(X \subset X') \ar[rr]_{(f, f')} \ar[rd] & & (Y \subset Y') \ar[ld] \\
& (B \subset B')
}
$$
을 생각하자. 이는 $S$ 위 대수공간의 두껍게 하기들의 도식이다.
$B \subset B'$이 유한 차수 두껍게 하기이고, $X'$이 $B'$ 위에서
평탄하며, $X = B \times_{B'} X'$이고
$Y = B \times_{B'} Y'$이라고 가정하자. 그러면
\begin{enumerate}
\item $f$가 표현 가능한 것과 $f'$이 표현 가능한 것은 동치이다.
\label{spaces-more-morphisms-item-representable}
\item $f$가 평탄한 것과 $f'$이 평탄한 것은 동치이다.
\label{spaces-more-morphisms-item-flat}
\item $f$가 동형사상인 것과 $f'$이 동형사상인 것은 동치이다.
\label{spaces-more-morphisms-item-isomorphism}
\item $f$가 열린 몰입인 것과 $f'$이 열린 몰입인 것은 동치이다.
\label{spaces-more-morphisms-item-open-immersion}
\item $f$가 준콤팩트인 것과 $f'$이 준콤팩트인 것은 동치이다.
\label{spaces-more-morphisms-item-quasi-compact}
\item $f$가 보편 닫힌인 것과 $f'$이 보편 닫힌인 것은 동치이다.
\label{spaces-more-morphisms-item-universally-closed}
\item $f$가 (준)분리인 것과 $f'$이 (준)분리인 것은 동치이다.
\label{spaces-more-morphisms-item-separated}
\item $f$가 단사상인 것과 $f'$이 단사상인 것은 동치이다.
\label{spaces-more-morphisms-item-monomorphism}
\item $f$가 전사인 것과 $f'$이 전사인 것은 동치이다.
\label{spaces-more-morphisms-item-surjective}
\item $f$가 보편 단사인 것과 $f'$이 보편 단사인 것은 동치이다.
\label{spaces-more-morphisms-item-universally-injective}
\item $f$가 아핀인 것과 $f'$이 아핀인 것은 동치이다.
\label{spaces-more-morphisms-item-affine}
\item
\label{spaces-more-morphisms-item-finite-type}
$f$가 국소 유한형인 것과 $f'$이 국소 유한형인 것은 동치이다.
\item $f$가 국소 준유한인 것과 $f'$이 국소 준유한인 것은 동치이다.
\label{spaces-more-morphisms-item-quasi-finite}
\item
\label{spaces-more-morphisms-item-finite-presentation}
$f$가 국소 유한 표시인 것과 $f'$이 국소 유한 표시인 것은 동치이다.
\item
\label{spaces-more-morphisms-item-relative-dimension-d}
$f$가 상대 차원 $d$인 국소 유한형인 것과 $f'$이 상대 차원 $d$인
국소 유한형인 것은 동치이다.
\item $f$가 보편 열린인 것과 $f'$이 보편 열린인 것은 동치이다.
\label{spaces-more-morphisms-item-universally-open}
\item $f$가 신토믹인 것과 $f'$이 신토믹인 것은 동치이다.
\label{spaces-more-morphisms-item-syntomic}
\item $f$가 매끄러운 것과 $f'$이 매끄러운 것은 동치이다.
\label{spaces-more-morphisms-item-smooth}
\item $f$가 비분기인 것과 $f'$이 비분기인 것은 동치이다.
\label{spaces-more-morphisms-item-unramified}
\item $f$가 에탈인 것과 $f'$이 에탈인 것은 동치이다.
\label{spaces-more-morphisms-item-etale}
\item $f$가 고유인 것과 $f'$이 고유인 것은 동치이다.
\label{spaces-more-morphisms-item-proper}
\item $f$가 정수적인 것과 $f'$이 정수적인 것은 동치이다.
\label{spaces-more-morphisms-item-integral}
\item $f$가 유한인 것과 $f'$이 유한인 것은 동치이다.
\label{spaces-more-morphisms-item-finite}
\item
\label{spaces-more-morphisms-item-finite-locally-free}
$f$가 유한 국소 자유(랭크 $d$)인 것과 $f'$이 유한 국소 자유
(랭크 $d$)인 것은 동치이다.
% P09-ERR-0572: frozen unresolved editorial placeholder is retained.
\item 여기에 더 추가할 것.
\end{enumerate}
\end{lemma}

\begin{proof}
경우 (\ref{spaces-more-morphisms-item-representable})는 Lemma
\ref{spaces-more-morphisms-lemma-thicken-property-morphisms}에서 따른다.

\medskip\noindent
스킴 $U'$과 전사 에탈 사상 $U' \to B'$을 택하자.
스킴 $V'$과 전사 에탈 사상
$V' \to U' \times_{B'} Y'$을 택하자.
스킴 $W'$과 전사 에탈 사상
$W' \to V' \times_{Y'} X'$을 택하자. $U, V, W$를
$U', V', W'$의 밑변환이라 하되, 이 밑변환은 $B \to B'$에 의해
취하자. 스킴의
두껍게 하기들의 도식
$$
\xymatrix{
(W \subset W') \ar[rr] \ar[rd] & & (V \subset V') \ar[ld] \\
& (U \subset U')
}
$$
을 생각하자. 원천과 타깃에 대해 에탈 국소적인 성질들 각각에
대해서는, 위 도식에 스킴의 두껍게 하기 사이의 사상에 관한 대응
결과를 적용하면 결론이 즉시 따른다. 따라서 경우
(\ref{spaces-more-morphisms-item-flat}), (\ref{spaces-more-morphisms-item-finite-type}),
(\ref{spaces-more-morphisms-item-quasi-finite}), (\ref{spaces-more-morphisms-item-finite-presentation}),
(\ref{spaces-more-morphisms-item-relative-dimension-d}), (\ref{spaces-more-morphisms-item-syntomic}),
(\ref{spaces-more-morphisms-item-smooth}), (\ref{spaces-more-morphisms-item-unramified}), (\ref{spaces-more-morphisms-item-etale})는
More on Morphisms, Lemma
\ref{more-morphisms-lemma-deform-property}의 대응하는 경우들에서
따른다.

\medskip\noindent
$X \to X'$과 $Y \to Y'$은 보편 위상동형이므로, 사상
$X \to Y$와 $X' \to Y'$의 위상에 관한 모든 물음은 같은 답을
갖는다. 따라서 경우 (\ref{spaces-more-morphisms-item-quasi-compact}),
(\ref{spaces-more-morphisms-item-universally-closed}),
(\ref{spaces-more-morphisms-item-surjective}), (\ref{spaces-more-morphisms-item-universally-injective}), and
(\ref{spaces-more-morphisms-item-universally-open})가 성립한다.

\medskip\noindent
남은 경우 각각에서는 함의
$f\text{가 }P \Rightarrow f'\text{이 }P$
만 증명한다. 다른 함의는 $P$가 밑변환 아래 안정적이라는 사실에서
따른다. Spaces, Lemma \ref{spaces-lemma-base-change-immersions}와
Morphisms of Spaces, Lemmas
\ref{spaces-morphisms-lemma-base-change-separated},
\ref{spaces-morphisms-lemma-base-change-monomorphism},
\ref{spaces-morphisms-lemma-base-change-affine},
\ref{spaces-morphisms-lemma-base-change-proper},
\ref{spaces-morphisms-lemma-base-change-integral}, and
\ref{spaces-morphisms-lemma-base-change-finite-locally-free}를 보라.

\medskip\noindent
경우 (\ref{spaces-more-morphisms-item-open-immersion}). $f$가 열린 몰입이라고 가정하자.
그러면 $f'$은 (\ref{spaces-more-morphisms-item-etale})에 의해 에탈이고
(\ref{spaces-more-morphisms-item-universally-injective})에 의해 보편 단사이므로
$f'$은 열린 몰입이다. Morphisms of Spaces, Lemma
\ref{spaces-morphisms-lemma-etale-universally-injective-open}을 보라.
먼저 (\ref{spaces-more-morphisms-item-representable})을 사용하여 스킴의 경우로
환원하면 이 보조정리를 쓰지 않을 수도 있다.

\medskip\noindent
경우 (\ref{spaces-more-morphisms-item-isomorphism}). 경우
(\ref{spaces-more-morphisms-item-open-immersion})과 (\ref{spaces-more-morphisms-item-surjective})에서 따른다.

\medskip\noindent
경우 (\ref{spaces-more-morphisms-item-separated}). Lemma
\ref{spaces-more-morphisms-lemma-thicken-property-morphisms}를 보라.

\medskip\noindent
경우 (\ref{spaces-more-morphisms-item-monomorphism}). $f$가 단사상이라고 가정하자.
대각 사상 $\Delta_{X'/Y'} : X' \to X' \times_{Y'} X'$을 생각하자.
$\Delta_{X'/Y'}$의 밑변환은, 이를 $B \to B'$에 의해 취할 때
$\Delta_{X/Y}$이고, 이는 가정에 의해 동형사상이다.
(\ref{spaces-more-morphisms-item-isomorphism})에 의해
$\Delta_{X'/Y'}$가 동형사상이라고 결론 내리며, 따라서
$f'$은 단사상이다.

\medskip\noindent
경우 (\ref{spaces-more-morphisms-item-affine}). Lemma
\ref{spaces-more-morphisms-lemma-thicken-property-morphisms}를 보라.

\medskip\noindent
경우 (\ref{spaces-more-morphisms-item-proper}). Lemma
\ref{spaces-more-morphisms-lemma-thicken-property-morphisms-cartesian}을 보라.

\medskip\noindent
경우 (\ref{spaces-more-morphisms-item-integral}). Lemma
\ref{spaces-more-morphisms-lemma-thicken-property-morphisms}를 보라.

\medskip\noindent
경우 (\ref{spaces-more-morphisms-item-finite}). Lemma
\ref{spaces-more-morphisms-lemma-thicken-property-morphisms-cartesian}을 보라.

\medskip\noindent
경우 (\ref{spaces-more-morphisms-item-finite-locally-free}). $f$가 유한 국소 자유라고
가정하자. (\ref{spaces-more-morphisms-item-finite})에 의해 $f'$은 유한이다.
(\ref{spaces-more-morphisms-item-flat})에 의해 $f'$은 평탄하다.
(\ref{spaces-more-morphisms-item-finite-presentation})에 의해 $f'$은 국소 유한
표시이다. 따라서 Morphisms of Spaces, Lemma
\ref{spaces-morphisms-lemma-finite-flat}에 의해 $f'$은 유한 국소
자유이다.
\end{proof}















\noindent
다음 보조정리는 “섬유별 평탄성 판정법”의 “국소 멱영” 판이다.
Section \ref{spaces-more-morphisms-section-criterion-flat-fibres}를 보라.

\begin{lemma}
\label{spaces-more-morphisms-lemma-flatness-morphism-thickenings-fp-over-ft}
$S$를 스킴이라 하자. 가환 도식
$$
\xymatrix{
(X \subset X') \ar[rr]_{(f, f')} \ar[rd] & & (Y \subset Y') \ar[ld] \\
& (B \subset B')
}
$$
을 생각하자. 이는 $S$ 위 대수공간의 두껍게 하기들의 도식이다.
다음을 가정하자.
\begin{enumerate}
\item $Y' \to B'$은 국소 유한형이다.
\item $X' \to B'$은 평탄하고 국소 유한 표시이다.
\item $f$는 평탄하다.
\item $X = B \times_{B'} X'$이고 $Y = B \times_{B'} Y'$이다.
\end{enumerate}
그러면 $f'$은 평탄하고, 모든 $y' \in |Y'|$에 대해 이 점이
$|f'|$의 상에 속하면 사상 $Y' \to B'$은 $y'$에서 평탄하다.
\end{lemma}

\begin{proof}
스킴 $U'$과 전사 에탈 사상 $U' \to B'$을 택하자.
스킴 $V'$과 전사 에탈 사상
$V' \to U' \times_{B'} Y'$을 택하자.
스킴 $W'$과 전사 에탈 사상
$W' \to V' \times_{Y'} X'$을 택하자. $U, V, W$를
$U', V', W'$의 밑변환이라 하되, 이 밑변환은 $B \to B'$에 의해
취하자. 그러면 $f'$의 평탄성은 $W' \to V'$의 평탄성과 동치이고,
$W \to V$가 평탄하다는 것이 주어졌다. 따라서 스킴의 두껍게
하기들의 도식
$$
\xymatrix{
(W \subset W') \ar[rr] \ar[rd] & & (V \subset V') \ar[ld] \\
& (U \subset U')
}
$$
에 스킴의 경우의 보조정리를 적용할 수 있다. More on Morphisms,
Lemma
\ref{more-morphisms-lemma-flatness-morphism-thickenings-fp-over-ft}를
보라. $Y'/B'$의 평탄성에 관한 주장도 $f'$의 상에 속하는 점들에서
같은 방식으로 따른다.
\end{proof}

\noindent
스킴 사상의 많은 성질은 위 보조정리에서와 같은 평탄 변형 아래
보존된다.

\begin{lemma}
\label{spaces-more-morphisms-lemma-deform-property-fp-over-ft}
$S$를 스킴이라 하자. 가환 도식
$$
\xymatrix{
(X \subset X') \ar[rr]_{(f, f')} \ar[rd] & & (Y \subset Y') \ar[ld] \\
& (B \subset B')
}
$$
을 생각하자. 이는 $S$ 위 대수공간의 두껍게 하기들의 도식이다.
$Y' \to B'$이 국소 유한형이고,
$X' \to B'$이 평탄하고 국소 유한 표시이며,
$X = B \times_{B'} X'$이고 $Y = B \times_{B'} Y'$이라고
가정하자. 그러면
\begin{enumerate}
\item $f$가 표현 가능한 것과 $f'$이 표현 가능한 것은 동치이다.
\label{spaces-more-morphisms-item-representable-fp-over-ft}
\item $f$가 평탄한 것과 $f'$이 평탄한 것은 동치이다.
\label{spaces-more-morphisms-item-flat-fp-over-ft}
\item $f$가 동형사상인 것과 $f'$이 동형사상인 것은 동치이다.
\label{spaces-more-morphisms-item-isomorphism-fp-over-ft}
\item $f$가 열린 몰입인 것과 $f'$이 열린 몰입인 것은 동치이다.
\label{spaces-more-morphisms-item-open-immersion-fp-over-ft}
\item $f$가 준콤팩트인 것과 $f'$이 준콤팩트인 것은 동치이다.
\label{spaces-more-morphisms-item-quasi-compact-fp-over-ft}
\item $f$가 보편 닫힌인 것과 $f'$이 보편 닫힌인 것은 동치이다.
\label{spaces-more-morphisms-item-universally-closed-fp-over-ft}
\item $f$가 (준)분리인 것과 $f'$이 (준)분리인 것은 동치이다.
\label{spaces-more-morphisms-item-separated-fp-over-ft}
\item $f$가 단사상인 것과 $f'$이 단사상인 것은 동치이다.
\label{spaces-more-morphisms-item-monomorphism-fp-over-ft}
\item $f$가 전사인 것과 $f'$이 전사인 것은 동치이다.
\label{spaces-more-morphisms-item-surjective-fp-over-ft}
\item $f$가 보편 단사인 것과 $f'$이 보편 단사인 것은 동치이다.
\label{spaces-more-morphisms-item-universally-injective-fp-over-ft}
\item $f$가 아핀인 것과 $f'$이 아핀인 것은 동치이다.
\label{spaces-more-morphisms-item-affine-fp-over-ft}
\item $f$가 국소 준유한인 것과 $f'$이 국소 준유한인 것은 동치이다.
\label{spaces-more-morphisms-item-quasi-finite-fp-over-ft}
\item
\label{spaces-more-morphisms-item-relative-dimension-d-fp-over-ft}
$f$가 상대 차원 $d$인 국소 유한형인 것과 $f'$이 상대 차원 $d$인
국소 유한형인 것은 동치이다.
\item $f$가 보편 열린인 것과 $f'$이 보편 열린인 것은 동치이다.
\label{spaces-more-morphisms-item-universally-open-fp-over-ft}
\item $f$가 신토믹인 것과 $f'$이 신토믹인 것은 동치이다.
\label{spaces-more-morphisms-item-syntomic-fp-over-ft}
\item $f$가 매끄러운 것과 $f'$이 매끄러운 것은 동치이다.
\label{spaces-more-morphisms-item-smooth-fp-over-ft}
\item $f$가 비분기인 것과 $f'$이 비분기인 것은 동치이다.
\label{spaces-more-morphisms-item-unramified-fp-over-ft}
\item $f$가 에탈인 것과 $f'$이 에탈인 것은 동치이다.
\label{spaces-more-morphisms-item-etale-fp-over-ft}
\item $f$가 고유인 것과 $f'$이 고유인 것은 동치이다.
\label{spaces-more-morphisms-item-proper-fp-over-ft}
\item $f$가 유한인 것과 $f'$이 유한인 것은 동치이다.
\label{spaces-more-morphisms-item-finite-fp-over-ft}
\item
\label{spaces-more-morphisms-item-finite-locally-free-fp-over-ft}
$f$가 유한 국소 자유(랭크 $d$)인 것과 $f'$이 유한 국소 자유
(랭크 $d$)인 것은 동치이다.
% P09-ERR-0573: frozen unresolved editorial placeholder is retained.
\item 여기에 더 추가할 것.
\end{enumerate}
\end{lemma}

\begin{proof}
경우 (\ref{spaces-more-morphisms-item-representable-fp-over-ft})는 Lemma
\ref{spaces-more-morphisms-lemma-thicken-property-morphisms}에서 따른다.

\medskip\noindent
스킴 $U'$과 전사 에탈 사상 $U' \to B'$을 택하자.
스킴 $V'$과 전사 에탈 사상
$V' \to U' \times_{B'} Y'$을 택하자.
스킴 $W'$과 전사 에탈 사상
$W' \to V' \times_{Y'} X'$을 택하자. $U, V, W$를
$U', V', W'$의 밑변환이라 하되, 이 밑변환은 $B \to B'$에 의해
취하자. 스킴의 두껍게 하기들의 도식
$$
\xymatrix{
(W \subset W') \ar[rr] \ar[rd] & & (V \subset V') \ar[ld] \\
& (U \subset U')
}
$$
을 생각하자. 원천과 타깃에 대해 에탈 국소적인 성질들 각각에
대해서는, 위 도식에 스킴의 두껍게 하기 사이의 사상에 관한 대응
결과를 적용하면 결론이 즉시 따른다. 따라서 경우
(\ref{spaces-more-morphisms-item-flat-fp-over-ft}),
(\ref{spaces-more-morphisms-item-quasi-finite-fp-over-ft}),
(\ref{spaces-more-morphisms-item-relative-dimension-d-fp-over-ft}),
(\ref{spaces-more-morphisms-item-syntomic-fp-over-ft}),
(\ref{spaces-more-morphisms-item-smooth-fp-over-ft}),
(\ref{spaces-more-morphisms-item-unramified-fp-over-ft}),
(\ref{spaces-more-morphisms-item-etale-fp-over-ft})는
More on Morphisms, Lemma
\ref{more-morphisms-lemma-deform-property-fp-over-ft}의 대응하는
경우들에서 따른다.

\medskip\noindent
$X \to X'$과 $Y \to Y'$은 보편 위상동형이므로, 사상
$X \to Y$와 $X' \to Y'$의 위상에 관한 모든 물음은 같은 답을
갖는다. 따라서 경우
(\ref{spaces-more-morphisms-item-quasi-compact-fp-over-ft}),
(\ref{spaces-more-morphisms-item-universally-closed-fp-over-ft}),
(\ref{spaces-more-morphisms-item-surjective-fp-over-ft}),
(\ref{spaces-more-morphisms-item-universally-injective-fp-over-ft}), and
(\ref{spaces-more-morphisms-item-universally-open-fp-over-ft})가 성립한다.

\medskip\noindent
남은 경우 각각에서는 함의
$f\text{가 }P \Rightarrow f'\text{이 }P$
만 증명한다. 다른 함의는 $P$가 밑변환 아래 안정적이라는 사실에서
따른다. Spaces, Lemma \ref{spaces-lemma-base-change-immersions}와
Morphisms of Spaces, Lemmas
\ref{spaces-morphisms-lemma-base-change-separated},
\ref{spaces-morphisms-lemma-base-change-monomorphism},
\ref{spaces-morphisms-lemma-base-change-affine},
\ref{spaces-morphisms-lemma-base-change-proper},
\ref{spaces-morphisms-lemma-base-change-integral}, and
\ref{spaces-morphisms-lemma-base-change-finite-locally-free}를 보라.

\medskip\noindent
경우 (\ref{spaces-more-morphisms-item-open-immersion-fp-over-ft}).
$f$가 열린 몰입이라고 가정하자. 그러면 $f'$은
(\ref{spaces-more-morphisms-item-etale-fp-over-ft})에 의해 에탈이고
(\ref{spaces-more-morphisms-item-universally-injective-fp-over-ft})에 의해 보편 단사이므로
$f'$은 열린 몰입이다. Morphisms of Spaces, Lemma
\ref{spaces-morphisms-lemma-etale-universally-injective-open}을 보라.
먼저 (\ref{spaces-more-morphisms-item-representable-fp-over-ft})을 사용하여 스킴의 경우로
환원하면 이 보조정리를 쓰지 않을 수도 있다.

\medskip\noindent
경우 (\ref{spaces-more-morphisms-item-isomorphism-fp-over-ft}). 경우
(\ref{spaces-more-morphisms-item-open-immersion-fp-over-ft})과
(\ref{spaces-more-morphisms-item-surjective-fp-over-ft})에서 따른다.

\medskip\noindent
경우 (\ref{spaces-more-morphisms-item-separated-fp-over-ft}). Lemma
\ref{spaces-more-morphisms-lemma-thicken-property-morphisms}를 보라.

\medskip\noindent
경우 (\ref{spaces-more-morphisms-item-monomorphism-fp-over-ft}). $f$가 단사상이라고
가정하자. 대각 사상
$\Delta_{X'/Y'} : X' \to X' \times_{Y'} X'$을 생각하자.
$\Delta_{X'/Y'}$의 밑변환은, 이를 $B \to B'$에 의해 취할 때
$\Delta_{X/Y}$이고, 이는 가정에 의해 동형사상이다.
(\ref{spaces-more-morphisms-item-isomorphism-fp-over-ft})에 의해
$\Delta_{X'/Y'}$가 동형사상이라고 결론 내리며, 따라서
$f'$은 단사상이다.

\medskip\noindent
경우 (\ref{spaces-more-morphisms-item-affine-fp-over-ft}). Lemma
\ref{spaces-more-morphisms-lemma-thicken-property-morphisms}를 보라.

\medskip\noindent
경우 (\ref{spaces-more-morphisms-item-proper-fp-over-ft}). Lemma
\ref{spaces-more-morphisms-lemma-properties-that-extend-over-thickenings}를 보라.

\medskip\noindent
경우 (\ref{spaces-more-morphisms-item-finite-fp-over-ft}). Lemma
\ref{spaces-more-morphisms-lemma-properties-that-extend-over-thickenings}를 보라.

\medskip\noindent
경우 (\ref{spaces-more-morphisms-item-finite-locally-free-fp-over-ft}).
$f$가 유한 국소 자유라고 가정하자.
(\ref{spaces-more-morphisms-item-finite-fp-over-ft})에 의해 $f'$은 유한이다.
(\ref{spaces-more-morphisms-item-flat-fp-over-ft})에 의해 $f'$은 평탄하다.
% P09-ERR-0574: frozen missing "of" in locally of finite presentation is recorded.
또한 Morphisms of Spaces, Lemma
\ref{spaces-morphisms-lemma-finite-presentation-permanence}에 의해
$f'$은 국소 유한 표시이다. 따라서 Morphisms of Spaces, Lemma
\ref{spaces-morphisms-lemma-finite-flat}에 의해 $f'$은 유한 국소
자유이다.
\end{proof}

\section{형식적으로 매끄러운 사상}
\label{spaces-more-morphisms-section-formally-smooth}

\noindent
이 절에서는 대수공간의 형식적으로 매끄러운 사상 $X \to Y$라는
개념을 도입한다. 이러한 사상은 $T$-값 점을 생각할 때, 이것이 $X$의
점이면 $T$의 무한소 두껍게 하기로 올라가며, 여기서 $T$가 아핀이라는
성질로 특징지어진다.
주요 결과는 형식적으로 매끄럽고 국소 유한 표시인 사상이
매끄럽다는 것이다. Lemma \ref{spaces-more-morphisms-lemma-smooth-formally-smooth}를 보라.
이 판정법은 야코비 판정법보다 사용하기 쉬운 경우가 많다.

\begin{definition}
\label{spaces-more-morphisms-definition-formally-smooth}
$S$를 스킴이라 하자. $f : X \to Y$를 $S$ 위 대수공간의 사상이라
하자. 이 사상이 {\it 형식적으로 매끄럽다}는 것은 Definition
\ref{spaces-more-morphisms-definition-formally-smooth-etale-unramified}에서처럼 함자의
변환으로서 형식적으로 매끄럽다는 뜻이다.
\end{definition}

\noindent
형식적으로 비분기인 사상과 형식적으로 에탈인 사상의 경우에는
$T'$이 아핀이라는 조건을 버릴 수 있었다. Lemmas
\ref{spaces-more-morphisms-lemma-formally-unramified-not-affine} and
\ref{spaces-more-morphisms-lemma-formally-etale-not-affine}를 보라. 형식적으로 매끄러운
사상의 경우에는 더 이상 그렇지 않다. 사실 조금 더 자연스러운
조건은 $T'$ 위 에탈 국소적으로 점선 화살표를 채울 수 있어야 한다는
것이다. Lemma \ref{spaces-more-morphisms-lemma-smooth-formally-smooth}의 증명을 분석하면
이 조건이 현재 형태의 정의와 동치임을 알 수 있다. 또한 $T'$ 위
fppf 국소적으로 점선 화살표의 존재를 요구하는 것으로도 충분하지만,
이를 증명하기는 조금 더 어렵다.

\medskip\noindent
Section \ref{spaces-more-morphisms-section-formally-smooth-etale-unramified}에서 함자의
형식적으로 매끄러운 변환이라는 더 일반적인 상황에서 증명한
결과들은 여기서 다시 진술하지 않는다.

\begin{lemma}
\label{spaces-more-morphisms-lemma-composition-formally-smooth}
형식적으로 매끄러운 사상들의 합성은 형식적으로 매끄럽다.
\end{lemma}

\begin{proof}
생략한다.
\end{proof}

\begin{lemma}
\label{spaces-more-morphisms-lemma-base-change-formally-smooth}
형식적으로 매끄러운 사상의 밑변환은 형식적으로 매끄럽다.
\end{lemma}

\begin{proof}
생략한다. 다만 대수적 판에 대해서는 Algebra, Lemma
\ref{algebra-lemma-base-change-fs}를 보라.
\end{proof}

\begin{lemma}
\label{spaces-more-morphisms-lemma-formally-etale-unramified-smooth}
$f : X \to S$를 스킴의 사상이라 하자.
그러면 $f$가 형식적으로 에탈인 것과 $f$가 형식적으로 매끄럽고
형식적으로 비분기인 것은 동치이다.
\end{lemma}

\begin{proof}
생략한다.
\end{proof}

\noindent
다음은 Lemma \ref{spaces-more-morphisms-lemma-formally-smooth}에 의해 대체될 보조정리이다.

\begin{lemma}
\label{spaces-more-morphisms-lemma-helper-formally-smooth}
$S$를 스킴이라 하자. 가환 도식
$$
\xymatrix{
U \ar[d] \ar[r]_\psi & V \ar[d] \\
X \ar[r]^f & Y
}
$$
을 생각하자. 이는 $S$ 위 대수공간의 사상들로 이루어진다.
수직 화살표들이 에탈이고 $f$가 형식적으로 매끄러우면
$\psi$도 형식적으로 매끄럽다.
\end{lemma}

\begin{proof}
Lemma \ref{spaces-more-morphisms-lemma-representable-property-formally-property}에 의해
사상 $U \to X$와 $V \to Y$는 형식적으로 에탈이다. Lemma
\ref{spaces-more-morphisms-lemma-composition-formally-smooth-etale-unramified}에 의해 합성
$U \to Y$는 형식적으로 매끄럽다. Lemma
\ref{spaces-more-morphisms-lemma-formally-permanence}에 의해
$\psi : U \to V$가 형식적으로 매끄러움을 알 수 있다.
\end{proof}

\noindent
다음 보조정리는 이 절의 주요 결과이다. Limits of Spaces,
Proposition
\ref{spaces-limits-proposition-characterize-locally-finite-presentation}과
결합하면, 대수공간의 사상 $f : X \to Y$가 매끄러운지 함자의 변환
$X \to Y$의 “단순한” 성질들로 판별할 수 있음을 함의한다.

\begin{lemma}[무한소 올림 판정법]
\label{spaces-more-morphisms-lemma-smooth-formally-smooth}
$S$를 스킴이라 하자.
$f : X \to Y$를 $S$ 위 대수공간의 사상이라 하자.
다음은 동치이다.
\begin{enumerate}
\item 사상 $f$는 매끄럽다.
\item 사상 $f$는 국소 유한 표시이고 형식적으로 매끄럽다.
\end{enumerate}
\end{lemma}

\begin{proof}
% P09-ERR-0575: frozen source uses codomain S here although the lemma has f : X -> Y.
$f : X \to S$가 국소 유한 표시이며 형식적으로 매끄럽다고 가정하자.
가환 도식
$$
\xymatrix{
U \ar[d] \ar[r]_\psi & V \ar[d] \\
X \ar[r]^f & Y
}
$$
을 생각하자. 여기서 $U$와 $V$는 스킴이고 수직 화살표들은 에탈
전사이다. Lemma \ref{spaces-more-morphisms-lemma-helper-formally-smooth}에 의해
$\psi : U \to V$가 형식적으로 매끄러움을 알 수 있다. Morphisms of
Spaces, Lemma
\ref{spaces-morphisms-lemma-finite-presentation-local}에 의해 사상
$\psi$는 국소 유한 표시이다. 따라서 스킴의 경우에 의해 사상
$\psi$는 매끄럽다. More on Morphisms, Lemma
\ref{more-morphisms-lemma-smooth-formally-smooth}를 보라. 그러므로
$f$는 매끄럽다. Morphisms of Spaces, Lemma
\ref{spaces-morphisms-lemma-smooth-local}을 보라.

\medskip\noindent
역으로 $f : X \to Y$가 매끄럽다고 가정하자. Definition
\ref{spaces-more-morphisms-definition-formally-smooth}에서와 같은 실선 가환 도식
$$
\xymatrix{
X \ar[d]_f & T \ar[d]^i \ar[l]^a \\
Y & T' \ar[l] \ar@{-->}[lu]
}
$$
을 생각하자. 점선 화살표가 존재함을 보여서 $f$가 형식적으로
매끄러움을 증명하겠다. $\mathcal{F}$를 $(T')_{spaces, \etale}$ 위의
Lemma \ref{spaces-more-morphisms-lemma-sheaf}의 집합 층이라 하되, Remark
\ref{spaces-more-morphisms-remark-special-case}에서 논한 특별한 경우를 취하자. 또한
$$
\mathcal{H} =
\SheafHom_{\mathcal{O}_T}(a^*\Omega_{X/Y}, \mathcal{C}_{T/T'})
$$
를 $\mathcal{O}_T$-가군의 층이라 하며, 이는 $T_{spaces, \etale}$ 위에
놓인다. 여기에
Lemma \ref{spaces-more-morphisms-lemma-action-sheaf}에서와 같은 작용
$\mathcal{H} \times \mathcal{F} \to \mathcal{F}$를 취한다. 이 작용
$\mathcal{H} \times \mathcal{F} \to \mathcal{F}$는
$\mathcal{F}$를 유사 $\mathcal{H}$-토서로 만든다. Cohomology on Sites,
Definition \ref{sites-cohomology-definition-torsor}를 보라. 목표는
$\mathcal{F}$가 자명한 $\mathcal{H}$-토서임을 보이는 것이다. 두 단계가
있다. (I) $\mathcal{F}$가 토서임을 보이려면 $\mathcal{F}$가 에탈
국소적으로 단면을 가짐을 보여야 한다. (II) $\mathcal{F}$가 자명한
토서임을 보이려면
$H^1(T_\etale, \mathcal{H}) = 0$임을 보이면 충분하다. Cohomology on
Sites, Lemma \ref{sites-cohomology-lemma-torsors-h1}을 보라.

\medskip\noindent
먼저 (I)을 증명한다. 이를 위해 가환 도식
$$
\xymatrix{
U \ar[d] \ar[r]_\psi & V \ar[d] \\
X \ar[r]^f & Y
}
$$
을 택하자. 여기서 $U$와 $V$는 스킴이고 수직 화살표들은 에탈
전사이다. $f$가 매끄럽다고 가정했으므로 $\psi$도 매끄럽고, 따라서
Lemma \ref{spaces-more-morphisms-lemma-representable-property-formally-property}에 의해
형식적으로 매끄럽다. 같은 보조정리에 의해 사상 $V \to Y$는
형식적으로 에탈이다. 그러므로 Lemma
\ref{spaces-more-morphisms-lemma-composition-formally-smooth-etale-unramified}에 의해 합성
$U \to Y$는 형식적으로 매끄럽다. 이제 Lemma
\ref{spaces-more-morphisms-lemma-etale-on-top}의 (4)에서 (I)이 따른다.

\medskip\noindent
마지막으로 (II)를 증명한다. Lemma
\ref{spaces-more-morphisms-lemma-finite-presentation-differentials}에 의해
% P09-ERR-0576: frozen source uses Omega_{X/S} twice where the argument needs Omega_{X/Y}.
$\Omega_{X/S}$는 유한 표시이다. 따라서
$a^*\Omega_{X/S}$는 유한 표시이다(Properties of Spaces, Section
\ref{spaces-properties-section-properties-modules}를 보라). 따라서 층
$\mathcal{H} =
\SheafHom_{\mathcal{O}_T}(a^*\Omega_{X/Y}, \mathcal{C}_{T/T'})$는
Properties of Spaces, Lemma
\ref{spaces-properties-lemma-properties-quasi-coherent}에 의해 준연접이다.
그러므로 Descent, Proposition
\ref{descent-proposition-same-cohomology-quasi-coherent} 및 Cohomology of
Schemes, Lemma
\ref{coherent-lemma-quasi-coherent-affine-cohomology-zero}에 의해
$$
H^1(T_{spaces, \etale}, \mathcal{H}) =
H^1(T_\etale, \mathcal{H}) =
H^1(T, \mathcal{H}) = 0
$$
을 얻으며, 이것이 원한 바이다.
\end{proof}

\noindent
매끄러운 사상은 강한 국소 올림 성질을 만족한다. Lemma
\ref{spaces-more-morphisms-lemma-smooth-strong-lift}를 보라. 그 보조정리에서 $T'$이
아핀이라고 가정할 때 에탈 덮개를 취해야 하는지는 알려져 있지 않다.
더 정확히 말해, 대수공간의 가환 도식
$$
\xymatrix{
X \ar[d] & T \ar[l] \ar[d] \\
Y & T' \ar[l] \ar@{..>}[lu]
}
$$
에서 $X \to Y$가 매끄럽고 $T \to T'$이 아핀 스킴들의 두껍게 하기일
때, 도식을 가환으로 만드는 점선 화살표가 항상 존재하는가?
% P09-ERR-0577: frozen source has the grammatical corruption "the does" in this question.
답이나 참고문헌을 알고 있다면
\href{mailto:stacks.project@gmail.com}{stacks.project@gmail.com}으로
이메일을 보내 주기 바란다.

\begin{lemma}
\label{spaces-more-morphisms-lemma-smooth-strong-lift}
$S$를 스킴이라 하자. 가환 도식
$$
\xymatrix{
X \ar[d] & T \ar[l] \ar[d] \\
Y & T' \ar[l]
}
$$
을 생각하자. 이는 $S$ 위 대수공간의 도식이며, 여기서 $X \to Y$는
매끄럽고 $T \to T'$은 두껍게 하기이다. 그러면 에탈 덮개
$\{T'_i \to T'\}$가 존재하여, 아래 도식에서 점선 화살표를 찾을 수
있다.
$$
\xymatrix{
X \ar[d] & T \ar[l] \ar[d] & T \times_{T'} T'_i \ar[l] \ar[d] \\
Y & T' \ar[l] & T'_i \ar[l] \ar@{..>}[llu]
}
$$
이 화살표는 모든 $i$에 대해 도식을 가환으로 만든다.
\end{lemma}

\begin{proof}
$\{Y_i \to Y\}$를 에탈 덮개로 택하되, 각 $Y_i$는 아핀이라 하자. $T'$을
유도된 에탈 덮개로 바꾸면 $Y$가 아핀이라고 가정할 수 있다.

\medskip\noindent
$Y$가 아핀이라고 가정하자. 에탈 덮개 $\{X_i \to X\}$를 택한다.
이로부터 $T$의 에탈 덮개를 얻는다. 이 $T$의 에탈 덮개는 $T'$의
에탈 덮개에서 온다(Theorem \ref{spaces-more-morphisms-theorem-topological-invariance} 및
Section \ref{spaces-more-morphisms-section-thickenings}의 논의를 보라). 따라서 $X$가
아핀이라고 가정할 수 있다.

\medskip\noindent
$X$와 $Y$가 아핀이라고 가정하자. $T'$에 에탈 덮개를 한 번 더
취하여 $T'$이 아핀이라고 가정할 수 있다. 이 경우 보조정리는
Algebra, Lemma \ref{algebra-lemma-smooth-strong-lift}에서 따른다.
\end{proof}







\noindent
형식적 매끄러움이 원천에서 에탈 국소적 성질임을 보이기 위해 조금 더
작업한다. 먼저 형식적으로 매끄러운 사상의 미분 층이 좋은 성질을
가짐을 보이자. 국소 사영 준연접 가군의 개념은 Properties of Spaces,
Section \ref{spaces-properties-section-locally-projective}에서 정의된다.

\begin{lemma}
\label{spaces-more-morphisms-lemma-formally-smooth-sheaf-differentials}
$S$를 스킴이라 하자. $f : X \to Y$를 $S$ 위 대수공간의 형식적으로
매끄러운 사상이라 하자. 그러면 $\Omega_{X/Y}$는 $X$ 위에서 국소
사영이다.
\end{lemma}

\begin{proof}
도식
$$
\xymatrix{
U \ar[d] \ar[r]_\psi & V \ar[d] \\
X \ar[r]^f & Y
}
$$
을 택하자. 여기서 $U$와 $V$는 아핀(!) 스킴이고 수직 화살표들은
에탈이다. Lemma \ref{spaces-more-morphisms-lemma-helper-formally-smooth}에 의해
$\psi : U \to V$는 형식적으로 매끄럽다. 따라서
$\Gamma(V, \mathcal{O}_V) \to \Gamma(U, \mathcal{O}_U)$는 형식적으로
매끄러운 환 준동형이다. More on Morphisms, Lemma
\ref{more-morphisms-lemma-affine-formally-smooth}를 보라. 따라서 Algebra,
Lemma \ref{algebra-lemma-characterize-formally-smooth-again}에 의해
$\Gamma(U, \mathcal{O}_U)$-가군
$\Omega_{\Gamma(U, \mathcal{O}_U)/\Gamma(V, \mathcal{O}_V)}$는
사영이다. 그러므로 $\Omega_{U/V}$는 국소 사영이다. Properties,
Section \ref{properties-section-locally-projective}을 보라.
$\Omega_{X/Y}|_U = \Omega_{U/V}$이므로 $\Omega_{X/Y}$도 국소
사영이다. (실제로 $X$를 덮으며 위와 같은 도식에 들어가는 아핀
$U$들의 에탈 덮개를 찾을 수 있다. 세부사항은 생략한다.)
\end{proof}

\begin{lemma}
\label{spaces-more-morphisms-lemma-h1-is-zero}
$T$를 아핀 스킴이라 하자. $\mathcal{F}$와 $\mathcal{G}$를 준연접
$\mathcal{O}_T$-가군이라 하되, 이들은 $T_\etale$ 위에 놓인다고 하자.
내부 준동형 층
$\mathcal{H} = \SheafHom_{\mathcal{O}_T}(\mathcal{F}, \mathcal{G})$를
$T_\etale$ 위에서 생각하자. $\mathcal{F}$가 국소 사영이면
$H^1(T_\etale, \mathcal{H}) = 0$이다.
\end{lemma}

\begin{proof}
대수공간 위 국소 사영 층의 정의(Properties of Spaces, Definition
\ref{spaces-properties-definition-locally-projective} 참조)에 의해
$\mathcal{F}_{Zar} = \mathcal{F}|_{T_{Zar}}$는 스킴 $T$ 위 국소 사영
층이다. 따라서 $\mathcal{F}_{Zar}$는 자유
$\mathcal{O}_{T_{Zar}}$-가군의 직합인자이다. 그러므로
$\mathcal{F} = (\mathcal{F}_{Zar})^a$임을 이용하면(Descent,
Proposition \ref{descent-proposition-equivalence-quasi-coherent} 참조),
$\mathcal{F}$가 자유 $\mathcal{O}_T$-가군의 직합인자임을 얻으며, 이
가군은 $T_\etale$ 위에 놓인다. 따라서
$\mathcal{F} = \bigoplus_{i \in I} \mathcal{O}_T$인 자유 가군의
경우로 가정할 수 있다. 이 경우
$\mathcal{H} = \prod_{i \in I} \mathcal{G}$는 준연접 가군들의
곱이다. Cohomology on Sites, Lemma
\ref{sites-cohomology-lemma-cohomology-products}에 의해 $H^1 = 0$이다.
실제로 아핀 스킴 위 준연접 층의 코호몰로지는 영이다. Descent,
Proposition \ref{descent-proposition-same-cohomology-quasi-coherent} 및
Cohomology of Schemes, Lemma
\ref{coherent-lemma-quasi-coherent-affine-cohomology-zero}를 보라.
\end{proof}

\begin{lemma}
\label{spaces-more-morphisms-lemma-formally-smooth}
$S$를 스킴이라 하자. $f : X \to Y$를 $S$ 위 대수공간의 사상이라
하자. 다음은 동치이다.
\begin{enumerate}
\item $f$는 형식적으로 매끄럽다.
\item 임의의 도식
$$
\xymatrix{
U \ar[d] \ar[r]_\psi & V \ar[d] \\
X \ar[r]^f & Y
}
$$
에서 $U$와 $V$가 스킴이고 수직 화살표들이 에탈이면, 스킴의 사상
$\psi$는 형식적으로 매끄럽다.
% P09-ERR-0578: frozen reference points to formally unramified rather than formally smooth.
(More on Morphisms, Definition
\ref{more-morphisms-definition-formally-unramified}의 의미에서).
\item 전사인 수직 화살표들을 갖는 이러한 도식 하나에 대해 사상
$\psi$가 형식적으로 매끄럽다.
\end{enumerate}
\end{lemma}

\begin{proof}
Lemma \ref{spaces-more-morphisms-lemma-helper-formally-smooth}에서 (1)이 (2)와 (3)을
함의함을 보았다. (3)을 가정하자. $f$가 형식적으로 매끄럽다는 증명은
Lemma \ref{spaces-more-morphisms-lemma-smooth-formally-smooth}의 (1) $\Rightarrow$ (2)의
증명과 완전히 같다.

\medskip\noindent
Definition \ref{spaces-more-morphisms-definition-formally-smooth}에서와 같은 실선 가환 도식
$$
\xymatrix{
X \ar[d]_f & T \ar[d]^i \ar[l]^a \\
Y & T' \ar[l] \ar@{-->}[lu]
}
$$
을 생각하자. 점선 화살표가 존재함을 보여 $f$가 형식적으로
매끄러움을 증명하겠다. $\mathcal{F}$를 $(T')_{spaces, \etale}$ 위의
Lemma \ref{spaces-more-morphisms-lemma-sheaf}의 집합 층이라 하되, Remark
\ref{spaces-more-morphisms-remark-special-case}에서 논한 특별한 경우를 취하자.
$$
\mathcal{H} =
\SheafHom_{\mathcal{O}_T}(a^*\Omega_{X/Y}, \mathcal{C}_{T/T'})
$$
를 $\mathcal{O}_T$-가군의 층이라 하며, 이는 $T_{spaces, \etale}$ 위에
놓인다. 여기에 Lemma \ref{spaces-more-morphisms-lemma-action-sheaf}에서와 같은 작용
$\mathcal{H} \times \mathcal{F} \to \mathcal{F}$를 취한다. 이 작용
$\mathcal{H} \times \mathcal{F} \to \mathcal{F}$는
$\mathcal{F}$를 유사 $\mathcal{H}$-토서로 만든다. Cohomology on Sites,
Definition \ref{sites-cohomology-definition-torsor}를 보라. 목표는
$\mathcal{F}$가 자명한 $\mathcal{H}$-토서임을 보이는 것이다. 두 단계가
있다. (I) $\mathcal{F}$가 토서임을 보이려면 $\mathcal{F}$가 에탈
국소적으로 단면을 가짐을 보여야 한다. (II) $\mathcal{F}$가 자명한
토서임을 보이려면
$H^1(T_\etale, \mathcal{H}) = 0$임을 보이면 충분하다. Cohomology on
Sites, Lemma \ref{sites-cohomology-lemma-torsors-h1}을 보라.

\medskip\noindent
먼저 (I)을 증명한다. 이를 위해 (3)을 가정했으므로 존재하는 도식
$$
\xymatrix{
U \ar[d] \ar[r]_\psi & V \ar[d] \\
X \ar[r]^f & Y
}
$$
을 생각하자. 여기서 $U$와 $V$는 스킴이고 수직 화살표들은 에탈
전사이며, $\psi$는 형식적으로 매끄럽다. Lemma
\ref{spaces-more-morphisms-lemma-representable-property-formally-property}에 의해 사상
$V \to Y$는 형식적으로 에탈이다. 따라서 Lemma
\ref{spaces-more-morphisms-lemma-composition-formally-smooth-etale-unramified}에 의해 합성
$U \to Y$는 형식적으로 매끄럽다. 이제 Lemma
\ref{spaces-more-morphisms-lemma-etale-on-top}의 (4)에서 (I)이 따른다.

\medskip\noindent
마지막으로 (II)를 증명한다. Lemma
\ref{spaces-more-morphisms-lemma-formally-smooth-sheaf-differentials}에 의해
% P09-ERR-0579: frozen source omits "is" before "locally projective".
$\Omega_{U/V}$가 국소 사영임을 알 수 있다. 따라서
$\Omega_{X/Y}$는 국소 사영이다. Descent on Spaces, Lemma
\ref{spaces-descent-lemma-locally-projective-descends}를 보라. 따라서
$a^*\Omega_{X/Y}$는 국소 사영이다. Properties of Spaces, Lemma
\ref{spaces-properties-lemma-locally-projective-pullback}을 보라. 그러므로
$$
H^1(T_\etale, \mathcal{H}) =
% P09-ERR-0580: frozen second H^1 expression lacks its closing parenthesis.
H^1(T_\etale,
\SheafHom_{\mathcal{O}_T}(a^*\Omega_{X/Y}, \mathcal{C}_{T/T'}) = 0
$$
이다. 이는 Lemma \ref{spaces-more-morphisms-lemma-h1-is-zero}에 의해 성립하며, 원한 바이다.
\end{proof}

\begin{lemma}
\label{spaces-more-morphisms-lemma-descending-property-formally-smooth}
성질 $\mathcal{P}(f) =$“$f$가 형식적으로 매끄럽다”는 밑에서 fpqc
국소적이다.
\end{lemma}

\begin{proof}
$f : X \to Y$를 스킴 $S$ 위 대수공간의 사상이라 하자. 지표집합
$I$와 도식들
$$
\xymatrix{
U_i \ar[d] \ar[r]_{\psi_i} & V_i \ar[d] \\
X \ar[r]^f & Y
}
$$
을 택하자. 여기서 수직 화살표들은 에탈이고 $U_i$, $V_i$는 아핀
스킴이다. 더 나아가 $\coprod U_i \to X$와
$\coprod V_i \to Y$가 전사라고 가정하자. Properties of Spaces,
Lemma \ref{spaces-properties-lemma-cover-by-union-affines}를 보라.
Lemma \ref{spaces-more-morphisms-lemma-formally-smooth}에 의해 $f$가 형식적으로 매끄러운
% P09-ERR-0581: frozen source has plural "are" after distributive "each".
것과 사상들 $\psi_i$가 각각 형식적으로 매끄러운 것은 동치이다.
따라서 아핀 스킴들의 사상인 경우로 환원된다. 이 경우 결과는 Algebra,
Lemma \ref{algebra-lemma-descent-formally-smooth}에서 따른다. 일부
세부사항은 생략한다.
\end{proof}

\begin{lemma}
\label{spaces-more-morphisms-lemma-triangle-differentials-formally-smooth}
$S$를 스킴이라 하자. $f : X \to Y$, $g : Y \to Z$를 $S$ 위
대수공간의 사상들이라 하자. $f$가 형식적으로 매끄럽다고 가정하자.
그러면
$$
0 \to f^*\Omega_{Y/Z} \to \Omega_{X/Z} \to \Omega_{X/Y} \to 0
$$
% P09-ERR-0582: frozen source omits "of" before the lemma reference.
라는 Lemma \ref{spaces-more-morphisms-lemma-triangle-differentials}의 열은 짧은 완전열이다.
\end{lemma}

\begin{proof}
스킴의 경우에서 따른다. More on Morphisms, Lemma
\ref{more-morphisms-lemma-triangle-differentials-formally-smooth}를 보라.
이는 에탈 국소화를 통해 적용한다. Lemmas
\ref{spaces-more-morphisms-lemma-formally-smooth} and \ref{spaces-more-morphisms-lemma-localize-differentials}를 보라.
\end{proof}

\begin{lemma}
\label{spaces-more-morphisms-lemma-differentials-formally-unramified-formally-smooth}
$S$를 스킴이라 하자. $B$를 $S$ 위 대수공간이라 하자.
$h : Z \to X$를 $B$ 위 대수공간의 형식적으로 비분기인 사상이라
하자. $Z$가 $B$ 위 형식적으로 매끄럽다고 가정하자. 그러면 Lemma
\ref{spaces-more-morphisms-lemma-universally-unramified-differentials-sequence}의 표준 완전열
$$
0 \to \mathcal{C}_{Z/X} \to h^*\Omega_{X/B} \to \Omega_{Z/B} \to 0
$$
은 짧은 완전열이다.
\end{lemma}

\begin{proof}
$Z \to Z'$을 $Z$의 보편 일차 두껍게 하기라 하되, 이는 $X$ 위에서
취한다. Lemma
\ref{spaces-more-morphisms-lemma-universally-unramified-differentials-sequence}의 증명에서
우리 열이 다음 열과 동일시됨을 알 수 있다.
$$
\mathcal{C}_{Z/Z'} \to \Omega_{Z'/B} \otimes \mathcal{O}_Z \to
\Omega_{Z/B} \to 0.
$$
% P09-ERR-0583: frozen source says Z -> S although the asserted formal smoothness is over B.
$Z \to S$가 형식적으로 매끄러우므로 $Z'$ 위 에탈 국소적으로
$Z' \to Z$를 찾을 수 있으며, 이는 $B$ 위에서 포함 사상
$Z \to Z'$의 왼쪽 역이다.
따라서 이 열은 에탈 국소적으로 분할된다. Lemma
\ref{spaces-more-morphisms-lemma-differentials-relative-immersion-section}을 보라.
\end{proof}

\begin{lemma}
\label{spaces-more-morphisms-lemma-two-unramified-morphisms-formally-smooth}
$S$를 스킴이라 하자. 다음 도식을 생각하자.
$$
\xymatrix{
Z \ar[r]_i \ar[rd]_j & X \ar[d]^f \\
& Y
}
$$
이는 $S$ 위 대수공간의 가환 도식이며, 여기서 $i$와 $j$는 형식적으로
비분기이고 $f$는 형식적으로
매끄럽다고 하자. 그러면 Lemma
\ref{spaces-more-morphisms-lemma-two-unramified-morphisms}의 표준 완전열
$$
0 \to
\mathcal{C}_{Z/Y} \to
\mathcal{C}_{Z/X} \to
i^*\Omega_{X/Y} \to 0
$$
은 완전하고 국소적으로 분할된다.
\end{lemma}

\begin{proof}
$Z \to Z'$을 $Z$의 보편 일차 두껍게 하기라 하되, 이는 $X$ 위에서
취한다. $Z \to Z''$을 $Z$의 보편 일차 두껍게 하기라 하되, 이는
$Y$ 위에서 취한다.
% P09-ERR-0584: frozen source cites the differential-sequence lemma instead of functoriality.
% P09-ERR-0585: frozen existential clause says "here is" rather than "there is".
Lemma
\ref{spaces-more-morphisms-lemma-universally-unramified-differentials-sequence}에 의해 표준 사상
$Z' \to Z''$이 존재하여 다음 가환 도식을 얻는다.
$$
\xymatrix{
Z \ar[r]_{i'} \ar[rd]_{j'} & Z' \ar[r]_a \ar[d]^k & X \ar[d]^f \\
& Z'' \ar[r]^b & Y
}
$$
위의 열은 다음 열과 동일시된다.
$$
\mathcal{C}_{Z/Z''} \to
\mathcal{C}_{Z/Z'} \to
(i')^*\Omega_{Z'/Z''} \to 0
$$
이는 형식적으로 비분기인 사상의 여법층에 관한 정의들을 통해 얻는다.
$U'' \to Z''$을 $U''$이 아핀인 에탈 사상이라 하자. $U \to Z$와
$U' \to Z'$을 각각 $Z$와 $Z'$ 위 에탈인 대응하는 아핀 스킴이라
하자. $f$가 형식적으로 매끄러우므로 사상 $h : U'' \to X$가
존재한다. 이 사상은 $i$와 $U$ 위에서 일치하고 $f \circ h$는
$b|_{U''}$와 같다.
$Z'$이 보편 일차 두껍게 하기이므로 유일한 사상
% P09-ERR-0586: frozen equality g = a \circ h is type-incompatible; likely a \circ g = h.
$g : U'' \to Z'$을 얻으며 $g = a \circ h$이다. $Z''$의 보편 성질에
의해 $k \circ g$는 포함 사상 $U'' \to Z''$이다. 따라서 $g$는
$k$의 왼쪽 역이다. 그림으로는
$$
\xymatrix{
U \ar[d] \ar[r] & Z' \ar[d]^k \\
U'' \ar[r] \ar[ru]^g & Z''
}
$$
이다. 따라서 $g$는 사상
$\mathcal{C}_{Z/Z'}|_U \to \mathcal{C}_{Z/Z''}|_U$를 유도하며, 이는
사상 $\mathcal{C}_{Z/Z''} \to \mathcal{C}_{Z/Z'}$의 왼쪽 역이다.
이 진술은 $U$ 위에서 성립한다.
\end{proof}












\section{뇌터 밑 위의 매끄러움}
\label{spaces-more-morphisms-section-smooth-Noetherian}

\noindent
이 절은 More on Morphisms, Section
\ref{more-morphisms-section-smooth-Noetherian}의 대응물이다.

\begin{lemma}
\label{spaces-more-morphisms-lemma-lifting-along-artinian-at-point}
$S$를 스킴이라 하자. $f : X \to Y$를 $S$ 위 대수공간의 사상이라
하자. $x \in |X|$라 하자. $Y$가 국소 뇌터이고 $f$가 국소 유한형이라고
가정하자. 다음은 동치이다.
\begin{enumerate}
\item $f$는 $x$에서 매끄럽다.
\item 임의의 실선 가환 도식
$$
\xymatrix{
X \ar[d]_f & \Spec(B) \ar[d]^i \ar[l]^-\alpha \\
Y & \Spec(B') \ar[l]_-{\beta} \ar@{-->}[lu]
}
$$
에서 $B' \to B$가 국소환들의 전사이고
$\Ker(B' \to B)$가 제곱 영이며, $\alpha$가 $\Spec(B)$의 닫힌점을
$x$로 보내면, 도식을 가환으로 만드는 점선 화살표가 존재한다.
\item (2)와 같되, $B' \to B$는 작은 확장들을 돈다(Algebra,
Definition \ref{algebra-definition-small-extension} 참조).
\end{enumerate}
\end{lemma}

\begin{proof}
조건 (1)은 열린 부분공간 $X' \subset X$가 존재하여
$X' \to Y$가 매끄럽다는 뜻이다. 따라서 Lemma
\ref{spaces-more-morphisms-lemma-smooth-formally-smooth}에 의해 (1)은 조건 (2)와 (3)을
함의한다. 조건 (2)가 조건 (3)을 함의함은 자명하다. (3)을 가정하자.
가환 도식
$$
\xymatrix{
X \ar[d] & U \ar[l] \ar[d] \\
Y & V \ar[l]
}
$$
을 택하자. 여기서 $U$와 $V$는 아핀이고 수평 화살표들은 에탈이며,
점 $u \in U$가 존재하여 $x$로 간다고 하자. 이제 도식
$$
\xymatrix{
X \ar[d] & U \ar[l] \ar[d] & \Spec(B) \ar[d]^i \ar[l]^-\alpha  \\
Y & V \ar[l] & \Spec(B') \ar[l]_-{\beta}
}
$$
을 생각하자. 이는 $u \in U \to V$에 대해 (3)에서와 같은 도식이다.
가정으로 얻는 화살표를 $\gamma : \Spec(B') \to X$라 하자. 이
화살표는 (3)이 $X$에 대해 성립한다는 데서 얻는다. $U \to X$가 에탈이고 따라서
형식적으로 에탈이므로(Lemma \ref{spaces-more-morphisms-lemma-etale-formally-etale}), 사상
$\gamma$는 $U$로의 유일한 올림을 가지며, 그 올림은 $\alpha$와
양립한다. 이제
$V \to Y$도 에탈이고 따라서 형식적으로 에탈이므로 이 올림은
$\beta$와 양립한다. 따라서 (3)은 $u \in U \to V$에 대해 성립하고,
More on Morphisms, Lemma
\ref{more-morphisms-lemma-lifting-along-artinian-at-point}에 의해
$U \to V$가 $u$에서 매끄럽다고 결론 내린다. 이로써
$X \to Y$가 $x$에서 매끄러움이 증명되어 증명이 끝난다.
\end{proof}

\noindent
때로는 유한형 점들을 “중심으로 하는” 작은 확장들에 대해서만 올림
판정법을 확인하면 된다는 사실이 유용하다(Morphisms of Spaces,
Section \ref{spaces-morphisms-section-points-finite-type} 참조).

\begin{lemma}
\label{spaces-more-morphisms-lemma-lifting-along-artinian}
$S$를 스킴이라 하자. $f : X \to Y$를 $S$ 위 대수공간의 사상이라
하자. $Y$가 국소 뇌터이고 $f$가 국소 유한형이라고 가정하자. 다음은
동치이다.
\begin{enumerate}
\item $f$는 매끄럽다.
\item 임의의 실선 가환 도식
$$
\xymatrix{
X \ar[d]_f & \Spec(B) \ar[d]^i \ar[l]^-\alpha \\
Y & \Spec(B') \ar[l]_-{\beta} \ar@{-->}[lu]
}
$$
에서 $B' \to B$가 아르틴 국소환들의 작은 확장이고 $\beta$가
유한형(!)이면, 도식을 가환으로 만드는 점선 화살표가 존재한다.
\end{enumerate}
\end{lemma}

\begin{proof}
$f$가 매끄러우면 무한소 올림 판정법(Lemma
\ref{spaces-more-morphisms-lemma-smooth-formally-smooth})에 의해 $f$는 형식적으로 매끄럽고
(2)가 성립한다.

\medskip\noindent
$f$가 매끄럽지 않다고 가정하자. 점들 $x \in X$ 가운데 $f$가
매끄럽지 않은 것들의 집합은 닫힌 부분집합 $T$를 이루며, 이는
$|X|$ 안에 있다. Morphisms of
Spaces, Lemma
\ref{spaces-morphisms-lemma-enough-finite-type-points}에 의해
$x \in T \subset X$이고 $x \in X_{\text{ft-pts}}$인 점이 존재한다.
가환 도식
$$
\xymatrix{
X \ar[d] & U \ar[l] \ar[d] & u \ar@{|->}[d] \\
Y & V \ar[l] & v
}
$$
을 택하자. 여기서 $U$와 $V$는 아핀이고 수평 화살표들은 에탈이며,
점 $u \in U$가 존재하여 $x$로 간다. 그러면 $u$는 $U$의 유한형
점이다. $U \to V$가 점 $u$에서 매끄럽지 않으므로 More on Morphisms,
Lemma \ref{more-morphisms-lemma-lifting-along-artinian-at-point}에 의해
도식
$$
\xymatrix{
X \ar[d] & U \ar[l] \ar[d] & \Spec(B) \ar[d]^i \ar[l]^-\alpha  \\
Y & V \ar[l] & \Spec(B') \ar[l]_-{\beta} \ar@{-->}[lu]
}
$$
이 존재한다. 여기서 $B' \to B$는 (아르틴) 국소환들의 작은 확장이고,
$B$의 잉여체는 $\kappa(v)$와 같으며, 점선 화살표는 존재하지 않는다.
$U \to V$가 유한형이므로 $v$는 $V$의 유한형 점이다. Morphisms,
Lemma \ref{morphisms-lemma-artinian-finite-type}에 의해 사상 $\beta$는
유한형이고, 따라서 합성
% P09-ERR-0587: frozen source uses Spec(B) although beta has source Spec(B').
$\Spec(B) \to Y$도 유한형이다. Lemma
\ref{spaces-more-morphisms-lemma-lifting-along-artinian-at-point}의 증명에서와 똑같이
논증하면($U \to X$와 $V \to Y$가 에탈이고 따라서 형식적으로
에탈임을 사용한다), 마지막 표시 도식의 바깥 직사각형에 들어맞는
화살표
% P09-ERR-0588: frozen source names the already-existing Spec(B) arrow, not the absent lift.
$\Spec(B) \to X$가 존재할 수 없음을 알 수 있다. 달리 말해 (2)는
성립하지 않으며, 증명이 끝난다.
\end{proof}

\noindent
다음은 유용한 응용이다.

\begin{lemma}
\label{spaces-more-morphisms-lemma-check-smoothness-on-infinitesimal-nbhds}
$S$를 스킴이라 하자. $f : X \to Y$를 $S$ 위 대수공간의 사상이라
하자. $f$가 국소 유한형이고 $Y$가 국소 뇌터라고 가정하자.
$Z \subset Y$를 닫힌 부분공간이라 하고, 그 $n$차 무한소 근방을
$Z_n \subset Y$라 하자. $X_n = Z_n \times_Y X$로 놓는다.
\begin{enumerate}
\item $X_n \to Z_n$이 모든 $n$에 대해 매끄러우면 $f$는
$f^{-1}(Z)$의 모든 점에서 매끄럽다.
\item $X_n \to Z_n$이 모든 $n$에 대해 에탈이면 $f$는
$f^{-1}(Z)$의 모든 점에서 에탈이다.
\end{enumerate}
\end{lemma}

\begin{proof}
$X_n \to Z_n$이 모든 $n$에 대해 매끄럽다고 가정하자.
$x \in X$를 $Z$의 한 점 위에 놓이는 점이라 하자. Lemma
\ref{spaces-more-morphisms-lemma-lifting-along-artinian-at-point}의 (3)에서와 같은 작은 확장
$B' \to B$ 및 사상 $\alpha$, $\beta$가 주어졌다고 하자. $B'$의
극대 아이디얼은 영멱이다($B'$이 아르틴이므로). 따라서 사상
$\beta$는 $Z_n$을 통해 분해되고 $\alpha$는 $X_n$을 통해 분해된다.
이는 적당한 $n$에 대해 성립한다. 이제 $X_n \to Z_n$의 올림 성질을 적용하여 도식의 원하는
점선 화살표를 얻는다. 이로써 (1)이 증명된다. (2)는 (1)과, 사상이
에탈인 것과 상대차원 $0$인 매끄러운 사상인 것이 동치라는 사실에서
따른다.
\end{proof}
\section{나이브 코탄젠트 복합체}
\label{spaces-more-morphisms-section-netherlander}

\noindent
이 절은 Modules on Sites, Section
\ref{sites-modules-section-netherlander}의 연속이며, 그 절은 다시
Algebra, Section \ref{algebra-section-netherlander}의 논의를 잇는다.

\begin{definition}
\label{spaces-more-morphisms-definition-netherlander}
$S$를 스킴이라 하자. $f : X \to Y$를 $S$ 위 대수공간의 사상이라
하자. {\it $f$의 나이브 코탄젠트 복합체}란 Modules on Sites,
Definition
\ref{sites-modules-definition-cotangent-complex-morphism-ringed-topoi}에서
환 달린 토포스의 사상 $f_{small}$에 대해 정의한 복합체이다. 이
사상은 $X$와 $Y$의 작은 에탈 사이트 사이에 놓인다. Properties of Spaces, Lemma
\ref{spaces-properties-lemma-morphism-ringed-topoi}를 보라. 표기는
$\NL_f$ 또는 $\NL_{X/Y}$이다.
\end{definition}

\noindent
다음 보조정리들은 이 정의가 환 준동형 및 스킴에 대한 정의와
양립하며, $\NL_{X/Y}$가 $D_\QCoh(\mathcal{O}_X)$의 대상임을 보인다.

\begin{lemma}
\label{spaces-more-morphisms-lemma-NL-etale-localization}
$S$를 스킴이라 하자. 가환 도식
$$
\xymatrix{
U \ar[d]_p \ar[r]_g & V \ar[d]^q \\
X \ar[r]^f & Y
}
$$
을 생각하자. 이는 $S$ 위 대수공간의 도식이며 $p$와 $q$는
에탈이라고 하자. 그러면 표준적 동일시
$\NL_{X/Y}|_{U_\etale} = \NL_{U/V}$가 존재하며, 이는
$D(\mathcal{O}_U)$에서의 동일시이다.
\end{lemma}

\begin{proof}
나이브 코탄젠트 복합체의 구성은 당김과 가환하며(Modules on Sites,
Lemma \ref{sites-modules-lemma-pullback-NL}), 다음 등식들이 성립한다.
$p_{small}^{-1}\mathcal{O}_X = \mathcal{O}_U$이고
$g_{small}^{-1}\mathcal{O}_{V_\etale} =
p_{small}^{-1}f_{small}^{-1}\mathcal{O}_{Y_\etale}$이다. 실제로
$q_{small}^{-1}\mathcal{O}_{Y_\etale} =
\mathcal{O}_{V_\etale}$이다. Properties of Spaces, Lemma
\ref{spaces-properties-lemma-etale-exact-pullback}을 보라. 정의를
추적하면 $\NL_{X/Y}|_{U_\etale} = \NL_{U/V}$를 얻는다.
\end{proof}

\begin{lemma}
\label{spaces-more-morphisms-lemma-NL-compare-spaces-schemes}
$S$를 스킴이라 하자. $f : X \to Y$를 $S$ 위 대수공간의 사상이라
하자.
% P09-ERR-0589: frozen source omits "are" before "representable".
$X$와 $Y$가 스킴 $X_0$와 $Y_0$로 표현된다고 가정하자. 그러면
표준적 동일시 $\NL_{X/Y} = \epsilon^*\NL_{X_0/Y_0}$가 존재하며,
이는 $D(\mathcal{O}_X)$에서의 동일시이다. 여기서
$\epsilon$은 Derived Categories of Spaces, Section
\ref{spaces-perfect-section-derived-quasi-coherent-etale}에서와 같고,
$\NL_{X_0/Y_0}$는 More on Morphisms, Definition
\ref{more-morphisms-definition-netherlander}에서와 같다.
\end{lemma}

\begin{proof}
$f_0 : X_0 \to Y_0$를 $f$에 대응하는 스킴의 사상이라 하자. 표준
사상
$\epsilon^{-1}f_0^{-1}\mathcal{O}_{Y_0} \to f_{small}^{-1}\mathcal{O}_Y$
이 존재하고, 이는
$\epsilon^\sharp : \epsilon^{-1}\mathcal{O}_{X_0} \to \mathcal{O}_X$와
양립한다. 그 이유는 다음 가환 도식이 있기 때문이다.
$$
\xymatrix{
X_{0, Zar} \ar[d]_{f_0} & X_\etale \ar[l]^\epsilon \ar[d]^f \\
Y_{0, Zar} & Y_\etale \ar[l]_\epsilon
}
$$
Derived Categories of Spaces, Remark
\ref{spaces-perfect-remark-match-total-direct-images}를 보라. 따라서
나이브 코탄젠트 복합체의 함자성으로부터 표준 사상
$$
\epsilon^{-1}\NL_{X_0/Y_0} =
\epsilon^{-1}\NL_{\mathcal{O}_{X_0}/f_0^{-1}\mathcal{O}_{Y_0}} =
\NL_{\epsilon^{-1}\mathcal{O}_{X_0}/\epsilon^{-1}f_0^{-1}\mathcal{O}_{Y_0}}
\to
\NL_{\mathcal{O}_X/f^{-1}_{small}\mathcal{O}_Y} = \NL_{X/Y}
$$
을 얻는다. 유도된 사상
$\epsilon^*\NL_{X_0/Y_0} \to \NL_{X/Y}$가
$D(\mathcal{O}_X)$에서 동형임을 보이기 위해 기하점의 줄기에서
확인해도 충분하다(Properties of Spaces, Theorem
\ref{spaces-properties-theorem-exactness-stalks}).
$\overline{x} : \Spec(k) \to X_0$를 $x \in X_0$ 위에 놓이는
기하점이라 하자. 또한 $\overline{y} = f \circ \overline{x}$가
$y \in Y_0$ 위에 놓인다고 하자. 그러면
$$
\NL_{X/Y, \overline{x}} =
\NL_{\mathcal{O}_{X, \overline{x}}/\mathcal{O}_{Y, \overline{y}}}
$$
이다. 이는 $\overline{x}$에서 줄기를 취하는 것이
$\overline{x} : \Spec(k) \to X$를 통한 역상을 취하는 것과 같고,
Modules on Sites, Lemma \ref{sites-modules-lemma-pullback-NL}을 적용할 수
있기 때문이다. 한편
$$
(\epsilon^*\NL_{X_0/Y_0})_{\overline{x}} =
\NL_{X_0/Y_0, x} \otimes_{\mathcal{O}_{X_0, x}}
\mathcal{O}_{X, \overline{x}} =
\NL_{\mathcal{O}_{X_0, x}/\mathcal{O}_{Y_0, y}}
\otimes_{\mathcal{O}_{X_0, x}} \mathcal{O}_{X, \overline{x}}
$$
이다. 일부 세부사항은 생략한다(힌트: 당김의 줄기는 상점에서의
줄기라는 사실, 즉 Sites, Lemma
\ref{sites-lemma-point-morphism-sites}와 가군에 대한 대응 결과인 Modules
on Sites, Lemma \ref{sites-modules-lemma-pullback-stalk}을 사용하라).
$\mathcal{O}_{X, \overline{x}}$는 $\mathcal{O}_{X_0, x}$의 엄밀
헨젤화이고, $\mathcal{O}_{Y, \overline{y}}$도 마찬가지이다(Properties
of Spaces, Lemma
\ref{spaces-properties-lemma-describe-etale-local-ring}). 따라서 결과는
More on Algebra, Lemma \ref{more-algebra-lemma-henselization-NL}에서
따른다.
\end{proof}

\begin{lemma}
\label{spaces-more-morphisms-lemma-netherlander-quasi-coherent}
$S$를 스킴이라 하자. $f : X \to Y$를 $S$ 위 대수공간의 사상이라
하자. 복합체 $\NL_{X/Y}$의 코호몰로지 층들은 준연접이고, 차수
$-1$, $0$ 바깥에서는 영이며, $\Omega_{X/Y}$와 같다. 마지막 등식은
차수 $0$에서 성립한다.
\end{lemma}

\begin{proof}
Modules on Sites, Section \ref{sites-modules-section-netherlander}의 나이브
코탄젠트 복합체 구성에 의해 $\NL_{X/Y}$는 차수 $-1$, $0$에 놓이는
복합체이고, 차수 $0$의 코호몰로지는 $\Omega_{X/Y}$이다(Section
\ref{spaces-more-morphisms-section-sheaf-differentials}의 $\Omega_{X/Y}$ 구성에 의한다).
미분 층은 준연접이다(Lemma
\ref{spaces-more-morphisms-lemma-module-differentials-quasi-coherent} 참조). 증명을 끝내려면
$H^{-1}(\NL_{X/Y})$가 준연접임을 보이면 충분하다. 이는 에탈
국소적으로 확인하고(Lemma \ref{spaces-more-morphisms-lemma-NL-etale-localization} 및
Properties of Spaces, Lemma
\ref{spaces-properties-lemma-characterize-quasi-coherent}에 의해 허용됨),
스킴의 경우로 환원한 뒤(Lemma
\ref{spaces-more-morphisms-lemma-NL-compare-spaces-schemes}), 마지막으로 스킴의 경우에 대한
결과(More on Morphisms, Lemma
\ref{more-morphisms-lemma-netherlander-quasi-coherent})를 적용하여 얻는다.
\end{proof}

\begin{lemma}
\label{spaces-more-morphisms-lemma-netherlander-fp}
$S$를 스킴이라 하자. $f : X \to Y$를 $S$ 위 대수공간의 사상이라
하자. $f$가 국소 유한 표시이면 $\NL_{X/Y}$는 $X$ 위 에탈 국소적으로
준동형인 복합체
$$
\ldots \to 0 \to \mathcal{F}^{-1} \to \mathcal{F}^0 \to 0 \to \ldots
$$
를 갖는다. 이 복합체는 준연접 $\mathcal{O}_X$-가군들로 이루어지고,
$\mathcal{F}^0$는 유한 표시이며 $\mathcal{F}^{-1}$은 유한형이다.
\end{lemma}

\begin{proof}
나이브 코탄젠트 복합체의 구성은 Lemma
\ref{spaces-more-morphisms-lemma-NL-etale-localization}에 의해 에탈 국소화와 가환한다.
Lemma \ref{spaces-more-morphisms-lemma-NL-compare-spaces-schemes}에 의해 스킴의 경우로
환원된다. 스킴의 경우에 대한 결과는 More on Morphisms, Lemma
\ref{more-morphisms-lemma-netherlander-fp}이다.
\end{proof}

\begin{lemma}
\label{spaces-more-morphisms-lemma-NL-formally-smooth}
$S$를 스킴이라 하자. $f : X \to Y$를 $S$ 위 대수공간의 사상이라
하자. 다음은 동치이다.
\begin{enumerate}
\item $f$는 형식적으로 매끄럽다.
\item $H^{-1}(\NL_{X/Y}) = 0$이고
$H^0(\NL_{X/Y}) = \Omega_{X/Y}$는 국소 사영이다.
\end{enumerate}
\end{lemma}

\begin{proof}
이는 Lemma \ref{spaces-more-morphisms-lemma-formally-smooth}, Lemma
\ref{spaces-more-morphisms-lemma-NL-etale-localization}, Lemma
\ref{spaces-more-morphisms-lemma-NL-compare-spaces-schemes}, 그리고 스킴의 경우인 More on
Morphisms, Lemma \ref{more-morphisms-lemma-NL-formally-smooth}에서 따른다.
\end{proof}

\begin{lemma}
\label{spaces-more-morphisms-lemma-NL-formally-etale}
% P09-ERR-0590: frozen statement says schemes although the proof is space-level.
$f : X \to Y$를 스킴의 사상이라 하자. 다음은 동치이다.
\begin{enumerate}
\item $f$는 형식적으로 에탈이다.
\item $H^{-1}(\NL_{X/Y}) = H^0(\NL_{X/Y}) = 0$이다.
\end{enumerate}
\end{lemma}

\begin{proof}
(1)을 가정하자. 형식적으로 에탈인 사상은 형식적으로 매끄러운
사상이다. 따라서 Lemma \ref{spaces-more-morphisms-lemma-NL-formally-smooth}에 의해
$H^{-1}(\NL_{X/Y}) = 0$이다. 한편
% P09-ERR-0591: frozen source has "if" in place of the copula "is".
형식적으로 에탈인 사상은 형식적으로 비분기이므로 Lemma
\ref{spaces-more-morphisms-lemma-characterize-formally-unramified}에 의해
$\Omega_{X/Y} = 0$이다. 역으로 (2)가 성립하면 Lemma
\ref{spaces-more-morphisms-lemma-NL-formally-smooth}에 의해 $f$는 형식적으로 매끄럽고,
Lemma \ref{spaces-more-morphisms-lemma-characterize-formally-unramified}에 의해 형식적으로
비분기이며, 따라서
% P09-ERR-0592: frozen source uses plural "Lemmas" for one citation.
Lemma \ref{spaces-more-morphisms-lemma-formally-etale-unramified-smooth}에 의해 형식적으로
에탈이다.
\end{proof}

\begin{lemma}
\label{spaces-more-morphisms-lemma-NL-smooth}
% P09-ERR-0593: frozen statement says schemes although its proof treats spaces.
$f : X \to Y$를 스킴의 사상이라 하자. 다음은 동치이다.
\begin{enumerate}
\item $f$는 매끄럽다.
\item $f$는 국소 유한 표시이고,
$H^{-1}(\NL_{X/Y}) = 0$이며, $H^0(\NL_{X/Y}) = \Omega_{X/Y}$는
유한 국소 자유이다.
\end{enumerate}
\end{lemma}

\begin{proof}
이는 Lemma \ref{spaces-more-morphisms-lemma-formally-smooth}, Lemma
\ref{spaces-more-morphisms-lemma-NL-etale-localization}, Lemma
\ref{spaces-more-morphisms-lemma-NL-compare-spaces-schemes}, 그리고 스킴의 경우인 More on
Morphisms, Lemma \ref{more-morphisms-lemma-NL-smooth}에서 따른다.
\end{proof}









\section{평탄 자취의 열림성}
\label{spaces-more-morphisms-section-open-flat}

\noindent
% P09-ERR-0594: frozen source omits the article in "is analogue of".
이 절은 More on Morphisms, Section
\ref{more-morphisms-section-open-flat}의 대응물이다. 대수공간 위
준연접 가군의 평탄성은 Morphisms of Spaces, Section
\ref{spaces-morphisms-section-flat-modules}에서 정의했음을 기억하라.

\begin{theorem}
\label{spaces-more-morphisms-theorem-openness-flatness}
$S$를 스킴이라 하자. $f : X \to Y$를 $S$ 위 대수공간의 사상이라
하자. $\mathcal{F}$를 $X$ 위 준연접 층이라 하자. $f$가 국소 유한
표시이고 $\mathcal{F}$가 국소 유한 표시인
$\mathcal{O}_X$-가군이라고 가정하자. 그러면
$$
\{x \in |X| : \mathcal{F}\text{ is flat over }Y\text{ at }x\}
$$
는 $|X|$에서 열린집합이다.
\end{theorem}

\begin{proof}
가환 도식
$$
\xymatrix{
U \ar[d]_p \ar[r]_\alpha &
V \ar[d]^q \\
X \ar[r]^a & Y
}
$$
을 택하자. Spaces, Lemma
\ref{spaces-lemma-lift-morphism-presentations}에서와 같이 $U$, $V$는
스킴이고 $p$, $q$는 전사 에탈이다. More on Morphisms, Theorem
\ref{more-morphisms-theorem-openness-flatness}에 의해 집합
$U' = \{u \in |U| : p^*\mathcal{F}\text{ is flat over }V\text{ at }u\}$
는 $U$에서 열린집합이다. Morphisms of Spaces, Definition
\ref{spaces-morphisms-definition-flat-module}에 의해 $U'$의 상은
$|X|$ 안에서 정리에 나타난 집합이다. 따라서 끝났다. 실제로 사상
$|U| \to |X|$는 열린 사상이다. Properties of Spaces, Lemma
\ref{spaces-properties-lemma-topology-points}를 보라.
\end{proof}

\begin{lemma}
\label{spaces-more-morphisms-lemma-flat-locus-base-change}
$S$를 스킴이라 하자. 다음 도식을 생각하자.
$$
\xymatrix{
X' \ar[r]_{g'} \ar[d]_{f'} & X \ar[d]^f \\
Y' \ar[r]^g & Y
}
$$
이는 $S$ 위 대수공간의 카르테시안 도식이다.
$\mathcal{F}$를 준연접 $\mathcal{O}_X$-가군이라 하자. $g$가
평탄하고 $f$가 국소 유한 표시이며 $\mathcal{F}$가 국소 유한
표시라고 가정하자. 그러면
$$
\{x' \in |X'| : (g')^*\mathcal{F}\text{ is flat over }Y'\text{ at }x'\}
$$
는 Theorem \ref{spaces-more-morphisms-theorem-openness-flatness}의 열린 부분집합을 연속 사상
$|g'| : |X'| \to |X|$로 역상 취한 것이다.
\end{lemma}

\begin{proof}
이는 Morphisms of Spaces, Lemma
\ref{spaces-morphisms-lemma-base-change-module-flat}에서 따른다.
\end{proof}












\section{섬유에 의한 평탄성 판정법}
\label{spaces-more-morphisms-section-criterion-flat-fibres}

\noindent
$S$를 스킴이라 하자. $S$ 위 대수공간의 가환 도식
$$
\xymatrix{
X \ar[rr]_f \ar[dr]_g & & Y \ar[dl]^h \\
& Z
}
$$
과 준연접 $\mathcal{O}_X$-가군 $\mathcal{F}$를 생각하자. 점
$x \in |X|$가 주어졌을 때 $\mathcal{F}$가 $Y$ 위에서 $x$에서
평탄한가라는 문제를 생각한다. $\mathcal{F}$가 $Z$ 위에서 $x$에서
평탄하면, 아래 정리는 이 문제가 $\mathcal{F}$를 $X \to Z$의
$g(x)$ 위 섬유로 제한한 것이 $Y \to Z$의 $g(x)$ 위 섬유 위에서
평탄한가라는 문제와 밀접한 관계가 있다고 말한다. 이를 의미 있게
만들기 위해 다음 예비 보조정리를 제시한다.

\begin{lemma}
\label{spaces-more-morphisms-lemma-flat-on-fibres-at-point}
위 상황에서 다음은 동치이다.
\begin{enumerate}
\item 기하점 $\overline{x}$를 택하되, 이는 $X$의 점이고 $x$ 위에
놓인다고 하자.
$\overline{y} = f \circ \overline{x}$ 및
$\overline{z} = g \circ \overline{x}$로 놓는다. 그러면 가군
$\mathcal{F}_{\overline{x}}/
\mathfrak m_{\overline{z}}\mathcal{F}_{\overline{x}}$는
$\mathcal{O}_{Y, \overline{y}}/
\mathfrak m_{\overline{z}}\mathcal{O}_{Y, \overline{y}}$ 위에서
평탄하다.
\item 사상 $x : \Spec(K) \to X$를 택하되, 이는 $x$의 동치류에
속한다고 하자.
$z = g \circ x$, $X_z = \Spec(K) \times_{z, Z} X$,
$Y_z = \Spec(K) \times_{z, Z} Y$로 놓고, $\mathcal{F}_z$를
$\mathcal{F}$의 $X_z$로의 당김이라 하자. 그러면 $\mathcal{F}_z$는
$x$에서 평탄하며, 이는 $Y_z$ 위에서의 평탄성이다(Morphisms of Spaces, Definition
\ref{spaces-morphisms-definition-flat-module}에서 정의된 의미에서).
\item 가환 도식
% P09-ERR-0595: frozen Xy-pic arrow label contains seven stray greater-than signs.
$$
\xymatrix{
& & & U \ar[llld]_a \ar[rr] \ar[dr] & & V \ar[llld]_>>>>>>>b \ar[dl] \\
X \ar[rr]_f \ar[dr]_g & & Y \ar[dl]^h &  & W \ar[llld]_c \\
& Z
}
$$
을 택하자. 여기서 $U, V, W$는 스킴이고 $a, b, c$는 에탈이며,
점 $u \in U$를 택한다. 이 점은 $x$로 간다. $w \in W$를 $u$의 상이라 하자.
$\mathcal{F}_w$를 $\mathcal{F}$의 당김이라 하되, 당김 대상은 섬유
$U_w$이고 이는 $U \to W$의 $w$에서의 섬유라 하자. 그러면 $\mathcal{F}_w$는
$V_w$ 위에서 $u$에서 평탄하다.
\end{enumerate}
\end{lemma}

\begin{proof}
(2)에서 사상 $x : \Spec(K) \to X$는 $K$-유리점을 정의하며, 이 점은
$X_z$에 놓이므로 진술은 의미가 있다. 더 나아가 (2)의 조건은
$\Spec(K) \to X$를 어떻게 택하는가에 의존하지 않는데, 그 선택은
$x$의 동치류에서 이루어진다
(세부사항은 생략한다. 다른 조건들이 이 선택에 의존하지 않으므로
아래 논증에서도 이 사실이 따른다). 또한 (3)과 같은 도식은 항상
다음처럼 택할 수 있다. 먼저 스킴 $W$와 전사 에탈 사상 $W \to Z$를
택하고, 다음으로 스킴 $V$와 전사 에탈 사상
$V \to W \times_Z Y$를 택하며, 마지막으로 스킴 $U$와 전사 에탈 사상
$U \to V \times_Y X$를 택한다. 이 선택들 아래 $U \to W$를 합성
$U \to V \to W$와 같게 놓는다. 또한 점 $u \in U$를 택할 수 있으며,
이 점은 $x$로 간다. 그 이유는 사상 $U \to X$가 전사이기 때문이다.

\medskip\noindent
(3)과 같은 도식과 (1)에서와 같은 기하점
$\overline{x} : \Spec(k) \to X$가 모두 주어졌다고 하자. Properties of
Spaces, Lemma
\ref{spaces-properties-lemma-geometric-lift-to-usual}에 의해 기하점
$\overline{u} : \Spec(k) \to U$를 택하되, 이는 $u$ 위에 놓이고
$\overline{x} = a \circ \overline{u}$가 되게 할 수 있다. 유도된
$\overline{v} : \Spec(k) \to V$ 및
$\overline{w} : \Spec(k) \to W$라 하자. 이들은 각각 $V$와 $W$의
유도된 기하점이다. 이 상황에서
$\mathcal{O}_{X, \overline{x}} = \mathcal{O}_{U, u}^{sh}$이고 $Y$와
$Z$에 대해서도 마찬가지이다. Properties of Spaces, Lemma
\ref{spaces-properties-lemma-describe-etale-local-ring}을 보라. 마찬가지로
$$
\mathcal{F}_{\overline{x}} =
(a^*\mathcal{F})_u \otimes_{\mathcal{O}_{U, u}}
\mathcal{O}_{U, u}^{sh}
$$
이다. Properties of Spaces, Lemma
\ref{spaces-properties-lemma-stalk-quasi-coherent}를 보라. $\mathcal{F}_w$의
$u$에서의 줄기는
$$
(\mathcal{F}_w)_u = (a^*\mathcal{F})_u/\mathfrak m_w(a^*\mathcal{F})_u
$$
이고, $V_w$의 $v$에서의 국소환은
$$
\mathcal{O}_{V_w, v} = \mathcal{O}_{V, v}/\mathfrak m_w\mathcal{O}_{V, v}.
$$
또한
$\mathfrak m_{\overline{z}} =
\mathfrak m_w \mathcal{O}_{Z, \overline{z}} =
\mathfrak m_w \mathcal{O}_{W, w}^{sh}$이므로
\begin{align*}
\mathcal{F}_{\overline{x}}/
\mathfrak m_{\overline{z}}\mathcal{F}_{\overline{x}} & =
(a^*\mathcal{F})_u \otimes_{\mathcal{O}_{U, u}}
\mathcal{O}_{X, \overline{x}}/
\mathfrak m_{\overline{z}}\mathcal{O}_{X, \overline{x}} \\
& =
(\mathcal{F}_w)_u \otimes_{\mathcal{O}_{U_w, u}}
\mathcal{O}_{U, u}^{sh}/\mathfrak m_w\mathcal{O}_{U, u}^{sh} \\
& = (\mathcal{F}_w)_u \otimes_{\mathcal{O}_{U_w, u}}
\mathcal{O}_{U_w, \overline{u}}^{sh} \\
& = (\mathcal{F}_w)_{\overline{u}}
\end{align*}
를 얻는다. 끝에서 두 번째 등식은 Algebra, Lemma
\ref{algebra-lemma-quotient-strict-henselization}에 의하고, 마지막 등식은
Properties of Spaces, Lemma
\ref{spaces-properties-lemma-stalk-quasi-coherent}에 의한다. $V$와 $Y$의
구조층에 같은 논증을 적용하면
$$
\mathcal{O}_{V_w, \overline{v}}^{sh} =
\mathcal{O}_{V, v}^{sh}/\mathfrak m_w \mathcal{O}_{V, v}^{sh} =
\mathcal{O}_{Y, \overline{y}}/
\mathfrak m_{\overline{z}}\mathcal{O}_{Y, \overline{y}}.
$$
이제 Morphisms of Spaces, Lemma
\ref{spaces-morphisms-lemma-flat-at-point}를 사용하면 (1)과 (3)이
동치임을 알 수 있다.

\medskip\noindent
마지막으로 (2)와 (3)의 동치를 증명한다. 체 $\tilde K$를 택하되,
이는 $K$의 확장이라 하자. 사상 $\tilde x : \Spec(\tilde K) \to U$를
택하되, 이는 $u$ 위에 놓인다고 하자
(이는 $u \times_{X, x} \Spec(K)$가 공집합이 아닌 스킴이므로
가능하다).
% P09-ERR-0596: frozen source says "Set ... be" rather than "Let ... be".
합성 $\tilde z : \Spec(\tilde K) \to U \to W$를 놓자. 다음 가환
도식을 얻는다.
% P09-ERR-0595: second frozen Xy-pic arrow label contains the same stray signs.
$$
\xymatrix{
& & & U_w \times_w \tilde z \ar[llld]_a \ar[rr] \ar[dr] & &
V_w \times_w \tilde z \ar[llld]_>>>>>>>b \ar[dl] \\
X_z \ar[rr]_f \ar[dr]_g & & Y_z \ar[dl]^h &  & \tilde z \ar[llld]_c \\
& z
}
$$
여기서 $z = \Spec(K)$이고 $w = \Spec(\kappa(w))$이다. 이제
$\mathcal{F}_w$와 $\mathcal{F}_z$를 $U_w \times_w \tilde z$로 당기면
같은 가군이 됨은 명백하다. 이로부터 가환 도식
$$
\xymatrix{
X_z \ar[d] & U_w \times_w \tilde z \ar[l] \ar[d] \ar[r] & U_w \ar[d] \\
Y_z & V_w \times_w \tilde z \ar[l] \ar[r] & V_w
}
$$
을 얻는다. 두 정사각형은 모두 카르테시안이고 아래쪽 수평 화살표들은
평탄하다. 아래 왼쪽 수평 화살표는 사상
$Y \times_Z \tilde z \to Y \times_Z z = Y_z$(평탄 사상의 밑변환),
에탈 사상 $V \times_Z \tilde z \to Y \times_Z \tilde z$, 그리고
에탈 사상 $V \times_W \tilde z \to V \times_Z \tilde z$의 합성이다.
따라서 Morphisms of Spaces, Lemma
\ref{spaces-morphisms-lemma-base-change-module-flat}에 의해
$$
\mathcal{F}_z\text{ flat at }x\text{ over }Y_z
\Leftrightarrow
\mathcal{F}|_{U_w \times_w \tilde z}
\text{ flat at }\tilde x\text{ over }V_w \times_w \tilde z
\Leftrightarrow
\mathcal{F}_w\text{ flat at }u\text{ over }V_w
$$
이며, 이것으로 끝난다.
\end{proof}

\begin{definition}
\label{spaces-more-morphisms-definition-module-flat-on-fibre}
$S$를 스킴이라 하자. $X \to Y \to Z$를 $S$ 위 대수공간의
사상들이라 하자. $\mathcal{F}$를 준연접 $\mathcal{O}_X$-가군이라
하자. $x \in |X|$를 점이라 하고 그 상을 $z \in |Z|$로 나타내자.
\begin{enumerate}
\item 우리는 {\it $\mathcal{F}$의 $z$ 위 섬유로의 제한이 $x$에서
평탄하며, 그 밑은 $Y$의 $z$ 위 섬유이다}라는 것은 Lemma
\ref{spaces-more-morphisms-lemma-flat-on-fibres-at-point}의 동치 조건들이 성립한다는 뜻이다.
\item 우리는 {\it $X$의 $z$ 위 섬유가 $x$에서 평탄하며, 그 밑은 $Y$의
$z$ 위 섬유이다}라는 것은 Lemma
\ref{spaces-more-morphisms-lemma-flat-on-fibres-at-point}의 동치 조건들이
$\mathcal{F} = \mathcal{O}_X$에 대해 성립한다는 뜻이다.
\item 우리는 {\it $X$의 $z$ 위 섬유가 $Y$의 $z$ 위 섬유 위에서
평탄하다}라는 것은, 모든 $x \in |X|$에 대해 이 점이 $z$ 위에 놓이면
$X$의 $z$ 위 섬유가 $x$에서 평탄하며 그 밑이 $Y$의 $z$ 위 섬유라는
뜻이다
% P09-ERR-0597: frozen third definition item lacks terminal punctuation.
\end{enumerate}
\end{definition}

\noindent
이 정의를 사용하면 판정법의 한 판을 다음과 같이 진술할 수 있다.
뇌터 판은 Section \ref{spaces-more-morphisms-section-flat-over-Noetherian}에 있다.

\begin{theorem}
\label{spaces-more-morphisms-theorem-criterion-flatness-fibre}
$S$를 스킴이라 하자. $f : X \to Y$와 $Y \to Z$를 $S$ 위
대수공간의 사상들이라 하자. $\mathcal{F}$를 준연접
$\mathcal{O}_X$-가군이라 하자. 다음을 가정한다.
\begin{enumerate}
\item $X$는 $Z$ 위 국소 유한 표시이다.
% P09-ERR-0598: frozen assumption omits the copula after F.
\item $\mathcal{F}$는 유한 표시인 $\mathcal{O}_X$-가군이다.
\item $Y$는 $Z$ 위 국소 유한형이다.
\end{enumerate}
$x \in |X|$라 하고, $y \in |Y|$와 $z \in |Z|$를 $x$의 상들이라
하자.
% P09-ERR-0599: frozen theorem has not introduced the barred geometric point.
$\mathcal{F}_{\overline{x}} \not = 0$이면 다음은 동치이다.
\begin{enumerate}
\item $\mathcal{F}$는 $Z$ 위에서 $x$에서 평탄하고, $\mathcal{F}$를
그 $z$ 위 섬유에 제한하면 $x$에서 평탄하며 그 밑은 $Y$의 $z$ 위
섬유이다.
\item $Y$는 $Z$ 위에서 $y$에서 평탄하고 $\mathcal{F}$는 $Y$
위에서 $x$에서 평탄하다.
\end{enumerate}
더 나아가 (1)과 (2)가 성립하는 점 $x$들의 집합은
$\text{Supp}(\mathcal{F})$에서 열린집합이다.
\end{theorem}

\begin{proof}
Lemma \ref{spaces-more-morphisms-lemma-flat-on-fibres-at-point}의 (3)에서와 같은 도식을
택하자. 정의들로부터 스킴의 사상들 $U \to V \to W$, 준연접 층
$a^*\mathcal{F}$, 그리고 점 $u \in U$에 대한 대응 정리로 환원된다.
따라서 정리는 스킴에 대한 대응 결과인 More on Morphisms, Theorem
\ref{more-morphisms-theorem-criterion-flatness-fibre}에서 따른다.
\end{proof}

\begin{lemma}
\label{spaces-more-morphisms-lemma-morphism-between-flat}
$S$를 스킴이라 하자.
% P09-ERR-0600: frozen opening uses a singular noun for two morphisms.
$f : X \to Y$와 $Y \to Z$를 $S$ 위 대수공간의 사상들이라 하자.
다음을 가정한다.
\begin{enumerate}
\item $X$는 $Z$ 위 국소 유한 표시이다.
\item $X$는 $Z$ 위에서 평탄하다.
\item 모든 $z \in |Z|$에 대해 $X$의 $z$ 위 섬유는
$Y$의 $z$ 위 섬유 위에서 평탄하다.
\item $Y$는 $Z$ 위 국소 유한형이다.
\end{enumerate}
그러면 $f$는 평탄하다. $f$가 전사이기도 하면 $Y$는 $Z$ 위에서
평탄하다.
\end{lemma}

\begin{proof}
이는 Theorem \ref{spaces-more-morphisms-theorem-criterion-flatness-fibre}의 특수한 경우이다.
\end{proof}

\begin{lemma}
\label{spaces-more-morphisms-lemma-base-change-criterion-flatness-fibre}
$S$를 스킴이라 하자. $f : X \to Y$와 $Y \to Z$를 $S$ 위
대수공간의 사상들이라 하자. $\mathcal{F}$를 준연접
$\mathcal{O}_X$-가군이라 하자.
다음을 가정한다.
\begin{enumerate}
\item $X$는 $Z$ 위 국소 유한 표시이다.
% P09-ERR-1119: frozen assumption omits the copula after F.
\item $\mathcal{F}$는 유한 표시인 $\mathcal{O}_X$-가군이다.
\item $\mathcal{F}$는 $Z$ 위에서 평탄하다.
\item $Y$는 $Z$ 위 국소 유한형이다.
\end{enumerate}
그러면 집합
$$
A = \{x \in |X| : \mathcal{F} \text{ flat at }x \text{ over }Y\}.
$$
% P09-ERR-0601: frozen display ends the sentence before its predicate.
는 $|X|$에서 열린집합이고 그 구성은 임의의 밑변환과 가환한다.
$Z' \to Z$가 대수공간의 사상이고, $A'$가
$X' = X \times_Z Z'$에서 $\mathcal{F}' = \mathcal{F} \times_Z Z'$가
$Y' = Y \times_Z Z'$ 위에서 평탄한 점들의 집합이면, $A'$는 $A$의
역상이고 이를 주는 연속사상은 $|X'| \to |X|$이다.
% P09-ERR-0602: frozen source writes an ill-typed fibre product for F'.
\end{lemma}

\begin{proof}
이를 증명하는 한 방법은 More on Morphisms, Lemma
\ref{more-morphisms-lemma-morphism-between-flat}의 증명을 대수공간의
범주로 옮기는 것이다. 대신 스킴의 경우로 환원하여 증명하겠다.
즉 Lemma \ref{spaces-more-morphisms-lemma-flat-on-fibres-at-point}의 (3)과 같은 도식에서
$a$, $b$, $c$가 전사가 되도록 택하자. 정의들로부터 스킴의 사상
$U \to V \to W$, 준연접 층 $a^*\mathcal{F}$, 그리고 점
$u \in U$에 대한 대응 정리로 환원된다. 유일하게 덧붙일 작은 점은
다음과 같다. 대수공간의 사상 $Z' \to Z$가 주어지면 스킴 $W'$와
전사 에탈 사상 $W' \to W \times_Z Z'$를 택한다. 그리고
$U' = W' \times_W U$, $V' = W' \times_W V$로 둔다.
$a', b', c'$를 각각 $U', V', W'$에서 $X', Y', Z'$로 가는
사상으로 쓴다. 이 경우 $A$, 각각 $A'$는
$U$, 각각 $U'$의 열린 부분집합들의 상이며, 이 열린 부분집합들은
$a^*\mathcal{F}$, 각각 $(a')^*\mathcal{F}'$에 대응한다.
이는 실제로 보조정리를 More on Morphisms, Lemma
\ref{more-morphisms-lemma-morphism-between-flat}로 환원한다.
\end{proof}

\begin{lemma}
\label{spaces-more-morphisms-lemma-base-change-flatness-fibres}
$S$를 스킴이라 하자.
% P09-ERR-0600: frozen opening again uses a singular noun for two morphisms.
$f : X \to Y$와 $Y \to Z$를 $S$ 위 대수공간의 사상들이라 하자.
다음을 가정한다.
\begin{enumerate}
\item $X$는 $Z$ 위 국소 유한 표시이다.
\item $X$는 $Z$ 위에서 평탄하다.
\item $Y$는 $Z$ 위 국소 유한형이다.
\end{enumerate}
그러면 집합
$$
\{x \in |X| : X\text{ flat at }x \text{ over }Y\}.
$$
% P09-ERR-0601: frozen display ends the sentence before its predicate.
는 $|X|$에서 열린집합이고 그 구성은 임의의 밑변환
$Z' \to Z$와 가환한다.
\end{lemma}

\begin{proof}
이는 Lemma \ref{spaces-more-morphisms-lemma-base-change-criterion-flatness-fibre}의
특수한 경우이다.
\end{proof}

\begin{lemma}
\label{spaces-more-morphisms-lemma-flat-and-free-at-point-fibre}
$S$를 스킴이라 하자. $f : X \to Y$를 $S$ 위 대수공간의 사상이라
하고 국소 유한 표시라고 하자. $\mathcal{F}$를 유한 표시인
$\mathcal{O}_X$-가군이라 하자. $x \in |X|$의 상을 $y \in |Y|$라
하자. $\mathcal{F}$가 $x$에서 $Y$ 위로 평탄하면 다음은 동치이다.
% P09-ERR-0603: frozen list-introducing clause lacks a colon.
\begin{enumerate}
\item $(\mathcal{F}_{\overline{y}})_{\overline{x}}$는 평탄한
$\mathcal{O}_{X_{\overline{y}}, \overline{x}}$-가군이다.
\item $(\mathcal{F}_{\overline{y}})_{\overline{x}}$는 자유
$\mathcal{O}_{X_{\overline{y}}, \overline{x}}$-가군이다.
\item $\mathcal{F}_{\overline{y}}$는 $\overline{x}$의 어떤 에탈
근방 위에서 유한 자유이며, 이 근방은 $X_{\overline{y}}$ 안에 있다.
\item $\mathcal{F}$는 $x$의 어떤 에탈 근방 위에서 유한 자유이며,
이 근방은 $X$ 안에 있다.
\end{enumerate}
여기서 $\overline{x}$는 $X$의 기하학적 점으로 $x$ 위에 놓이고
$\overline{y} = f \circ \overline{x}$이다.
\end{lemma}

\begin{proof}
가환 도식
$$
\xymatrix{
U \ar[d] \ar[r] & V \ar[d] \\
X \ar[r] & Y
}
$$
에서 $U$와 $V$가 스킴이고 수직 화살표들이 에탈이며, 점
$u \in U$가 존재하여 그 상이 $x$가 되도록 택하자.
$v \in V$를 $u$의 상이라 하자.
$\text{id} : X \to X$를 $Y$ 위에서 보아 Lemma
\ref{spaces-more-morphisms-lemma-flat-on-fibres-at-point}를 적용하면 (1)은
“$\mathcal{F}|_{U_v}$가 $U_v$ 위에서 $u$에서 평탄하다”라는
조건으로 옮겨진다. 달리 말해, (1)은
$(\mathcal{F}|_{U_v})_u$가 평탄한
$\mathcal{O}_{U_v, u}$-가군인 것과 동치이다. 스킴의 경우인
More on Morphisms, Lemma
\ref{more-morphisms-lemma-flat-and-free-at-point-fibre}에 의해
이는 $\mathcal{F}|_U$가 $u$의 어떤 열린 근방 위에서 유한 자유임을
함의한다. 이로써 (1)이 (4)를 함의함을 알 수 있다.
(4) $\Rightarrow$ (3)과 (2) $\Rightarrow$ (1)은 즉시 따른다.
(3) $\Rightarrow$ (2)에는 국소환과 줄기의 기술인
Properties of Spaces, Lemmas
\ref{spaces-properties-lemma-describe-etale-local-ring}와
\ref{spaces-properties-lemma-stalk-quasi-coherent}를 사용한다.
\end{proof}

\begin{lemma}
\label{spaces-more-morphisms-lemma-finite-free-open}
$S$를 스킴이라 하자. $f : X \to Y$를 $S$ 위 대수공간의 사상이라
하고 국소 유한 표시라고 하자. $\mathcal{F}$를 유한 표시
$\mathcal{O}_X$-가군이라 하고 $Y$ 위에서 평탄하다 하자. 그러면 집합
$$
\{x \in |X| : \mathcal{F}\text{ free in an \'etale neighbourhood of }x\}
$$
은 $|X|$에서 열린집합이고 그 구성은 임의의 밑변환
$Y' \to Y$와 가환한다.
\end{lemma}

\begin{proof}
열림성은 자명하다. $Y' \to Y$를 대수공간의 사상이라 하고
$X' = Y' \times_Y X$로 두며, $x' \in |X'|$를 $x \in |X|$ 위에
놓이는 점이라 하자. Lemma \ref{spaces-more-morphisms-lemma-flat-and-free-at-point-fibre}에
의해 $x$가 우리 집합에 속하는 것은
$(\mathcal{F}_{\overline{y}})_{\overline{x}}$가 평탄한
$\mathcal{O}_{X_{\overline{y}}, \overline{x}}$-가군인 것과 동치이다.
마찬가지로 $x'$가 당김 $\mathcal{F}'$에 대한 대응 집합에 속한다고
하자. 여기서 이는 $\mathcal{F}$를 $X'$로 당긴 것이다.
그러면 그 조건은
$(\mathcal{F}'_{\overline{y}'})_{\overline{x}'}$가 평탄한
$\mathcal{O}_{X'_{\overline{y}'}, \overline{x}'}$-가군인 것과
동치이다(기호의 뜻은 명백하다). 이 두 주장은 사상
$\text{id} : X \to X$를 $Y$ 위에서 보아 Lemma
\ref{spaces-more-morphisms-lemma-flat-on-fibres-at-point}를 적용하면 동치이다.
따라서 밑변환에 관한 주장이 성립한다.
\end{proof}






\section{뇌터 밑 위의 평탄성}
\label{spaces-more-morphisms-section-flat-over-Noetherian}

\noindent
뇌터 경우의 “섬유에 의한 평탄성 판정법”은 다음과 같다.

\begin{theorem}
\label{spaces-more-morphisms-theorem-criterion-flatness-fibre-Noetherian}
$S$를 스킴이라 하자.
$f : X \to Y$와 $Y \to Z$를 $S$ 위 대수공간의 사상들이라 하자.
$\mathcal{F}$를 준연접 $\mathcal{O}_X$-가군이라 하자.
다음을 가정한다.
\begin{enumerate}
% P09-ERR-0604: frozen hypotheses omit copulas here and in the next item.
\item $X$, $Y$, $Z$는 국소 뇌터이다.
\item $\mathcal{F}$는 연접 $\mathcal{O}_X$-가군이다.
\end{enumerate}
$x \in |X|$라 하고, $y \in |Y|$와 $z \in |Z|$를 $x$의 상들이라
하자.
% P09-ERR-0605: frozen theorem has not introduced the barred geometric point.
$\mathcal{F}_{\overline{x}} \not = 0$이면 다음은 동치이다.
\begin{enumerate}
\item $\mathcal{F}$는 $Z$ 위에서 $x$에서 평탄하고,
$\mathcal{F}$를 $z$ 위 섬유에 제한하면 $x$에서 평탄하며 그 밑은
$Y$의 $z$ 위 섬유이다.
\item $Y$는 $Z$ 위에서 $y$에서 평탄하고 $\mathcal{F}$는
$Y$ 위에서 $x$에서 평탄하다.
\end{enumerate}
\end{theorem}

\begin{proof}
Lemma \ref{spaces-more-morphisms-lemma-flat-on-fibres-at-point}의 (3)과 같은 도식을 택하자.
정의들로부터 스킴의 사상들 $U \to V \to W$, 준연접 층
$a^*\mathcal{F}$, 그리고 점 $u \in U$에 대한 대응 정리로
환원된다. 따라서 정리는 스킴에 대한 대응 결과인
More on Morphisms, Theorem
\ref{more-morphisms-theorem-criterion-flatness-fibre-Noetherian}에서
따른다.
\end{proof}

\begin{lemma}
\label{spaces-more-morphisms-lemma-morphism-between-flat-Noetherian}
$S$를 스킴이라 하자.
% P09-ERR-0606: frozen opening uses a singular noun for two morphisms.
$f : X \to Y$와 $Y \to Z$를 $S$ 위 대수공간의 사상들이라 하자.
다음을 가정한다.
\begin{enumerate}
% P09-ERR-0604: frozen hypothesis again omits the copula.
\item $X$, $Y$, $Z$는 국소 뇌터이다.
\item $X$는 $Z$ 위에서 평탄하다.
\item 모든 $z \in |Z|$에 대해 $X$의 $z$ 위 섬유는
$Y$의 $z$ 위 섬유 위에서 평탄하다.
\end{enumerate}
그러면 $f$는 평탄하다. $f$가 전사이기도 하면 $Y$는 $Z$ 위에서
평탄하다.
\end{lemma}

\begin{proof}
이는 Theorem
\ref{spaces-more-morphisms-theorem-criterion-flatness-fibre-Noetherian}의 특수한 경우이다.
\end{proof}

\noindent
매끄러움을 확인할 때와 마찬가지로 밑이 뇌터이면 아르틴 환 위에서
평탄성을 확인하는 것으로 충분하다. 다음은 그 한 예이다.

\begin{lemma}
\label{spaces-more-morphisms-lemma-flatness-over-Noetherian-ring}
$A$를 뇌터 환이라 하고 $I \subset A$를 아이디얼이라 하자.
$X$를 $S = \Spec(A)$ 위 국소 유한 표시인 대수공간이라 하자.
$n \geq 1$에 대해 $S_n = \Spec(A/I^n)$ 및
$X_n = S_n \times_S X$로 둔다.
% P09-ERR-0607: frozen English omits an article before coherent module.
$\mathcal{F}$를 연접 $\mathcal{O}_X$-가군이라 하자.
% P09-ERR-0608: frozen source says the pullback is to X, not X_n.
모든 $n \geq 1$에 대해 $\mathcal{F}_n$을 $\mathcal{F}$를 $X$로
당긴 것으로 두고 이것이 $S_n$ 위에서 평탄하면,
% P09-ERR-0609: frozen conclusion says flat over X although the proof uses S.
$\mathcal{F}$가 $X$ 위에서 평탄한 (열린) 자취는 $V(I)$의 역상을
포함하며, 그 역상은 $X \to S$에 의한다.
\end{lemma}

\begin{proof}
$\mathcal{F}$가 $S$ 위에서 평탄한 자취는 Theorem
\ref{spaces-more-morphisms-theorem-openness-flatness}에 의해 $|X|$에서 열린집합이다.
명제는 $X$를 에탈 덮개의 성원들로 바꾸어도 달라지지 않으므로
$X$가 아핀 스킴이라고 가정할 수 있다. 이 경우 결과는
Algebra, Lemma \ref{algebra-lemma-flat-module-powers}에서 즉시 따른다.
일부 세부사항은 생략한다.
\end{proof}







\section{정규화 다시 보기}
\label{spaces-more-morphisms-section-normalization}

\noindent
정규화는 매끄러운 밑변환과 가환한다.

\begin{lemma}
\label{spaces-more-morphisms-lemma-integral-closure-smooth-pullback}
$S$를 스킴이라 하자. $f : Y \to X$를 $S$ 위 대수공간의 매끄러운
사상이라 하자. $\mathcal{A}$를 준연접
$\mathcal{O}_X$-대수층이라 하자. $\mathcal{O}_Y$를
$f^*\mathcal{A}$ 안에서 정수적으로 닫아 얻는 폐포는
$f^*\mathcal{A}'$와 같다. 여기서 $\mathcal{A}' \subset \mathcal{A}$는
$\mathcal{O}_X$를 $\mathcal{A}$ 안에서 정수적으로 닫아 얻는다.
\end{lemma}

\begin{proof}
정수적 폐포의 구성, 즉 Morphisms of Spaces, Definition
\ref{spaces-morphisms-definition-integral-closure}에 의해 이는
$X$와 $Y$가 아핀인 경우로 즉시 환원된다. 이 경우 결과는
Algebra, Lemma
\ref{algebra-lemma-integral-closure-commutes-smooth}이다.
\end{proof}

\begin{lemma}[정규화는 매끄러운 밑변환과 가환한다]
\label{spaces-more-morphisms-lemma-normalization-smooth-localization}
$S$를 스킴이라 하자.
$$
\xymatrix{
Y_2 \ar[r] \ar[d] & Y_1 \ar[d]^f \\
X_2 \ar[r]^\varphi & X_1
}
$$
를 $S$ 위 대수공간의 섬유곱 정사각형이라 하자. $f$가 준콤팩트이자
준분리이고 $\varphi$가 매끄럽다고 가정하자.
$Y_i \to X_i' \to X_i$를 $X_i$의 정규화라 하며, 정규화를 취하는
곳은 $Y_i$이다.
그러면 $X_2' \cong X_2 \times_{X_1} X_1'$이다.
\end{lemma}

\begin{proof}
% P09-ERR-0610: frozen base-changed factorization ends in X_1 twice.
$Y_1 \to X_1' \to X_1$의 분해를 $X_2$로 밑변환하면
$Y_2 \to X_2 \times_{X_1} X_1' \to X_1$이라는 분해가 되고,
$X_2 \times_{X_1} X_1' \to X_1$은 정수적이다
(Morphisms of Spaces, Lemma
\ref{spaces-morphisms-lemma-base-change-integral}).
따라서 Morphisms of Spaces, Lemma
\ref{spaces-morphisms-lemma-characterize-normalization}의
보편 성질에 의해 사상
$h : X_2' \to X_2 \times_{X_1} X_1'$을 얻는다.
$X_2'$는 $\mathcal{O}_{X_2}$를
$f_{2, *}\mathcal{O}_{Y_2}$ 안에서 정수적으로 닫아 얻는 상대
스펙트럼임을 주목하자.
% P09-ERR-0611: frozen source uses O_{X_2} for an algebra on X_1.
$\mathcal{A}' \subset f_{1, *}\mathcal{O}_{Y_1}$가
$\mathcal{O}_{X_2}$의 정수적 폐포를 나타내면,
$X_2 \times_{X_1} X_1'$은 $\varphi^*\mathcal{A}'$의 상대
스펙트럼이다. 상대 스펙트럼의 구성은 임의의 밑변환과 가환하기
때문이다. Cohomology of Spaces, Lemma
\ref{spaces-cohomology-lemma-flat-base-change-cohomology}에 의해
$f_{2, *}\mathcal{O}_{Y_2} = \varphi^*f_{1, *}\mathcal{O}_{Y_1}$임을
안다. 따라서 결과는 Lemma
\ref{spaces-more-morphisms-lemma-integral-closure-smooth-pullback}에서 따른다.
\end{proof}









\section{코언-매콜리 사상}
\label{spaces-more-morphisms-section-CM}

\noindent
이는 More on Morphisms, Section
\ref{more-morphisms-section-CM}의 대수공간 판이다.

\begin{lemma}
\label{spaces-more-morphisms-lemma-CM-local-ring-fibre}
스킴의 싹들 사이의 사상에 대한 성질
\begin{align*}
& \mathcal{P}((X, x) \to (S, s)) = \\
& \text{섬유의 국소환 }
\mathcal{O}_{X_s, x}
\text{이 뇌터이고 코언-매콜리이다}
\end{align*}
는 원천과 목표에서 에탈 국소적이다(Descent, Definition
\ref{descent-definition-local-source-target-at-point}).
\end{lemma}

\begin{proof}
Descent, Definition
\ref{descent-definition-local-source-target-at-point}와 같은 도식이
주어지면 섬유들의 에탈 사상 $U'_{v'} \to U_v$를 얻으며, 이는
$u'$를 $u$로 보낸다. Descent, Lemma
\ref{descent-lemma-etale-on-fiber}를 보라. 따라서 국소환
$\mathcal{O}_{U'_{v'}, u'}$와 $\mathcal{O}_{U_v, u}$의 강 헨젤화는
같다. More on Algebra, Lemma
\ref{more-algebra-lemma-henselization-CM}에 의해 결론을 얻는다.
\end{proof}

\begin{definition}
\label{spaces-more-morphisms-definition-CM}
$S$를 스킴이라 하자.
$f : X \to Y$를 $S$ 위 대수공간의 사상이라 하자.
$f$의 섬유들이 국소 뇌터라고 가정하자
(Divisors on Spaces, Definition
\ref{spaces-divisors-definition-locally-Noetherian-fibre}).
\begin{enumerate}
\item $x \in |X|$라 하고 $y = f(x)$라 하자. $f$가
{\it $x$에서 코언-매콜리이다}라는 것은 $f$가 $x$에서 평탄하고,
Morphisms of Spaces, Lemma
\ref{spaces-morphisms-lemma-local-source-target-at-point}의 동치
조건들이 Lemma \ref{spaces-more-morphisms-lemma-CM-local-ring-fibre}에서 기술한 성질
$\mathcal{P}$에 대해 성립한다는 뜻이다.
\item $f$가 {\it 코언-매콜리 사상}이라는 것은 $f$가 $X$의 모든
점에서 코언-매콜리라는 뜻이다.
\end{enumerate}
\end{definition}

\noindent
풀어 쓰면 다음과 같다.

\begin{lemma}
\label{spaces-more-morphisms-lemma-CM}
$S$를 스킴이라 하자. $f : X \to Y$를 $S$ 위 대수공간의 사상이라
하자. $f$의 섬유들이 국소 뇌터라고 가정하자. 다음은 동치이다.
% P09-ERR-0612: frozen equivalence heading lacks list punctuation.
\begin{enumerate}
\item $f$는 코언-매콜리이다.
\item $f$는 평탄하고, 어떤 전사 에탈 사상 $V \to Y$에 대해,
여기서 $V$는 스킴이며, $X_V \to V$의 섬유들이 코언-매콜리
대수공간이다.
\item $f$는 평탄하고, 임의의 에탈 사상 $V \to Y$에 대해,
여기서 $V$는 스킴이며, $X_V \to V$의 섬유들이 코언-매콜리
대수공간이다.
\end{enumerate}
$x \in |X|$의 상을 $y \in |Y|$라 하면 다음은 동치이다.
% P09-ERR-0612: frozen second equivalence heading lacks list punctuation.
% P09-ERR-0613: frozen statement does not choose the barred geometric points.
\begin{enumerate}
\item[(a)] $f$는 $x$에서 코언-매콜리이다.
\item[(b)] $\mathcal{O}_{Y, \overline{y}} \to
\mathcal{O}_{X, \overline{x}}$는 평탄하고
$$
\mathcal{O}_{X, \overline{x}}/
\mathfrak m_{\overline{y}}\mathcal{O}_{X, \overline{x}}
$$
는 코언-매콜리이다.
\end{enumerate}
\end{lemma}

\begin{proof}
$V \to Y$를 스킴 $V$에서 오는 에탈 사상이라 하자. 스킴 $U$와
전사 에탈 사상 $U \to X \times_Y V$를 택한다. 가환 도식
$$
\xymatrix{
U \ar[d] \ar[r] & V \ar[d] \\
X \ar[r] & Y
}
$$
을 생각하자. $u \in U$의 상들을 $x \in |X|$, $y \in |Y|$,
$v \in V$라 하자. 그러면 $f$가 $x$에서 코언-매콜리인 것과
$U \to V$가 $u$에서 코언-매콜리인 것은 동치이다(정의에 의해).
또한 사상 $U_v \to X_v = (X_V)_v$는 전사 에탈이다. 따라서 스킴
$U_v$가 코언-매콜리인 것과 대수공간 $X_v$가 코언-매콜리인 것은
동치이다. 그러므로 (1), (2), (3)의 동치는 스킴 사상들에 대한
대응 동치에서 형식적으로 따른다. More on Morphisms, Lemma
\ref{more-morphisms-lemma-CM}을 보라.

\medskip\noindent
(a)와 (b)의 동치를 증명하자. 평탄성에 대한 대응 동치는
Morphisms of Spaces, Lemma
\ref{spaces-morphisms-lemma-flat-at-point-etale-local-rings}이다.
따라서 동치를 증명할 때 $f$가 $x$에서 평탄하다고 가정할 수 있다.
위와 같은 도식과 $x, y, u, v$를 생각하자. 그러면
$$
\mathcal{O}_{Y, \overline{y}} \to \mathcal{O}_{X, \overline{x}}
$$
는 국소환의 강 헨젤화 사이의 사상
$$
\mathcal{O}_{V, v}^{sh} \to \mathcal{O}_{U, u}^{sh}
$$
과 같다. Properties of Spaces, Lemma
\ref{spaces-properties-lemma-describe-etale-local-ring}를 보라.
따라서 Algebra, Lemma
\ref{algebra-lemma-quotient-strict-henselization}에 의해
$$
\mathcal{O}_{X, \overline{x}}/
\mathfrak m_{\overline{y}}\mathcal{O}_{X, \overline{x}} =
(\mathcal{O}_{U, u}/\mathfrak m_v \mathcal{O}_{U, u})^{sh}
$$
이다. 그러므로 뇌터 국소환
$\mathcal{O}_{U, u}/\mathfrak m_v \mathcal{O}_{U, u}$이
코언-매콜리인 것과 그 강 헨젤화가 코언-매콜리인 것이 동치임을
보이면 된다. 이는 More on Algebra, Lemma
\ref{more-algebra-lemma-henselization-CM}이다.
\end{proof}

\begin{lemma}
\label{spaces-more-morphisms-lemma-composition-CM}
$S$를 스킴이라 하자.
$f : X \to Y$와 $g : Y \to Z$를 $S$ 위 대수공간의 사상들이라
하자. $f$, $g$, $g \circ f$의 섬유들이 국소 뇌터라고 가정하자.
$x \in |X|$의 상들을 $y \in |Y|$, $z \in |Z|$라 하자.
\begin{enumerate}
\item $f$가 $x$에서 코언-매콜리이고 $g$가 $f(x)$에서
코언-매콜리이면 $g \circ f$는 $x$에서 코언-매콜리이다.
\item $f$와 $g$가 코언-매콜리이면 $g \circ f$는
코언-매콜리이다.
\item $g \circ f$가 $x$에서 코언-매콜리이고 $f$가 $x$에서
평탄하면, $f$는 $x$에서 코언-매콜리이고 $g$는 $f(x)$에서
코언-매콜리이다.
% P09-ERR-0614: frozen hypothesis uses the ill-typed composite f o g.
\item $f \circ g$가 코언-매콜리이고 $f$가 평탄하면, $f$는
코언-매콜리이고 $g$는 $f$의 상에 속하는 모든 점에서
코언-매콜리이다.
\end{enumerate}
\end{lemma}

\begin{proof}
에탈 국소적으로 작업하면 스킴에 대한 대응 결과인
More on Morphisms, Lemma
\ref{more-morphisms-lemma-composition-CM}에서 따른다.
다른 방법으로 Lemma \ref{spaces-more-morphisms-lemma-CM}의 (a)와 (b)의 동치를 사용할
수 있다. 즉 뇌터 국소환들의 국소 준동형
$$
\mathcal{O}_{Y, \overline{y}}/
\mathfrak m_{\overline{z}}\mathcal{O}_{Y, \overline{y}}
\longrightarrow
\mathcal{O}_{X, \overline{x}}/
\mathfrak m_{\overline{z}}\mathcal{O}_{X, \overline{x}}
$$
을 생각한다. 그 섬유는
$$
\mathcal{O}_{X, \overline{x}}/
\mathfrak m_{\overline{y}}\mathcal{O}_{X, \overline{x}}
$$
이고, Algebra, Lemma \ref{algebra-lemma-CM-goes-up}을 사용한다.
\end{proof}

\begin{lemma}
\label{spaces-more-morphisms-lemma-flat-morphism-from-CM}
$S$를 스킴이라 하자.
$f : X \to Y$를 $S$ 위 국소 뇌터 대수공간들의 평탄한 사상이라
하자. $X$가 코언-매콜리이면 $f$는 코언-매콜리이고
% P09-ERR-0615: frozen conclusion uses f(xbar) without choosing xbar.
$\mathcal{O}_{Y, f(\overline{x})}$는 모든 $x \in |X|$에 대해
코언-매콜리이다.
\end{lemma}

\begin{proof}
Lemma \ref{spaces-more-morphisms-lemma-CM}을 사용하여 대수로 옮기면(Lemma
\ref{spaces-more-morphisms-lemma-composition-CM}의 증명과 비교하라), 이는
Algebra, Lemma \ref{algebra-lemma-CM-goes-up}에서 따른다.
\end{proof}

\begin{lemma}
\label{spaces-more-morphisms-lemma-base-change-CM}
$S$를 스킴이라 하자.
$f : X \to Y$를 $S$ 위 대수공간의 사상이라 하자.
$f$의 섬유들이 국소 뇌터라고 가정하자.
$Y' \to Y$가 국소 유한형이라고 하자. $f' : X' \to Y'$를
$f$의 밑변환이라 하자.
$x' \in |X'|$의 상을 $x \in |X|$라 하자.
\begin{enumerate}
\item $f$가 $x$에서 코언-매콜리이면
$f' : X' \to Y'$는 $x'$에서 코언-매콜리이다.
\item $f$가 $x$에서 평탄하고 $f'$가 $x'$에서 코언-매콜리이면
$f$는 $x$에서 코언-매콜리이다.
\item $Y' \to Y$가 $f'(x')$에서 평탄하고 $f'$가 $x'$에서
코언-매콜리이면 $f$는 $x$에서 코언-매콜리이다.
\end{enumerate}
\end{lemma}

\begin{proof}
% P09-ERR-0616: frozen source calls y and y' one image collectively.
$y \in |Y|$와 $y' \in |Y'|$를 $x'$의 상으로 나타내자.
전사 에탈 사상 $V \to Y$를 택하되 $V$는 스킴이라 하자.
전사 에탈 사상 $U \to X \times_Y V$를 택하되 $U$는 스킴이라 하자.
% P09-ERR-0617: frozen source misspells surjective.
전사 에탈 사상 $V' \to Y' \times_Y V$를 택하되 $V'$는 스킴이라
하자. 그러면 $U' = U \times_V V'$는 스킴이고 전사 에탈 사상
$U' \to X'$가 함께 주어진다. $u' \in U'$를 택하여 그 상이
$x'$가 되게 하자.
$u \in U$를 $u'$의 상이라 하자. 이제 보조정리는
$U \to V$와 그 밑변환 $U' \to V'$ 및 점 $u'$, $u$에 대한
보조정리에서 따른다(정의에서 따른다). 따라서 보조정리는 스킴의
경우인 More on Morphisms, Lemma
\ref{more-morphisms-lemma-base-change-CM}에서 따른다.
\end{proof}

\begin{lemma}
\label{spaces-more-morphisms-lemma-flat-finite-presentation-CM-open}
$S$를 스킴이라 하자.
$f : X \to Y$를 $S$ 위 대수공간의 사상이라 하고 평탄하며
국소 유한 표시라고 하자. 다음과 같이 두자.
$$
W = \{x \in |X| : f\text{ is Cohen-Macaulay at }x\}
$$
그러면 $W$는 $|X|$에서 열린집합이고 $W$의 구성은 $f$의 임의의
밑변환과 가환한다. 임의의 사상 $g : Y' \to Y$에 대해 밑변환
$f' : X' \to Y'$를 생각하되 이는 $f$에서 얻고, 사영
$g' : X' \to X$도 생각하자.
% P09-ERR-0618: frozen source duplicates W' in the base-change identity.
그러면 대응 집합 $W'$를 사상 $f'$에 대해 취했을 때 이는
$W' = (g')^{-1}(W)$와 같다.
\end{lemma}

\begin{proof}
가환 도식
$$
\xymatrix{
U \ar[d] \ar[r] & V \ar[d] \\
X \ar[r] & Y
}
$$
에서 수직 화살표들이 에탈이고 $U$와 $V$가 스킴이 되도록 택하자.
$u \in U$의 상을 $x \in |X|$라 하자. 그러면 $f$가 $x$에서
코언-매콜리인 것과 $U \to V$가 $u$에서 코언-매콜리인 것은
동치이다(정의에 의해). 따라서 사상 $U \to V$의 경우로 환원된다.
More on Morphisms, Lemma
\ref{more-morphisms-lemma-flat-finite-presentation-CM-open}을 보라.
\end{proof}

\begin{lemma}
\label{spaces-more-morphisms-lemma-lfp-CM-relative-dimension}
$S$를 스킴이라 하자. $f : X \to Y$를 $S$ 위 대수공간의 사상이라
하자. $f$가 국소 유한 표시이고 코언-매콜리라고 가정하자.
% P09-ERR-0619: frozen statement says subschemes although X is a space.
열리고 닫힌 부분스킴 $X_d \subset X$들이 존재하여
$X = \coprod_{d \geq 0} X_d$이고
$f|_{X_d} : X_d \to Y$의 상대 차원은 $d$이다.
\end{lemma}

\begin{proof}
가환 도식
$$
\xymatrix{
U \ar[d] \ar[r] & V \ar[d] \\
X \ar[r] & Y
}
$$
에서 수직 화살표들이 에탈이고 $U$와 $V$가 스킴이 되도록 택하자.
그러면 $U \to V$는 국소 유한 표시이고 코언-매콜리이다
(우리 정의들에서 즉시 따른다). 따라서
$U = \coprod_{d \geq 0} U_d$라는 열리고 닫힌 부분스킴들로의 분해가
있으며,
% P09-ERR-0620: frozen proof restricts f to U_d although f has domain X.
$f|_{U_d} : U_d \to V$의 상대 차원은 $d$이다. Morphisms, Lemma
\ref{morphisms-lemma-flat-finite-presentation-CM-fibres-relative-dimension}를
보라. $u \in U$의 상을 $x \in |X|$라 하자. 그러면 $f$의
상대 차원이 $d$이고 이를 보는 점이 $x$인 것과 $U \to V$의
상대 차원이 $d$이고 이를 보는 점이 $u$인 것은 동치이다
(우리 정의들에서 따른다). 이로부터 $U_d$가
어떤 부분집합 $X_d \subset |X|$의 역상임을 알 수 있고, 이 부분집합은
반드시 열리고 닫혀 있다. $X_d$로 $X$의 대응하는 열리고 닫힌
대수부분공간을 나타내면 보조정리가 성립한다.
\end{proof}







\section{고렌슈타인 사상}
\label{spaces-more-morphisms-section-gorenstein}

\noindent
이는 Duality for Schemes, Section
\ref{duality-section-gorenstein-morphisms}의 대수공간 판이다.

\begin{lemma}
\label{spaces-more-morphisms-lemma-gorenstein-local-ring-fibre}
스킴의 싹들 사이의 사상에 대한 성질
\begin{align*}
& \mathcal{P}((X, x) \to (S, s)) = \\
& \text{섬유의 국소환 }
\mathcal{O}_{X_s, x}
\text{이 뇌터이고 고렌슈타인이다}
\end{align*}
는 원천과 목표에서 에탈 국소적이다(Descent, Definition
\ref{descent-definition-local-source-target-at-point}).
\end{lemma}

\begin{proof}
Descent, Definition
\ref{descent-definition-local-source-target-at-point}과 같은 도식이
주어지면 섬유들의 에탈 사상 $U'_{v'} \to U_v$를 얻으며, 이는
$u'$를 $u$로 보낸다. Descent, Lemma
\ref{descent-lemma-etale-on-fiber}를 보라. 따라서
$\mathcal{O}_{U_v, u} \to \mathcal{O}_{U'_{v'}, u'}$는 에탈 환
사상의 국소화이다. 그러므로 첫째 환이 뇌터인 것과 둘째 환이
뇌터인 것은 동치이다. More on Algebra, Lemma
\ref{more-algebra-lemma-Noetherian-etale-extension}을 보라. 이제
$$
\mathcal{O}_{U'_{v'}, u'}/\mathfrak m_u \mathcal{O}_{U'_{v'}, u'}
= \kappa(u')
$$
는 고렌슈타인 환이므로(Algebra, Lemma
\ref{algebra-lemma-etale-at-prime}),
$\mathcal{O}_{U_v, u}$가 고렌슈타인인 것과
$\mathcal{O}_{U'_{v'}, u'}$가 고렌슈타인인 것은
Dualizing Complexes, Lemma
\ref{dualizing-lemma-flat-under-gorenstein}에 의해 동치이다.
\end{proof}

\begin{definition}
\label{spaces-more-morphisms-definition-gorenstein}
$S$를 스킴이라 하자.
$f : X \to Y$를 $S$ 위 대수공간의 사상이라 하자.
$f$의 섬유들이 국소 뇌터라고 가정하자
(Divisors on Spaces, Definition
\ref{spaces-divisors-definition-locally-Noetherian-fibre}).
\begin{enumerate}
\item $x \in |X|$라 하고 $y = f(x)$라 하자. $f$가
{\it $x$에서 고렌슈타인이다}라는 것은 $f$가 $x$에서 평탄하고,
Morphisms of Spaces, Lemma
\ref{spaces-morphisms-lemma-local-source-target-at-point}의 동치
조건들이 Lemma \ref{spaces-more-morphisms-lemma-gorenstein-local-ring-fibre}에서 기술한
성질 $\mathcal{P}$에 대해 성립한다는 뜻이다.
\item $f$가 {\it 고렌슈타인 사상}이라는 것은 $f$가 $X$의 모든
점에서 고렌슈타인이라는 뜻이다.
\end{enumerate}
\end{definition}

\noindent
풀어 쓰면 다음과 같다.

\begin{lemma}
\label{spaces-more-morphisms-lemma-gorenstein}
$S$를 스킴이라 하자. $f : X \to Y$를 $S$ 위 대수공간의 사상이라
하자. $f$의 섬유들이 국소 뇌터라고 가정하자. 다음은 동치이다.
% P09-ERR-0621: frozen equivalence heading lacks list punctuation.
\begin{enumerate}
\item $f$는 고렌슈타인이다.
\item $f$는 평탄하고, 어떤 전사 에탈 사상 $V \to Y$에 대해,
여기서 $V$는 스킴이며, $X_V \to V$의 섬유들이 고렌슈타인
대수공간이다.
\item $f$는 평탄하고, 임의의 에탈 사상 $V \to Y$에 대해,
여기서 $V$는 스킴이며, $X_V \to V$의 섬유들이 고렌슈타인
대수공간이다.
\end{enumerate}
$x \in |X|$의 상을 $y \in |Y|$라 하면 다음은 동치이다.
% P09-ERR-0621: frozen second equivalence heading lacks list punctuation.
% P09-ERR-0622: frozen statement does not choose the barred geometric points.
\begin{enumerate}
\item[(a)] $f$는 $x$에서 고렌슈타인이다.
\item[(b)] $\mathcal{O}_{Y, \overline{y}} \to
\mathcal{O}_{X, \overline{x}}$는 평탄하고
$$
\mathcal{O}_{X, \overline{x}}/
\mathfrak m_{\overline{y}}\mathcal{O}_{X, \overline{x}}
$$
는 고렌슈타인이다.
\end{enumerate}
\end{lemma}

\begin{proof}
$V \to Y$를 스킴 $V$에서 오는 에탈 사상이라 하자. 스킴 $U$와
전사 에탈 사상 $U \to X \times_Y V$를 택한다. 가환 도식
$$
\xymatrix{
U \ar[d] \ar[r] & V \ar[d] \\
X \ar[r] & Y
}
$$
을 생각하자. $u \in U$의 상들을 $x \in |X|$, $y \in |Y|$,
$v \in V$라 하자. 그러면 $f$가 $x$에서 고렌슈타인인 것과
$U \to V$가 $u$에서 고렌슈타인인 것은 동치이다(정의에 의해).
또한 사상 $U_v \to X_v = (X_V)_v$는 전사 에탈이다. 따라서 스킴
$U_v$가 고렌슈타인인 것과 대수공간 $X_v$가 고렌슈타인인 것은
동치이다. 그러므로 (1), (2), (3)의 동치는 스킴 사상들에 대한
대응 동치에서 형식적으로 따른다. Duality for Schemes, Lemma
\ref{duality-lemma-gorenstein}을 보라.

\medskip\noindent
(a)와 (b)의 동치를 증명하자. 평탄성에 대한 대응 동치는
Morphisms of Spaces, Lemma
\ref{spaces-morphisms-lemma-flat-at-point-etale-local-rings}이다.
따라서 동치를 증명할 때 $f$가 $x$에서 평탄하다고 가정할 수 있다.
위와 같은 도식과 $x, y, u, v$를 생각하자. 그러면
$$
\mathcal{O}_{Y, \overline{y}} \to \mathcal{O}_{X, \overline{x}}
$$
는 국소환의 강 헨젤화 사이의 사상
$$
\mathcal{O}_{V, v}^{sh} \to \mathcal{O}_{U, u}^{sh}
$$
과 같다. Properties of Spaces, Lemma
\ref{spaces-properties-lemma-describe-etale-local-ring}를 보라.
따라서 Algebra, Lemma
\ref{algebra-lemma-quotient-strict-henselization}에 의해
$$
\mathcal{O}_{X, \overline{x}}/
\mathfrak m_{\overline{y}}\mathcal{O}_{X, \overline{x}} =
(\mathcal{O}_{U, u}/\mathfrak m_v \mathcal{O}_{U, u})^{sh}
$$
이다. 그러므로 뇌터 국소환
$\mathcal{O}_{U, u}/\mathfrak m_v \mathcal{O}_{U, u}$이
고렌슈타인인 것과 그 강 헨젤화가 고렌슈타인인 것이 동치임을
보이면 된다. 이는 Dualizing Complexes, Lemma
\ref{dualizing-lemma-completion-henselization-dualizing}와
자기 자신 위의 쌍대화 복합체인 뇌터 국소환이라는 고렌슈타인
국소환의 정의에서 즉시 따른다.
\end{proof}

\begin{lemma}
\label{spaces-more-morphisms-lemma-composition-gorenstein}
$S$를 스킴이라 하자.
$f : X \to Y$와 $g : Y \to Z$를 $S$ 위 대수공간의 사상들이라
하자. $f$, $g$, $g \circ f$의 섬유들이 국소 뇌터라고 가정하자.
$x \in |X|$의 상들을 $y \in |Y|$, $z \in |Z|$라 하자.
\begin{enumerate}
\item $f$가 $x$에서 고렌슈타인이고 $g$가 $f(x)$에서
고렌슈타인이면 $g \circ f$는 $x$에서 고렌슈타인이다.
\item $f$와 $g$가 고렌슈타인이면 $g \circ f$는
고렌슈타인이다.
\item $g \circ f$가 $x$에서 고렌슈타인이고 $f$가 $x$에서
평탄하면, $f$는 $x$에서 고렌슈타인이고 $g$는 $f(x)$에서
고렌슈타인이다.
% P09-ERR-0623: frozen hypothesis uses the ill-typed composite f o g.
\item $f \circ g$가 고렌슈타인이고 $f$가 평탄하면, $f$는
고렌슈타인이고 $g$는 $f$의 상에 속하는 모든 점에서
고렌슈타인이다.
\end{enumerate}
\end{lemma}

\begin{proof}
에탈 국소적으로 작업하면 스킴에 대한 대응 결과인
Duality for Schemes, Lemma
\ref{duality-lemma-composition-gorenstein}에서 따른다.
다른 방법으로 Lemma \ref{spaces-more-morphisms-lemma-gorenstein}의 (a)와 (b)의 동치를
사용할 수 있다. 즉 뇌터 국소환들의 국소 준동형
$$
\mathcal{O}_{Y, \overline{y}}/
\mathfrak m_{\overline{z}}\mathcal{O}_{Y, \overline{y}}
\longrightarrow
\mathcal{O}_{X, \overline{x}}/
\mathfrak m_{\overline{z}}\mathcal{O}_{X, \overline{x}}
$$
을 생각한다. 그 섬유는
$$
\mathcal{O}_{X, \overline{x}}/
\mathfrak m_{\overline{y}}\mathcal{O}_{X, \overline{x}}
$$
이고, Dualizing Complexes, Lemma
\ref{dualizing-lemma-flat-under-gorenstein}을 사용한다.
\end{proof}

\begin{lemma}
\label{spaces-more-morphisms-lemma-flat-morphism-from-gorenstein}
$S$를 스킴이라 하자.
$f : X \to Y$를 $S$ 위 국소 뇌터 대수공간들의 평탄한 사상이라
하자. $X$가 고렌슈타인이면 $f$는 고렌슈타인이고
% P09-ERR-0624: frozen conclusion uses f(xbar) without choosing xbar.
$\mathcal{O}_{Y, f(\overline{x})}$는 모든 $x \in |X|$에 대해
고렌슈타인이다.
\end{lemma}

\begin{proof}
Lemma \ref{spaces-more-morphisms-lemma-gorenstein}을 사용하여 대수로 옮기면(Lemma
\ref{spaces-more-morphisms-lemma-composition-gorenstein}의 증명과 비교하라), 이는
Dualizing Complexes, Lemma
\ref{dualizing-lemma-flat-under-gorenstein}에서 따른다.
\end{proof}

\begin{lemma}
\label{spaces-more-morphisms-lemma-base-change-gorenstein}
$S$를 스킴이라 하자.
$f : X \to Y$를 $S$ 위 대수공간의 사상이라 하자.
$f$의 섬유들이 국소 뇌터라고 가정하자.
$Y' \to Y$가 국소 유한형이라고 하자. $f' : X' \to Y'$를
$f$의 밑변환이라 하자.
$x' \in |X'|$의 상을 $x \in |X|$라 하자.
\begin{enumerate}
\item $f$가 $x$에서 고렌슈타인이면
$f' : X' \to Y'$는 $x'$에서 고렌슈타인이다.
\item $f$가 $x$에서 평탄하고 $f'$가 $x'$에서 고렌슈타인이면
$f$는 $x$에서 고렌슈타인이다.
\item $Y' \to Y$가 $f'(x')$에서 평탄하고 $f'$가 $x'$에서
고렌슈타인이면 $f$는 $x$에서 고렌슈타인이다.
\end{enumerate}
\end{lemma}

\begin{proof}
% P09-ERR-0625: frozen source calls y and y' one image collectively.
$y \in |Y|$와 $y' \in |Y'|$를 $x'$의 상으로 나타내자.
전사 에탈 사상 $V \to Y$를 택하되 $V$는 스킴이라 하자.
전사 에탈 사상 $U \to X \times_Y V$를 택하되 $U$는 스킴이라 하자.
% P09-ERR-0626: frozen source misspells surjective.
전사 에탈 사상 $V' \to Y' \times_Y V$를 택하되 $V'$는 스킴이라
하자. 그러면 $U' = U \times_V V'$는 스킴이고 전사 에탈 사상
$U' \to X'$가 함께 주어진다. $u' \in U'$를 택하여 그 상이
$x'$가 되게 하자. $u \in U$를 $u'$의 상이라 하자. 이제 보조정리는
$U \to V$와 그 밑변환 $U' \to V'$ 및 점 $u'$, $u$에 대한
보조정리에서 따른다(정의에서 따른다). 따라서 보조정리는 스킴의
경우인 Duality for Schemes, Lemma
\ref{duality-lemma-base-change-gorenstein}에서 따른다.
\end{proof}

\begin{lemma}
\label{spaces-more-morphisms-lemma-flat-finite-presentation-gorenstein-open}
$S$를 스킴이라 하자.
$f : X \to Y$를 $S$ 위 대수공간의 사상이라 하고 평탄하며
국소 유한 표시라고 하자. 다음과 같이 두자.
$$
W = \{x \in |X| : f\text{ is Gorenstein at }x\}
$$
그러면 $W$는 $|X|$에서 열린집합이고 $W$의 구성은 $f$의 임의의
밑변환과 가환한다. 임의의 사상 $g : Y' \to Y$에 대해 밑변환
$f' : X' \to Y'$를 생각하되 이는 $f$에서 얻고, 사영
$g' : X' \to X$도 생각하자.
% P09-ERR-0627: frozen source duplicates W' in the base-change identity.
그러면 대응 집합 $W'$를 사상 $f'$에 대해 취했을 때 이는
$W' = (g')^{-1}(W)$와 같다.
\end{lemma}

\begin{proof}
% P09-ERR-0628: frozen proof leaves the diagram's etale/scheme hypotheses unstated.
가환 도식
$$
\xymatrix{
U \ar[d] \ar[r] & V \ar[d] \\
X \ar[r] & Y
}
$$
을 택하자. $u \in U$의 상을 $x \in |X|$라 하자. 그러면 $f$가
$x$에서 고렌슈타인인 것과 $U \to V$가 $u$에서 고렌슈타인인 것은
동치이다(정의에 의해). 따라서 스킴 사상 $U \to V$의 경우로
환원된다. 열림성은 Duality for Schemes, Lemma
\ref{duality-lemma-flat-finite-presentation-characterize-gorenstein}에서
증명되었고 밑변환과의 호환성은 Duality for Schemes, Lemma
\ref{duality-lemma-flat-lft-base-change-gorenstein}에서 증명되었다.
\end{proof}











\section{코언-매콜리 사상의 절단}
\label{spaces-more-morphisms-section-slice}

\noindent
$S$를 스킴이라 하자. $X$를 $S$ 위 대수공간이라 하자.
$f_1, \ldots, f_r \in \Gamma(X, \mathcal{O}_X)$라 하자. 이 경우
$V(f_1, \ldots, f_r)$로 {\it $X$의 닫힌 부분공간, 즉
$f_1, \ldots, f_r$가 잘라내는 부분공간}을 나타낸다. 더 정확히
말하면, $V(f_1, \ldots, f_r)$를 $X$의 닫힌 부분공간으로 정의하되
이는 $f_1, \ldots, f_r$가 생성하는 준연접 아이디얼층에 대응한다.
Morphisms of Spaces, Lemma
\ref{spaces-morphisms-lemma-closed-immersion-ideals}를 보라.
다른 방법으로 표시 $X = U/R$을 택하고 닫힌 부분스킴
$Z \subset U$를 생각할 수 있으며, 이는 $f_1|U, \ldots, f_r|_U$가
잘라낸다. $Z$가 $R$-불변인 닫힌 부분스킴임은
명백하고(Groupoids, Definition
\ref{groupoids-definition-invariant-open}을 보라),
$V(f_1, \ldots, f_r) = Z/R_Z$로 둘 수 있다.

\begin{lemma}
\label{spaces-more-morphisms-lemma-slice}
$S$를 스킴이라 하자. 카르테시안 도식
$$
\xymatrix{
X \ar[d] & F \ar[l]^p \ar[d] \\
Y & \Spec(k) \ar[l]
}
$$
을 생각하자. 여기서 $X \to Y$는 $S$ 위 대수공간의 사상으로
평탄하고 국소 유한 표시이며, $k$는 $S$ 위 체이다.
$f_1, \ldots, f_r \in \Gamma(X, \mathcal{O}_X)$라 하고
% P09-ERR-0629: frozen statement does not choose the barred geometric point.
$z \in |F|$에서 $f_1, \ldots, f_r$의 상들이 국소환
$\mathcal{O}_{F, \overline{z}}$에서 정칙열을 이룬다고 하자.
그러면 $X$를 $p(z)$를 포함하는 열린 부분공간으로 바꾼 뒤 사상
$$
V(f_1, \ldots, f_r) \longrightarrow Y
$$
은 평탄하고 국소 유한 표시이다.
\end{lemma}

\begin{proof}
$Z = V(f_1, \ldots, f_r)$로 두자. $Z \to X$가 국소 유한 표시임은
명백하므로 합성 $Z \to Y$도 국소 유한 표시이다. Morphisms of Spaces,
Lemma \ref{spaces-morphisms-lemma-composition-finite-presentation}를
보라. 따라서 $Z \to Y$가 $p(z)$의 어떤 근방에서 평탄함을 보이면
충분하다. $k'/k$를 체의 확대라 하자.
% P09-ERR-0630: frozen source attributes surjectivity directly to F'.
$F' = F \times_{\Spec(k)} \Spec(k')$은 $F$ 위에서 전사이고 평탄하다.
따라서 점 $z' \in |F'|$를 찾아 그것이 $z$로 가게 할 수 있고
국소환 사상
$$
\mathcal{O}_{F, \overline{z}} \to
\mathcal{O}_{F', \overline{z}'}
$$
은 평탄하다. Morphisms of Spaces, Lemma
\ref{spaces-morphisms-lemma-flat-at-point-etale-local-rings}를 보라.
따라서 $f_1, \ldots, f_r$의 $\mathcal{O}_{F', \overline{z}'}$에서의
상들도 정칙열을 이룬다. Algebra, Lemma
\ref{algebra-lemma-flat-increases-depth}를 보라. 그러므로 증명 중에
$k$를 확대체로 바꾸어도 된다. 특히 $z \in |F|$가 절단
$z : \Spec(k) \to F$에서 온다고 가정할 수 있으며, 이 절단이 놓이는
구조사상은 $F \to \Spec(k)$이다.

\medskip\noindent
스킴 $V$와 전사 에탈 사상 $V \to Y$를 택하자. 스킴 $U$와 전사
에탈 사상 $U \to X \times_Y V$를 택하자. 필요하면 $k$를 한 번 더
확대하여 $\Spec(k) \to F \to X$가 $U$를 통해 분해된다고 가정할 수
있다($U \to X$가 전사이므로). 그러한 분해를
$u : \Spec(k) \to U$라 하고 $v \in V$를 $u$의 상이라 하자. 사상들
$$
U_v \times_{\Spec(\kappa(v))} \Spec(k) =
U \times_V \Spec(k) \to U \times_Y \Spec(k) \to F
$$
은 에탈임을 주목하자. 첫째는 $V \to V \times_Y V$의 밑변환이고,
둘째는 $U \to X$의 밑변환이기 때문이다. 또한 구성에 의해 점
% P09-ERR-0631: frozen source splits the compound adjective leftmost.
$u : \Spec(k) \to U$는 가장 왼쪽 공간의 점을 주며 오른쪽의
$z$로 간다. 따라서 원소 $f_1, \ldots, f_r$는 다음 사상의 오른쪽
국소환에서 정칙열로 간다.
% P09-ERR-0632: frozen source has a misplaced parenthesis in the subscript.
$$
\mathcal{O}_{U_v, u}
\longrightarrow
\mathcal{O}_{U_v \times_{\Spec(\kappa(v)} \Spec(k), \overline{u}}
=
\mathcal{O}_{U \times_V \Spec(k), \overline{u}}.
$$
그러나 표시된 화살표는 평탄하다(More on Flatness, Lemma
\ref{flat-lemma-flat-up-down-henselization}와 Morphisms of Spaces,
Lemma \ref{spaces-morphisms-lemma-flat-at-point-etale-local-rings}를
함께 사용한다). 따라서 Algebra, Lemma
\ref{algebra-lemma-flat-increases-depth}에 의해
% P09-ERR-0633: frozen plural subject takes a singular English verb.
$f_1, \ldots, f_r$가 $\mathcal{O}_{U_v, u}$에서 정칙열로 감을
알 수 있다. More on Morphisms, Lemma
\ref{more-morphisms-lemma-slice-given-elements}에 의해 스킴의 사상
$$
V(f_1, \ldots, f_r) \times_X U =
V(f_1|_U, \ldots, f_r|_U) \to V
$$
은 어떤 열린 근방 $U'$에서 평탄하며 이 근방은 $u$를 포함한다.
$X' \subset X$를
$|U'| \to |X|$의 상에 대응하는 열린 부분공간이라 하자
(Properties of Spaces, Lemmas
\ref{spaces-properties-lemma-topology-points}와
\ref{spaces-properties-lemma-open-subspaces}를 보라). 그러면
$V(f_1, \ldots, f_r) \cap X' \to Y$는 평탄하다
(Morphisms of Spaces, Definition
\ref{spaces-morphisms-definition-flat}을 보라). 실제로 가환 도식
$$
\xymatrix{
V(f_1, \ldots, f_r) \times_X U' \ar[d]_a \ar[r] & V \ar[d]^b \\
V(f_1, \ldots, f_r) \cap X' \ar[r] & Y
}
$$
에서 $a, b$는 에탈이고 $a$는 전사이다.
\end{proof}






\section{축약 섬유}
\label{spaces-more-morphisms-section-reduced}

\noindent
이 절은 More on Morphisms, Section
\ref{more-morphisms-section-reduced}의 대수공간 판이다.

\begin{lemma}
\label{spaces-more-morphisms-lemma-geometrically-reduced-fibre}
$S$를 스킴이라 하자. $f : X \to Y$를 $S$ 위 대수공간의 사상이라
하자. $y \in |Y|$라 하자. 다음은 동치이다.
% P09-ERR-0634: frozen equivalence heading lacks list punctuation.
\begin{enumerate}
\item 어떤 사상 $\Spec(k) \to Y$가 $y$의 동치류에 속하고
대수공간 $X_k$가 $k$ 위에서 기하학적으로 축약이다.
\item 모든 사상 $\Spec(k) \to Y$가 $y$의 동치류에 속할 때
대수공간 $X_k$가 $k$ 위에서 기하학적으로 축약이다.
\item 모든 사상 $\Spec(k) \to Y$가 $y$의 동치류에 속할 때
대수공간 $X_k$가 축약이다.
\end{enumerate}
\end{lemma}

\begin{proof}
이는 Spaces over Fields, Lemma
\ref{spaces-over-fields-lemma-geometrically-reduced-upstairs}와
% P09-ERR-0635: frozen proof cites the equivalence relation defining |X|.
$|X|$를 정의하는 동치관계의 정의에서 즉시 따른다. Properties of
Spaces, Section \ref{spaces-properties-section-points}에 그 정의가 있다.
\end{proof}

\begin{definition}
\label{spaces-more-morphisms-definition-geometrically-reduced-fibre}
$S$를 스킴이라 하자. $f : X \to Y$를 $S$ 위 대수공간의 사상이라
하자. $y \in |Y|$라 하자. Lemma
\ref{spaces-more-morphisms-lemma-geometrically-reduced-fibre}의 동치 조건들이 성립하면
{\it $f : X \to Y$의 섬유가 $y$에서 기하학적으로 축약이다}라고
한다.
\end{definition}

\noindent
필수적인 보조정리들은 다음과 같다.

\begin{lemma}
\label{spaces-more-morphisms-lemma-base-change-fibres-geometrically-reduced}
$S$를 스킴이라 하자. $f : X \to Y$와 $g : Y' \to Y$를 $S$ 위
대수공간의 사상들이라 하자. $f' : X' \to Y'$로 $f$의 밑변환을
나타내되, 이 밑변환은 $g$에 의한다. 그러면
\begin{align*}
\{y' \in |Y'| :
\text{the fibre of }f' : X' \to Y'\text{ at }y'
\text{ is geometrically reduced}\} \\
= g^{-1}(\{y \in |Y| :
\text{the fibre of }f : X \to Y\text{ at }y
\text{ is geometrically reduced}\}).
\end{align*}
\end{lemma}

\begin{proof}
$y' \in |Y'|$에 대해 사상 $\Spec(k) \to Y'$를 택하여 이것이
$y'$의 동치류에 속하게 하자. 그러면 $g(y')$는 합성
$\Spec(k) \to Y' \to Y$로 표시된다. 따라서
$X' \times_{Y'} \Spec(k) = X \times_Y \Spec(k)$이고 결과는
정의에서 따른다.
\end{proof}

\begin{lemma}
\label{spaces-more-morphisms-lemma-geometrically-reduced-constructible}
$S$를 스킴이라 하자. $f : X \to Y$를 $S$ 위 대수공간의 사상이라
하고 준콤팩트이며 국소 유한 표시라고 하자. 그러면 집합
$$
E = \{y \in |Y| : \text{the fibre of }f : X \to Y\text{ at }y
\text{ is geometrically reduced}\}
$$
은 에탈 국소적으로 구성가능하다.
\end{lemma}

\begin{proof}
아핀 스킴 $V$와 에탈 사상 $V \to Y$를 택하자. 명제의 뜻은
$E$의 $|V|$에서의 역상이 구성가능하다는 것이다. Lemma
\ref{spaces-more-morphisms-lemma-base-change-fibres-geometrically-reduced}에 의해 $Y$를
$V$로 바꿀 수 있다. 즉 $Y$가 아핀 스킴이라고 가정할 수 있다.
그러면 $X$는 준콤팩트이다. 아핀 스킴 $U$와 전사 에탈 사상
$U \to X$를 택하자. 사상 $\Spec(k) \to Y$에 대해 섬유 사이의
사상 $U_k \to X_k$는 전사 에탈이다. 따라서 $U_k$가 $k$ 위에서
기하학적으로 축약인 것과 $X_k$가 $k$ 위에서 기하학적으로 축약인
것은 동치이다. Spaces over Fields, Lemma
\ref{spaces-over-fields-lemma-geometrically-reduced-etale-local}을
보라. 그러므로 집합 $E$는 $X \to Y$에 대해 취하나, 집합 $E$는
$U \to Y$에 대해 취하나 같다. 이로써 보조정리는 스킴의 경우인
More on Morphisms,
Lemma \ref{more-morphisms-lemma-geometrically-reduced-constructible}에서
따른다.
\end{proof}

\begin{lemma}
\label{spaces-more-morphisms-lemma-proper-flat-over-dvr-reduced-fibre}
$X$를 이산 값매김환 $R$ 위 대수공간이라 하고 구조사상
$X \to \Spec(R)$이 고유하고 평탄하다고 하자. 특수섬유가 축약이면
$X$와 일반섬유 $X_\eta$는 모두 축약이다.
\end{lemma}

\begin{proof}
$U \to X$를 에탈 사상이라 하되 $U$는 아핀 스킴이라 하자.
그러면 $U$는 $R$ 위 유한형이다. $u \in U$를 특수섬유에 속하는
점이라 하자. 국소환 $A = \mathcal{O}_{U, u}$는 $R$ 위 본질적으로
유한형이므로 뇌터이다. $\pi \in R$를 균등화원이라 하자.
$X$가 $R$ 위에서 평탄하므로 $\pi \in \mathfrak m_A$는 $A$ 위
영인자가 아니고, $X$의 특수섬유가 축약이므로 $A/\pi A$는
축약이다. $a \in A$, $a \not = 0$이면 어떤 $n \geq 0$과 원소
$a' \in A$가 존재하여 $a = \pi^n a'$이고 $a' \not \in \pi A$이다.
이는 크룰 교차 정리(Algebra, Lemma
\ref{algebra-lemma-intersect-powers-ideal-module-zero})에서 따른다.
$a$가 멱영이면 $a'$도 멱영인데, 이는 $\pi$가 영인자가 아니기
때문이다.
그러나 $a'$는 축약환 $A/\pi A$의 영이 아닌 원소로 가므로 이는
불가능하다. 따라서 $A$는 축약이다. 그러므로 $u$의 어떤 열린
근방이 축약이고 이 근방은 $U$ 안에 있다.
% P09-ERR-0636: frozen proof explicitly omits the finite-nilradical detail.
(작은 세부사항은 생략한다. $U$가 뇌터임을 사용하라.)
따라서 에탈 사상 $U \to X$를 택하여 $U$가 축약 스킴이고
$X$의 특수섬유의 모든 점이 그 상에 속하게 할 수 있다.
$X$가 $R$ 위에서 고유하므로 $U \to X$는 전사이다. 따라서 $X$는
축약이다. $U \to \Spec(R)$의 일반섬유도 축약이므로(아핀 조각에서는
국소화를 취하여 계산된다), 일반섬유에 대해서도 같은 결론을 얻는다.
\end{proof}

\begin{lemma}
\label{spaces-more-morphisms-lemma-geometrically-reduced-open}
$S$를 스킴이라 하자. $f : X \to Y$를 $S$ 위 대수공간의 사상이라
하자. $f$가 평탄하고 고유하며 유한 표시이면 집합
$$
E = \{y \in |Y| : \text{the fibre of }f : X \to Y\text{ at }y
\text{ is geometrically reduced}\}
$$
은 $|Y|$에서 열린집합이다.
\end{lemma}

\begin{proof}
Lemma \ref{spaces-more-morphisms-lemma-base-change-fibres-geometrically-reduced}에 의해
$E$의 구성은 밑변환과 가환한다. $|Y|$의 부분집합이 열린지
확인할 때 $Y$를 에탈 덮개의 성원들로 바꿀 수 있다. 따라서
$Y$가 아핀이라고 가정할 수 있다. 그러면 $Y$는
$\mathbf{Z}$ 위 유한형 아핀 스킴들의 여과 여극한이다. 따라서
$X \to Y$가 $X_0 \to Y_0$의 밑변환이라고 가정할 수 있다.
여기서 $Y_0$는 유한형 $\mathbf{Z}$-대수의 스펙트럼이고
$X_0 \to Y_0$는 평탄하고 고유하다. Limits of Spaces, Lemmas
\ref{spaces-limits-lemma-descend-finite-presentation},
\ref{spaces-limits-lemma-descend-flat}, 그리고
\ref{spaces-limits-lemma-eventually-proper}를 보라. $E$의 구성은
밑변환과 가환하므로(위 참조), 밑이 뇌터라고 가정할 수 있다.

\medskip\noindent
$Y$가 뇌터라고 가정하자. 집합은 Lemma
\ref{spaces-more-morphisms-lemma-geometrically-reduced-constructible}에 의해
구성가능하다. 따라서 집합이 일반화에 대해 안정함을 보이면 충분하다
(Topology, Lemma \ref{topology-lemma-characterize-closed-Noetherian}).
Properties, Lemma
\ref{properties-lemma-locally-Noetherian-specialization-dvr}에 의해
$Y = \Spec(R)$이고 $R$가 이산 값매김환이며 닫힌 섬유 $X_y$가
기하학적으로 축약인 경우로 환원된다. 보일 것은 일반섬유
$X_\eta$가 기하학적으로 축약이라는 것이다.

\medskip\noindent
그렇지 않으면 어떤 유한 확대 $L$이 존재하며, 이는 $R$의 분수체를
확대하고 $X_L$은 축약이 아니다. Spaces over Fields, Lemmas
\ref{spaces-over-fields-lemma-perfect-reduced}(표수가 영인 경우)와
\ref{spaces-over-fields-lemma-geometrically-reduced-positive-characteristic}
(양의 표수인 경우)를 보라. 이산 값매김환 $R' \subset L$이 존재하며
그 분수체가 $L$이고 $R$를 지배한다. Algebra, Lemma
\ref{algebra-lemma-integral-closure-Dedekind}를 보라.
$R$를 $R'$로 바꾸면 Lemma
\ref{spaces-more-morphisms-lemma-proper-flat-over-dvr-reduced-fibre}로 환원된다.
\end{proof}







\section{섬유의 연결성분}
\label{spaces-more-morphisms-section-connected}

\noindent
이 절은 More on Morphisms, Section
\ref{more-morphisms-section-connected}의 대수공간 판이다.

\begin{lemma}
\label{spaces-more-morphisms-lemma-base-change-fibres-nr-geometrically-connected-components}
$S$를 스킴이라 하자.
$f : X \to Y$를 $S$ 위 대수공간의 사상이라 하자.
$$
n_{X/Y} : |Y| \to \{0, 1, 2, 3, \ldots, \infty\}
$$
를 다음 함수라 하자.
% P09-ERR-0637: frozen source uses y in Y rather than y in |Y|.
$y \in Y$에 대응시키는 값은 $X_k$의 연결성분 개수이다. 여기서
$\Spec(k) \to Y$는 $y$의 동치류에 속하고 $k$는 대수적으로 닫힌
체이다. 이는 잘 정의되며 사상 $g : Y' \to Y$가 주어지면
% P09-ERR-0645: frozen composition uses g instead of the underlying |g|.
$$
n_{X'/Y'} = n_{X/Y} \circ g
$$
이다. 여기서 $X' \to Y'$는 $f$의 밑변환이다.
\end{lemma}

\begin{proof}
% P09-ERR-0638: frozen proof omits bars in both point memberships.
$y' \in Y'$의 상이 $y \in Y$라고 하자. 사상
$\Spec(k') \to Y'$가 $y'$의 동치류에 속하고 $k'$가 대수적으로
닫혔다고 하자. 그러면 가환 도식
$$
\xymatrix{
\Spec(K) \ar[r] \ar[rd] &
\Spec(k') \ar[r] & Y' \ar[d] \\
& \Spec(k) \ar[r] & Y
}
$$
을 택할 수 있다. 여기서 $K$는 대수적으로 닫힌 체이다. 스킴의 사상들
$$
\xymatrix{
X'_{k'} & (X'_{k'})_K = (X_k)_K \ar[l] \ar[r] & X_k
}
$$
은 연결성분들 사이의 전단사를 유도하므로 결과가 따른다.
Spaces over Fields, Lemma
\ref{spaces-over-fields-lemma-separably-closed-field-connected-components}를
보라. 이를 사용하여 함수가 잘 정의됨을 보일 때는 $Y' = Y$로
취한다.
\end{proof}







\section{섬유의 차원}
\label{spaces-more-morphisms-section-dimension}

\noindent
이 절은 More on Morphisms, Section
\ref{more-morphisms-section-dimension}의 대수공간 판이다.

\begin{lemma}
\label{spaces-more-morphisms-lemma-dimension-fibre}
$S$를 스킴이라 하자. $f : X \to Y$를 $S$ 위 대수공간의 유한형
사상이라 하자. $y \in |Y|$라 하자. 다음 양들은 같다.
% P09-ERR-0639: frozen list-introducing clause lacks punctuation.
\begin{enumerate}
\item $d = -\infty$이며, 이는 $y$가 $|f|$의 상에 속하지 않을
때이다. 그렇지 않으면 최소 정수 $d$로서, $f$의 상대 차원이
$\leq d$인 조건이 모든 $x \in |X|$에서 성립하고 이 점은 $y$로 간다.
\item 대수공간 $X_k = \Spec(k) \times_Y X$의 차원. 여기서 임의의
사상 $\Spec(k) \to Y$는 $y$를 정의하는 동치류에 속한다.
\end{enumerate}
\end{lemma}

\begin{proof}
이 명제를 해석하려면 Morphisms of Spaces, Definition
\ref{spaces-morphisms-definition-dimension-fibre},
Properties of Spaces, Definition
\ref{spaces-properties-definition-dimension}, 그리고
Properties of Spaces, Definition
\ref{spaces-properties-definition-dimension-at-point}를 참고해야 한다.
고정된 사상 $\Spec(k) \to Y$에 대해 (1)과 (2)의 수가 같음을
보이겠다. 에탈 사상 $V \to Y$를 택하되 $V$는 아핀 스킴이고,
점 $v \in V$를 택하여 그 상이 $y$가 되게 하자.
$V \times_Y \Spec(k) \to \Spec(k)$는 전사 에탈이므로(Properties of
Spaces, Lemma \ref{spaces-properties-lemma-points-cartesian}),
유한 분리 확대 $k'/k$와(Morphisms, Lemma
\ref{morphisms-lemma-etale-over-field}) 가환 도식
$$
\xymatrix{
\Spec(k') \ar[r] \ar[d] & V \ar[d] \\
\Spec(k) \ar[r] & Y
}
$$
을 찾을 수 있다. $X \to Y$를 $V \times_Y X \to V$로 바꾸고
$X_k$를 $X_{k'} = \Spec(k') \times_V (V \times_Y X)$로 바꾸어도
문제의 차원들은 달라지지 않는다. Properties of Spaces, Lemma
\ref{spaces-properties-lemma-dimension-decent-invariant-under-etale}와
Morphisms of Spaces, Lemma
\ref{spaces-morphisms-lemma-dimension-fibre-after-base-change}에 의한다.
따라서 $Y$가 아핀 스킴이라고 가정할 수 있다. 이 경우
$k = \kappa(y)$라고 가정할 수 있다. 실제로
$X_{\kappa(y)}$와 $X_k$의 차원은 앞서 든 Morphisms of Spaces,
Lemma \ref{spaces-morphisms-lemma-dimension-fibre-after-base-change}와
다음 사실에 의해 같다. 대수공간 $Z$가 체 $K$ 위에 있을 때,
$Z$의 상대 차원을 점 $z \in |Z|$에서 보면 정의에 의해
$\dim_z(Z)$와 같다.
$Y$가 아핀이고 $k = \kappa(y)$라고 가정하자. 그러면
% P09-ERR-0640: frozen English joins two clauses without a connector.
$X$는 준콤팩트이므로 아핀 스킴 $U$와
% P09-ERR-0641: frozen English uses the wrong indefinite article.
전사 에탈 사상 $U \to X$를 택할 수 있다. 그러면
% P09-ERR-0642: frozen maximum omits its indexing domain.
$$
\dim(X_k) = \dim(U_k) = \max \dim_u(U_k)
$$
는 정의에 의해 (1)의 수와 같다.
\end{proof}

\begin{lemma}
\label{spaces-more-morphisms-lemma-base-change-dimension-fibres}
$S$를 스킴이라 하자. $f : X \to Y$를 $S$ 위 대수공간의 유한형
사상이라 하자.
$$
n_{X/Y} : |Y| \to \{-\infty, 0, 1, 2, 3, \ldots\}
$$
를 $y \in |Y|$에 Lemma \ref{spaces-more-morphisms-lemma-dimension-fibre}에서 논한 정수를
대응시키는 함수라 하자. $g : Y' \to Y$가 사상이면
$$
n_{X'/Y'} = n_{X/Y} \circ |g|
$$
이다. 여기서 $X' \to Y'$는 $f$의 밑변환이다.
\end{lemma}

\begin{proof}
이는 Lemma \ref{spaces-more-morphisms-lemma-dimension-fibre}에서 즉시 따른다.
\end{proof}

\begin{lemma}
\label{spaces-more-morphisms-lemma-dimension-fibres-flat}
$S$를 스킴이라 하자.
$f : X \to Y$를 $S$ 위 대수공간의 평탄한 유한 표시 사상이라 하자.
$n_{X/Y}$를 Lemma \ref{spaces-more-morphisms-lemma-base-change-dimension-fibres}에서 도입한,
$Y$ 위에서 $f$의 섬유 차원을 주는 함수라 하자. 그러면
$n_{X/Y}$는 아래쪽 반연속이다.
\end{lemma}

\begin{proof}
전사 에탈 사상 $V \to Y$를 택하되 $V$는 스킴이라 하자.
% P09-ERR-0643: frozen proof uses undefined g and |g| for this morphism.
합성 $n_{X/Y} \circ |g|$가 아래쪽 반연속임을 보일 수 있으면
$|g|$가 열린사상이므로 보조정리가 따른다. 따라서 $Y$가 스킴이라고
가정할 수 있다.
% P09-ERR-0644: frozen proof localizes V rather than the current base Y.
국소적으로 작업하여 $V$가 아핀 스킴이라고 가정할 수 있다.
그러면 아핀 스킴 $U$와 전사 에탈 사상 $U \to X$를 택할 수 있다.
이때 $n_{X/Y} = n_{U/Y}$이다. 따라서 $X$와 $Y$가 모두
스킴이라고 가정할 수 있다. 이 경우 보조정리는
More on Morphisms, Lemma
\ref{more-morphisms-lemma-dimension-fibres-flat}에서 따른다.
\end{proof}

\begin{lemma}
\label{spaces-more-morphisms-lemma-dimension-fibres-proper}
$S$를 스킴이라 하자.
$f : X \to Y$를 $S$ 위 대수공간의 고유 사상이라 하자.
$n_{X/Y}$를 Lemma \ref{spaces-more-morphisms-lemma-base-change-dimension-fibres}에서 도입한,
$Y$ 위에서 $f$의 섬유 차원을 주는 함수라 하자. 그러면
$n_{X/Y}$는 위쪽 반연속이다.
\end{lemma}

\begin{proof}
$$
Z_d = \{x \in |X| :
\text{the fibre of }f\text{ at }x\text{ has dimension }> d\}
$$
로 두자. 그러면 $Z_d$는 Morphisms of Spaces, Lemma
\ref{spaces-morphisms-lemma-openness-bounded-dimension-fibres}에 의해
$|X|$의 닫힌 부분집합이다. $f$가 고유하므로 $f(Z_d)$는 $|Y|$에서
닫혀 있다. $y \in f(Z_d) \Leftrightarrow n_{X/Y}(y) > d$이므로
보조정리가 성립한다.
\end{proof}

\begin{lemma}
\label{spaces-more-morphisms-lemma-dimension-fibres-proper-flat}
$S$를 스킴이라 하자. $f : X \to Y$를 $S$ 위 대수공간의 고유하고
평탄한 유한 표시 사상이라 하자. $n_{X/Y}$를 Lemma
\ref{spaces-more-morphisms-lemma-base-change-dimension-fibres}에서 도입한,
$Y$ 위에서 $f$의 섬유 차원을 주는 함수라 하자. 그러면
$n_{X/Y}$는 국소 상수이다.
\end{lemma}

\begin{proof}
Lemmas \ref{spaces-more-morphisms-lemma-dimension-fibres-flat}과
\ref{spaces-more-morphisms-lemma-dimension-fibres-proper}의 즉각적인 귀결이다.
\end{proof}





\section{카테나리 대수공간}
\label{spaces-more-morphisms-section-catenary}

\noindent
이 절에서는 Decent Spaces, Section
\ref{decent-spaces-section-catenary}에서 시작한 논의를 이어간다.
다음 보조정리의 증명에는 아래 보조정리를 사용한다.

\begin{lemma}
\label{spaces-more-morphisms-lemma-construct-glueing}
$S$를 스킴이라 하자. $f : X \to Y$를 $S$ 위 대수공간의
정수적 사상이라 하자. $y \in |Y|$를 닫힌 몰입
$y : \Spec(k) \to Y$로 나타낼 수 있는 점이라 하자. 그러면 분해
$X \to X' \to Y$, 즉 $f$의 분해로서 다음을 만족하는 것이 존재한다.
\begin{enumerate}
\item $X' \to Y$는 정수적이다.
\item $X \to X'$은 $X' \setminus X'_y$ 위에서 동형이다.
\item $X'_y$에는 유일한 점 $x'$이 있고 $\kappa(x') = k$이다.
\end{enumerate}
더욱이 $f$가 유한이고 $Y$가 국소 뇌터이면 $X' \to Y$는 유한이다.
\end{lemma}

\begin{proof}
Morphisms of Spaces, Lemma \ref{spaces-morphisms-lemma-pushforward}에 의해
층 $f_*\mathcal{O}_X$, $(X_y \to Y)_*\mathcal{O}_{X_y}$,
$y_*\mathcal{O}_{\Spec(k)}$는 준연접 $\mathcal{O}_Y$-대수의 층이다.
사상
$$
% P09-ERR-0647: frozen source uses f_* O_Y instead of f_* O_X.
f_*\mathcal{O}_Y \longrightarrow
(X_y \to Y)_*\mathcal{O}_{X_y} \longleftarrow
y_*\mathcal{O}_{\Spec(k)}
$$
을 생각하자. 그 섬유곱은 준연접 $\mathcal{O}_Y$-대수층
$\mathcal{A}'$이고, $X' \to Y$를 $\mathcal{A}'$의 $Y$ 위
상대 스펙트럼으로 정의할 수 있다. Morphisms, Lemma
\ref{morphisms-lemma-affine-equivalence-algebras}를 보라.
이 구성은 임의의 밑변환과 가환한다. 특히 열린 부분공간
$|Y| \setminus \{y\}$ 위에서 사상 $X \to X'$이 동형이고,
$|Y| \setminus \{y\}$ 위에서 사상 $X' \to Y$가 정수적임은 명백하다.
이제 $y$의 작은 근방에서 명제들을 증명하면 된다.
아핀 스킴 $V = \Spec(R)$와 에탈 사상 $\varphi : V \to Y$로서
$y$가 $\varphi$의 상에 속하는 것을 택하자. 그러면 $V_y$는 $V$의
닫힌 부분스킴이고 $k$ 위 에탈이므로, 유한 개의 점들로 이루어지며
각 점의 잉여체는 $k$ 위에서 분리적이다
(Decent Spaces, Remark \ref{decent-spaces-remark-recall} 참조).
$V$를 줄인 뒤, 유일한 닫힌점 $v = \Spec(l) \to V$가 존재하여
$y$로 가고 $l/k$가 유한 분리라고 가정할 수 있다.
$V \times_Y X = \Spec(C)$로 쓸 수 있고 $R \to C$는 정수적 환 사상이다.
위의 밑변환과의 양립성에 의해
% P09-ERR-0648: frozen source uses undefined U and Y' in this base change.
$U \times_X Y' = \Spec(C')$이며, 여기서
$$
C' = C \times_{C \otimes_R l} l
$$
이다. $R \to l$이 전사이므로 $C \to C \otimes_R l$도 전사이고,
이는 More on Algebra, Situation
\ref{more-algebra-situation-module-over-fibre-product}에서 연구한 종류의
섬유곱임을 알 수 있다
($A, A', B, B'$는 각각 $C \otimes_R l, C, l, C'$에 대응한다).
$C'$은 $R$-부분대수, 즉 $C$의 부분대수이므로 $R$ 위에서 정수적이다.
이로써 (1)을 증명했다. 끝으로 More on Algebra, Lemma
\ref{more-algebra-lemma-points-of-fibre-product}에 의해
% P09-ERR-0649: frozen source repeats the ill-typed V x_X Y'.
$V \times_X Y' = \Spec(C')$에는 유일한 점 $y''$이 있으며,
이는 $v$ 위에 놓이고 잉여체가 $l$이다
(이는 명백한 전사 $C' \to l$에 대응한다).
따라서
% P09-ERR-0650: frozen source uses X_y where the constructed fibre is X'_y.
$X_y \times_{\Spec(k)} \Spec(l)$에는 잉여체가 $l$인 유일한 점이 있다.
$l/k$가 유한 분리이므로, 이는 $X'_y$에 잉여체가 $k$인 유일한 점이
있음을 뜻한다. 즉 (3)이 성립한다.

\medskip\noindent
마지막 명제를 증명하자. $Y$가 국소 뇌터이면 $R$은 뇌터 환이고,
$f$가 유한이면 $R \to C$는 유한임을 관찰하자. 그러면
More on Algebra, Lemma
\ref{more-algebra-lemma-fibre-product-finite-type}에 의해 $C'$은
유한형 $R$-대수이다. 이로써 $X' \to Y$가 유한임을 증명했다.
\end{proof}

\begin{lemma}
\label{spaces-more-morphisms-lemma-universally-catenary-dimension-function}
$S$를 스킴이라 하자. $B$를 $S$ 위의 대수공간이라 하자.
$\delta : |B| \to \mathbf{Z}$를 함수라 하자. $B$가 decent이고
국소 뇌터이며 보편적으로 카테나리이고, $\delta$가 차원함수라고
가정하자. $X$가 $B$ 위의 decent 대수공간이고 그 구조사상
$f : X \to B$가 국소 유한형이면, 다음 규칙으로
$\delta_X : |X| \to \mathbf{Z}$를 정의한다.
$$
% P09-ERR-0651: frozen source runs "degreeof" together.
\delta_X(x) = \delta(f(x)) + \text{초월차수 }x/f(x)
$$
(Morphisms of Spaces, Definition
\ref{spaces-morphisms-definition-dimension-fibre}).
그러면 $\delta_X$는 차원함수이다.
\end{lemma}

\begin{proof}
문제는 $B$ 위에서 국소적이다. 따라서 $B$가 준콤팩트라고 가정할 수 있다.
Decent Spaces, Lemma
\ref{decent-spaces-lemma-locally-Noetherian-decent-quasi-separated}에 의해
$B$는 준분리이다. Limits of Spaces, Proposition
\ref{spaces-limits-proposition-there-is-a-scheme-finite-over}에 의해
% P09-ERR-0652: frozen source targets X although the finite cover must target B.
유한 전사 사상 $\pi : Y \to X$를 택할 수 있고, $Y$는 스킴이다.
주장: $\delta_Y$는 차원함수이다.

\medskip\noindent
이 주장은 보조정리를 함의한다. 보조정리에서와 같은 $X \to B$에 대해
$Z = Y \times_B X$로 놓고 사영을 $p : Z \to Y$와
$q : Z \to X$로 쓰자. 그러면
$$
% P09-ERR-0651: second frozen "degreeof" occurrence.
\delta_Z(z) = \delta_Y(p(z)) + \text{초월차수 }z/p(z)
$$
이고 $\delta_Z(z) = \delta_X(q(z))$이다. 이는 Morphisms of Spaces,
Lemma \ref{spaces-morphisms-lemma-dimension-fibre-at-a-point-additive}와
유한 사상에 대해서는 이 초월차수들이 영이라는 사실에서 따른다.
Decent Spaces, Lemma
\ref{decent-spaces-lemma-scheme-with-dimension-function} 및 위 주장에 의해
$\delta_Z$가 차원함수임을 얻는다. 이어서 Decent Spaces, Lemma
\ref{decent-spaces-lemma-check-dimension-function-finite-cover}에 의해
$\delta_X$가 차원함수임을 얻는다.

\medskip\noindent
주장을 증명하자. 특수화 $y \leadsto y'$, $y \not = y'$을 생각하자.
이는 뇌터 스킴 $Y$의 점들 사이의 특수화이다. 그러면
$\delta_Y(y) > \delta_Y(y')$이다. 실제로 $Y$의 섬유 안에 있는
점들 사이에는 특수화가 없다
(Decent Spaces, Lemma
\ref{decent-spaces-lemma-conditions-on-fibre-and-qf} 참조). 이를 특수화의 사슬에 적용하면
$$
\delta_Y(y) - \delta_Y(y') \geq
\text{여차원}(\overline{\{y'\}}, \overline{\{y\}})
$$
을 얻는다. 등호를 보이는 것이 목표이다. Properties, Lemma
\ref{properties-lemma-locally-Noetherian-closed-point}에 의해 특수화
$y' \leadsto y_0$를 택할 수 있다.
$$
\delta_Y(y) - \delta_Y(y_0) =
\text{여차원}(\overline{\{y_0\}}, \overline{\{y\}})
$$
을 보이면 충분하다. 실제로 이 등식은 $y \leadsto y'$ 및
$y' \leadsto y_0$ 모두에 대한 등식을 함의한다.

\medskip\noindent
$y = y_c \leadsto y_{c - 1} \leadsto \ldots \leadsto y_0$를
$Y$ 안의 특수화의 극대사슬로 택하자.
$b = \pi(y)$ 및 $b_0 = \pi(y_0)$로 놓는다. 특수화의 극대사슬
$b = b_e \leadsto b_{e - 1} \leadsto \ldots \leadsto b_0$를
$|B|$ 안에서 택하자.
$e = c$임을 보여야 한다. $\pi$가 닫힌사상이므로
(Morphisms of Spaces, Lemma \ref{spaces-morphisms-lemma-finite-proper}),
특수화의 열
$y = y'_e \leadsto y'_{e - 1} \leadsto \ldots \leadsto y'_0$로서
$b = b_e \leadsto b_{e - 1} \leadsto \ldots \leadsto b_0$로 가는 것을
찾을 수 있다. $y'_e \leadsto y'_{e - 1} \leadsto \ldots \leadsto y'_0$
또한 극대사슬임을 관찰하자. $y_0 = y'_0$이면 $Y$가 카테나리이므로
원하는 대로 $e = c$를 얻는다. 다음 단락에서는 약간의 기교로 이 경우로
귀착시키고 같은 방식으로 결론을 얻는다.

\medskip\noindent
$\pi$가 닫힌사상이므로 $b_0$는 $|B|$의 닫힌점이다.
Decent Spaces, Lemma \ref{decent-spaces-lemma-decent-space-closed-point}에 의해
$b_0$를 닫힌 몰입 $b_0 : \Spec(k) \to B$로 나타낼 수 있다.
Lemma \ref{spaces-more-morphisms-lemma-construct-glueing}에 의해 다음 분해를 얻을 수 있다.
$$
% P09-ERR-0653: frozen source factors through X rather than B.
Y \to Y' \to X
$$
여기서 $\pi' : Y' \to X$는 유한이고, $Y \to Y'$는 $y_0$와
$y'_0$를 같은 점으로 보내며 $b_0$의 역상 밖에서는 동형인 사상이다.
% P09-ERR-0654: frozen source has singular/plural verb disagreement.
(물론 $Y'$은 스킴이 아니지만, 이는 이어지는 논증에 문제가 되지 않는다.)
특수화의 두 극대사슬
$y_c \leadsto y_{c - 1} \leadsto \ldots \leadsto y_0$와
$y'_e \leadsto y'_{e - 1} \leadsto \ldots \leadsto y'_0$는
$Y'$ 안에서 같은 시작점과 끝점을 갖는 특수화의 극대사슬들로 감을 알 수 있다.
$B$가 보편적으로 카테나리이므로 $|Y'|$도 카테나리이고,
앞과 같이 결론을 얻는다.
\end{proof}





\section{사상의 에탈 국소화}
\label{spaces-more-morphisms-section-etale-localization}

\noindent
이 절은 More on Morphisms, Section
\ref{more-morphisms-section-etale-localization}의 대응물이다.

\begin{lemma}
\label{spaces-more-morphisms-lemma-etale-splits-off-quasi-finite-part}
$S$를 스킴이라 하자. $f : X \to Y$를 $S$ 위 대수공간의 사상이라 하자.
$y \in |Y|$라 하자. $x_1, \ldots, x_n \in |X|$를 택하자.
이 점들은 $y$로 간다.
다음을 가정하자.
\begin{enumerate}
\item $f$는 국소 유한형이다.
\item $f$는 분리이다.
\item $f$는 $x_1, \ldots, x_n$에서 준유한이다.
\item $f$는 준콤팩트이거나 $Y$가 decent이다.
\end{enumerate}
그러면 점이 표시된 대수공간의 에탈 사상
$(U, u) \to (Y, y)$와 열린닫힌 부분공간들로의 분해
$$
U \times_Y X = W \amalg V
$$
가 존재하여, 사상 $V \to U$가 유한이고,
$|V| \to |U|$의 $u$ 위 섬유에 있는 모든 점은 어떤 $x_i$로 가며,
$|W| \to |U|$의 $u$ 위 섬유에는 어떤 $x_i$로 가는 점도 없다.
\end{lemma}

\begin{proof}
점이 표시된 대수공간의 에탈 사상 $(U, u) \to (Y, y)$를 택하고,
$w \in |U \times_Y X|$들 중 $u \in |U|$로 가면서
어떤 $x_i \in |X|$로 가는 것들의 집합을 생각하자. Decent Spaces, Lemma
\ref{decent-spaces-lemma-qf-and-qc-finite-fibre}
($f$가 유한형인 경우) 또는 Decent Spaces, Lemma
\ref{decent-spaces-lemma-decent-finite-fibre}
($Y$가 decent인 경우)에 의해 이 집합은 유한이다.
따라서 $f$를 밑변환 $U \times_Y X \to U$로 바꾸고,
$x_1, \ldots, x_n$을 이 $w$들의 집합으로 바꿀 수 있다.
특히 $Y$가 아핀 스킴이라고 가정할 수 있고 그렇게 하며,
그러면 $X$는 분리 대수공간이다.

\medskip\noindent
아핀 스킴 $Z$와 에탈 사상 $Z \to X$를 택하여
$x_1, \ldots, x_n$이 $|Z| \to |X|$의 상에 속하게 하자.
$|Z| \to |X|$의 섬유들은 유한하다. Properties of Spaces, Lemma
\ref{spaces-properties-lemma-finite-fibres-presentation}
(또는 Decent Spaces, Section
\ref{decent-spaces-section-reasonable-decent}의 더 일반적인 논의)을 보라.
$\{z_1, \ldots, z_m\} \subset |Z|$를
$\{x_1, \ldots, x_n\}$의 역상이라 하자. More on Morphisms, Lemma
\ref{more-morphisms-lemma-etale-splits-off-quasi-finite-part-technical}에 의해
에탈 사상 $(U, u) \to (Y, y)$가 존재하여
$U \times_Y Z = Z_1 \amalg Z_2$이고 $Z_1 \to U$는 유한이며,
% P09-ERR-0655: frozen source indexes this fibre by y instead of u.
$(Z_1)_y = \{z_1, \ldots, z_m\}$이다. $U$가 아핀이라고 가정할 수 있고,
따라서 $Z_1$도 아핀이다.

\medskip\noindent
$f$가 분리이므로 $V$, 즉 $Z_1 \to X$의 상은 열린 동시에 닫혀 있다
% P09-ERR-0656: frozen source uses X instead of the base-changed U x_Y X here and below.
(Morphisms of Spaces, Lemma
\ref{spaces-morphisms-lemma-universally-closed-permanence}).
$W = X \setminus V$로 놓으면 보조정리와 같은 분해를 얻는다.
증명을 끝내려면 $V \to U$가 유한임을 보여야 한다.
$Z_1 \to V$가 전사 에탈이므로 $V$는 $Z_1$의 몫이고,
그 에탈 동치관계는 $R = Z_1 \times_V Z_1$이다. Spaces, Lemma
\ref{spaces-lemma-space-presentation}를 보라.
$f$가 분리이므로 $V \to U$는 분리이고 $R$은
$Z_1 \times_U Z_1$ 안에서 닫혀 있다. $Z_1 \to U$가 유한이므로
사영들 $s, t : R \to Z_1$은 유한이다. 따라서 Groupoids, Proposition
\ref{groupoids-proposition-finite-flat-equivalence}에 의해 $V$는 아핀 스킴이다.
Morphisms, Lemma \ref{morphisms-lemma-image-proper-is-proper}에 의해
$V \to U$가 고유임을 얻고, Morphisms, Lemma
\ref{morphisms-lemma-finite-proper}에 의해 $V \to U$가 유한임을 얻는다.
이로써 증명이 끝난다.
\end{proof}

\begin{lemma}
\label{spaces-more-morphisms-lemma-etale-splits-off-just-one-quasi-finite-part}
$S$를 스킴이라 하자. $f : X \to Y$를 $S$ 위 대수공간의 사상이라 하자.
$x \in |X|$의 상을 $y \in |Y|$라 하자. 다음을 가정하자.
\begin{enumerate}
\item $f$는 국소 유한형이다.
\item $f$는 분리이다.
\item $f$는 $x$에서 준유한이다.
\end{enumerate}
그러면 점이 표시된 대수공간의 에탈 사상
$(U, u) \to (Y, y)$와 열린닫힌 부분공간들로의 분해
$$
U \times_Y X = W \amalg V
$$
가 존재하여, 사상 $V \to U$가 유한이고, 점 $v \in |V|$가 존재한다.
이 점은 $x$로 가는데 이는 $|X|$에서이며,
$u$로 가는데 이는 $|U|$에서이다.
\end{lemma}

\begin{proof}
스킴 $U$, 점 $u \in U$, 그리고 에탈 사상 $U \to Y$를 택하자.
이 사상은 $u$를 $y$로 보낸다. 점 $x' \in |U \times_Y X|$이 존재하며,
이 점은 $x$로 가는데 이는 $|X|$에서이고, $u$로 가는데 이는 $|U|$에서이다
(Properties of Spaces,
Lemma \ref{spaces-properties-lemma-points-cartesian}).
증명을 끝내기 위해 Lemma \ref{spaces-more-morphisms-lemma-etale-splits-off-quasi-finite-part}를
사상 $U \times_Y X \to U$와 점 $x'$에 적용하자.
$U$가 스킴이고 따라서 $u$가 단사상 $\Spec(\kappa(u)) \to U$에서
오기 때문에 이를 적용할 수 있다.
\end{proof}










\section{자리스키의 주정리}
\label{spaces-more-morphisms-section-structure-quasi-finite}

\noindent
이 절에서는 앞 절의 결과들을 적용하여 대수공간의 사상에 대한
자리스키의 주정리를 증명한다. 이 절은 More on Morphisms, Section
\ref{more-morphisms-section-application-etale-neighbourhoods}의 대응물이다.

\begin{lemma}
\label{spaces-more-morphisms-lemma-finite-type-separated}
$S$를 스킴이라 하자. $f : X \to Y$를 $S$ 위 대수공간의
유한형 분리 사상이라 하자. $Y'$을 $Y$의 $X$ 안에서의 정규화라 하자.
그림은 다음과 같다.
$$
\xymatrix{
X \ar[rd]_f \ar[rr]_{f'} & & Y' \ar[ld]^\nu \\
& Y &
}
$$
그러면 열린 부분공간 $U' \subset Y'$가 존재하여 다음이 성립한다.
\begin{enumerate}
\item $(f')^{-1}(U') \to U'$는 동형이다.
\item $(f')^{-1}(U') \subset X$는 $f$가 준유한인 점들의 집합이다.
\end{enumerate}
\end{lemma}

\begin{proof}
Morphisms of Spaces, Lemma
\ref{spaces-morphisms-lemma-locally-finite-type-quasi-finite-part}에 의해,
열린 부분공간 $U \subset X$가 있고, 이는 $|X|$에서 $f$가 준유한인
점들에 대응한다.
다음을 증명해야 한다.
\begin{enumerate}
\item[(a)] $|U| \to |Y'|$의 상은 $|U'|$이다. 여기서 $U'$은
$Y'$의 어떤 열린 부분공간이다.
\item[(b)] $U = f^{-1}(U')$이다. % P09-ERR-0657: frozen source uses f rather than f'.
\item[(c)] $U \to U'$는 동형이다.
\end{enumerate}
$U$의 형성은 임의의 밑변환과 가환하고
(Morphisms of Spaces, Lemma
\ref{spaces-morphisms-lemma-locally-finite-type-quasi-finite-part}),
정규화 $Y'$의 형성은 매끄러운 밑변환과 가환하며
(Lemma \ref{spaces-more-morphisms-lemma-normalization-smooth-localization}),
에탈 사상은 열린사상이고, “동형인 것”은 밑 위 fpqc 국소적이므로
(Descent on Spaces, Lemma
\ref{spaces-descent-lemma-descending-property-isomorphism}),
(a), (b), (c)를 $Y$ 위에서 에탈 국소적으로 증명하면 충분하다
(일부 세부사항은 생략한다). % P09-ERR-0658: frozen source omits the descent details.
따라서 $Y$가 아핀 스킴이라고 가정할 수 있다. 이는 $Y'$도
(아핀) 스킴임을 함의한다.

\medskip\noindent
$x \in |U|$라 하자. 주장: 열린 근방
$f'(x) \in V \subset Y'$로서 $(f')^{-1}V \to V$가 동형인 것이 존재한다.
먼저 이 주장이 보조정리를 함의함을 증명하자. 실제로
$(f')^{-1}V \cong V$는 스킴이고 ($Y'$의 열린 부분으로서),
$Y$ 위에서 국소 유한형이며 ($X$의 열린 부분공간으로서),
$v \in V$에 대해 잉여체 확대 $\kappa(v)/\kappa(\nu(v))$는 대수적이다
($V \subset Y'$이고 $Y'$은 $Y$ 위에서 정수적이므로).
따라서 $V \to Y$의 섬유들은 이산적이고 (Morphisms, Lemma
\ref{morphisms-lemma-algebraic-residue-field-extension-closed-point-fibre}),
$(f')^{-1}V \to Y$는 국소 준유한이다
(Morphisms, Lemma \ref{morphisms-lemma-locally-quasi-finite-fibres}).
이는 $(f')^{-1}V \subset U$ 및
% P09-ERR-0659: frozen source uses U' before constructing it.
$V \subset U'$을 함의한다. $x$가 임의였으므로 (a), (b), (c)가 성립한다.

\medskip\noindent
$y = f(x) \in |Y|$로 놓자. 점이 표시된 스킴의 에탈 사상
$(T, t) \to (Y, y)$를 택하자. 아래첨자 ${}_T$로 $T$로의 밑변환을 나타낸다.
% P09-ERR-0660: frozen source writes a topological point as z in X_T.
$z \in X_T$를 섬유 $X_t$의 점이라 하자. 이 점은 $x$ 위에 놓인다.
$U_T \subset X_T$는 $f_T$가 준유한인 점들의 집합임을 관찰하자.
Morphisms of Spaces, Lemma
\ref{spaces-morphisms-lemma-locally-finite-type-quasi-finite-part}를 보라.
또한
$$
X_T \xrightarrow{f'_T} Y'_T \xrightarrow{\nu_T} T
$$
는 $T$의 $X_T$ 안에서의 정규화이다. Lemma
\ref{spaces-more-morphisms-lemma-normalization-smooth-localization}을 보라.
% P09-ERR-0660: frozen source writes a topological point as z in U_T.
주장이 $z \in U_T \subset X_T \to Y'_T \to T$에 대해 성립한다고 하자.
즉, 열린 근방 $f'_T(z) \in V' \subset Y'_T$로서
$(f'_T)^{-1}V' \to V'$가 동형인 것을 찾을 수 있다고 하자.
사상 $Y'_T \to Y'$은 에탈이므로 $V \subset Y'$, 즉 $V'$의 상은 열려 있다.
$f'(x) \in V$임을 관찰하자. 이는 $f'_T(z) \in V'$이기 때문이다. 또한
$$
\xymatrix{
(f'_T)^{-1}V' \ar[r] \ar[d] & (f')^{-1}(V) \ar[d] \\
V' \ar[r] & V
}
$$
는 섬유곱 정사각형이다 ($Y'_T \times_{Y'} X = X_T$이므로).
왼쪽 수직 화살표가 동형이고 $\{V' \to V\}$가 에탈 덮개이므로,
Descent on Spaces, Lemma
\ref{spaces-descent-lemma-descending-property-isomorphism}에 의해
오른쪽 수직 화살표가 동형이라고 결론내린다.
달리 말하면 주장은
% P09-ERR-0660: frozen source writes a topological point as x in U.
$x \in U \subset X \to Y' \to Y$에 대해 성립한다.

\medskip\noindent
앞 단락의 결과에 의해 $x \in |U|$에 대한 주장을 증명하기 위해
$Y$를 에탈 근방 $T$, 즉 $y = f(x)$의 근방으로 바꾸고, $x$를
% P09-ERR-0660: frozen source omits the underlying-space bars for the replacement point.
$x$ 위에 놓이는 임의의 점으로 바꿀 수 있다. 이 점은 $T \times_Y X$ 안에 있다.
따라서 열린닫힌 부분공간들로의 분해
$$
X = V \amalg W
$$
가 존재하여 $V \to Y$가 유한이고 $x \in V$라고 가정할 수 있다.
Lemma \ref{spaces-more-morphisms-lemma-etale-splits-off-quasi-finite-part}를 보라.
$X$가 $V$와 $W$의 서로소 합이고 이는 $Y$ 위의 합이며,
$V \to Y$가 유한이므로, $Y$의 $X$ 안에서의 정규화는 다음 사상임을 알 수 있다.
$$
% P09-ERR-0661: frozen factorization ends in S rather than Y.
X = V \amalg W \longrightarrow V \amalg W' \longrightarrow S
$$
여기서 $W'$은 $Y$의 $W$ 안에서의 정규화이다. Morphisms of Spaces,
Lemmas \ref{spaces-morphisms-lemma-normalization-in-disjoint-union},
\ref{spaces-morphisms-lemma-finite-integral},
\ref{spaces-morphisms-lemma-normalization-in-integral}을 보라.
주장이 따르고 증명이 끝난다.
\end{proof}

\noindent
다음 보조정리는 Morphisms of Spaces, Lemma
\ref{spaces-morphisms-lemma-quasi-finite-separated-quasi-affine}와 중복된다.
같은 보조정리를 두 번 싣는 이유는 증명들이 다소 다르기 때문이다.
아래 증명은 위에서 제시한 대수공간의 비표현가능 사상에 대한
자리스키 주정리에 의존하는 반면, Morphisms of Spaces, Lemma
\ref{spaces-morphisms-lemma-quasi-finite-separated-quasi-affine}의 증명은
Morphisms of Spaces, Proposition
\ref{spaces-morphisms-proposition-locally-quasi-finite-separated-over-scheme}을
사용하여 스킴 사상의 경우로 귀착한다.

\begin{lemma}
\label{spaces-more-morphisms-lemma-quasi-finite-separated-quasi-affine}
$S$를 스킴이라 하자. $f : X \to Y$를 $S$ 위 대수공간의 사상이라 하자.
$f$가 준유한이고 분리라고 가정하자. $Y'$을 $Y$의 $X$ 안에서의
정규화라 하자. 그림은 다음과 같다.
$$
\xymatrix{
X \ar[rd]_f \ar[rr]_{f'} & & Y' \ar[ld]^\nu \\
& Y &
}
$$
그러면 $f'$은 준콤팩트 열린 몰입이고 $\nu$는 정수적이다.
특히 $f$는 준아핀이다.
\end{lemma}

\begin{proof}
Lemma \ref{spaces-more-morphisms-lemma-finite-type-separated}에서 따른다. 실제로 그 보조정리에 의해
열린 부분공간 $U' \subset Y'$가 존재하여
$(f')^{-1}(U') = X$이고 (!) $X \to U'$는 동형이다!
달리 말하면 $f'$은 열린 몰입이다. $f'$이 준콤팩트임을 관찰하자.
실제로 $f$는 준콤팩트이고 $\nu : Y' \to Y$는 분리이다
(Morphisms of Spaces, Lemma
\ref{spaces-morphisms-lemma-quasi-compact-permanence}).
따라서 모든 아핀 스킴 $Z$와 사상 $Z \to Y$에 대해 섬유곱
$Z \times_Y X$는 아핀 스킴 $Z \times_Y Y'$의 준콤팩트 열린 부분스킴이다.
그러므로 정의에 의해 $f$는 준아핀이다.
\end{proof}

\begin{lemma}[자리스키의 주정리]
\label{spaces-more-morphisms-lemma-quasi-finite-separated-pass-through-finite}
$S$를 스킴이라 하자. $f : X \to Y$를 $S$ 위 대수공간의 사상이라 하자.
$f$가 준유한이고 분리이며, $Y$가 준콤팩트이고 준분리라고 가정하자.
그러면 다음 분해가 존재한다.
$$
\xymatrix{
X \ar[rd]_f \ar[rr]_j & & T \ar[ld]^\pi \\
& Y &
}
$$
여기서 $j$는 준콤팩트 열린 몰입이고 $\pi$는 유한이다.
\end{lemma}

\begin{proof}
$X \to Y' \to Y$를 Lemma
\ref{spaces-more-morphisms-lemma-quasi-finite-separated-quasi-affine}의 결론에 나오는 것으로 잡자.
Limits of Spaces, Lemma
\ref{spaces-limits-lemma-integral-algebra-directed-colimit-finite}에 의해
$$
\nu_*\mathcal{O}_{Y'} = \colim_{i \in I} \mathcal{A}_i
$$
를 유한 준연접
% P09-ERR-0662: frozen source calls these O_X-algebras although they live on Y.
$\mathcal{O}_X$-대수 $\mathcal{A}_i \subset \nu_*\mathcal{O}_{Y'}$들의
유향 여극한으로 쓸 수 있다. 그러면
$\pi_i : T_i = \underline{\Spec}_Y(\mathcal{A}_i) \to Y$는 각 $i$에 대해
유한 사상이다.
전이사상 $T_{i'} \to T_i$가 아핀이고 $Y' = \lim T_i$임을 관찰하자.

\medskip\noindent
Limits of Spaces, Lemma \ref{spaces-limits-lemma-descend-opens}에 의해
어떤 $i$와 준콤팩트 열린 부분공간 $U_i \subset T_i$가 존재하여,
$Y'$ 안에서 그 역상이 $f'(X)$와 같다. $i' \geq i$에 대해
$U_{i'}$를 $U_i$의 $T_{i'}$ 안에서의 역상이라 하자. 그러면
$X \cong f'(X) = \lim_{i' \geq i} U_{i'}$이다. Limits of Spaces, Lemma
\ref{spaces-limits-lemma-directed-inverse-system-has-limit}를 보라.
Limits of Spaces, Lemma
\ref{spaces-limits-lemma-finite-type-eventually-closed}에 의해
$X \to U_{i'}$가 어떤 $i' \geq i$에 대해 닫힌 몰입임을 알 수 있다.
(사실 $X \cong U_{i'}$는 충분히 큰 $i'$에 대해 성립하지만, 이것은 필요 없다.)
따라서 $X \to T_{i'}$는 몰입이다. Morphisms of Spaces, Lemma
\ref{spaces-morphisms-lemma-factor-the-other-way}에 의해 이를
$X \to T \to T_{i'}$로 분해할 수 있다. 첫째 화살표는 열린 몰입이고
둘째 화살표는 닫힌 몰입이다. 이로써 증명이 끝난다.
\end{proof}

\begin{lemma}
\label{spaces-more-morphisms-lemma-quasi-finite-separated-pass-through-finite-addendum}
Lemma \ref{spaces-more-morphisms-lemma-quasi-finite-separated-pass-through-finite}와 같은 표기와
가정을 사용하자. 더욱이 $f$가 국소 유한 표시라고 가정하자.
그러면 $T$가 $Y$ 위에서 유한이고 유한 표시가 되도록 분해를 택할 수 있다.
\end{lemma}

\begin{proof}
Limits of Spaces, Lemma
\ref{spaces-limits-lemma-finite-in-finite-and-finite-presentation}에 의해
$T = \lim T_i$로 쓸 수 있다. 여기서 모든 $T_i$는 $Y$ 위에서 유한이고
유한 표시이며, 전이사상 $T_{i'} \to T_i$는 닫힌 몰입이다.
Limits of Spaces, Lemma \ref{spaces-limits-lemma-descend-opens}에 의해
어떤 $i$와 열린 부분스킴
% P09-ERR-0663: frozen source says open subscheme although T_i is an algebraic space.
$U_i \subset T_i$가 존재하여, $T$ 안에서 그 역상은 $X$이다.
Limits of Spaces, Lemma
\ref{spaces-limits-lemma-finite-type-eventually-closed}에 의해
$X \cong U_i$가 충분히 큰 $i$에 대해 성립함을 알 수 있다.
$T$를 $T_i$로 바꾸면 증명이 끝난다.
\end{proof}












\section{자리스키 주정리의 응용, I}
\label{spaces-more-morphisms-section-applications-zmt}

\noindent
첫째 응용은 유한 사상을 섬유가 유한인 고유 사상으로 특징짓는 것이다.

\begin{lemma}
\label{spaces-more-morphisms-lemma-characterize-finite}
$S$를 스킴이라 하자. $f : X \to Y$를 $S$ 위 대수공간의 사상이라 하자.
다음은 동치이다.
\begin{enumerate}
\item $f$는 유한이다.
\item $f$는 고유이고 국소 준유한이다.
\item $f$는 고유이고, $|X_k|$는 모든 사상 $\Spec(k) \to Y$에 대해
이산공간이다. 여기서 $k$는 체이다.
\item $f$는 보편 닫힌사상이고 분리이며 국소 유한형이고,
$|X_k|$는 모든 사상 $\Spec(k) \to Y$에 대해 이산공간이다.
여기서 $k$는 체이다.
\end{enumerate}
\end{lemma}

\begin{proof}
Morphisms of Spaces, Lemmas \ref{spaces-morphisms-lemma-finite-proper},
\ref{spaces-morphisms-lemma-finite-quasi-finite}에 의해
(1) $\Rightarrow$ (2)이다. Morphisms of Spaces, Lemma
\ref{spaces-morphisms-lemma-locally-quasi-finite}에 의해
(2) $\Rightarrow$ (3)이다. 정의에 의해 (3)은 (4)를 함의한다.

\medskip\noindent
(4)를 가정하자. $f$가 보편 닫힌사상이므로 준콤팩트이다
(Morphisms of Spaces, Lemma
\ref{spaces-morphisms-lemma-universally-closed-quasi-compact}).
점 $y$를 $|Y|$에서 택하자. $y$를 사상 $\Spec(k) \to Y$로 나타내자.
$|X_k|$는 준콤팩트 이산공간이므로 유한 이산공간이다.
사상 $|X_k| \to |X|$는 $|X| \to |Y|$의 $y$ 위 섬유로 전사한다
(Properties of Spaces, Lemma \ref{spaces-properties-lemma-points-cartesian}).
Morphisms of Spaces, Lemma
\ref{spaces-morphisms-lemma-quasi-finite-at-point}에 의해
$X \to Y$는 $|X| \to |Y|$의 $y$ 위 섬유의 모든 점에서 준유한이다.
점이 표시된 기본 에탈 근방 $(U, u) \to (Y, y)$와 분해
$X_U = V \amalg W$를 Lemma
\ref{spaces-more-morphisms-lemma-etale-splits-off-quasi-finite-part}에서와 같이 택하되,
$|X|$에서 $y$ 위에 놓이는 모든 점에 맞추어 택하자.
$W_u = \emptyset$임을 관찰하자. 실제로 $|X| \to |Y|$의
$y$ 위 섬유에 있는 모든 점을 사용했다. $f$가 보편 닫힌사상이므로
$|W|$의 $|U|$에서의 상은 $u$를 포함하지 않는 닫힌집합이다.
$U$를 줄인 뒤 $W = \emptyset$이라고 가정할 수 있다.
달리 말하면 $X_U = V$는 $U$ 위에서 유한이다.
$y \in |Y|$가 임의였으므로, 에탈 사상의 족 $\{U_i \to Y\}$가 존재하고
그 상들은 $Y$를 덮으며 밑변환 $X_{U_i} \to U_i$는 유한이다.
Morphisms of Spaces, Lemma \ref{spaces-morphisms-lemma-integral-local}에 의해
$f$가 유한이라고 결론내린다.
\end{proof}

\noindent
그 결과 다음과 같은 유용한 명제를 얻는다.

\begin{lemma}
\label{spaces-more-morphisms-lemma-proper-finite-fibre-finite-in-neighbourhood}
$S$를 스킴이라 하자. $f : X \to Y$를 $S$ 위 대수공간의 사상이라 하자.
$y \in |Y|$라 하자. 다음을 가정하자.
\begin{enumerate}
\item $f$는 고유이다.
\item $f$는 모든 $x \in |X|$에서 준유한이고, 이 점들은 $y$ 위에 놓인다
(Decent Spaces, Lemma \ref{decent-spaces-lemma-conditions-on-fibre-and-qf}).
\end{enumerate}
그러면 열린집합 $V \subset Y$가 존재하여 $y$의 근방이고,
$f|_{f^{-1}(V)} : f^{-1}(V) \to V$가 유한이다.
\end{lemma}

\begin{proof}
Morphisms of Spaces, Lemma
\ref{spaces-morphisms-lemma-locally-finite-type-quasi-finite-part}에 의해
$f$가 준유한인 점들의 집합은 열린집합 $U \subset X$이다.
$Z = X \setminus U$로 놓자. 그러면 $y \not \in f(Z)$이다.
$f$가 고유이므로 $f(Z) \subset Y$는 닫힌집합이다.
임의의 열린집합 $V \subset Y$로서 $y$의 근방이고
% P09-ERR-0664: frozen source intersects subspaces of different ambient spaces.
$Z \cap V = \emptyset$인 것을 택하자. 그러면
$f^{-1}(V) \to V$는 국소 준유한이고 고유이다.
따라서 Lemma \ref{spaces-more-morphisms-lemma-characterize-finite}에 의해
$f^{-1}(V) \to V$는 유한이다.
\end{proof}

\begin{lemma}
\label{spaces-more-morphisms-lemma-flat-proper-family-cannot-collapse-fibre}
\begin{slogan}
고유 족의 한 섬유를 붕괴시키면 근방의 섬유들도 붕괴한다.
\end{slogan}
$S$를 스킴이라 하자. 다음 가환도를 생각하자.
$$
\xymatrix{
X \ar[rr]_h \ar[rd]_f & & Y \ar[ld]^g \\
& B
}
$$
이는 $S$ 위 대수공간의 사상들로 이루어진다.
% P09-ERR-0665: frozen source uses singular “morphism” for this diagram.
% P09-ERR-0666: frozen source writes a topological point as b in B.
$b \in B$라 하고, 사상 $\Spec(k) \to B$를 택하자. 이 사상은
$b$의 동치류에 속한다.
다음을 가정하자.
\begin{enumerate}
\item $X \to B$는 고유 사상이다.
\item $Y \to B$는 분리이고 국소 유한형이다.
\item 다음 중 하나가 성립한다.
\begin{enumerate}
\item $|X_k| \to |Y_k|$의 상은 유한이다.
\item $|f|^{-1}(\{b\})$의 $|Y|$ 안에서의 상은 유한이고 $B$는 decent이다.
\end{enumerate}
\end{enumerate}
그러면 열린 부분공간 $B' \subset B$가 존재하여 $b$를 포함하고,
$X_{B'} \to Y_{B'}$는 닫힌 부분공간 $Z \subset Y_{B'}$를 통해
분해되며 이 부분공간은 $B'$ 위에서 유한이다.
\end{lemma}

\begin{proof}
$Z \subset Y$를 $h$의 스킴 이론적 상이라 하자. Morphisms of Spaces,
Section \ref{spaces-morphisms-section-scheme-theoretic-image}를 보라.
Morphisms of Spaces, Lemma
\ref{spaces-morphisms-lemma-scheme-theoretic-image-is-proper}에 의해
$X \to Z$는 전사이고 $Z \to B$는 고유이다. 따라서
$$
\{x \in |X|\text{ 중 밑의 점이 }b\} \to
\{z \in |Z|\text{ 중 밑의 점이 }b\}
$$
및 $|X_k| \to |Z_k|$는 전사이다. (3)(a) 또는 (3)(b) 어느 쪽이든
% P09-ERR-0667: frozen source omits “at” before the fibre points.
$Z \to B$가 $|Z|$의 모든 점에서 준유한임을 함의하며, 이 점들은
$b$ 위에 놓인다.
이는 Decent Spaces, Lemma
\ref{decent-spaces-lemma-conditions-on-fibre-and-qf}에 의한 것이다.
따라서 Lemma \ref{spaces-more-morphisms-lemma-proper-finite-fibre-finite-in-neighbourhood}에 의해
$Z \to B$는 $b$의 어떤 열린 근방에서 유한이다.
\end{proof}








\section{슈타인 분해}
\label{spaces-more-morphisms-section-stein-factorization}

\noindent
슈타인 분해는 고유 사상 $f : X \to S$가
$f_*\mathcal{O}_X = \mathcal{O}_S$를 만족하면 그 섬유들이 연결이라는 명제이다.

\begin{lemma}
\label{spaces-more-morphisms-lemma-stein-universally-closed}
$S$를 스킴이라 하자. $f : X \to Y$를 $S$ 위 대수공간의
보편 닫힌 준분리 사상이라 하자. 다음 분해가 존재한다.
$$
\xymatrix{
X \ar[rr]_{f'} \ar[rd]_f & & Y' \ar[dl]^\pi \\
& Y &
}
$$
이 분해는 다음 성질들을 갖는다.
\begin{enumerate}
\item 사상 $f'$은 보편 닫힌사상이고, 준콤팩트이고, 준분리이고, 전사이다.
\item 사상 $\pi : Y' \to Y$는 정수적이다.
\item $f'_*\mathcal{O}_X = \mathcal{O}_{Y'}$이다.
\item $Y' = \underline{\Spec}_Y(f_*\mathcal{O}_X)$이다.
\item $Y'$은 Morphisms of Spaces, Definition
\ref{spaces-morphisms-definition-normalization-X-in-Y}에서 정의한
$Y$의 $X$ 안에서의 정규화이다.
\end{enumerate}
분해 $f = \pi \circ f'$의 형성은 평탄한 밑변환과 가환한다.
\end{lemma}

\begin{proof}
Morphisms of Spaces, Lemma
\ref{spaces-morphisms-lemma-universally-closed-quasi-compact}에 의해
사상 $f$는 준콤팩트이다. $Y'$을 $Y$의 $X$ 안에서의 정규화로 정의하면
(5)와 (2)는 자동으로 성립한다. Morphisms of Spaces, Lemma
\ref{spaces-morphisms-lemma-normalization-in-universally-closed}에 의해
(4)가 성립한다. 사상 $f'$은 Morphisms of Spaces, Lemma
\ref{spaces-morphisms-lemma-universally-closed-permanence}에 의해
보편 닫힌사상이다. 이는 Morphisms of Spaces, Lemma
\ref{spaces-morphisms-lemma-quasi-compact-permanence}에 의해 준콤팩트이고,
Morphisms of Spaces, Lemma
\ref{spaces-morphisms-lemma-compose-after-separated}에 의해 준분리이다.

\medskip\noindent
남은 명제들을 보이기 위해 밑 $Y$가 아핀이라고 가정할 수 있다
(정규화를 취하는 것은 에탈 국소화와 가환하기 때문이다).
$Y = \Spec(R)$이라 하자. 그러면 $Y' = \Spec(A)$이고
$A = \Gamma(X, \mathcal{O}_X)$는 정수적 $R$-대수이다.
따라서 $f'_*\mathcal{O}_X$가 $\mathcal{O}_{Y'}$임은 명백하다
($f'_*\mathcal{O}_X$는 Morphisms of Spaces, Lemma
\ref{spaces-morphisms-lemma-pushforward}에 의해 준연접이고,
따라서 $\widetilde{A}$와 같기 때문이다). 이로써 (3)을 증명했다.

\medskip\noindent
$f'$이 전사임을 보이자. 위에서 보았듯이 $f'$은 보편 닫힌사상이므로,
$f'$의 상은 닫힌 부분집합 $V(I) \subset Y' = \Spec(A)$이다.
$h \in I$를 택하자. 그러면
% P09-ERR-0684: frozen source uses f-sharp although h lives on Y'.
$h|_X = f^\sharp(h)$는 $X$의 구조층의 대역절단이고 모든 점에서 사라진다.
$X$가 준콤팩트이므로 이는 $h|_X$가 멱영 절단임을 뜻한다. 즉
% P09-ERR-0668: frozen restriction notation omits the subscript marker.
$h^n|X = 0$인 $n > 0$가 있다. 그런데 $A = \Gamma(X, \mathcal{O}_X)$이므로
$h^n = 0$이다. 달리 말하면
% P09-ERR-0669: frozen source says Jacobson radical where nilradical is needed.
$I$는 $A$의 야코브슨 근기에 포함되고, 따라서 원하는 대로
$V(I) = Y'$이라고 결론내린다.
\end{proof}

\begin{lemma}
\label{spaces-more-morphisms-lemma-stein-universally-closed-residue-fields}
Lemma \ref{spaces-more-morphisms-lemma-stein-universally-closed}에서 $f$가 국소 유한형이고
$Y$가 아핀이라고 더 가정하자. 그러면
% P09-ERR-0670: frozen source writes the topological point as y in Y.
$y \in Y$에 대해 섬유
$\pi^{-1}(\{y\}) = \{y_1, \ldots, y_n\}$은 유한이고, 체 확대
$\kappa(y_i)/\kappa(y)$들은 유한이다.
\end{lemma}

\begin{proof}
$\pi^{-1}(\{y\})$의 점들 사이에는 특수화가 없음을 상기하자.
Algebra, Lemma \ref{algebra-lemma-integral-no-inclusion}을 보라.
$f'$이 전사이므로 $|X_y| \to \pi^{-1}(\{y\})$가 전사임을 얻는다.
$X_y$는 체 위의 유한형 준분리 대수공간임을 관찰하자
(준콤팩트성은 인용한 보조정리의 증명에서 보였다).
따라서 $|X_y|$는 뇌터 위상공간이다
(Morphisms of Spaces, Lemma
\ref{spaces-morphisms-lemma-finite-presentation-noetherian}).
% P09-ERR-0671: frozen source explicitly omits the finite-quotient argument.
이제 생략된 위상적 논증으로 $\pi^{-1}(\{y\})$가 유한임을 알 수 있다.
각 $i$에 대해 유한형 점 $x_i \in |X_y|$를 택할 수 있고 이는 $y_i$로 간다
(Morphisms of Spaces, Lemma
\ref{spaces-morphisms-lemma-enough-finite-type-points}).
$\kappa(y_i)/\kappa(y)$가 유한이라고 결론내리자.
$x_i$는 유한형 사상 $\Spec(k_i) \to X_y$로 나타낼 수 있고
(유한형 점의 정의에 의해), 따라서
% P09-ERR-0672: frozen source identifies y with Spec(kappa(y)).
$\Spec(k_i) \to y = \Spec(\kappa(y))$도 유한형이다
(유한형 사상들의 합성이므로). 그러므로 $k_i/\kappa(y)$는 유한이다
(Morphisms, Lemma \ref{morphisms-lemma-point-finite-type}).
\end{proof}

\noindent
$f : X \to Y$를 대수공간의 사상이라 하고
$\overline{y} : \Spec(k) \to Y$를 기하점이라 하자. 그러면
{\it $f$의 $\overline{y}$ 위 섬유}는 대수공간
$X_{\overline{y}} = X \times_{Y, \overline{y}} \Spec(k)$이고 이는 $k$ 위에 있다.
$Y$가 스킴이고
% P09-ERR-0673: frozen source writes the topological point as y in Y.
$y \in Y$가 점이면, 보통과 같이 섬유를
$X_y = X \times_Y \Spec(\kappa(y))$로 쓴다.

\begin{lemma}
\label{spaces-more-morphisms-lemma-characterize-geometrically-connected-fibres}
$S$를 스킴이라 하자. $f : X \to Y$를 $S$ 위 대수공간의 사상이라 하자.
$\overline{y}$를 $Y$의 기하점이라 하자. 그러면 $X_{\overline{y}}$가
연결인 것과, 모든 에탈 근방
$(V, \overline{v}) \to (Y, \overline{y})$에 대해, 여기서 $V$는 스킴이고, 밑변환
$X_V \to V$의 섬유 $X_v$가 연결인 것은 동치이다.
\end{lemma}

\begin{proof}
$\overline{y}$의 에탈 근방들의 범주는 여쌍여과이고 스킴들로 이루어진
공종 집합을 포함하므로
(Properties of Spaces, Lemma \ref{spaces-properties-lemma-cofinal-etale}),
$Y$를 이 근방들 가운데 하나로 바꾸고 $Y$가 스킴이라고 가정할 수 있다.
$y \in Y$를 $\overline{y}$에 대응하는 점이라 하자.
그러면 $X_y$가 $\kappa(y)$ 위에서 기하학적으로 연결인 것과
$X_{\overline{y}}$가 연결인 것은 동치이고, 이는 다시
$(X_y)_{k'}$가 연결인 것과 동치이다. 여기서 $k'$는 $\kappa(y)$의
임의의 유한 분리 확대이다. Spaces over Fields, Section
\ref{spaces-over-fields-section-geometrically-connected}, 특히 Lemma
\ref{spaces-over-fields-lemma-characterize-geometrically-disconnected}를 보라.
More on Morphisms, Lemma
\ref{more-morphisms-lemma-realize-prescribed-residue-field-extension-etale}에 의해
아핀 에탈 근방 $(V, v) \to (Y, y)$가 존재하여
% P09-ERR-0674: frozen source suddenly uses undefined s and u.
$\kappa(s) \subset \kappa(u)$가 주어진 임의의 유한 분리 확대
$\kappa(s) \subset k'$와 동일시된다. 보조정리가 따른다.
\end{proof}

\begin{theorem}[슈타인 분해; 뇌터 경우]
\label{spaces-more-morphisms-theorem-stein-factorization-Noetherian}
$S$를 스킴이라 하자. $f : X \to Y$를 $S$ 위 대수공간의 고유 사상이라 하고,
$Y$가 국소 뇌터라고 하자. 다음 분해가 존재한다.
$$
\xymatrix{
X \ar[rr]_{f'} \ar[rd]_f & & Y' \ar[dl]^\pi \\
& Y &
}
$$
이는 다음 성질들을 갖는다.
\begin{enumerate}
\item 사상 $f'$은 고유이고 기하학적으로 연결된 섬유들을 갖는다.
\item 사상 $\pi : Y' \to Y$는 유한이다.
\item $f'_*\mathcal{O}_X = \mathcal{O}_{Y'}$이다.
\item $Y' = \underline{\Spec}_Y(f_*\mathcal{O}_X)$이다.
\item $Y'$은 $Y$의 $X$ 안에서의 정규화이다. Morphisms, Definition
\ref{morphisms-definition-normalization-X-in-Y}를 보라.
\end{enumerate}
\end{theorem}

\begin{proof}
$f = \pi \circ f'$을 Lemma \ref{spaces-more-morphisms-lemma-stein-universally-closed}의 분해라 하자.
Lemma \ref{spaces-more-morphisms-lemma-stein-universally-closed}의 결론들뿐 아니라 $f'$이 분리이고
(Morphisms of Spaces, Lemma
\ref{spaces-morphisms-lemma-compose-after-separated}) 유한형임도 알 수 있다
(Morphisms of Spaces, Lemma
\ref{spaces-morphisms-lemma-permanence-finite-type}).
따라서 $f'$은 고유이다. Cohomology of Spaces, Lemma
\ref{spaces-cohomology-lemma-proper-pushforward-coherent}에 의해
$f_*\mathcal{O}_X$가 연접 $\mathcal{O}_Y$-가군임을 알 수 있다.
따라서 $\pi$가 유한임, 즉 (2)가 성립함을 알 수 있다.

\medskip\noindent
이로써 가장 흥미로운 명제, 즉 $f'$의 기하 섬유들이 연결이라는 것을
제외하고 모두 증명했다. 위 논의에서 $Y$를 $Y'$으로 바꿀 수 있음은 명백하다.
그러면 $Y$는 국소 뇌터이고, $f : X \to Y$는 고유이며,
$f_*\mathcal{O}_X = \mathcal{O}_Y$이다. $\overline{y}$를 $Y$의 기하점이라 하자.
여기서 형식함수 정리, 더 정확히 Cohomology of Spaces, Lemma
\ref{spaces-cohomology-lemma-formal-functions-stalk}를 적용한다.
그 정리는 다음을 말한다.
$$
\mathcal{O}^\wedge_{Y, \overline{y}} =
\lim_n H^0(X_n, \mathcal{O}_{X_n})
$$
여기서
$X_n = \Spec(\mathcal{O}_{Y, \overline{y}}/\mathfrak m_{\overline{y}}^n) \times_Y X$이다.
$X_1 = X_{\overline{y}} \to X_n$은 (유한 차수) 두껍게 하기이고,
따라서 $X_n$의 바탕 위상공간은 $X_{\overline{y}}$의 것과 같다.
그러므로 $X_{\overline{y}} = T_1 \amalg T_2$가 공집합이 아닌
열린닫힌 부분공간들의 서로소 합이면, 마찬가지로
$X_n = T_{1, n} \amalg T_{2, n}$이며 이는 모든 $n$에 대해 성립한다. 이는 다시
$H^0(X_n, \mathcal{O}_{X_n})$가 비자명 멱등원 $e_{1, n}$을 포함함을 뜻한다.
이 멱등원은 항등적으로 $1$인 함수인데 이는 $T_{1, n}$ 위에서이고,
항등적으로 $0$인데 이는 $T_{2, n}$ 위에서이다. $e_{1, n + 1}$은
$e_{1, n}$으로 제한되며 이는 $X_n$ 위에서임이 명백하다. 따라서
$e_1 = \lim e_{1, n}$은 그 극한의 비자명 멱등원이다.
이는 $\mathcal{O}^\wedge_{Y, \overline{y}}$가 국소환이라는 사실과 모순이다.
따라서 가정이 틀렸고, 원하는 대로 $X_{\overline{y}}$는 연결이다.
\end{proof}

\begin{theorem}[슈타인 분해; 일반 경우]
\label{spaces-more-morphisms-theorem-stein-factorization-general}
$S$를 스킴이라 하자. $f : X \to Y$를 $S$ 위 대수공간의 고유 사상이라 하자.
다음 분해가 존재한다.
$$
\xymatrix{
X \ar[rr]_{f'} \ar[rd]_f & & Y' \ar[dl]^\pi \\
& Y &
}
$$
이는 다음 성질들을 갖는다.
\begin{enumerate}
\item 사상 $f'$은 고유이고 기하학적으로 연결된 섬유들을 갖는다.
\item 사상 $\pi : Y' \to Y$는 정수적이다.
\item $f'_*\mathcal{O}_X = \mathcal{O}_{Y'}$이다.
\item $Y' = \underline{\Spec}_Y(f_*\mathcal{O}_X)$이다.
\item $Y'$은 $Y$의 $X$ 안에서의 정규화이다
(Morphisms of Spaces, Definition
\ref{spaces-morphisms-definition-normalization-X-in-Y}).
\end{enumerate}
\end{theorem}

\begin{proof}
Lemma \ref{spaces-more-morphisms-lemma-stein-universally-closed}를 적용하여 사상
$f' : X \to Y'$을 얻을 수 있다. Lemma
\ref{spaces-more-morphisms-lemma-stein-universally-closed}의 결론들뿐 아니라
$f'$이 분리이고 (Morphisms of Spaces, Lemma
\ref{spaces-morphisms-lemma-compose-after-separated}) 유한형임도 알 수 있다
(Morphisms of Spaces, Lemma
\ref{spaces-morphisms-lemma-permanence-finite-type}).
따라서 $f'$은 고유이다. 이제 $f'$이 기하학적으로 연결된 섬유들을
갖는다는 명제를 제외한 모든 명제를 증명했다.

\medskip\noindent
위 논의에서 $Y$를 $Y'$으로 바꿀 수 있음은 명백하다.
그러면 $f : X \to Y$는 고유이고 $f_*\mathcal{O}_X = \mathcal{O}_Y$이다.
이 조건들은 평탄한 밑변환 아래 보존됨을 관찰하자
(Morphisms of Spaces, Lemma \ref{spaces-morphisms-lemma-base-change-proper} 및
Cohomology of Spaces, Lemma
\ref{spaces-cohomology-lemma-flat-base-change-cohomology}).
$\overline{y}$를 $Y$의 기하점이라 하자. Lemma
\ref{spaces-more-morphisms-lemma-characterize-geometrically-connected-fibres}와 방금 한 관찰에 의해
$Y$가 스킴이고,
% P09-ERR-0675: frozen source writes the topological point as y in Y.
$y \in Y$가 점이며, $f : X \to Y$가 $Y$ 위의 고유 대수공간이고,
$f_*\mathcal{O}_X = \mathcal{O}_Y$이며, 섬유 $X_y$가 연결임을 보여야 하는
경우로 귀착된다. $Y$를 $y$의 아핀 근방으로 바꾸어
$Y = \Spec(R)$이 아핀이라고 가정할 수 있다. 그러면
$f_*\mathcal{O}_X = \mathcal{O}_Y$라는 것은 환 사상
$R \to \Gamma(X, \mathcal{O}_X)$가 전단사라는 뜻이다.

\medskip\noindent
Limits of Spaces, Lemma
\ref{spaces-limits-lemma-proper-limit-of-proper-finite-presentation-noetherian}에 의해
$(X \to Y) = \lim (X_i \to Y_i)$로 쓸 수 있다. 여기서 $X_i \to Y_i$는
고유이고 유한 표시이며 $Y_i$는 뇌터이다. 충분히 큰 $i$에 대해
$Y_i$는 아핀이다 (Limits of Spaces, Lemma
\ref{spaces-limits-lemma-limit-is-affine}).
$Y_i = \Spec(R_i)$라 하자. $R'_i = \Gamma(X_i, \mathcal{O}_{X_i})$로 놓자.
환 사상 $R_i \to R_i' \to R$가 있음을 관찰하자. 실제로 첫째 사상은
$X_i$가 $R_i$ 위 대수공간이기 때문에 있고, 둘째 사상은
$X \to X_i$ 및 $R = \Gamma(X, \mathcal{O}_X)$가 있기 때문에 있다.
Limits of Spaces, Lemma \ref{spaces-limits-lemma-descend-section}에 의해
$R = \colim R'_i$임을 관찰하자. 그러면
$$
\xymatrix{
X \ar[d]  \ar[r] & X_i \ar[d] \\
Y \ar[r] & Y'_i \ar[r] & Y_i
}
$$
는 가환하고 $Y'_i = \Spec(R'_i)$이다.
% P09-ERR-0676: frozen source writes the topological point as y'_i in Y'_i.
$y'_i \in Y'_i$를 $y$의 상이라 하자.
$X_y = \lim X_{i, y'_i}$이다. 실제로 $X = \lim X_i$,
$Y = \lim Y'_i$, 그리고 $\kappa(y) = \colim \kappa(y'_i)$이다.
이제 $X_y = U \amalg V$라 하고 $U$와 $V$가 열린 동시에 닫혀 있다고 하자.
그러면 $U, V$는 열린집합 $U_i, V_i$의 역상이고, 이 열린집합들은
$X_{i, y'_i}$ 안에 있다
(Limits of Spaces, Lemma \ref{spaces-limits-lemma-descend-opens}).
Theorem \ref{spaces-more-morphisms-theorem-stein-factorization-Noetherian}에 의해
$X_i \to Y'_i$의 섬유들은 연결이고, 따라서 $U$ 또는 $V$ 중 하나는 공집합이다.
이로써 증명이 끝난다.
\end{proof}

\noindent
응용 하나를 제시한다.

\begin{lemma}
\label{spaces-more-morphisms-lemma-geometrically-connected-fibres-towards-normal}
$S$를 스킴이라 하자. $f : X \to Y$를 $S$ 위 대수공간의 사상이라 하자.
다음을 가정하자.
\begin{enumerate}
\item $f$는 고유이다.
\item $Y$는 정수적이고 (Spaces over Fields, Definition
\ref{spaces-over-fields-definition-integral-algebraic-space}) 일반점은 $\xi$이다.
\item $Y$는 정규이다.
\item $X$는 축약이다.
\item $|X|$의 기약성분의 모든 일반점은 $\xi$로 간다.
\item $H^0(X_\xi, \mathcal{O}) = \kappa(\xi)$이다.
\end{enumerate}
그러면 $f_*\mathcal{O}_X = \mathcal{O}_Y$이고 $f$는
기하학적으로 연결된 섬유들을 갖는다.
\end{lemma}

\begin{proof}
Theorem \ref{spaces-more-morphisms-theorem-stein-factorization-general}을 적용하여 분해
$X \to Y' \to Y$를 얻는다. $Y' = Y$임을 보이면 충분하다.
$Y' \times_Y V \to V$가 동형임을 보이면 충분하다. 여기서
$V \to Y$는 에탈 사상이고 $V$는 아핀 정수적 스킴이다. Spaces over Fields,
Lemma \ref{spaces-over-fields-lemma-normal-integral-cover-by-affines}를 보라.
$Y'$의 형성은 에탈 밑변환과 가환한다. Morphisms of Spaces, Lemma
\ref{spaces-morphisms-lemma-properties-normalization}을 보라.
$X \times_Y V$의 일반점들은 $X$의 일반점들 위에 놓이고
(Decent Spaces, Lemma \ref{decent-spaces-lemma-decent-generic-points}),
따라서 가정 (5)에 의해 $V$의 일반점으로 간다. 더욱이 조건 (6)은
$V \to Y$에 의한 밑변환 아래 보존된다. 예를 들어 평탄한 밑변환
(Cohomology of Spaces, Lemma
\ref{spaces-cohomology-lemma-flat-base-change-cohomology})을 사용하라.
따라서 $Y$가 정규 정수적 아핀 스킴인 경우에 보조정리를 증명하면 충분하다.

\medskip\noindent
$Y$가 정규 정수적 아핀 스킴이라고 가정하자. Morphisms, Lemma
\ref{morphisms-lemma-finite-birational-over-normal}을 적용하여
$Y' \to Y$가 동형임을 보이겠다. 실제로 $Y'$도 축약이다. 이는 $X$가 축약이기 때문이다
(Morphisms of Spaces, Lemma
\ref{spaces-morphisms-lemma-normalization-in-reduced}).
위에서 인용한 정리에 의해 사상 $Y' \to Y$는 정수적이다.
$Y$가 decent이고 $X \to Y$가 분리이므로 $X$도 decent임을 알 수 있다.
이를 위해 Decent Spaces, Lemmas
\ref{decent-spaces-lemma-properties-trivial-implications},
\ref{decent-spaces-lemma-property-over-property}를 사용하라.
가정 (5), Morphisms of Spaces, Lemma
\ref{spaces-morphisms-lemma-normalization-generic}, Decent Spaces, Lemma
\ref{decent-spaces-lemma-decent-generic-points}에 의해 $|Y'|$의
기약성분의 모든 일반점은 $\xi$로 간다. 한편 $Y'$은
$f_*\mathcal{O}_X$의 상대 스펙트럼이므로 스킴 이론적 섬유
$Y'_\xi$는 $H^0(X_\xi, \mathcal{O})$의 스펙트럼이고, 이는 가정에 의해
$\kappa(\xi)$와 같다. 따라서 $Y'$은 정수적 스킴이고 그 함수체는
$Y$의 함수체와 같다. 이로써 증명이 끝난다.
\end{proof}

\noindent
또 다른 응용을 제시한다.

\begin{lemma}
\label{spaces-more-morphisms-lemma-proper-flat-nr-geom-conn-comps-lower-semicontinuous}
$S$를 스킴이라 하자.
% P09-ERR-0677: frozen source leaves the displayed morphism unnamed but later uses f.
$X \to Y$를 $S$ 위 대수공간의 사상이라 하자.
$f$가 고유이고 평탄하고 유한 표시이면, 함수
$n_{X/Y} : |Y| \to \mathbf{Z}$는 $f$의 섬유들의 기하학적
연결성분 개수를 세며
(Lemma \ref{spaces-more-morphisms-lemma-base-change-fibres-nr-geometrically-connected-components})
아래쪽 반연속이다.
\end{lemma}

\begin{proof}
문제는 $Y$ 위에서 에탈 국소적이므로 $Y$가 아핀 스킴이라고 가정할 수 있고
그렇게 한다. $y \in Y$라 하자.
% P09-ERR-0678: frozen source uses n_{X/S} instead of n_{X/Y}.
$n = n_{X/S}(y)$로 놓자. $n < \infty$임을 관찰하자. 실제로
$X \to Y$의 $y$에서의 기하 섬유는 체 위의 고유 대수공간이므로
뇌터이고, 따라서 연결성분이 유한 개이다. 열린 근방 $V$로서 $y$의 근방이고
% P09-ERR-0679: frozen source uses n_{X/S}|_V instead of n_{X/Y}|_V.
$n_{X/S}|_V \geq n$인 것을 찾아야 한다.
$X \to Y' \to Y$를 Theorem
\ref{spaces-more-morphisms-theorem-stein-factorization-general}에서와 같은 슈타인 분해라 하자.
Lemma \ref{spaces-more-morphisms-lemma-stein-universally-closed-residue-fields}에 의해
유한 개의 점 $y'_1, \ldots, y'_m \in Y'$이 있고 이들은 $y$ 위에 놓이며,
확대 $\kappa(y'_i)/\kappa(y)$는 유한이다. More on Morphisms, Lemma
\ref{more-morphisms-lemma-etale-makes-integral-split}에 의하면 $Y$를
$y$의 에탈 근방으로 바꾼 뒤, 스킴으로서
$Y' = V_1 \amalg \ldots \amalg V_m$이고 $y'_i \in V_i$이며
$\kappa(y'_i)/\kappa(y)$가 순수 비분리라고 가정할 수 있다.
그러면 대수공간 $X_{y_i'}$들은 $\kappa(y)$ 위에서 기하학적으로 연결이고,
따라서 $m = n$이다. 대수공간들
$X_i = (f')^{-1}(V_i)$, $i = 1, \ldots, n$은 $Y$ 위에서 평탄하고
유한 표시이다. 따라서 $X_i \to Y$의 상은 열려 있다
(Morphisms of Spaces, Lemma \ref{spaces-morphisms-lemma-fppf-open}).
그러므로 $y$의 한 근방에서 $n_{X/Y}$는 적어도 $n$이다.
\end{proof}

\begin{lemma}
\label{spaces-more-morphisms-lemma-proper-flat-geom-red}
$S$를 스킴이라 하자. $f : X \to Y$를 $S$ 위 대수공간의 사상이라 하자.
다음을 가정하자.
\begin{enumerate}
\item $f$는 고유이고 평탄하고 유한 표시이다.
\item $f$의 기하 섬유들은 축약이다.
\end{enumerate}
% P09-ERR-0680: frozen source names the Y-relative function n_{X/S}.
$n_{X/S} : |Y| \to \mathbf{Z}$를 생각하자. 이 함수는 $f$의 섬유들의
기하학적 연결성분 개수를 세며
(Lemma \ref{spaces-more-morphisms-lemma-base-change-fibres-nr-geometrically-connected-components})
국소 상수이다.
\end{lemma}

\begin{proof}
Lemma \ref{spaces-more-morphisms-lemma-proper-flat-nr-geom-conn-comps-lower-semicontinuous}에 의해
% P09-ERR-0681: frozen source misspells “semi-continuous”.
함수 $n_{X/Y}$는 아래쪽 반연속이다. 따라서 위쪽 반연속임을 보이면 충분하다.
이를 위해 $Y$ 위에서 에탈 국소적으로 작업할 수 있고, 따라서
$Y$가 아핀 스킴이라고 가정할 수 있다. $y \in Y$에 대해
$\kappa(y)$-대수
$$
A = H^0(X_y, \mathcal{O}_{X_y})
$$
를 생각하자. Spaces over Fields, Lemma
\ref{spaces-over-fields-lemma-proper-geometrically-reduced-global-sections}와
$X_y$가 기하학적으로 축약이라는 사실에 의해 $A$는
$\kappa(y)$의 유한 분리 확대들의 유한 곱이다. 따라서
$A \otimes_{\kappa(y)} \kappa(\overline{y})$는
$\beta_0(y) = \dim_{\kappa(y)} A$개의 $\kappa(\overline{y})$ 복사본으로
이루어진 곱이다.
그러므로 $X_{\overline{y}}$에는 연결성분이 $\beta_0(y)$개 있다.
% P09-ERR-0682: frozen source writes n_{X/S}=beta_0 instead of n_{X/Y}=beta_0.
$n_{X/S} = \beta_0$이고, 이는 $Y$ 위의 함수들로서의 등식이다.
따라서 Derived Categories of Spaces, Lemma
\ref{spaces-perfect-lemma-jump-loci-geometric}에 의해 $n_{X/Y}$는
위쪽 반연속이다. 이로써 증명이 끝난다.
\end{proof}

\begin{lemma}
\label{spaces-more-morphisms-lemma-stein-factorization-etale}
$S$를 스킴이라 하자. $f : X \to Y$를 $S$ 위 대수공간의 고유 사상이라 하자.
$X \to Y' \to Y$를 $f$의 슈타인 분해라 하자
(Theorem \ref{spaces-more-morphisms-theorem-stein-factorization-general}).
$f$가 유한 표시이고 평탄하며 기하학적으로 축약인 섬유들을 가지면
(Definition \ref{spaces-more-morphisms-definition-geometrically-reduced-fibre}),
$Y' \to Y$는 유한 에탈이다.
\end{lemma}

\begin{proof}
슈타인 분해의 형성은 평탄한 밑변환과 가환한다. Lemma
\ref{spaces-more-morphisms-lemma-stein-universally-closed}를 보라. 따라서 $Y$ 위에서
에탈 국소적으로 작업하여 $Y$가 아핀 스킴이라고 가정할 수 있다.
그러면 $Y'$은 아핀 스킴이고 $Y' \to Y$는 정수적이다.

\medskip\noindent
$y \in Y$라 하자. $n$을 기하 섬유 $X_{\overline{y}}$의 연결성분 개수라 놓자.
$n < \infty$이다. 실제로 $X \to Y$의 $y$에서의 기하 섬유는 체 위의
고유 대수공간이므로 뇌터이고, 따라서 연결성분이 유한 개이다.
Lemma \ref{spaces-more-morphisms-lemma-stein-universally-closed-residue-fields}에 의해
유한 개의 점 $y'_1, \ldots, y'_m \in Y'$이 있고 이들은 $y$ 위에 놓이며,
각 $i$에 대해 유한형 점 $x_i \in |X_y|$를 택할 수 있다.
% P09-ERR-0683: frozen source omits a conjunction after y'_i.
이 점은 $y'_i$로 가고 확대 $\kappa(y'_i)/\kappa(y)$는 유한이다.
따라서 More on Morphisms, Lemma
\ref{more-morphisms-lemma-etale-makes-integral-split}에 의하면 $Y$를
$y$의 에탈 근방으로 바꾼 뒤, 스킴으로서
$Y' = V_1 \amalg \ldots \amalg V_m$이고 $y'_i \in V_i$이며
$\kappa(y'_i)/\kappa(y)$가 순수 비분리라고 가정할 수 있다.
이 경우 대수공간 $X_{y_i'}$들은 $\kappa(y)$ 위에서
기하학적으로 연결이고, 따라서 $m = n$이다. 대수공간들
$X_i = (f')^{-1}(V_i)$, $i = 1, \ldots, n$은 $Y$ 위에서 고유이고,
평탄하고, 유한 표시이며, 기하학적으로 축약인 섬유들을 갖는다.
각 사상 $X_i \to Y$에 대해 보조정리를 증명하면 충분하다.
이는 $X_{\overline{y}}$가 연결인 경우로 우리를 귀착시킨다.

\medskip\noindent
$X_{\overline{y}}$가 연결이라고 가정하자. Lemma
\ref{spaces-more-morphisms-lemma-proper-flat-geom-red}에 의해 $X \to Y$는 $y$의 한 근방에서
기하학적으로 연결된 섬유들을 갖는다. 따라서 $X \to Y$의 섬유들이
기하학적으로 연결이라고 가정할 수 있다. 그러면 Derived Categories of Spaces,
Lemma \ref{spaces-perfect-lemma-proper-flat-geom-red-connected}에 의해
$f_*\mathcal{O}_X = \mathcal{O}_Y$이고, 이것으로 증명이 끝난다.
\end{proof}

\noindent
다음 보조정리의 증명은 스킴에 대한 슈타인 분해를 사용한다.
이것이 이 보조정리가 이 절에 들어온 이유이다.

\begin{lemma}
\label{spaces-more-morphisms-lemma-split-off-proper-part-henselian}
$(A, I)$를 헨젤 쌍이라 하자. $X$를 $A$ 위에서 분리이고 유한형인
대수공간이라 하자. $X_0 = X \times_{\Spec(A)} \Spec(A/I)$로 놓자.
$Y \subset X_0$를 열린닫힌 부분공간으로서
$Y \to \Spec(A/I)$가 고유인 것이라 하자. 그러면 열린닫힌 부분공간
$W \subset X$가 존재한다. 이는 $A$ 위에서 고유이고
$W \times_{\Spec(A)} \Spec(A/I) = Y$이다.
\end{lemma}

\begin{proof}
$T \mapsto T_0$로 $\Spec(A/I) \to \Spec(A)$에 의한 밑변환을 나타내겠다.
차우 보조정리의 약한 판
(Cohomology of Spaces, Lemma \ref{spaces-cohomology-lemma-weak-chow}의 형태)에 의해
전사 고유 사상 $\varphi : X' \to X$가 존재하여 $X'$은
$\mathbf{P}^n_A$ 안으로의 몰입을 갖는다. $Y' = \varphi^{-1}(Y)$로 놓자.
이는 $X'_0$의 열린닫힌 부분스킴이다. More on Morphisms, Lemma
\ref{more-morphisms-lemma-split-off-proper-part-henselian}에 의해
보조정리는 $(X', Y')$에 대해 성립한다. $W' \subset X'$을 $A$ 위에서
고유인 열린닫힌 부분스킴으로서 $Y' = W'_0$인 것이라 하자.
Morphisms of Spaces, Lemma
\ref{spaces-morphisms-lemma-universally-closed-permanence}에 의해
$Q_1 = \varphi(|W'|) \subset |X|$ 및
$Q_2 = \varphi(|X' \setminus W'|) \subset |X|$는 닫힌 부분집합이다.
또한 Morphisms of Spaces, Lemma
\ref{spaces-morphisms-lemma-image-proper-is-proper}에 의해 $Q_1$ 위의
임의의 닫힌 부분공간 구조는 $A$ 위에서 고유이다.
$Q_1 \cap Q_2$의 $\Spec(A)$ 안에서의 상은 닫혀 있다.
$(A, I)$가 헨젤이므로 $Q_1 \cap Q_2$가 공집합이 아니면
$Q_1 \cap Q_2$에는 $\Spec(A/I)$ 위에 놓이는 점이 있음을 알 수 있다.
$W'_0 = Y' = \varphi^{-1}(Y)$이므로 이는 불가능하다.
따라서 $Q_1$은 $|X|$ 안에서 열리고 닫혀 있다.
$W \subset X$를 이에 대응하는 열린닫힌 부분공간이라 하자.
그러면 $W$는 $A$ 위에서 고유이고 $W_0 = Y$이다.
\end{proof}










\section{열린 부분공간에서 성질의 연장}
\label{spaces-more-morphisms-section-extending-properties}

\noindent
이 절에서는 다음 형태의 결과 몇 가지를 모은다. $f : X \to Y$가
대수공간의 평탄 사상이고, $f$가 $Y$의 조밀한 열린 부분공간 위에서
어떤 성질을 만족하면, $f$는 $Y$ 전체 위에서 같은 성질을 만족한다.

\begin{lemma}
\label{spaces-more-morphisms-lemma-flat-finite-type-finitely-presented-over-dense-open}
$S$를 스킴이라 하자. $f : X \to Y$를 $S$ 위의 대수공간들의 사상이라 하자.
$\mathcal{F}$를 준연접 $\mathcal{O}_X$-가군이라 하자.
$V \subset Y$를 열린 부분공간이라 하자. 다음을 가정하자.
\begin{enumerate}
\item $f$는 국소 유한 표시이다.
\item $\mathcal{F}$는 유한형이고 $Y$ 위에서 평탄하다.
\item $V \to Y$는 준콤팩트이고 스킴론적으로 조밀하다.
\item $\mathcal{F}|_{f^{-1}V}$는 유한 표시이다.
\end{enumerate}
그러면 $\mathcal{F}$는 유한 표시이다.
\end{lemma}

\begin{proof}
$\mathcal{F}$를 $X$ 위로 전사 에탈인 스킴에 당긴 것이 유한 표시임을
보이면 충분하다. 따라서 $X$가 스킴이라고 가정해도 된다.
% P09-ERR-0685: frozen source omits the modifier relating the scheme cover to X.
마찬가지로 $Y$를 $Y$ 위로 전사 에탈인 스킴으로 바꾸어도 된다.
% P09-ERR-0686: frozen source has a duplicated conjunction in the Y-cover phrase.
이 스킴에서 $V$의 역상은 스킴론적으로 조밀하다. 다음을 보라.
Morphisms of Spaces, Section
\ref{spaces-morphisms-section-scheme-theoretic-closure}.
이로써 스킴의 경우로 환원되며, 이는 More on Flatness, Lemma
\ref{flat-lemma-flat-finite-type-finitely-presented-over-dense-open}이다.
\end{proof}

\begin{lemma}
\label{spaces-more-morphisms-lemma-flat-finite-type-finitely-presented-over-dense-open-X}
$S$를 스킴이라 하자. $f : X \to Y$를 $S$ 위의 대수공간들의 사상이라 하자.
$V \subset Y$를 열린 부분공간이라 하자. 다음을 가정하자.
\begin{enumerate}
\item $f$는 국소 유한형이고 평탄하다.
\item $V \to Y$는 준콤팩트이고 스킴론적으로 조밀하다.
\item $f|_{f^{-1}V} : f^{-1}V \to V$는 국소 유한 표시이다.
\end{enumerate}
그러면 $f$는 국소 유한 표시이다.
% P09-ERR-0687: frozen source duplicates “of” in the conclusion.
\end{lemma}

\begin{proof}
증명은 Lemma
\ref{spaces-more-morphisms-lemma-flat-finite-type-finitely-presented-over-dense-open}의
증명과 동일하되, More on Flatness, Lemma
\ref{flat-lemma-flat-finite-type-finitely-presented-over-dense-open-X}를 사용한다.
\end{proof}

\begin{lemma}
\label{spaces-more-morphisms-lemma-flat-fp-dimension-over-dense-open}
$S$를 스킴이라 하자. $f : X \to Y$를 $S$ 위의 대수공간들의 사상으로서
평탄하고 국소 유한형인 것이라 하자. $V \subset Y$를 열린 부분공간으로서
$|V| \subset |Y|$가 조밀하고 $X_V \to V$의 상대 차원이 $\leq d$인
것이라 하자. 또한 다음 중 하나가 성립한다고 하자.
\begin{enumerate}
\item $f$는 국소 유한 표시이다.
\item $V \to Y$는 준콤팩트이다.
\end{enumerate}
그러면 $f : X \to Y$의 상대 차원은 $\leq d$이다.
\end{lemma}

\begin{proof}
$Y$를 그 축약으로 바꾸어도 되므로 $Y$가 축약이라고 가정해도 된다.
그러면 $V$는 $Y$ 안에서 스킴론적으로 조밀하다. 다음을 보라.
Morphisms of Spaces, Lemma
\ref{spaces-morphisms-lemma-quasi-compact-immersion}.
정의에 의해 상대 차원이 $\leq d$라는 성질은 에탈 덮개에서 확인할 수 있다.
% P09-ERR-0688: frozen source uses plural “Sections” for one section reference.
다음을 보라. Morphisms of Spaces, Section
\ref{spaces-morphisms-section-relative-dimension}.
따라서 $f$의 상대 차원이 $\leq d$임을 보이면 충분하며, 이를 위해
$X$를 $X$ 위로 전사 에탈인 스킴으로 바꾸어도 된다.
마찬가지로 $Y$를 $Y$ 위로 전사 에탈인 스킴으로 바꾸어도 된다.
% P09-ERR-0689: frozen source repeats the malformed Y-cover phrase.
이 스킴에서 $V$의 역상은 스킴론적으로 조밀하다. 다음을 보라.
Morphisms of Spaces, Section
\ref{spaces-morphisms-section-scheme-theoretic-closure}.
스킴의 스킴론적으로 조밀한 열린 부분공간은 특히 조밀하므로
스킴의 경우로 환원되며, 이는 More on Flatness, Lemma
\ref{flat-lemma-flat-finite-presentation-dimension-over-dense-open}이다.
\end{proof}

\begin{lemma}
\label{spaces-more-morphisms-lemma-proper-flat-finite-over-dense-open}
$S$를 스킴이라 하자. $f : X \to Y$를 $S$ 위의 대수공간들의 사상으로서
평탄하고 고유인 것이라 하자. $V \to Y$를 열린 부분공간으로서
% P09-ERR-0690: frozen source conflates an open subspace with its morphism.
$|V| \subset |Y|$가 조밀하고 $X_V \to V$가 유한인 것이라 하자.
또한 $f$가 국소 유한 표시이거나 $V \to Y$가 준콤팩트라고 하자.
% P09-ERR-0691: frozen source has an ungrammatical alternative-hypothesis phrase.
그러면 $f$는 유한이다.
\end{lemma}

\begin{proof}
Lemma \ref{spaces-more-morphisms-lemma-flat-fp-dimension-over-dense-open}에 의해 $f$의 섬유들은
차원이 영이다. Morphisms of Spaces, Lemma
\ref{spaces-morphisms-lemma-locally-quasi-finite-rel-dimension-0}에 의해
이는 $f$가 국소 준유한임을 뜻한다. Morphisms of Spaces, Lemma
\ref{spaces-morphisms-lemma-locally-quasi-finite-separated-representable}에 의해
이는 $f$가 표현 가능임을 뜻한다. $f$가 유한인지 여부는 $Y$ 위에서
에탈 국소적으로 확인할 수 있으므로 $Y$가 스킴이라고 가정해도 된다.
$f$가 표현 가능이므로 스킴의 경우로 환원되며, 이는
More on Flatness, Lemma \ref{flat-lemma-proper-flat-finite-over-dense-open}이다.
\end{proof}

\begin{lemma}
\label{spaces-more-morphisms-lemma-zariski}
$S$를 스킴이라 하자. $f : X \to Y$를 $S$ 위의 대수공간들의 사상이라 하자.
$V \subset Y$를 열린 부분공간이라 하자. 다음을 가정하자.
\begin{enumerate}
\item $f$는 분리이고 국소 유한형이며 평탄하다.
\item $f^{-1}(V) \to V$는 동형이다.
\item $V \to Y$는 준콤팩트이고 스킴론적으로 조밀하다.
\end{enumerate}
그러면 $f$는 열린 몰입이다.
\end{lemma}

\begin{proof}
Lemma \ref{spaces-more-morphisms-lemma-flat-finite-type-finitely-presented-over-dense-open-X}를
적용하면 $f$는 국소 유한 표시이다. Lemma
\ref{spaces-more-morphisms-lemma-flat-fp-dimension-over-dense-open}를 적용하면 $f$의 상대 차원은
$\leq 0$이다. Morphisms of Spaces, Lemma
\ref{spaces-morphisms-lemma-locally-quasi-finite-rel-dimension-0}에 의해
이는 $f$가 국소 준유한임을 뜻한다. Morphisms of Spaces, Lemma
\ref{spaces-morphisms-lemma-locally-quasi-finite-separated-representable}에 의해
이는 $f$가 표현 가능임을 뜻한다. Descent on Spaces, Lemma
\ref{spaces-descent-lemma-descending-property-open-immersion}에 의해
$f$가 열린 몰입인지 여부는 $Y$ 위에서 에탈 국소적으로 확인할 수 있다.
따라서 $Y$가 스킴이라고 가정해도 된다. $f$가 표현 가능이므로
스킴의 경우로 환원되며, 이는 More on Flatness, Lemma
\ref{flat-lemma-zariski}이다.
\end{proof}







\section{블로업과 평탄성}
\label{spaces-more-morphisms-section-blowup-flat}

\noindent
More on Flatness, Section \ref{flat-section-blowup-flat}의 작업을
되풀이하는 대신 More on Flatness, Lemma \ref{flat-lemma-push-ideal}의
유사한 결과를 증명한다. 이는 적절한 블로업을 찾는 문제가 흔히
밑공간 위에서 에탈 국소적임을 말해 준다.

\begin{lemma}
\label{spaces-more-morphisms-lemma-push-ideal}
$S$를 스킴이라 하자. $X$를 $S$ 위의 준콤팩트이고 준분리인
대수공간이라 하자. $\varphi : W \to X$를 준콤팩트이고 분리인
에탈 사상이라 하자. $U \subset X$를 준콤팩트 열린 부분공간이라 하자.
$\mathcal{I} \subset \mathcal{O}_W$를 유한형 준연접 아이디얼 층으로서
$V(\mathcal{I}) \cap \varphi^{-1}(U) = \emptyset$인 것이라 하자.
그러면 유한형 준연접 아이디얼 층
$\mathcal{J} \subset \mathcal{O}_X$가 존재하여 다음을 만족한다.
\begin{enumerate}
\item $V(\mathcal{J}) \cap U = \emptyset$이다.
\item $\varphi^{-1}(\mathcal{J})\mathcal{O}_W = \mathcal{I} \mathcal{I}'$이며,
여기서 $\mathcal{I}' \subset \mathcal{O}_W$는 어떤 유한형 준연접
아이디얼이다.
\end{enumerate}
\end{lemma}

\begin{proof}
$W \to Y \to X$라는 분해를 택하자. 여기서 $j : W \to Y$는
준콤팩트 열린 몰입이고 $\pi : Y \to X$는 유한 표시인 유한 사상이다
(Lemma \ref{spaces-more-morphisms-lemma-quasi-finite-separated-pass-through-finite-addendum}).
$V = j(W) \cup \pi^{-1}(U) \subset Y$로 놓자. $\mathcal{I}$는
$W \cong j(W)$ 위에서 $\mathcal{O}_{\pi^{-1}(U)}$와 서로 붙어서 유한형
준연접 아이디얼 층 $\mathcal{I}_1 \subset \mathcal{O}_V$을 이룬다.
Limits of Spaces, Lemma \ref{spaces-limits-lemma-extend}에 의해
유한형 준연접 아이디얼 층 $\mathcal{I}_2 \subset \mathcal{O}_Y$가
존재하여 $\mathcal{I}_2|_V = \mathcal{I}_1$이다. 달리 말해,
$\mathcal{I}_2 \subset \mathcal{O}_Y$는 유한형 준연접 아이디얼 층으로서
$V(\mathcal{I}_2)$가 $\pi^{-1}(U)$와 서로소이고
$j^{-1}\mathcal{I}_2 = \mathcal{I}$이다. 이에 대응하는 닫힌 몰입을
$i : Z \to Y$로 나타내자.
% P09-ERR-0692: frozen source omits the preposition in the naming construction.
이는 유한 표시이다(Morphisms of Spaces, Lemma
\ref{spaces-morphisms-lemma-closed-immersion-finite-presentation}).
특히 합성 $\tau = \pi \circ i : Z \to X$는 유한이고 유한 표시이다
(Morphisms of Spaces, Lemmas
\ref{spaces-morphisms-lemma-composition-finite-presentation} and
\ref{spaces-morphisms-lemma-composition-integral}).

\medskip\noindent
$\mathcal{F} = \tau_*\mathcal{O}_Z$로 놓고 이를 준연접
$\mathcal{O}_X$-가군으로 보자. Descent on Spaces, Lemma
\ref{spaces-descent-lemma-finite-finitely-presented-module}에 의해
$\mathcal{F}$는 유한 표시 $\mathcal{O}_X$-가군이다.
$\mathcal{J} = \text{Fit}_0(\mathcal{F})$로 놓자. (환 달린 토포스 위의
Fitting 가군에 대한 참고문헌을 여기에 넣어야 한다.)
% P09-ERR-0693: frozen source renders this editorial placeholder.
이는 $X$ 위의 유한형 준연접 아이디얼 층이다($\mathcal{F}$가 유한
표시인 까닭이며, More on Algebra, Lemma
\ref{more-algebra-lemma-fitting-ideal-basics}를 보라).
$|\tau|(|Z|) \cap |U| = \emptyset$인 것은 아이디얼 $\mathcal{I}_2$의
선택에 의한 것이고, 영 가군의 제$0$ Fitting 아이디얼은
구조층과 같으므로 보조정리의 (1)이 성립한다. (2)를 보이기 위해,
Fitting 아이디얼을 취하는 것이 밑변환과 가환하므로
$\varphi^{-1}(\mathcal{J})\mathcal{O}_W = \text{Fit}_0(\varphi^*\mathcal{F})$임을
주목하자. 한편 $\varphi : W \to X$가 분리이고 에탈이므로
$(1, j) : W \to W \times_X Y$는 열린닫힌 몰입이다. 따라서
$W \times_Y Z = V(\mathcal{I}) \amalg Z'$이며, 여기서
$\tau' : Z' \to W$는 유한이고 유한 표시인 대수공간들의 사상이다.
% P09-ERR-0694: frozen source uses W times_Y Z twice where base change requires X.
따라서 다음을 얻는다.
\begin{align*}
\text{Fit}_0(\varphi^*\mathcal{F}) & =
\text{Fit}_0((W \times_Y Z \to W)_*\mathcal{O}_{W \times_Y Z}) \\
& =
\text{Fit}_0(\mathcal{O}_W/\mathcal{I}) \cdot
\text{Fit}_0(\tau'_*\mathcal{O}_{Z'}) \\
& =
\mathcal{I} \cdot \text{Fit}_0(\tau'_*\mathcal{O}_{Z'})
\end{align*}
여기서 두 번째 등식은 More on Algebra, Lemma
\ref{more-algebra-lemma-fitting-ideal-basics}를 환 달린 토포스 위의 층으로
옮겨 적용한 결과이다.
% P09-ERR-0695: frozen source's justification is a sentence fragment.
$\mathcal{I}' = \text{Fit}_0(\tau'_*\mathcal{O}_{Z'})$로 놓으면
보조정리의 증명이 끝난다.
\end{proof}

\begin{theorem}
\label{spaces-more-morphisms-theorem-flatten-module}
$S$를 스킴이라 하자. $B$를 $S$ 위의 준콤팩트이고 준분리인
대수공간이라 하자. $X$를 $B$ 위의 대수공간이라 하자.
$\mathcal{F}$를 $X$ 위의 준연접 가군이라 하자.
$U \subset B$를 준콤팩트 열린 부분공간이라 하자. 다음을 가정하자.
\begin{enumerate}
\item $X$는 준콤팩트이다.
\item $X$는 $B$ 위에서 국소 유한 표시이다.
\item $\mathcal{F}$는 유한형 가군이다.
\item $\mathcal{F}_U$는 유한 표시이다.
\item $\mathcal{F}_U$는 $U$ 위에서 평탄하다.
\end{enumerate}
그러면 $U$-허용 블로업 $B' \to B$가 존재하여 진변환
$\mathcal{F}'$, 즉 $\mathcal{F}$의 진변환은 유한 표시
$\mathcal{O}_{X \times_B B'}$-가군이고 $B'$ 위에서 평탄하다.
\end{theorem}

\begin{proof}
아핀 스킴 $V$와 전사 에탈 사상 $V \to X$를 택하자. 진변환은
에탈 국소화와 가환하므로(Divisors on Spaces, Lemma
\ref{spaces-divisors-lemma-strict-transform-local}), 결과를 $X$ 대신
$V$에 대해 증명하면 충분하다. 따라서 보조정리의 가정들에 더하여
$X \to B$가 표현 가능이라고 가정해도 된다.
% P09-ERR-0696: frozen source says “hypotheses of the lemma” inside this theorem.

\medskip\noindent
$X \to B$가 표현 가능이라고 가정하자. 아핀 스킴 $W$와 전사 에탈 사상
$\varphi : W \to B$를 택하자. $X \times_B W$는 스킴임을 주목하자.
스킴의 경우(More on Flatness, Theorem
\ref{flat-theorem-flatten-module})에 의해 유한형 준연접 아이디얼 층
$\mathcal{I} \subset \mathcal{O}_W$를 다음과 같이 찾을 수 있다.
(a) $|V(\mathcal{I})| \cap |\varphi^{-1}(U)| = \emptyset$이고, (b)
$\mathcal{F}|_{X \times_B W}$의, 블로업 $W' \to W$에 관한 진변환은
(이 블로업은 $\mathcal{I}$에서 취한다) $W'$ 위에서 평탄한
유한 표시 가군이다. Lemma \ref{spaces-more-morphisms-lemma-push-ideal}에 나오는 것과 같은
유한형 아이디얼 층 $\mathcal{J} \subset \mathcal{O}_B$를 택하자.
$B' \to B$를 $\mathcal{J}$의 블로업이라 하자. 이 블로업이 원하는
것임을 보이겠다. 실제로 $B' \to B$가 $U$-허용임은 아이디얼
$\mathcal{J}$의 선택에 의해 분명하다. 더욱이 밑변환
$B' \times_B W \to W$는 $W$의
$\varphi^{-1}\mathcal{J} = \mathcal{I}\mathcal{I}'$에서의 블로업이다
(평탄 밑변환과 블로업의 호환성은 Divisors on Spaces, Lemma
\ref{spaces-divisors-lemma-flat-base-change-blowing-up}를 보라).
따라서 다음 분해가 존재한다.
$$
W \times_B B' \to W' \to W
$$
여기서 첫 번째 사상도 블로업이다. 다음을 보라. Divisors on Spaces, Lemma
\ref{spaces-divisors-lemma-blowing-up-two-ideals}.
$\mathcal{F}'$은 $B' \times_B X$ 위에 놓이며, 이를
$W \times_B B' \times_B X$에 제한한 것은
$\mathcal{F}|_{X \times_B W}$의 진변환이다(Divisors on Spaces, Lemma
\ref{spaces-divisors-lemma-strict-transform-local}). 따라서 이는 위에 표시한
두 블로업에 따라 $\mathcal{F}|_{X \times_B W}$의 진변환을 두 번
연달아 취한 것이다(Divisors on Spaces, Lemma
\ref{spaces-divisors-lemma-strict-transform-composition-blowups}).
첫 번째 블로업 뒤에는 구성에 의해 우리 층이 이미 밑공간 위에서 평탄하고
유한 표시이다. 따라서 두 번째 진변환 뒤에도 그러하다. 실제로 이는
당김이다(Divisors on Spaces, Lemma
\ref{spaces-divisors-lemma-strict-transform-flat}).
이로써 $\mathcal{F}'$을 $B' \times_B X$의 에탈 덮개에 제한한 것이
원하는 성질들을 가짐을 알 수 있고, 정리가 증명되었다.
\end{proof}







\section{응용}
\label{spaces-more-morphisms-section-applications-flattening-by-blowing-up}

\noindent
이 절에서는 블로업에 의한 평탄화 결과를 응용한다.

\begin{lemma}
\label{spaces-more-morphisms-lemma-flat-after-blowing-up}
$S$를 스킴이라 하자. $f : X \to B$를 $S$ 위의 대수공간들의
사상이라 하자. $U \subset B$를 열린 부분공간이라 하자. 다음을 가정하자.
\begin{enumerate}
\item $B$는 준콤팩트이고 준분리이다.
\item $U$는 준콤팩트이다.
\item $f : X \to B$는 유한형이고 준분리이다.
\item $f^{-1}(U) \to U$는 평탄하고 국소 유한 표시이다.
\end{enumerate}
그러면 $U$-허용 블로업 $B' \to B$가 존재하여 진변환
$X'$, 즉 $X$의 진변환은 $B'$ 위에서 평탄하고 유한 표시이다.
\end{lemma}

\begin{proof}
$B' \to B$를 $U$-허용 블로업이라 하자. $X$의 진변환은 $B'$ 위에서
준콤팩트이고 준분리임을 주목하자. 이는 $X$가 $B$ 위에서 준콤팩트이고
준분리이기 때문이다. 따라서 진변환이 평탄하고 국소 유한 표시가 되게 하는
$U$-허용 블로업을 찾는 일만 고려하면 된다. $X$가 $B$ 위에서 국소 유한
표시가 아니므로 Theorem \ref{spaces-more-morphisms-theorem-flatten-module}을 곧바로 적용할 수 없다.

\medskip\noindent
아핀 스킴 $V$와 전사 에탈 사상 $V \to X$를 택하자. (이는 $X$가
준콤팩트 공간 $B$ 위의 유한형 공간이라 준콤팩트이므로 가능하다.)
그러면 사상 $V \to B$에 대해 결과를 보이면 충분하다. 진변환은 에탈
국소화와 가환하기 때문이다. 다음을 보라. Divisors on Spaces, Lemma
\ref{spaces-divisors-lemma-strict-transform-local}.
따라서 $X \to B$가 유한형일 뿐 아니라 분리라고 가정해도 된다.
이 경우 닫힌 몰입 $i : X \to Y$를 찾을 수 있으며 $Y \to B$는
분리이고 유한 표시이다. 다음을 보라. Limits of Spaces, Proposition
\ref{spaces-limits-proposition-separated-closed-in-finite-presentation}.

\medskip\noindent
Theorem \ref{spaces-more-morphisms-theorem-flatten-module}을
$\mathcal{F} = i_*\mathcal{O}_X$에 적용하자. 여기서 이 가군은
$Y/B$ 위에 있다. 그러면 $U$-허용 블로업
$B' \to B$를 찾을 수 있어서 $\mathcal{F}$의 진변환은 $B'$ 위에서
평탄하고 유한 표시이다. $X'$을 $X$의 진변환이라 하되 블로업
$B' \to B$ 아래에서 취하자. $i' : X' \to Y \times_B B'$를 유도된
사상이라 하자. 진변환을 취하는 것은 아핀 사상에 따른 앞으로 보내기와
가환하므로(Divisors on Spaces, Lemma
\ref{spaces-divisors-lemma-strict-transform-affine}),
$i'_*\mathcal{O}_{X'}$은 $B'$ 위에서 평탄하고
$\mathcal{O}_{Y \times_B B'}$-가군으로서 유한 표시이다.
따라서 $X' \to B'$는 평탄하고 국소 유한 표시이다.
앞의 관찰에 의해 이것으로 보조정리가 증명된다.
\end{proof}

\begin{lemma}
\label{spaces-more-morphisms-lemma-finite-after-blowing-up}
$S$를 스킴이라 하자. $f : X \to B$를 $S$ 위의 대수공간들의
사상이라 하자. $U \subset B$를 열린 부분공간이라 하자. 다음을 가정하자.
\begin{enumerate}
\item $B$는 준콤팩트이고 준분리이다.
\item $U$는 준콤팩트이다.
\item $f : X \to B$는 고유이다.
\item $f^{-1}(U) \to U$는 유한 국소 자유이다.
\end{enumerate}
그러면 $U$-허용 블로업 $B' \to B$가 존재하여 진변환
$X'$, 즉 $X$의 진변환은 $B'$ 위에서 유한 국소 자유이다.
\end{lemma}

\begin{proof}
Lemma \ref{spaces-more-morphisms-lemma-flat-after-blowing-up}에 의해 $X \to B$가 평탄하고
유한 표시라고 가정해도 된다. 필요하면 $B$를 $U$-허용 블로업으로
바꾸어 $U \subset B$가 스킴론적으로 조밀하다고 가정해도 된다.
그러면 Lemma \ref{spaces-more-morphisms-lemma-proper-flat-finite-over-dense-open}에 의해
$f$는 유한이다. 따라서 Morphisms of Spaces, Lemma
\ref{spaces-morphisms-lemma-finite-flat}에 의해 $f$는 유한 국소 자유이다.
\end{proof}

\begin{lemma}
\label{spaces-more-morphisms-lemma-isomorphism-after-blowing-up}
$S$를 스킴이라 하자. $f : X \to B$를 $S$ 위의 대수공간들의
사상이라 하자. $U \subset B$를 열린 부분공간이라 하자. 다음을 가정하자.
\begin{enumerate}
\item $B$는 준콤팩트이고 준분리이다.
\item $U$는 준콤팩트이다.
\item $f : X \to B$는 고유이다.
\item $f^{-1}(U) \to U$는 동형이다.
\end{enumerate}
그러면 $U$-허용 블로업 $B' \to B$가 존재하여 진변환
$X'$, 즉 $X$의 진변환은 $B'$에 동형으로 사상된다.
\end{lemma}

\begin{proof}
Lemma \ref{spaces-more-morphisms-lemma-flat-after-blowing-up}에 의해 $X \to B$가 평탄하고
유한 표시라고 가정해도 된다. 필요하면 $B$를 $U$-허용 블로업으로
바꾸어 $U \subset B$가 스킴론적으로 조밀하다고 가정해도 된다.
그러면 Lemma \ref{spaces-more-morphisms-lemma-proper-flat-finite-over-dense-open}에 의해
$f$는 유한이고 Lemma \ref{spaces-more-morphisms-lemma-zariski}에 의해 열린 몰입이다.
따라서 $f$는 상이 닫혀 있고 조밀한 열린 부분공간 $U$를 포함하는
열린 몰입이므로, $f$는 동형이다.
\end{proof}

\begin{lemma}
\label{spaces-more-morphisms-lemma-zariski-after-blowup}
$S$를 스킴이라 하자. $f : X \to B$를 $S$ 위의 대수공간들의
사상이라 하자. $U \subset B$를 열린 부분공간이라 하자. 다음을 가정하자.
\begin{enumerate}
\item $B$는 준콤팩트이고 준분리이다.
% P09-ERR-0698: frozen source omits the copula in this hypothesis.
\item $U$는 준콤팩트이다.
\item $f$는 유한형이다.
\item $f^{-1}(U) \to U$는 동형이다.
\end{enumerate}
그러면 $U$-허용 블로업 $B' \to B$가 존재하여 $U$는 $B'$ 안에서
스킴론적으로 조밀하고, 진변환 $X'$, 즉 $X$의 진변환은 $B'$의
열린 부분공간에 동형으로 사상된다.
\end{lemma}

\begin{proof}
이 보조정리는 Lemma \ref{spaces-more-morphisms-lemma-isomorphism-after-blowing-up}의 일반화이다.
$U$-허용 블로업들의 합성은 $U$-허용이므로(Divisors on Spaces, Lemma
\ref{spaces-divisors-lemma-composition-admissible-blowups}), 단계별로
진행할 수 있다. 유한형 준연접 아이디얼 층
$\mathcal{I} \subset \mathcal{O}_B$를 택하여
$|B| \setminus |U| = |V(\mathcal{I})|$이 되게 하자. $B$를
$B$의 블로업으로 바꾸되 $\mathcal{I}$에서 취하고, $X$를 $X$의
진변환으로 바꾸자. 이 교체 뒤에 $B \setminus U$는 유효 카르티에
약수 $D$의 지지집합이다(Divisors on Spaces, Lemma
\ref{spaces-divisors-lemma-blowing-up-gives-effective-Cartier-divisor}).
특히 $U$는 $B$ 안에서 스킴론적으로 조밀하다(Divisors on Spaces, Lemma
\ref{spaces-divisors-lemma-complement-effective-Cartier-divisor}).
다음으로 $U$-허용 블로업을 하나 더 취하여 $X \to B$가 평탄하고
유한 표시인 상황에 이른다. 다음을 보라. Lemma
\ref{spaces-more-morphisms-lemma-flat-after-blowing-up}. $U$는 여전히 $B$ 안에서
스킴론적으로 조밀함을 주목하자. 따라서 Lemma \ref{spaces-more-morphisms-lemma-zariski}에 의해
$X \to B$는 열린 몰입이다.
\end{proof}

\noindent
다음 보조정리는 수정이 블로업에 의해 지배될 수 있음을 말한다.

\begin{lemma}
\label{spaces-more-morphisms-lemma-dominate-modification-by-blowup}
$S$를 스킴이라 하자. $f : X \to B$를 $S$ 위의 대수공간들의
사상이라 하자. $U \subset B$를 열린 부분공간이라 하자. 다음을 가정하자.
\begin{enumerate}
\item $B$는 준콤팩트이고 준분리이다.
\item $U$는 준콤팩트이다.
\item $f : X \to B$는 고유이다.
\item $f^{-1}(U) \to U$는 동형이다.
% P09-ERR-0699: frozen source says “us an isomorphism”.
\end{enumerate}
그러면 $U$-허용 블로업 $B' \to B$가 존재하여 $X$를 지배한다.
즉, 이 블로업 사상의 분해 $B' \to X \to B$가 존재한다.
\end{lemma}

\begin{proof}
Lemma \ref{spaces-more-morphisms-lemma-isomorphism-after-blowing-up}에 의해 $U$-허용 블로업
$B' \to B$를 찾아서 진변환 $X'$이 $B'$에 동형으로 사상되게 할 수 있다.
그러면 $B' = X' \to X$를 이 분해로 사용할 수 있다.
\end{proof}

\begin{lemma}
\label{spaces-more-morphisms-lemma-get-section-after-blowup}
$S$를 스킴이라 하자. $X$, $Y$를 $S$ 위의 대수공간들이라 하자.
$U \subset W \subset Y$를 열린 부분공간들이라 하자.
$f : X \to W$와 $s : U \to X$를 사상들로서
$f \circ s = \text{id}_U$인 것들이라 하자. 다음을 가정하자.
\begin{enumerate}
\item $f$는 고유이다.
\item $Y$는 준콤팩트이고 준분리이다.
\item $U$와 $W$는 준콤팩트이다.
\end{enumerate}
그러면 $U$-허용 블로업 $b : Y' \to Y$와 사상
$s' : b^{-1}(W) \to X$가 존재하며, 이는 $s$를 연장하고
$f \circ s' = b|_{b^{-1}(W)}$이다.
\end{lemma}

\begin{proof}
$X$를 $s$의 스킴론적 상으로 바꾸어도 되며 그렇게 하자.
그러면 $X \to W$는 $U$ 위에서 동형이다. 다음을 보라.
Morphisms of Spaces, Lemma
\ref{spaces-morphisms-lemma-scheme-theoretic-image-of-partial-section}.
Lemma \ref{spaces-more-morphisms-lemma-dominate-modification-by-blowup}에 의해
$U$-허용 블로업 $W' \to W$와 $W' \to X$가 존재하며, 후자는
$s$의 연장이다.
Divisors on Spaces, Lemma
\ref{spaces-divisors-lemma-extend-admissible-blowups}를 적용하여
$W' \to W$를 $U$-허용 블로업으로 연장하되 $Y$ 위에서 취하면
증명이 끝난다.
\end{proof}















\section{차우 보조정리}
\label{spaces-more-morphisms-section-chow}

\noindent
이 절에서는 차우 보조정리를 증명한다
(Lemma \ref{spaces-more-morphisms-lemma-chow-noetherian-separated}).
쉽게 증명되면서 많은 응용에 충분한 차우 보조정리의 약한 판은
Cohomology of Spaces, Section
\ref{spaces-cohomology-section-weak-chow}를 참조하기 바란다.

\medskip\noindent
대수공간의 사영 사상을 아직 정의하지 않았으므로, 보조정리들의 명제
(예를 들어 Lemma \ref{spaces-more-morphisms-lemma-blowup-to-find-embedding})에는 사영 사상 대신
표현 가능한 고유 사상이 등장한다.

\begin{lemma}
\label{spaces-more-morphisms-lemma-find-common-blowups}
$S$를 스킴이라 하자. $Y$를 $S$ 위의 준콤팩트이고 준분리인
대수공간이라 하자. $U \to X_1$과 $U \to X_2$를 $Y$ 위 대수공간들의
열린 몰입이라 하고, $U$, $X_1$, $X_2$가 $Y$ 위에서 유한형이고
분리라고 가정하자.
% P09-ERR-0700: frozen source omits the copula in this coordinated hypothesis.
그러면 다음 가환 그림이 존재한다.
$$
\xymatrix{
X_1' \ar[d] \ar[r] & X & X_2' \ar[l] \ar[d] \\
X_1 & U \ar[l] \ar[lu] \ar[u] \ar[ru] \ar[r] & X_2
}
$$
이는 $Y$ 위 대수공간들의 그림이며, $X_i' \to X_i$는
$U$-허용 블로업이고, $X_i' \to X$는 열린 몰입이며,
$X$는 $Y$ 위에서 분리이고 유한형이다.
\end{lemma}

\begin{proof}
증명 전체에서 모든 대수공간은 $Y$ 위에서 분리이고 유한형이다.
% P09-ERR-0701: frozen source omits the conjunction between these properties.
특히 이 대수공간들은 준콤팩트이고 준분리이며, 이들 사이의 사상은
준콤팩트이고 분리이다. 다음을 보라. Morphisms of Spaces, Sections
\ref{spaces-morphisms-section-separation-axioms} and
\ref{spaces-morphisms-section-quasi-compact}.
이러한 공간들의 몰입 $U \to W$가 $Y$ 위에 있으면, 스킴론적 상
$Z$, 즉 $U$의 $W$ 안에서의 스킴론적 상은 $W$의 닫힌 부분공간이고,
$U \to Z$는 열린 몰입이며,
$U \subset Z$는 스킴론적으로 조밀하고 $|U| \subset |Z|$는 조밀하다는
사실을 사용하겠다. 다음을 보라. Morphisms of Spaces, Lemma
\ref{spaces-morphisms-lemma-quasi-compact-immersion}.

\medskip\noindent
$X_{12} \subset X_1 \times_Y X_2$를 $U \to X_1 \times_Y X_2$의
스킴론적 상이라 하자. 사영 $p_i : X_{12} \to X_i$는 다음에 의해
동형 $p_i^{-1}(U) \to U$를 유도한다. Morphisms of Spaces, Lemma
\ref{spaces-morphisms-lemma-scheme-theoretic-image-of-partial-section}.
$U$-허용 블로업 $X_i^i \to X_i$를 택하여 진변환 $X_{12}^i$, 즉
$X_{12}$의 진변환이 $X_i^i$의 열린 부분공간에 동형이 되게 하자.
다음을 보라.
Lemma \ref{spaces-more-morphisms-lemma-zariski-after-blowup}.
$\mathcal{I}_i \subset \mathcal{O}_{X_i}$를 이에 대응하는 유한형
준연접 아이디얼 층이라 하자. $X_{12}^i \to X_{12}$는
$p_i^{-1}\mathcal{I}_i \mathcal{O}_{X_{12}}$에서의 블로업임을 상기하자.
다음을 보라. Divisors on Spaces, Lemma
\ref{spaces-divisors-lemma-strict-transform}.
$X_{12}'$을 $X_{12}$의
$p_1^{-1}\mathcal{I}_1 p_2^{-1}\mathcal{I}_2 \mathcal{O}_{X_{12}}$에서의
블로업이라 하자. 이것이 뜻하는 바는 Divisors on Spaces, Lemma
\ref{spaces-divisors-lemma-blowing-up-two-ideals}를 보라.
다음 가환 그림을 얻는다.
$$
\xymatrix{
X_{12}' \ar[d] \ar[r] & X_{12}^2 \ar[d] \\
X_{12}^1 \ar[r] & X_{12}
}
$$
여기서 모든 사상은 $U$-허용 블로업이다. $X_{12}^i \subset X_i^i$가
열린 부분공간이므로, $U$-허용 블로업 $X_i' \to X_i^i$를 택할 수 있으며
그 제한은 $X_{12}' \to X_{12}^i$이다. 다음을 보라.
Divisors on Spaces, Lemma
\ref{spaces-divisors-lemma-extend-admissible-blowups}.
그러면 $X_{12}' \subset X_i'$은 열린 부분공간이고 다음 그림은 가환이다.
$$
\xymatrix{
X_{12}' \ar[d] \ar[r] & X_i' \ar[d] \\
X_{12}^i \ar[r] & X_i^i
}
$$
세로 화살표들은 블로업이고 가로 화살표들은 열린 몰입이다.
% P09-ERR-0702: frozen source twice uses the malformed plural “blowing ups”.
$X'_{12} \to X_1' \times_Y X_2'$은 몰입이자 고유 사상임을 주목하자.
실제로 $X'_{12} \to X_{12}$는 고유이고,
$X_{12} \to X_1 \times_Y X_2$는 닫힌 몰입이며,
$X_1' \times_Y X_2' \to X_1 \times_Y X_2$는 분리라는 사실과
Morphisms of Spaces, Lemma
\ref{spaces-morphisms-lemma-universally-closed-permanence}를 사용한다.
따라서 $X'_{12} \to  X_1' \times_Y X_2'$은 닫힌 몰입이다.
$X$를 $X_1'$과 $X_2'$을 공통 열린 부분공간 $X_{12}'$을 따라 붙인
것으로 정의하면, $X \to Y$는 유한형이고 분리이다\footnote{대각 사상
$X \to X \times_Y X$의 닫힘은 네 열린 부분
$X'_i \times_Y X'_j$, 즉 $X \times_Y X$의 이 열린 부분들 위에서
확인할 수 있고, 거기서는 명백하기 때문이다.}.
$U$-허용 블로업들의 합성은 $U$-허용 블로업이므로
(Divisors on Spaces, Lemma
\ref{spaces-divisors-lemma-composition-admissible-blowups}),
보조정리가 증명되었다.
\end{proof}

\begin{lemma}
\label{spaces-more-morphisms-lemma-blowup-to-find-embedding}
$S$를 스킴이라 하자. $f : X \to Y$를 $S$ 위의 대수공간들의
사상이라 하자. $U \subset X$를 열린 부분공간이라 하자. 다음을 가정하자.
\begin{enumerate}
\item $U$는 준콤팩트이다.
\item $Y$는 준콤팩트이고 준분리이다.
\item 몰입 $U \to \mathbf{P}^n_Y$가 $Y$ 위에 존재한다.
\item $f$는 유한형이고 분리이다.
\end{enumerate}
그러면 다음 가환 그림이 존재한다.
$$
\xymatrix{
& U \ar[ld] \ar[d] \ar[rd] \ar[rrd] \\
X \ar[rd] & X' \ar[l] \ar[d] \ar[r] & Z' \ar[ld] \ar[r] & Z \ar[ld] \\
& Y & \mathbf{P}^n_Y \ar[l]
}
$$
여기서 $U$에서 나가는 화살표들은 열린 몰입이고,
$X' \to X$는 $U$-허용 블로업이며,
$X' \to Z'$은 열린 몰입이고,
$Z' \to Y$는 대수공간들의 표현 가능한 고유 사상이다.
더 정확히는 $Z' \to Z$가 $U$-허용 블로업이고
$Z \to \mathbf{P}^n_Y$는 닫힌 몰입이다.
\end{lemma}

\begin{proof}
$Z \subset \mathbf{P}^n_Y$를 몰입 $U \to \mathbf{P}^n_Y$의
스킴론적 상이라 하자. $U \to \mathbf{P}^n_Y$가 준콤팩트이므로
$U \subset Z$는 (스킴론적으로) 조밀한 열린 부분공간이다
(Morphisms of Spaces, Lemma
\ref{spaces-morphisms-lemma-quasi-compact-immersion}).
Lemma \ref{spaces-more-morphisms-lemma-find-common-blowups}를 적용하여 다음 그림을 얻자.
$$
\xymatrix{
X' \ar[d] \ar[r] & \overline{X}' & Z' \ar[l] \ar[d] \\
X & U \ar[l] \ar[lu] \ar[u] \ar[ru] \ar[r] & Z
}
$$
이 그림은 그 보조정리의 명제에 열거된 성질들을 가진다.
$X' \to X$와 $Z' \to Z$가 $U$-허용 블로업이므로 $U$는
$X'$과 $Z'$ 모두의 스킴론적으로 조밀한 열린 부분공간이다. 다음을 보라.
Divisors on Spaces, Lemmas
\ref{spaces-divisors-lemma-blowing-up-gives-effective-Cartier-divisor} and
\ref{spaces-divisors-lemma-complement-effective-Cartier-divisor}.
$Z' \to Z \to Y$가 고유이므로 $Z' \subset \overline{X}'$은
닫힌 부분공간이다. 다음을 보라. Morphisms of Spaces, Lemma
\ref{spaces-morphisms-lemma-universally-closed-permanence}.
따라서 스킴론적으로 $X' \subset Z'$이고, 그러므로 $X'$은
$Z'$의 열린 부분공간이다(작은 세부사항은 생략한다).
% P09-ERR-0703: frozen source explicitly omits the decisive local argument.
이것으로 보조정리가 증명된다.
\end{proof}

\begin{lemma}
\label{spaces-more-morphisms-lemma-chow-noetherian}
$S$를 스킴이라 하자. $f : X \to Y$를 $S$ 위의 대수공간들의
사상이라 하자. $f$가 분리이고 유한형이며 $Y$가 뇌터라고 가정하자.
% P09-ERR-0704: frozen source omits both copulas in this hypothesis.
그러면 조밀한 열린 부분공간 $U \subset X$와 다음 가환 그림이 존재한다.
$$
\xymatrix{
& U \ar[ld] \ar[d] \ar[rd] \ar[rrd] \\
X \ar[rd] & X' \ar[l] \ar[d] \ar[r] & Z' \ar[ld] \ar[r] & Z \ar[ld] \\
& Y & \mathbf{P}^n_Y \ar[l]
}
$$
여기서 $U$에서 나가는 화살표들은 열린 몰입이고,
$X' \to X$는 $U$-허용 블로업이며,
$X' \to Z'$은 열린 몰입이고,
$Z' \to Y$는 대수공간들의 표현 가능한 고유 사상이다.
더 정확히는 $Z' \to Z$가 $U$-허용 블로업이고
$Z \to \mathbf{P}^n_Y$는 닫힌 몰입이다.
\end{lemma}

\begin{proof}
Limits of Spaces, Lemma
\ref{spaces-limits-lemma-embedding-into-affine-over-qs}에 의해
조밀한 열린 부분공간 $U \subset X$와 몰입
$U \to \mathbf{A}^n_Y$가 $Y$ 위에 존재한다. 열린 몰입
$\mathbf{A}^n_Y \to \mathbf{P}^n_Y$와 합성하면 Lemma
\ref{spaces-more-morphisms-lemma-blowup-to-find-embedding}의 상황을 얻고, 결과가 따른다.
\end{proof}

\begin{remark}
\label{spaces-more-morphisms-remark-chow-Noetherian}
Lemmas \ref{spaces-more-morphisms-lemma-blowup-to-find-embedding} and
\ref{spaces-more-morphisms-lemma-chow-noetherian}에서 사상 $g : Z' \to Y$는 사영 사상들의
합성이다. 아마도 대수공간에 대한 Morphisms, Lemma
\ref{morphisms-lemma-ample-composition}의 유사한 결과에 의해
$g$-풍부한 가역층이 $Z'$ 위에 존재할 것이다. 이것이 필요해지면
여기에서 명제를 진술하고 증명하겠다.
\end{remark}

\noindent
다음 결과는 \cite[IV Theorem 3.1]{Kn}이다. 몰입
$X' \to \mathbf{P}^n_Y$는 준콤팩트이므로
$X' \to Z' \to \mathbf{P}^n_Y$로 분해할 수 있으며, 첫 사상은
열린 몰입이고 둘째 사상은 닫힌 몰입이다(Morphisms of Spaces, Lemma
\ref{spaces-morphisms-lemma-quasi-compact-immersion}).

\begin{lemma}[차우 보조정리]
\label{spaces-more-morphisms-lemma-chow-noetherian-separated}
\begin{reference}
\cite[IV Theorem 3.1]{Kn}
\end{reference}
$S$를 스킴이라 하자. $f : X \to Y$를 $S$ 위의 대수공간들의
사상이라 하자. $f$가 분리이고 유한형이며 $Y$가 분리이고 뇌터라고
가정하자.
% P09-ERR-0705: frozen source omits the copula and conjunction for f.
그러면 다음 가환 그림이 존재한다.
$$
\xymatrix{
X \ar[rd] & X' \ar[l] \ar[d] \ar[r] & \mathbf{P}^n_Y \ar[ld] \\
& Y
}
$$
여기서 $X' \to X$는 $U$-허용 블로업이고, 이때
$U \subset X$는 어떤 조밀한 열린 부분공간이며, 사상
$X' \to \mathbf{P}^n_Y$는 몰입이다.
\end{lemma}

\begin{proof}
증명의 첫 문단에서는 보조정리를 $Y$가 $\Spec(\mathbf{Z})$ 위에서
유한형인 경우로 환원한다. 밑스킴 $S$를 $\Spec(\mathbf{Z})$로
바꾸어도 되며 그렇게 하자. $Y = \lim Y_i$를
$\Spec(\mathbf{Z})$ 위 유한형인 분리 대수공간들의 유향 극한으로 쓸 수 있다.
다음을 보라. Limits of Spaces, Proposition
\ref{spaces-limits-proposition-approximate} and Lemma
\ref{spaces-limits-lemma-descend-separated}.
충분히 큰 모든 $i$에 대해 분리인 유한형 사상 $X_i \to Y_i$를 찾아서
$X = Y \times_{Y_i} X_i$가 되게 할 수 있다. 다음을 보라.
Limits of Spaces, Lemmas
\ref{spaces-limits-lemma-descend-finite-presentation} and
\ref{spaces-limits-lemma-descend-separated-morphism}.
$\eta_1, \ldots, \eta_n$을 $|X|$의 기약 성분들의 일반점들이라 하자.
($X$는 뇌터 대수공간 $Y$ 위의 분리인 유한형 대수공간이므로 뇌터이고,
따라서 $|X|$는 뇌터 위상공간이다.)
Limits of Spaces, Lemma \ref{spaces-limits-lemma-topology-limit}에 의해
$\eta_1, \ldots, \eta_n$의 $|X_i|$ 안에서의 상들은 충분히 큰 $i$에
대해 서로 다르다. $X_i$를 사상 $X \to X_i$의 스킴론적 상으로
바꾸어도 된다. 이 사상은 준콤팩트이며, 사실 아핀이다.
이렇게 바꾼 뒤에 $\eta_1, \ldots, \eta_n$의 $|X_i|$ 안에서의 상들은
$|X_i|$의 기약 성분들의 일반점들이다. 다음을 보라.
Morphisms of Spaces, Lemma
\ref{spaces-morphisms-lemma-quasi-compact-scheme-theoretic-image}.
이제 다음과 같은 그림을 찾을 수 있다고 하자.
$$
\xymatrix{
X_i \ar[rd] & X_i' \ar[l] \ar[d] \ar[r] & \mathbf{P}^n_{Y_i} \ar[ld] \\
& Y
}
$$
% P09-ERR-0706: frozen source uses Y below objects and arrows over Y_i.
여기서 $X_i' \to X_i$는 $U_i$-허용 블로업이고, 이때
$U_i \subset X_i$는 어떤 조밀한 열린 부분공간이며, 사상
$X_i' \to \mathbf{P}^n_{Y_i}$는 몰입이다. 그러면
진변환 $X' \to X$, 즉 $X$의 $X_i' \to X_i$에 관한 진변환은
$U$-허용 블로업이다. 여기서 $U \subset X$는 $U_i$의 $X$ 안에서의
역상이다. 지표 $i$를 주의 깊게 택했으므로
$\eta_1, \ldots, \eta_n \in |U|$이고 $U \subset X$는 조밀하다.
더욱이 $X' \to \mathbf{P}^n_Y$는 몰입이다. 실제로 $X'$은
$X_i' \times_{X_i} X = X_i' \times_{Y_i} Y$ 안에서 닫혀 있고,
후자는 $\mathbf{P}^n_Y$ 안으로의 몰입을 갖는다. 이로써 다음 문단의
상황으로 환원했다.

\medskip\noindent
$Y$가 $\Spec(\mathbf{Z})$ 위에서 분리이고 유한형이라고 가정하자.
그러면 $X \to \Spec(\mathbf{Z})$도 분리이고 유한형이다.
Lemma \ref{spaces-more-morphisms-lemma-chow-noetherian}을 $X \to \Spec(\mathbf{Z})$에 적용하여
조밀한 열린 부분공간 $U \subset X$와 다음 가환 그림을 얻자.
$$
\xymatrix{
& U \ar[ld] \ar[d] \ar[rd] \ar[rrd] \\
X \ar[rd] & X' \ar[l] \ar[d] \ar[r] & Z' \ar[ld] \ar[r] & Z \ar[ld] \\
& \Spec(\mathbf{Z}) & \mathbf{P}^n_\mathbf{Z} \ar[l]
}
$$
이 그림은 보조정리에 열거된 모든 성질을 가진다.
$Z$에는 풍부한 가역층, 즉
$\mathcal{O}_{\mathbf{P}^n}(1)|_Z$가 있음을 주목하자. 따라서
Morphisms, Lemma
\ref{morphisms-lemma-projective-over-quasi-projective-is-H-projective}에 의해
$Z' \to Z$는 H-사영 사상이다. 그러므로 Morphisms, Lemma
\ref{morphisms-lemma-H-projective-composition}에 의해
$Z' \to \Spec(\mathbf{Z})$는 H-사영이다. 따라서 닫힌 몰입
$Z' \to \mathbf{P}^m_{\Spec(\mathbf{Z})}$가 어떤 $m \geq 0$에 대해
존재한다.
이에 따라 대각 사상
$$
X' \to Y \times \mathbf{P}^m_\mathbf{Z} = \mathbf{P}^m_Y
$$
은 몰입이며, 이는 $\mathbf{P}^m_\mathbf{Z}$로의 사영과 합성한 것이
몰입이기 때문이다. 이것으로 증명이 끝난다.
\end{proof}






\section{차우 보조정리의 변형들}
\label{spaces-more-morphisms-section-chows-lemma}

\noindent
이 절에서는 비뇌터 대수공간들 사이의 사상을 다루는 차우 보조정리의
몇 가지 변형을 증명한다. 뇌터 판들은 Lemma
\ref{spaces-more-morphisms-lemma-chow-noetherian}과 Lemma
\ref{spaces-more-morphisms-lemma-chow-noetherian-separated}이다.

\begin{lemma}
\label{spaces-more-morphisms-lemma-chow-finite-type}
$S$를 스킴이라 하자. $Y$를 $S$ 위의 준콤팩트이고 준분리인
대수공간이라 하자. $f : X \to Y$를 유한형인 분리 사상이라 하자.
그러면 다음 가환 그림이 존재한다.
$$
\xymatrix{
X \ar[rd] & X' \ar[l] \ar[d] \ar[r] & \overline{X}' \ar[ld] \\
& Y
}
$$
여기서 $X' \to X$는 고유 전사이고,
$X' \to \overline{X}'$은 열린 몰입이며,
$\overline{X}' \to Y$는 대수공간들의 표현 가능한 고유 사상이다.
% P09-ERR-0707: frozen source omits the article before “morphism”.
\end{lemma}

\begin{proof}
Limits of Spaces, Proposition
\ref{spaces-limits-proposition-separated-closed-in-finite-presentation}에 의해
닫힌 몰입 $X \to X_1$을 찾을 수 있으며, 여기서 $X_1$은 $Y$ 위에서
분리이고 유한 표시이다. $X_1 \to Y$에 대해 주장을 증명하면 $X$에
대한 결과가 따름은 분명하다. 따라서 $X$가 $Y$ 위에서 유한 표시라고
가정해도 된다.

\medskip\noindent
밑스킴 $S$를 $\Spec(\mathbf{Z})$로 바꾸어도 되며 그렇게 하자.
$Y = \lim_i Y_i$를 $\Spec(\mathbf{Z})$ 위 유한형인 준분리
대수공간들의 유향 극한으로 쓰자. 다음을 보라. Limits of Spaces,
Proposition \ref{spaces-limits-proposition-approximate}.
Limits of Spaces, Lemma
\ref{spaces-limits-lemma-descend-finite-presentation}에 의해 지표
$i \in I$와 유한 표시인 스킴 $X_i \to Y_i$를 찾아서
% P09-ERR-0708: frozen source calls this descended algebraic space a scheme.
$X = Y \times_{Y_i} X_i$가 되게 할 수 있다. Limits of Spaces, Lemma
\ref{spaces-limits-lemma-descend-separated-morphism}에 의해
$X_i \to Y_i$가 분리라고 가정해도 된다. $X_i$를 $Y_i$ 위에서
다룬 주장이 성립하면 $X$에 대한 주장도 성립함은 분명하다.
$X_i \to Y_i$의 경우는 Lemma \ref{spaces-more-morphisms-lemma-chow-noetherian}에서 다룬다.
\end{proof}

\begin{lemma}
\label{spaces-more-morphisms-lemma-chow-finite-type-separated}
$S$를 스킴이라 하자. $f : X \to Y$를 $S$ 위의 대수공간들의
사상이라 하자. $f$가 분리이고 유한형이며 $Y$가 분리이고 준콤팩트라고
가정하자.
% P09-ERR-0709: frozen source omits the copula and conjunction for f.
그러면 다음 가환 그림이 존재한다.
$$
\xymatrix{
X \ar[rd] & X' \ar[l] \ar[d] \ar[r] & \mathbf{P}^n_Y \ar[ld] \\
& Y
}
$$
여기서 $X' \to X$는 고유 전사 사상이고,
% P09-ERR-0710: frozen source omits the article before “proper surjective morphism”.
$X' \to \mathbf{P}^n_Y$는 몰입이다.
\end{lemma}

\begin{proof}
Limits of Spaces, Proposition
\ref{spaces-limits-proposition-separated-closed-in-finite-presentation}에 의해
닫힌 몰입 $X \to X_1$을 찾을 수 있으며, 여기서 $X_1$은 $Y$ 위에서
분리이고 유한 표시이다. $X_1 \to Y$에 대해 주장을 증명하면 $X$에
대한 결과가 따름은 분명하다. 따라서 $X$가 $Y$ 위에서 유한 표시라고
가정해도 된다.

\medskip\noindent
밑스킴 $S$를 $\Spec(\mathbf{Z})$로 바꾸어도 되며 그렇게 하자.
$Y = \lim_i Y_i$를 $\Spec(\mathbf{Z})$ 위 유한형인 준분리
대수공간들의 유향 극한으로 쓰자. 다음을 보라. Limits of Spaces,
Proposition \ref{spaces-limits-proposition-approximate}.
Limits of Spaces, Lemma \ref{spaces-limits-lemma-descend-separated}에 의해
$Y_i$가 모든 $i$에 대해 분리라고 가정해도 된다.
Limits of Spaces, Lemma
\ref{spaces-limits-lemma-descend-finite-presentation}에 의해 지표
$i \in I$와 유한 표시인 스킴 $X_i \to Y_i$를 찾아서
% P09-ERR-0711: frozen source again calls the descended algebraic space a scheme.
$X = Y \times_{Y_i} X_i$가 되게 할 수 있다. Limits of Spaces, Lemma
\ref{spaces-limits-lemma-descend-separated-morphism}에 의해
$X_i \to Y_i$가 분리라고 가정해도 된다. $X_i$를 $Y_i$ 위에서
다룬 주장이 성립하면 $X$에 대한 주장도 성립함은 분명하다.
$X_i \to Y_i$의 경우는 Lemma
\ref{spaces-more-morphisms-lemma-chow-noetherian-separated}에서 다룬다.
\end{proof}






\section{그로텐디크 존재 정리}
\label{spaces-more-morphisms-section-existence-proper}

\noindent
이 절에서는 대수공간에 대한 그로텐디크 존재 정리를 논의한다.
``형식적 대수공간'' 이론을 전개하는 대신, ``뇌터 아딕 형식 공간 위의
연접 가군'' 개념을 대신할 약간의 언어를 잠시 전개한다.

\medskip\noindent
$S$를 스킴이라 하자. $X$를 $S$ 위의 뇌터 대수공간이라 하자.
$\mathcal{I} \subset \mathcal{O}_X$를 준연접 아이디얼 층이라 하자.
이하에서는 다음을 만족하는 역계 $(\mathcal{F}_n)$, 즉 연접
$\mathcal{O}_X$-가군들의 역계를 고려한다.
\begin{enumerate}
\item $\mathcal{F}_n$은 $\mathcal{I}^n$에 의해 소멸한다.
\item 전이 사상은 동형
$\mathcal{F}_{n + 1}/\mathcal{I}^n\mathcal{F}_{n + 1} \to \mathcal{F}_n$을
유도한다.
\end{enumerate}
이러한 역계들의 사상
$\alpha : (\mathcal{F}_n) \to (\mathcal{G}_n)$은 단순히 서로 호환되는
사상계 $\alpha_n : \mathcal{F}_n \to \mathcal{G}_n$이다.
이 역계들의 범주를 $\textit{Coh}(X, \mathcal{I})$로 나타내자.
이 계들에 관한 이론을 전개할 것인데, 이는 스킴의 경우에 해당하는 결과들과
평행하다.
% P09-ERR-0712: frozen source combines “parallel” with an extraneous “to”.
다음을 보라. Cohomology of Schemes, Sections
\ref{coherent-section-existence},
\ref{coherent-section-existence-proper},
\ref{coherent-section-existence-proper-support}, and
\ref{coherent-section-algebraization}.

\medskip\noindent
함자성. $f : X \to Y$를 스킴 $S$ 위의 뇌터 대수공간들의 사상이라 하고,
$\mathcal{J} \subset \mathcal{O}_Y$를 준연접 아이디얼 층이라 하자.
$\mathcal{I} = f^{-1}\mathcal{J}\mathcal{O}_X$로 놓자. 이 상황에는 함자
$$
f^* : \textit{Coh}(Y, \mathcal{J}) \longrightarrow \textit{Coh}(X, \mathcal{I})
$$
가 있으며, 이는 $(\mathcal{G}_n)$을 $(f^*\mathcal{G}_n)$으로 보낸다.
다음과 비교하라. Cohomology of Schemes, Lemma
\ref{coherent-lemma-inverse-systems-pullback}.
$f$가 에탈이면 이를 단순히 계를 $X$에 제한한 것으로 생각할 수 있다.
다음을 보라. Properties of Spaces,
Equation \ref{spaces-properties-equation-restrict-modules}.

\medskip\noindent
에탈 하강. $S$를 스킴이라 하자. $U_0 \to X$를 뇌터 대수공간들의
전사 에탈 사상이라 하자. 다음과 같이 놓자.
$U_1 = U_0 \times_X U_0$ 및
$U_2 = U_0 \times_X U_0 \times_X U_0$.
$\mathcal{I} \subset \mathcal{O}_{X}$를 준연접 아이디얼 층이라 하자.
$\mathcal{I}_i = \mathcal{I}|_{U_i}$로 놓자. 이 상황에서 다음 범주들의
그림을 얻는다.
$$
\xymatrix{
\textit{Coh}(X, \mathcal{I}) \ar[r] &
\textit{Coh}(U_0, \mathcal{I}_0) \ar@<0.5ex>[r] \ar@<-0.5ex>[r] &
\textit{Coh}(U_1, \mathcal{I}_1) \ar@<1ex>[r] \ar[r] \ar@<-1ex>[r] &
\textit{Coh}(U_2, \mathcal{I}_2)
}
$$
첫 화살표는 $\textit{Coh}(X, \mathcal{I})$를 이 그림 오른쪽 부분의
호모토피 극한으로 나타낸다.
% P09-ERR-0713: frozen source begins this sentence with the typo “an”.
더 정확히 말해, {\it 하강 데이터}, 즉 쌍
$((\mathcal{G}_n), \varphi)$가 주어졌다고 하자. 여기서
$(\mathcal{G}_n)$은 $\textit{Coh}(U_0, \mathcal{I}_0)$의 대상이고,
$\varphi : \text{pr}_0^*(\mathcal{G}_n) \to \text{pr}_1^*(\mathcal{G}_n)$는
$\textit{Coh}(U_1, \mathcal{I}_1)$ 안의 동형으로서
$\textit{Coh}(U_2, \mathcal{I}_2)$ 안에서 코사이클 조건을 만족한다.
그러면 유일한 대상 $(\mathcal{F}_n)$이
$\textit{Coh}(X, \mathcal{I})$에 존재하여, 이에 딸린 표준 하강 데이터가
$((\mathcal{G}_n), \varphi)$와 동형이다. 다음과 비교하라.
Descent on Spaces, Definition
\ref{spaces-descent-definition-descent-datum-effective-quasi-coherent}.
이 명제의 증명은 하강 데이터 $(\mathcal{G}_n, \varphi_n)$에
Descent on Spaces, Proposition
\ref{spaces-descent-proposition-fpqc-descent-quasi-coherent}을 적용하면
곧바로 따르며, 여기서 $n$을 변화시킨다.

\begin{lemma}
\label{spaces-more-morphisms-lemma-inverse-systems-abelian}
$S$를 스킴이라 하자. $X$를 $S$ 위의 뇌터 대수공간이라 하고,
$\mathcal{I} \subset \mathcal{O}_X$를 준연접 아이디얼 층이라 하자.
\begin{enumerate}
\item 범주 $\textit{Coh}(X, \mathcal{I})$는 아벨 범주이다.
\item $\textit{Coh}(X, \mathcal{I})$ 안의 완전성은 에탈 국소적으로
확인할 수 있다.
\item 뇌터 대수공간들의 임의의 평탄 사상 $f : X' \to X$에 대해 함자
$f^* : \textit{Coh}(X, \mathcal{I}) \to
\textit{Coh}(X', f^{-1}\mathcal{I}\mathcal{O}_{X'})$는 완전하다.
\end{enumerate}
\end{lemma}

\begin{proof}
(1)의 증명. 아핀 스킴 $U_0$와 전사 에탈 사상 $U_0 \to X$를 택하자.
위의 에탈 하강 논의에서처럼 다음과 같이 놓자.
$U_1 = U_0 \times_X U_0$ 및
$U_2 = U_0 \times_X U_0 \times_X U_0$.
범주 $\textit{Coh}(U_i, \mathcal{I}_i)$는 아벨 범주이고
(Cohomology of Schemes, Lemma
\ref{coherent-lemma-inverse-systems-abelian}), 당김 함자
$\textit{Coh}(U_0, \mathcal{I}_0) \to \textit{Coh}(U_1, \mathcal{I}_1)$와
$\textit{Coh}(U_1, \mathcal{I}_1) \to \textit{Coh}(U_2, \mathcal{I}_2)$는
완전 함자이다(Cohomology of Schemes, Lemma
\ref{coherent-lemma-inverse-systems-pullback}). 그러면
$\textit{Coh}(X, \mathcal{I})$를 하강 데이터들의 범주로 나타낸 서술에서
보조정리가 형식적으로 따른다. 몇몇 세부사항은 생략한다.
% P09-ERR-0714: frozen source explicitly omits the kernel/cokernel details.
다음 증명과 비교하라. Groupoids, Lemma \ref{groupoids-lemma-abelian}.

\medskip\noindent
(2)는 앞 문단의 논의에서 곧바로 따른다. (3)의 상황에서는 다음 가환
그림을 택하자.
$$
\xymatrix{
U' \ar[d] \ar[r] & U \ar[d] \\
X' \ar[r] & X
}
$$
여기서 $U'$과 $U$는 아핀 스킴이고 세로 사상들은 전사 에탈이다.
그러면 $U' \to U$는 뇌터 스킴들의 평탄 사상이다(Morphisms of Spaces,
Lemma \ref{spaces-morphisms-lemma-flat-local}). 따라서 당김 함자
$\textit{Coh}(U, \mathcal{I}\mathcal{O}_U) \to
\textit{Coh}(U', \mathcal{I}\mathcal{O}_{U'})$는 Cohomology of Schemes,
Lemma \ref{coherent-lemma-inverse-systems-pullback}에 의해 완전하다.
$\textit{Coh}(X, \mathcal{O}_X)$ 안의 완전성은 $U$ 위에서 확인할 수 있고,
% P09-ERR-0715: frozen source uses O_X here instead of the ideal I.
$X', U'$에 대해서도 마찬가지이므로 주장이 따른다.
\end{proof}

\begin{lemma}
\label{spaces-more-morphisms-lemma-inverse-systems-surjective}
$S$를 스킴이라 하자. $X$를 $S$ 위의 뇌터 대수공간이라 하고,
$\mathcal{I} \subset \mathcal{O}_X$를 준연접 아이디얼 층이라 하자.
사상 $(\mathcal{F}_n) \to (\mathcal{G}_n)$이
$\textit{Coh}(X, \mathcal{I})$ 안에서 전사일 필요충분조건은
$\mathcal{F}_1 \to \mathcal{G}_1$이 전사인 것이다.
\end{lemma}

\begin{proof}
Lemma \ref{spaces-more-morphisms-lemma-inverse-systems-abelian}에 의해 $X$의 아핀 에탈 덮개에서
확인할 수 있다. 따라서 스킴의 경우로 환원되며, 이는
Cohomology of Schemes, Lemma
\ref{coherent-lemma-inverse-systems-surjective}이다.
\end{proof}

\noindent
$S$를 스킴이라 하자. $X$를 $S$ 위의 뇌터 대수공간이라 하고,
$\mathcal{I} \subset \mathcal{O}_X$를 준연접 아이디얼 층이라 하자.
다음 함자가 존재한다.
\begin{equation}
\label{spaces-more-morphisms-equation-completion-functor}
\textit{Coh}(\mathcal{O}_X) \longrightarrow \textit{Coh}(X, \mathcal{I}), \quad
\mathcal{F} \longmapsto \mathcal{F}^\wedge
\end{equation}
이는 연접 $\mathcal{O}_X$-가군 $\mathcal{F}$에
$\mathcal{F}^\wedge = (\mathcal{F}/\mathcal{I}^n\mathcal{F})$, 즉
$\textit{Coh}(X, \mathcal{I})$의 대상을 대응시킨다.

\begin{lemma}
\label{spaces-more-morphisms-lemma-exact}
함자 (\ref{spaces-more-morphisms-equation-completion-functor})는 완전하다.
\end{lemma}

\begin{proof}
Lemma \ref{spaces-more-morphisms-lemma-inverse-systems-abelian}에 의해 이를 $X$ 위에서 에탈
국소적으로 확인하면 충분하다. 따라서 스킴의 경우로 환원되며, 이는
Cohomology of Schemes, Lemma \ref{coherent-lemma-exact}이다.
\end{proof}

\begin{lemma}
\label{spaces-more-morphisms-lemma-completion-internal-hom}
$S$를 스킴이라 하자. $X$를 $S$ 위의 뇌터 대수공간이라 하고,
$\mathcal{I} \subset \mathcal{O}_X$를 준연접 아이디얼 층이라 하자.
$\mathcal{F}$, $\mathcal{G}$를 연접 $\mathcal{O}_X$-가군들이라 하자.
$\mathcal{H} = \SheafHom_{\mathcal{O}_X}(\mathcal{F}, \mathcal{G})$로 놓자.
그러면
$$
\lim H^0(X, \mathcal{H}/\mathcal{I}^n\mathcal{H}) =
\Mor_{\textit{Coh}(X, \mathcal{I})}
(\mathcal{F}^\wedge, \mathcal{G}^\wedge).
$$
\end{lemma}

\begin{proof}
$\mathcal{H}$는 $X_\etale$ 위의 층이고
$\textit{Coh}(X, \mathcal{I})$의 대상들에 대해 에탈 하강이 있으므로,
이를 에탈 국소적으로 증명하면 충분하다. 따라서 스킴의 경우로 환원되며,
이는 Cohomology of Schemes, Lemma
\ref{coherent-lemma-completion-internal-hom}이다.
\end{proof}

\noindent
이 절의 나머지 전체에서 집중할 상황을 도입한다.

\begin{situation}
\label{spaces-more-morphisms-situation-existence}
여기서 $A$는 아이디얼 $I$에 관해 완비인 뇌터 환이다.
또한 $f : X \to \Spec(A)$는 대수공간들의 유한형 분리 사상이고
$\mathcal{I} = I\mathcal{O}_X$이다.
\end{situation}

\noindent
이 상황에서
$$
\textit{Coh}_{\text{지지가 위에서 고유인 } A}(\mathcal{O}_X)
$$
를 $\textit{Coh}(\mathcal{O}_X)$의 충만한 부분범주로 정의한다.
% P09-ERR-0716: frozen source has a malformed “denote ... be” construction.
그 대상은 연접 $\mathcal{O}_X$-가군 중 지지가 $\Spec(A)$ 위에서
고유인 것들, 또는 동치로 스킴론적 지지가 $\Spec(A)$ 위에서 고유인
것들이다.
다음을 보라. Derived Categories of Spaces, Lemma
\ref{spaces-perfect-lemma-module-support-proper-over-base}.
마찬가지로
$$
\textit{Coh}_{\text{지지가 위에서 고유인 } A}(X, \mathcal{I})
$$
를 $\textit{Coh}(X, \mathcal{I})$의 충만한 부분범주로 정의한다.
그 대상은 $(\mathcal{F}_n)$ 중에서 $\mathcal{F}_1$의 지지가
$\Spec(A)$ 위에서 고유인 것들이다. 몫가군의 지지는 원래 가군의
지지에 포함되므로, (\ref{spaces-more-morphisms-equation-completion-functor})는 함자
\begin{equation}
\label{spaces-more-morphisms-equation-completion-functor-proper-over-A}
\textit{Coh}_{\text{지지가 위에서 고유인 }A}(\mathcal{O}_X)
\longrightarrow
\textit{Coh}_{\text{지지가 위에서 고유인 }A}(X, \mathcal{I})
\end{equation}
를 유도한다. 첫 결과는 이 함자가 충실충만하다는 것이다.

\begin{lemma}
\label{spaces-more-morphisms-lemma-fully-faithful}
Situation \ref{spaces-more-morphisms-situation-existence}의 상황이라 하자.
$\mathcal{F}$, $\mathcal{G}$를 연접 $\mathcal{O}_X$-가군들이라 하자.
$\mathcal{F}$와 $\mathcal{G}$의 지지들의 교집합이 $\Spec(A)$ 위에서
고유라고 가정하자. 그러면 (\ref{spaces-more-morphisms-equation-completion-functor})에서 오는 사상
$$
\Mor_{\textit{Coh}(\mathcal{O}_X)}(\mathcal{F}, \mathcal{G})
\longrightarrow
\Mor_{\textit{Coh}(X, \mathcal{I})}
(\mathcal{F}^\wedge, \mathcal{G}^\wedge)
$$
은 전단사이다. 특히 (\ref{spaces-more-morphisms-equation-completion-functor-proper-over-A})는
충실충만하다.
\end{lemma}

\begin{proof}
$\mathcal{H} = \SheafHom_{\mathcal{O}_X}(\mathcal{G}, \mathcal{F})$로 놓자.
% P09-ERR-0717: frozen source reverses F and G in both internal-Hom formulas.
이는 연접 $\mathcal{O}_X$-가군이다. 실제로 $X$ 위로 에탈인 스킴들에
제한하면 다음에 의해 연접이다.
% P09-ERR-0718: frozen source uses the wrong preposition in this restriction phrase.
Modules, Lemma \ref{modules-lemma-internal-hom-locally-kernel-direct-sum}.
Lemma \ref{spaces-more-morphisms-lemma-completion-internal-hom}에 의해 사상
$$
\lim_n H^0(X, \mathcal{H}/\mathcal{I}^n\mathcal{H})
\to
\Mor_{\textit{Coh}(X, \mathcal{I})}
(\mathcal{G}^\wedge, \mathcal{F}^\wedge)
$$
은 전단사이다. $i : Z \to X$를 $\mathcal{H}$의 스킴론적 지지라 하자.
$Z$는 닫힌 부분공간이고, $|Z|$는 $\mathcal{F}$와 $\mathcal{G}$의
지지들의 교집합에 포함됨이 분명하다. 따라서 가정에 의해
$Z \to \Spec(A)$는 고유이다(Derived Categories of Spaces, Section
\ref{spaces-perfect-section-proper-over-base}를 보라).
$\mathcal{H} = i_*\mathcal{H}'$로 쓰자. 여기서 연접
$\mathcal{O}_Z$-가군을 $\mathcal{H}'$로 나타낸다. 다음 등식이 성립한다.
$i_*(\mathcal{H}'/I^n\mathcal{H}') = \mathcal{H}/I^n\mathcal{H}$.
따라서 다음을 얻는다.
\begin{align*}
\lim_n H^0(X, \mathcal{H}/\mathcal{I}^n\mathcal{H})
& =
\lim_n H^0(Z, \mathcal{H}'/\mathcal{I}^n\mathcal{H}') \\
& =
H^0(Z, \mathcal{H}') \\
& =
H^0(X, \mathcal{H}) \\
& =
\Mor_{\textit{Coh}(\mathcal{O}_X)}(\mathcal{F}, \mathcal{G})
\end{align*}
두 번째 등식은 형식함수 정리에 의한다(Cohomology of Spaces, Lemma
\ref{spaces-cohomology-lemma-spell-out-theorem-formal-functions}).
이것으로 보조정리가 증명된다.
\end{proof}

\begin{remark}
\label{spaces-more-morphisms-remark-inverse-systems-kernel-cokernel-annihilated-by}
$S$를 스킴이라 하자. $X$를 $S$ 위의 뇌터 대수공간이라 하고,
$\mathcal{I}, \mathcal{K} \subset \mathcal{O}_X$를 준연접 아이디얼
층들이라 하자. $\alpha : (\mathcal{F}_n) \to (\mathcal{G}_n)$을
$\textit{Coh}(X, \mathcal{I})$의 사상이라 하자.
아핀 스킴 $U = \Spec(A)$와 전사 에탈 사상 $U \to X$가 주어졌을 때,
$I, K \subset A$를 제한 $\mathcal{I}|_U, \mathcal{K}|_U$에 대응하는
아이디얼들이라 하자. $\alpha_U : M \to N$을 유한
$A^\wedge$-가군들 사이의 사상으로서, Cohomology of Schemes, Lemma
\ref{coherent-lemma-inverse-systems-affine}를 통해 $\alpha|_U$에 대응하는
것이라 하자.
% P09-ERR-0719: frozen source's naming sentence does not state the map relation.
다음 조건들이 동치임을 주장한다.
\begin{enumerate}
\item 정수 $t \geq 1$이 존재하여,
$\Ker(\alpha_n)$과 $\Coker(\alpha_n)$은 $\mathcal{K}^t$에 의해
모든 $n \geq 1$에 대해 소멸한다.
\item 위와 같은 임의의(또는 어떤) 아핀 열린 부분공간
$\Spec(A) = U \subset X$에 대해,
% P09-ERR-0720: frozen source turns the etale cover into an open inclusion.
$\Ker(\alpha_U)$와 $\Coker(\alpha_U)$는 $K^t$에 의해 소멸하며,
여기서 어떤 정수 $t \geq 1$을 택할 수 있다.
\end{enumerate}
이 동치 조건들이 성립하면 $\alpha$를 {\it 핵과 여핵이
$\mathcal{K}$의 한 거듭제곱에 의해 소멸하는 사상}이라 부르겠다.
동치성은 Cohomology of Schemes, Remark
\ref{coherent-remark-inverse-systems-kernel-cokernel-annihilated-by}를 보라.
\end{remark}

\begin{lemma}
\label{spaces-more-morphisms-lemma-existence-easy}
$S$를 스킴이라 하자. $X$를 $S$ 위의 뇌터 대수공간이라 하고,
$\mathcal{I} \subset \mathcal{O}_X$를 준연접 아이디얼 층이라 하자.
$\mathcal{G}$를 연접 $\mathcal{O}_X$-가군,
$(\mathcal{F}_n)$을 $\textit{Coh}(X, \mathcal{I})$의 대상,
$\alpha : (\mathcal{F}_n) \to \mathcal{G}^\wedge$를 그 핵과 여핵이
$\mathcal{I}$의 한 거듭제곱에 의해 소멸하는 사상이라 하자.
그러면 다음을 만족하는 삼중항 $(\mathcal{F}, a, \beta)$가 유일한
동형을 제외하고 유일하게 존재한다.
\begin{enumerate}
\item $\mathcal{F}$는 연접 $\mathcal{O}_X$-가군이다.
\item $a : \mathcal{F} \to \mathcal{G}$는 $\mathcal{O}_X$-가군 사상이고,
그 핵과 여핵은 $\mathcal{I}$의 한 거듭제곱에 의해 소멸한다.
\item $\beta : (\mathcal{F}_n) \to \mathcal{F}^\wedge$는 동형이다.
\item $\alpha = a^\wedge \circ \beta$이다.
\end{enumerate}
\end{lemma}

\begin{proof}
$\textit{Coh}(X, \mathcal{I})$와 $\textit{Coh}(\mathcal{O}_X)$의 대상들에
대한 유일성과 에탈 하강으로부터 $(\mathcal{F}, a, \beta)$를 $X$ 위에서
에탈 국소적으로 구성하면 충분하다는 결론이 따른다.
% P09-ERR-0721: frozen source gives this compound subject a singular verb.
따라서 스킴의 경우로 환원되며, 이는 Cohomology of Schemes, Lemma
\ref{coherent-lemma-existence-easy}이다.
\end{proof}

\begin{lemma}
\label{spaces-more-morphisms-lemma-existence-tricky}
Situation \ref{spaces-more-morphisms-situation-existence}의 상황이라 하자.
$\mathcal{K} \subset \mathcal{O}_X$를 준연접 아이디얼 층이라 하자.
$X_e \subset X$를 $\mathcal{K}^e$가 정의하는 닫힌 부분공간이라 하자.
$\mathcal{I}_e = \mathcal{I}\mathcal{O}_{X_e}$로 놓자.
$(\mathcal{F}_n)$을
$\textit{Coh}_{\text{지지가 위에서 고유인 } A}(X, \mathcal{I})$의
대상이라 하자. 다음을 가정하자.
\begin{enumerate}
\item 함자
$\textit{Coh}_{\text{지지가 위에서 고유인 } A}(\mathcal{O}_{X_e})
\to \textit{Coh}_{\text{지지가 위에서 고유인 } A}(X_e, \mathcal{I}_e)$는
모든 $e \geq 1$에 대해 범주의 동치이다.
\item 대상 $\mathcal{H}$가
$\textit{Coh}_{\text{지지가 위에서 고유인 } A}(\mathcal{O}_X)$에 존재하고,
사상
$\alpha : (\mathcal{F}_n) \to \mathcal{H}^\wedge$가 존재하여,
그 핵과 여핵은 $\mathcal{K}$의 한 거듭제곱에 의해 소멸한다.
\end{enumerate}
그러면 $(\mathcal{F}_n)$은
(\ref{spaces-more-morphisms-equation-completion-functor-proper-over-A})의 본질적 상에 속한다.
\end{lemma}

\begin{proof}
이 증명에서는 닫힌 몰입 $i : Z \to X$에 대해 함자 $i_*$가
$Z$ 위의 연접 가군들의 범주와, $X$ 위의 연접 가군 중 $Z$의
아이디얼 층에 의해 소멸하는 것들의 범주 사이의 동치를 준다는 사실을 더 언급하지
않고 사용하겠다. 다음을 보라. Cohomology of Spaces, Lemma
\ref{spaces-cohomology-lemma-i-star-equivalence}.
특히
$$
\textit{Coh}_{\text{지지가 위에서 고유인 } A}(\mathcal{O}_{X_e})
\subset
\textit{Coh}_{\text{지지가 위에서 고유인 } A}(\mathcal{O}_X)
$$
를 $\mathcal{K}^e$에 의해 소멸하는 가군들로 이루어진 충만한 부분범주로
보고,
% P09-ERR-0722: frozen source says “subcategory of consisting of”.
$$
\textit{Coh}_{\text{지지가 위에서 고유인 } A}(X_e, \mathcal{I}_e)
\subset
\textit{Coh}_{\text{지지가 위에서 고유인 } A}(X, \mathcal{I})
$$
를 $\mathcal{K}^e$에 의해 소멸하는 대상들로 이루어진 충만한 부분범주로
본다. 또한 (1)은 이 두 범주가 완비화 함자
(\ref{spaces-more-morphisms-equation-completion-functor-proper-over-A}) 아래에서 동치임을 말한다.

\medskip\noindent
이 동치를 적용하면 연접 $\mathcal{O}_X$-가군 $\mathcal{G}_e$를 얻으며,
이는 $\mathcal{K}^e$에 의해 소멸하고 계
$(\mathcal{F}_n/\mathcal{K}^e\mathcal{F}_n)$에 대응한다. 이 계는
$\textit{Coh}_{\text{지지가 위에서 고유인 } A}(X, \mathcal{I})$의
대상이다. 사상들
$\mathcal{F}_n/\mathcal{K}^{e + 1}\mathcal{F}_n \to
\mathcal{F}_n/\mathcal{K}^e\mathcal{F}_n$은 표준 사상
$\mathcal{G}_{e + 1} \to \mathcal{G}_e$에 대응하고, 후자는 동형
$\mathcal{G}_{e + 1}/\mathcal{K}^e\mathcal{G}_{e + 1} \to \mathcal{G}_e$를
유도한다.
따라서 대상 $(\mathcal{G}_e)$, 즉 범주
$\textit{Coh}_{\text{지지가 위에서 고유인 } A}(X, \mathcal{K})$의
대상을 얻는다. 사상 $\alpha$는 다음 사상계를 유도한다.
$$
\mathcal{F}_n/\mathcal{K}^e\mathcal{F}_n
\longrightarrow
\mathcal{H}/(\mathcal{I}^n + \mathcal{K}^e)\mathcal{H}
$$
따라서 다시 범주의 동치를 사용하면 사상
$\mathcal{G}_e \to \mathcal{H}/\mathcal{K}^e\mathcal{H}$를 얻는다.
가정 (2)에 의해 존재하는 정수 $t \geq 1$을 택하여
$\mathcal{K}^t$가 모든 사상
$\mathcal{F}_n \to \mathcal{H}/\mathcal{I}^n\mathcal{H}$의 핵과 여핵을
소멸하게 하자. 그러면 $\mathcal{K}^{2t}$는 사상
$\mathcal{F}_n/\mathcal{K}^e\mathcal{F}_n \to
\mathcal{H}/(\mathcal{I}^n + \mathcal{K}^e)\mathcal{H}$의 핵과 여핵을
소멸한다(세부사항은 생략한다. Cohomology of Schemes, Remark
\ref{coherent-remark-inverse-systems-kernel-cokernel-annihilated-by}를 보라).
따라서 $\mathcal{K}^{4t}$가 다음 사상들의 핵과 여핵을 소멸한다고
결론 내린다.
$$
\mathcal{G}_e
\longrightarrow
\mathcal{H}/\mathcal{K}^e\mathcal{H},
$$
(세부사항은 생략한다. Cohomology of Schemes, Remark
\ref{coherent-remark-inverse-systems-kernel-cokernel-annihilated-by}를 보라.)
% P09-ERR-0723: frozen source explicitly omits both quantitative diagram chases.
Lemma \ref{spaces-more-morphisms-lemma-existence-easy}를 적용하여 연접 $\mathcal{O}_X$-가군
$\mathcal{F}$, 사상 $a : \mathcal{F} \to \mathcal{H}$, 그리고 동형
$\beta : (\mathcal{G}_e) \to (\mathcal{F}/\mathcal{K}^e\mathcal{F})$를
얻는다. 이 동형은 $\textit{Coh}(X, \mathcal{K})$ 안에 있다.
거슬러 올라가면, 주어진 $n$에 대해 삼중항
$(\mathcal{F}/\mathcal{I}^n\mathcal{F}, a \bmod \mathcal{I}^n, \beta
\bmod \mathcal{I}^n)$은 사상
$\alpha_n \bmod \mathcal{K}^e :
(\mathcal{F}_n/\mathcal{K}^e\mathcal{F}_n) \to
(\mathcal{H}/(\mathcal{I}^n + \mathcal{K}^e)\mathcal{H})$에 대한
$\textit{Coh}(X, \mathcal{K})$ 안의
보조정리의 삼중항이다. 따라서 Lemma \ref{spaces-more-morphisms-lemma-existence-easy}의
유일성은 눈에 보이는 모든 사상과 호환되는 표준 동형
$\mathcal{F}/\mathcal{I}^n\mathcal{F} \to \mathcal{F}_n$을 준다.

\medskip\noindent
보조정리의 증명을 끝내려면 $\mathcal{F}$의 지지가 $A$ 위에서
고유임을 보여야 한다. 구성에 의해
$a : \mathcal{F} \to \mathcal{H}$의 핵은 $\mathcal{K}$의 한 거듭제곱에
의해 소멸한다. 따라서 이 핵의 지지는 $\mathcal{G}_1$의 지지에 포함된다.
$\mathcal{G}_1$은
$\textit{Coh}_{\text{지지가 위에서 고유인 } A}(\mathcal{O}_{X_1})$의
대상이므로 이것이 $A$ 위에서 고유임을 알 수 있다. $\mathcal{H}$의
지지가 $A$ 위에서 고유라는 사실과 합치면, Derived Categories of Spaces,
Lemma \ref{spaces-perfect-lemma-union-closed-proper-over-base}에 의해
$\mathcal{F}$의 지지가 $A$ 위에서 고유라고 결론 내린다.
\end{proof}

\begin{lemma}
\label{spaces-more-morphisms-lemma-inverse-systems-push-pull}
$S$를 스킴이라 하자. $f : X \to Y$를 $S$ 위의 뇌터 대수공간들 사이의
표현 가능한 고유 사상이라 하자. $\mathcal{J}, \mathcal{K} \subset \mathcal{O}_Y$를
준연접 아이디얼 층이라 하자.
$f$가 $V = Y \setminus V(\mathcal{K})$ 위에서 동형이라고 가정하자.
$\mathcal{I} = f^{-1}\mathcal{J} \mathcal{O}_X$로 놓자.
$(\mathcal{G}_n)$을 $\textit{Coh}(Y, \mathcal{J})$의 대상으로,
$\mathcal{F}$를 연접 $\mathcal{O}_X$-가군으로 하고,
$\beta : (f^*\mathcal{G}_n)  \to \mathcal{F}^\wedge$를
$\textit{Coh}(X, \mathcal{I})$ 안의 동형이라 하자. 그러면
$$
\alpha :
(\mathcal{G}_n)
\longrightarrow
(f_*\mathcal{F})^\wedge
$$
와 같은 $\textit{Coh}(Y, \mathcal{J})$ 안의 사상이 존재하며, 그 핵과 여핵은
$\mathcal{K}$의 한 거듭제곱에 의해 소멸한다.
\end{lemma}

\begin{proof}
$f$가 고유 사상이므로 $f_*\mathcal{F}$는 연접 $\mathcal{O}_Y$-가군이다
(Cohomology of Spaces, Lemma
\ref{spaces-cohomology-lemma-proper-pushforward-coherent}).
따라서 보조정리의 명제는 잘 정의된다. 다음 합성들을 생각하자.
$$
\gamma_n : \mathcal{G}_n \to
f_*f^*\mathcal{G}_n \to
f_*(\mathcal{F}/\mathcal{I}^n\mathcal{F}).
$$
여기서 첫 사상은 수반 사상이고 둘째 사상은 $f_*\beta_n$이다.
보조정리와 같은 유일한 $\alpha$가 존재하여 다음 합성들이
$$
\mathcal{G}_n \xrightarrow{\alpha_n}
f_*\mathcal{F}/\mathcal{J}^nf_*\mathcal{F} \to
f_*(\mathcal{F}/\mathcal{I}^n\mathcal{F})
$$
$\gamma_n$과 같고, 이것이 모든 $n$에 대해 성립한다고 주장한다.
유일성과 $\textit{Coh}(Y, \mathcal{J})$에 대한 에탈 강하에 의해
$Y$ 위에서 이를 에탈 국소적으로 증명하면 충분하다. 따라서 $Y$가 뇌터 환의
스펙트럼이라고 가정할 수 있다. $f$가 표현 가능하므로 $X$ 역시 스킴이다.
그러므로 스킴의 경우로 환원되며, Cohomology of Schemes, Lemma
\ref{coherent-lemma-inverse-systems-push-pull}의 증명을 보라.
\end{proof}

\begin{theorem}[그로텐디크 존재 정리]
\label{spaces-more-morphisms-theorem-grothendieck-existence}
Situation \ref{spaces-more-morphisms-situation-existence}에서 함자
(\ref{spaces-more-morphisms-equation-completion-functor-proper-over-A})는 동치이다.
\end{theorem}

\begin{proof}
이 정리의 증명에서는 Cohomology of Spaces, Lemma
\ref{spaces-cohomology-lemma-i-star-equivalence}의 범주 동치를 더 언급하지
않고 사용하겠다.
Lemma \ref{spaces-more-morphisms-lemma-fully-faithful}에 의해 이 함자는 충실충만하다.
따라서 본질적으로 전사임을 증명하면 된다.

\medskip\noindent
$\Xi$를 다음 성질을 가진 준연접 아이디얼 층
$\mathcal{K} \subset \mathcal{O}_X$의 모음이라 하자. 모든 대상
$(\mathcal{F}_n)$이
$\textit{Coh}_{\text{지지가 위에서 고유인 }A}(X, \mathcal{I})$에 속하고
$\mathcal{K}$에 의해 소멸하면 그 대상에 대해 명제가 성립한다.
$(0)$이 $\Xi$에 속함을 보이고자 한다. 그렇지 않다면 $X$가 뇌터이므로
극대 준연접 아이디얼 층 $\mathcal{K}$가 존재하며 이는 $\Xi$에 속하지 않는다.
Cohomology of Spaces, Lemma
\ref{spaces-cohomology-lemma-acc-coherent}를 보라.
% P09-ERR-0724: frozen source calls the closed subspace of an algebraic space a closed subscheme.
$X$를 $X$의 닫힌 부분스킴, 즉 $\mathcal{K}$에 대응하는 것으로 대체한 뒤에는,
0이 아닌 모든 $\mathcal{K}$가 $\Xi$에 속한다고 가정할 수 있다.
$(\mathcal{F}_n)$을
$\textit{Coh}_{\text{지지가 위에서 고유인 }A}(X, \mathcal{I})$의 대상으로
하자. 이 대상이 본질적 상에 속함을 보이면 정리의 증명이 끝난다.

\medskip\noindent
Chow's lemma (Lemma \ref{spaces-more-morphisms-lemma-chow-noetherian-separated})를 적용하여,
고유 전사 사상 $f : Y \to X$를 택하되 조밀한 열린집합 $U \subset X$
위에서 동형이고 $Y$가 $A$ 위에서 H-준사영이 되게 하자.
$Y$는 스킴이고 $f$는 표현 가능함에 유의하자.
열린 몰입 $j : Y \to Y'$를 택하되 $Y'$가 $A$ 위에서 사영이 되게 하자.
Morphisms, Lemma \ref{morphisms-lemma-H-quasi-projective-open-H-projective}를 보라.
$T_n$을 $\mathcal{F}_n$의 스킴론적 지지라 하자.
$|T_n| = |T_1|$이므로 $T_n$은 $A$ 위에서 모든 $n$에 대해 고유하다
(Morphisms of Spaces, Lemma
\ref{spaces-morphisms-lemma-image-proper-is-proper}).
그러면 $f^*\mathcal{F}_n$은 닫힌 부분스킴 $f^{-1}T_n$에 지지되며,
이 부분스킴은 $A$ 위에서 고유하다. 이는 Morphisms of Spaces, Lemma
\ref{spaces-morphisms-lemma-composition-proper}와 $f$의 고유성에서 따른다.
특히 합성 $f^{-1}T_n \to Y \to Y'$는 닫힌 사상이다
(Morphisms, Lemma \ref{morphisms-lemma-image-proper-scheme-closed}).
$T'_n \subset Y'$를 이에 대응하는 닫힌 부분스킴이라 하자. 이는 열린
부분스킴 $Y$에 포함되고 $f^{-1}T_n$과 같으며, 그 등식은 $Y$의 닫힌
부분스킴으로서의 등식이다.
$\mathcal{F}_n'$을 연접 $\mathcal{O}_{Y'}$-가군으로 다음과 같이 정하자.
이는 $f^*\mathcal{F}_n$을 $Y'$ 위의 연접 가군으로 본 것에 대응하며,
그 보기는 닫힌 몰입 $f^{-1}T_n = T'_n \subset Y'$를 통해 이루어진다.
그러면 $(\mathcal{F}_n')$은
$\textit{Coh}(Y', I\mathcal{O}_{Y'})$의 대상이다.
그로텐디크 존재 정리의 사영 경우
(Cohomology of Schemes, Lemma \ref{coherent-lemma-existence-projective})에 의해
연접 $\mathcal{O}_{Y'}$-가군 $\mathcal{F}'$과 동형
$(\mathcal{F}')^\wedge \cong (\mathcal{F}'_n)$이 존재한다. 이 동형은
$\textit{Coh}(Y', I\mathcal{O}_{Y'})$ 안에 있다.
$Z' \subset Y'$를 $\mathcal{F}'$의 스킴론적 지지라 하자.
$\mathcal{F}'/I\mathcal{F}' = \mathcal{F}'_1$이므로 집합론적으로
$Z' \cap V(I\mathcal{O}_{Y'}) = T'_1$이다.
구조 사상 $p' : Y' \to \Spec(A)$는 고유이므로
$p'(Z' \cap (Y' \setminus Y))$는 $\Spec(A)$에서 닫혀 있다.
이 집합이 비어 있지 않다면 $V(I)$의 한 점을 포함할 것이다. 이는
$I$가 $A$의 야코브슨 근기에 포함되기 때문이다
(Algebra, Lemma \ref{algebra-lemma-radical-completion}).
그러나 위에서 $Z' \cap (p')^{-1}V(I) = T'_1 \subset Y$임을 보았으므로
$Z' \subset Y$라고 결론 내린다. 따라서 $\mathcal{F}'|_Y$는 $Y$의 닫힌
부분스킴에 지지되고, 이 부분스킴은 $A$ 위에서 고유하다.

\medskip\noindent
준연접 아이디얼 층 $\mathcal{K}$를, 축약된 여집합 $X \setminus U$를
잘라내는 것으로 정하자. Cohomology of Spaces, Lemma
\ref{spaces-cohomology-lemma-proper-pushforward-coherent}에 의해
% P09-ERR-0725: frozen source applies f_* to F' on Y' rather than explicitly to its restriction to Y.
$\mathcal{O}_X$-가군 $\mathcal{H} = f_*\mathcal{F}'$은 연접이고,
Lemma \ref{spaces-more-morphisms-lemma-inverse-systems-push-pull}에 의해 사상
$\alpha : (\mathcal{F}_n) \to \mathcal{H}^\wedge$가
$\textit{Coh}_{\text{지지가 위에서 고유인 } A}(X, \mathcal{I})$ 안에
존재하며, 그 핵과 여핵은 $\mathcal{K}$의 한 거듭제곱에 의해 소멸한다.
$Z_0 \subset X$를 $\mathcal{H}$의 스킴론적 지지라 하자.
$|Z_0| \subset f(|Z'|)$임은 명백하다. 따라서
$Z_0 \to \Spec(A)$는 고유하다
(Morphisms of Spaces, Lemma
\ref{spaces-morphisms-lemma-image-proper-is-proper}).
그러므로 $\mathcal{H}$는
$\textit{Coh}_{\text{지지가 위에서 고유인 } A}(\mathcal{O}_X)$의 대상이다.
아이디얼 층 $\mathcal{K}^e$가 각각 $\Xi$의 원소이므로
Lemma \ref{spaces-more-morphisms-lemma-existence-tricky}의 가정들이 만족되고, 결론을 얻는다.
\end{proof}

\begin{remark}[그로텐디크 존재 정리의 구체적 서술]
\label{spaces-more-morphisms-remark-reformulate-existence-theorem}
$A$를 아이디얼 $I$에 관해 완비인 뇌터 환이라 하자.
$S = \Spec(A)$ 및 $S_n = \Spec(A/I^n)$로 쓰자.
$X \to S$를 분리이고 유한형인 대수공간의 사상이라 하자.
$n \geq 1$에 대해 $X_n = X \times_S S_n$으로 놓자. 그림은 다음과 같다.
$$
\xymatrix{
X_1 \ar[r]_{i_1} \ar[d] & X_2 \ar[r]_{i_2} \ar[d] & X_3 \ar[r] \ar[d] &
\ldots & X \ar[d] \\
S_1 \ar[r] & S_2 \ar[r] & S_3 \ar[r] & \ldots & S
}
$$
이 상황에서 다음 조건을 만족하는 계 $(\mathcal{F}_n, \varphi_n)$을 생각한다.
\begin{enumerate}
\item $\mathcal{F}_n$은 연접 $\mathcal{O}_{X_n}$-가군이다.
\item $\varphi_n : i_n^*\mathcal{F}_{n + 1} \to \mathcal{F}_n$은
동형이다.
\item $\text{Supp}(\mathcal{F}_1)$은 $S_1$ 위에서 고유하다.
\end{enumerate}
Theorem \ref{spaces-more-morphisms-theorem-grothendieck-existence}는 완비화 함자
$$
\begin{matrix}
\text{연접 }\mathcal{O}_X\text{-가군 }\mathcal{F} \\
\text{지지가 위에서 고유인 }A
\end{matrix}
\quad
\longrightarrow
\quad
\begin{matrix}
\text{계 }(\mathcal{F}_n) \\
\text{위와 같음}
\end{matrix}
$$
가 범주의 동치임을 말한다. 특히 $X$가 $A$ 위에서 고유한 경우에는
지지에 관한 조건을 생략할 수 있다.
\end{remark}







\section{그로텐디크 대수화 정리}
\label{spaces-more-morphisms-section-algebraization}

\noindent
이 절은 Cohomology of Schemes, Section
\ref{coherent-section-algebraization}에 대응한다. 그러나 풍부한 가역 층을 갖춘
고유 대수공간의 변형을 대수화하는 결과는 이 절에 들어 있지 않다. 풍부한
가역 층을 갖춘 고유 대수공간은 이미 스킴이기 때문이다. 상대 차원이 $1$인
고유 대수공간에 대해서는 대수화 결과가 있다. 첫 결과는 그로텐디크 존재 정리를
닫힌 부분스킴과 유한 사상의 언어로 옮긴 것이다.

\begin{lemma}
\label{spaces-more-morphisms-lemma-algebraize-formal-closed-subscheme}
$A$를 아이디얼 $I$에 관해 완비인 뇌터 환이라 하자.
$S = \Spec(A)$ 및 $S_n = \Spec(A/I^n)$로 쓰자.
$X \to S$를 분리이고 유한형인 대수공간의 사상이라 하자.
$n \geq 1$에 대해 $X_n = X \times_S S_n$으로 놓자.
대수공간들의 다음 가환 그림이 주어졌다고 하자.
$$
\xymatrix{
Z_1 \ar[r] \ar[d] & Z_2 \ar[r] \ar[d] & Z_3 \ar[r] \ar[d] & \ldots \\
X_1 \ar[r]^{i_1} & X_2 \ar[r]^{i_2} & X_3 \ar[r] & \ldots
}
$$
각 정사각형은 카르테시안이라고 하자. 다음을 가정한다.
\begin{enumerate}
\item $Z_1 \to X_1$은 닫힌 몰입이다.
\item $Z_1 \to S_1$은 고유이다.
\end{enumerate}
그러면 대수공간의 닫힌 몰입 $Z \to X$가 존재하여,
$Z_n = Z \times_S S_n$이 모든 $n \geq 1$에 대해 성립한다.
또한 $Z$는 $S$ 위에서 고유하다.
\end{lemma}

\begin{proof}
수직 사상들을 $j_n : Z_n \to X_n$이라 쓰자. 명제의 정사각형들이
카르테시안이므로 $j_n$을 $X_1$로 밑변환한 것은 $j_1$이다.
따라서 Limits of Spaces, Lemma
\ref{spaces-limits-lemma-check-closed-infinitesimally}에 의해 $j_n$은 닫힌
몰입이다. $\mathcal{F}_n = j_{n, *}\mathcal{O}_{Z_n}$으로 놓으면
$j_n^\sharp$는 전사 $\mathcal{O}_{X_n} \to \mathcal{F}_n$이다.
정사각형들이 카르테시안이라는 사실을 다시 사용하면
$\mathcal{F}_{n + 1}$을 $X_n$으로 당겨온 것은 $\mathcal{F}_n$임을 알 수 있다.
따라서 Remark \ref{spaces-more-morphisms-remark-reformulate-existence-theorem}에서 다시 서술한
그로텐디크 존재 정리에 의해 사상 $\mathcal{O}_X \to \mathcal{F}$가
존재한다. 이는 연접 $\mathcal{O}_X$-가군들의 사상이며, $X_n$에 제한하면
$\mathcal{O}_{X_n} \to \mathcal{F}_n$을 되찾는다. 또한 $\mathcal{F}$의
지지는 $S$ 위에서 고유하다. 완비화 함자가 완전하므로
(Lemma \ref{spaces-more-morphisms-lemma-exact}) $\mathcal{O}_X \to \mathcal{F}$는 전사이다.
따라서 $\mathcal{F} = \mathcal{O}_X/\mathcal{J}$이고, 여기서
$\mathcal{J}$는 어떤 준연접 아이디얼 층이다.
$Z = V(\mathcal{J})$로 놓으면 증명이 끝난다.
\end{proof}

\begin{lemma}
\label{spaces-more-morphisms-lemma-algebraize-formal-algebraic-space-finite-over-proper}
$A$를 아이디얼 $I$에 관해 완비인 뇌터 환이라 하자.
$S = \Spec(A)$ 및 $S_n = \Spec(A/I^n)$로 쓰자.
$X \to S$를 분리이고 유한형인 대수공간의 사상이라 하자.
$n \geq 1$에 대해 $X_n = X \times_S S_n$으로 놓자.
대수공간들의 다음 가환 그림이 주어졌다고 하자.
$$
\xymatrix{
Y_1 \ar[r] \ar[d] & Y_2 \ar[r] \ar[d] & Y_3 \ar[r] \ar[d] & \ldots \\
X_1 \ar[r]^{i_1} & X_2 \ar[r]^{i_2} & X_3 \ar[r] & \ldots
}
$$
각 정사각형은 카르테시안이라고 하자. 다음을 가정한다.
\begin{enumerate}
\item $Y_1 \to X_1$은 유한 사상이다.
\item $Y_1 \to S_1$은 고유이다.
\end{enumerate}
그러면 대수공간의 유한 사상 $Y \to X$가 존재하여,
$Y_n = Y \times_S S_n$이 모든 $n \geq 1$에 대해 성립한다.
또한 $Y$는 $S$ 위에서 고유하다.
\end{lemma}

\begin{proof}
수직 사상들을 $f_n : Y_n \to X_n$이라 쓰자. 명제의 정사각형들이
카르테시안이므로 $f_n$을 $X_1$로 밑변환한 것은 $f_1$이다.
따라서 Lemma \ref{spaces-more-morphisms-lemma-thicken-property-morphisms-cartesian}에 의해
$f_n$은 유한 사상이다. $\mathcal{F}_n = f_{n, *}\mathcal{O}_{Y_n}$으로 놓자.
정사각형들이 카르테시안이라는 사실을 사용하면 $\mathcal{F}_{n + 1}$을
$X_n$으로 당겨온 것은 $\mathcal{F}_n$임을 알 수 있다. 따라서 Remark
\ref{spaces-more-morphisms-remark-reformulate-existence-theorem}에서 다시 서술한 그로텐디크
존재 정리에 의해 연접 $\mathcal{O}_X$-가군 $\mathcal{F}$가 존재하며,
이를 $X_n$에 제한하면 $\mathcal{F}_n$을 되찾는다. 또한 $\mathcal{F}$의
지지는 $S$ 위에서 고유하다. 완비화 함자가 충실충만하므로
(Theorem \ref{spaces-more-morphisms-theorem-grothendieck-existence}) 곱셈 사상들
$\mathcal{F}_n \otimes_{\mathcal{O}_{X_n}} \mathcal{F}_n \to
\mathcal{F}_n$은 서로 맞물려 $\mathcal{F}$ 위의 대수 구조를 준다.
$Y = \underline{\Spec}_X(\mathcal{F})$로 놓으면 증명이 끝난다.
\end{proof}

\begin{lemma}
\label{spaces-more-morphisms-lemma-algebraize-morphism}
$A$를 아이디얼 $I$에 관해 완비인 뇌터 환이라 하자.
$S = \Spec(A)$ 및 $S_n = \Spec(A/I^n)$로 쓰자. $X$, $Y$를 $S$ 위의
대수공간이라 하자. $n \geq 1$에 대해 $X_n = X \times_S S_n$ 및
$Y_n = Y \times_S S_n$으로 놓자. 다음과 같은 서로 호환되는 가환 그림들의
계를 주었다고 하자.
$$
\xymatrix{
& & X_{n + 1} \ar[rd] \ar[rr]_{g_{n + 1}} & & Y_{n + 1} \ar[ld] \\
X_n \ar[rru] \ar[rd] \ar[rr]_{g_n} & & Y_n \ar[rru] \ar[ld] & S_{n + 1} \\
& S_n \ar[rru]
}
$$
다음을 가정한다.
\begin{enumerate}
\item $X \to S$는 고유이다.
\item $Y \to S$는 분리 유한형이다.
\end{enumerate}
그러면 유일한 사상 $g : X \to Y$가 존재하며, 이는 $S$ 위 대수공간의
사상이고
$g_n$은 $g$를 $S_n$으로 밑변환한 것이다.
\end{lemma}

\begin{proof}
사상들 $(1, g_n) : X_n \to X_n \times_S Y_n$은 닫힌 몰입이다. 이는
$Y_n \to S_n$이 분리이기 때문이다
(Morphisms of Spaces, Lemma \ref{spaces-morphisms-lemma-section-immersion}).
따라서 Lemma \ref{spaces-more-morphisms-lemma-algebraize-formal-closed-subscheme}에 의해
닫힌 부분공간 $Z \subset X \times_S Y$가 존재하며 이는 $S$ 위에서 고유하고,
이를 $S_n$으로 밑변환하면 $X_n \subset X_n \times_S Y_n$을 되찾는다.
첫 사영 $p : Z \to X$는 고유 사상이다. 실제로 $Z$는 $S$ 위에서
고유이며, Morphisms of Spaces, Lemma
\ref{spaces-morphisms-lemma-universally-closed-permanence}를 보라.
이를 $S_n$으로 밑변환하면 모든 $n$에 대해 동형이다. 특히
$p : Z \to X$는 $Z$의 한 열린 부분공간 위에서 준유한이며, 이 부분공간은
$Z_0$의 모든 점을 포함한다.
예를 들어 Morphisms of Spaces, Lemma
\ref{spaces-morphisms-lemma-locally-finite-type-quasi-finite-part}를 적용하라.
$Z$가 $S$ 위에서 고유이므로 이 열린 근방은 $Z$ 전체이다.
자리츠키 주정리에 의해 $p : Z \to X$가 유한이라고 결론 내린다.
예를 들어 Lemma \ref{spaces-more-morphisms-lemma-quasi-finite-separated-pass-through-finite}를
적용하고, $Z$가 $X$ 위에서 고유임을 사용하여 그 몰입이 닫힌 몰입임을 보라.
Theorem \ref{spaces-more-morphisms-theorem-grothendieck-existence}의 동치를 적용하면
$p_*\mathcal{O}_Z = \mathcal{O}_X$임을 알 수 있다. 이는 $I^n$으로 나눈 뒤
모든 $n$에 대해 성립하기 때문이다. 따라서 $p$는 동형이고 사상 $g$를
얻는데, 이는 합성 $X \cong Z \to Y$이다.
유일성의 증명은 생략한다.
\end{proof}

\begin{remark}
\label{spaces-more-morphisms-remark-weaken-separation-axioms-question}
그로텐디크 대수화 정리(Lemma \ref{spaces-more-morphisms-lemma-algebraize-morphism}의 형태)에서
표적에 더 약한 분리 공리를 두어도 되는지 물을 수 있다. 더 정확히 말하자.
$A$, $I$, $S$, $S_n$, $X$, $Y$, $X_n$, $Y_n$, $g_n$을 Lemma
\ref{spaces-more-morphisms-lemma-algebraize-morphism}의 명제에서와 같이 두고 다음을 가정하자.
\begin{enumerate}
\item $X \to S$는 고유이다.
\item $Y \to S$는 국소 유한형이다.
\end{enumerate}
$g : X \to Y$라는 대수공간의 사상이 존재하여, 이는 $S$ 위의 사상이고
$g_n$이 $g$를
$S_n$으로 밑변환한 것이 되는가? 일반적으로 그 답을 알지 못한다. 답을 안다면
\href{mailto:stacks.project@gmail.com}{stacks.project@gmail.com}으로 이메일을
보내 달라. $Y \to S$가 분리이면 보조정리에 의해 결과가 성립한다. 실제로
$X$가 $S$ 위에서 유한형인 경우로 즉시 환원되는데, 이는 $g_1$의 상을
포함하는 준콤팩트 열린집합을 택하면 된다. $Y \to S$가 준분리라고만
가정해도 결과는 성립한다.
먼저 앞에서처럼 $Y$가 준콤팩트이면서 준분리라고 가정할 수 있다.
% P09-ERR-0726: frozen source has the malformed coordination “either ... or from ...”.
그러면 $(g_n)$을 대수화하기 위해 \cite{Bhatt-Algebraize} 또는
\cite{Hall-Rydh-coherent}의 결과를 사용할 수 있다. 첫 참고문헌을 적용할
때에는 다음을 사용한다.
$$
D_{perf}(X) \to \lim D_{perf}(X_n)
\xrightarrow{\lim Lg_n^*}
\lim D_{perf}(Y_n) = D_{perf}(Y)
$$
마지막 단계에서는 고유 대수공간 $Y$의 유도 범주에 대한 그로텐디크 존재
결과를 사용하며, 여기서 밑은 $R$이다. 다음과 비교하라.
Flatness on Spaces, Remark \ref{spaces-flat-remark-correct-generality}.
% P09-ERR-0727: frozen source's displayed pullback direction and stated morphism Y -> X oppose the desired X -> Y.
인용한 논문은 이 화살이 원하는 사상 $Y \to X$를 결정함을 보인다.
둘째 참고문헌을 적용할 때에는 연접 가군들에 대해 같은 논증을 사용한다.
$$
\textit{Coh}(\mathcal{O}_X) \to \lim \textit{Coh}(\mathcal{O}_{X_n})
\xrightarrow{\lim g_n^*}
\lim \textit{Coh}(\mathcal{O}_{Y_n}) = \textit{Coh}(\mathcal{O}_Y)
$$
마지막 등식은 그로텐디크 존재 정리
(Theorem \ref{spaces-more-morphisms-theorem-grothendieck-existence})의 귀결이다.
% P09-ERR-0728: frozen source switches the base symbol from A to R without defining R here.
둘째 참고문헌은 이 함자가 사상 $Y \to X$에 대응한다고 말하며, 이 사상은
$R$ 위의 사상이다.
이 일반화가 필요해지면 여기서 결과를 정확히 서술하고 주의 깊게 증명할 것이다.
\end{remark}

\begin{lemma}
\label{spaces-more-morphisms-lemma-formal-algebraic-space-proper-reldim-1}
$(A, \mathfrak m, \kappa)$를 완비 국소 뇌터 환이라 하자.
$S = \Spec(A)$ 및 $S_n = \Spec(A/\mathfrak m^n)$으로 놓자.
다음 가환 그림을 생각하자.
$$
\xymatrix{
X_1 \ar[r]_{i_1} \ar[d] & X_2 \ar[r]_{i_2} \ar[d] & X_3 \ar[r] \ar[d] &
\ldots \\
S_1 \ar[r] & S_2 \ar[r] & S_3 \ar[r] & \ldots
}
$$
이는 정사각형들이 카르테시안인 대수공간의 그림이다.
$\dim(X_1) \leq 1$이면 스킴의 사영 사상 $X \to S$와 동형들
$X_n \cong X \times_S S_n$이 존재하고, 이 동형들은 $i_n$과 호환된다.
\end{lemma}

\begin{proof}
Spaces over Fields, Lemma
\ref{spaces-over-fields-lemma-codim-1-point-in-schematic-locus}에 의해 대수공간
$X_1$은 스킴이다. 따라서 $X_1$은 차원 $\leq 1$인, $\kappa$ 위의 고유
스킴이다. Varieties, Lemma
\ref{varieties-lemma-dim-1-proper-projective}에 의해 $X_1$은 $\kappa$ 위에서
H-사영이다. $\mathcal{L}_1$을 $X_1$ 위의 풍부한 가역 층이라 하자.

\medskip\noindent
$\mathcal{L}_1$이 가역 층들의 호환되는 계 $\{\mathcal{L}_n\}$으로
$\{X_n\}$ 위에 올라감을 보이겠다. Lemma \ref{spaces-more-morphisms-lemma-thickening-scheme}에 의해
$X_n$ 역시 스킴임에 유의하자. $X_1 \to X_n$은 바탕 위상공간의 위상동형을
유도한다. 증명의 나머지 부분에서는 $X_n$ 위의 층과 $X_1$ 위의 층을
구별하지 않는다. $\mathcal{L}_n$을 $X_n$으로 올렸다고 하자. 다음 완전열
$$
1 \to
(1 + \mathfrak m^n\mathcal{O}_{X_{n + 1}})^* \to
\mathcal{O}_{X_{n + 1}}^* \to \mathcal{O}_{X_n}^* \to 1
$$
을 생각하자. 이는 $X_{n + 1}$ 위의 층들의 완전열이다.
$\mathcal{L}_n$의 류는
$H^1(X_n, \mathcal{O}_{X_n}^*)$에 놓인다(Cohomology, Lemma
\ref{cohomology-lemma-h1-invertible}를 보라). 이 류를
$H^1(X_{n + 1}, \mathcal{O}_{X_{n + 1}}^*)$의 원소로 올릴 수 있을 필요충분조건은
$H^2(X_{n + 1}, (1 + \mathfrak m^n\mathcal{O}_{X_{n + 1}})^*)$ 안의 장애물이
0인 것이다. $X_1$은 차원 $\leq 1$인 뇌터 스킴이므로 이 코호몰로지 군은
사라진다(Cohomology, Proposition
\ref{cohomology-proposition-vanishing-Noetherian}).

\medskip\noindent
그로텐디크 대수화 정리
(Cohomology of Schemes, Theorem \ref{coherent-theorem-algebraization})에 의해
스킴의 사영 사상 $X \to S = \Spec(A)$와 호환되는 동형들의 계
$X_n = S_n \times_S X$를 얻는다.
\end{proof}

\begin{lemma}
\label{spaces-more-morphisms-lemma-projective-over-complete}
$(A, \mathfrak m, \kappa)$를 완비 뇌터 국소환이라 하자.
$X$를 $\Spec(A)$ 위의 대수공간이라 하자.
$X \to \Spec(A)$가 고유이고 $\dim(X_\kappa) \leq 1$이면,
$X$는 스킴이며 $A$ 위에서 사영이다.
\end{lemma}

\begin{proof}
$X_n = X \times_{\Spec(A)} \Spec(A/\mathfrak m^n)$으로 놓자.
Lemma \ref{spaces-more-morphisms-lemma-formal-algebraic-space-proper-reldim-1}에 의해 사영 사상
$Y \to \Spec(A)$와 호환되는 동형들
$Y \times_{\Spec(A)} \Spec(A/\mathfrak m^n) \cong
X \times_{\Spec(A)} \Spec(A/\mathfrak m^n)$이 존재한다.
Lemma \ref{spaces-more-morphisms-lemma-algebraize-morphism}에 의해 $X \cong Y$이고 증명이 끝난다.
\end{proof}











\section{정칙 몰입}
\label{spaces-more-morphisms-section-regular-immersions}

\noindent
이 절은 대수공간의 사상에 관해 Divisors, Section
\ref{divisors-section-regular-immersions}에 대응한다. 독자는 먼저 그 절에서
스킴의 정칙 몰입을 읽어 보기를 권한다.

\medskip\noindent
Divisors, Section \ref{divisors-section-regular-immersions}에서는 스킴의
사상에 대해 네 종류의 정칙 몰입을 정의했다. 우리가 아는 한 이 가운데 세
종류만 에탈 위상에서 표적에 관해 국소적이며, 보통의 정칙 몰입은 늘 그렇듯
그렇지 않다. 이 때문에 대수공간의 사상에 대해서는 정칙 몰입을 실제로 정의할
수 없다. 이러한 불편 때문에 그로텐디크와 그의 학파는 정칙 몰입의 원래 개념을
코쥘 정칙 몰입으로 대체하였다. \cite[Exposee VII, Definition 1.4]{SGA6}을
보라. 그러나 코쥘 정칙, $H_1$-정칙, 준정칙 몰입은 정의할 수 있다.
또 한 가지 주의할 점은 코쥘 정칙 몰입이 임의의 밑변환으로 보존되지 않으므로
이를 정의할 때 Morphisms of Spaces, Section
\ref{spaces-morphisms-section-representable}의 전략을 사용할 수 없다는 것이다.
마찬가지로 코쥘 정칙 몰입은 원천에 관해 에탈 국소적이지 않으므로 Morphisms
of Spaces, Lemma \ref{spaces-morphisms-lemma-local-source-target}도 사용할 수
없다. 대신 이 보조정리를 다음 것으로 바꾼다.

\begin{lemma}
\label{spaces-more-morphisms-lemma-representable-etale-local-target}
$\mathcal{P}$를 표적에 관해 에탈 국소적인 스킴 사상의 성질이라 하자.
$S$를 스킴이라 하자. $f : X \to Y$를 $S$ 위 대수공간의 표현 가능한
사상이라 하자. 다음 가환 그림들을 생각하자.
$$
\xymatrix{
X \times_Y V \ar[d] \ar[r] & V \ar[d] \\
X \ar[r]^f & Y
}
$$
여기서 $V$는 스킴이고 $V \to Y$는 에탈이다. 다음 조건들은 동치이다.
\begin{enumerate}
\item 위와 같은 모든 그림에서 사영 $X \times_Y V \to V$가 성질
$\mathcal{P}$를 갖는다.
\item 위와 같은 어떤 그림에서 $V \to Y$가 전사이고, 사영
$X \times_Y V \to V$가 성질 $\mathcal{P}$를 갖는다.
\end{enumerate}
$X$와 $Y$가 표현 가능하면, 이는 $f$가 스킴의 사상으로서 성질
$\mathcal{P}$를 갖는 것과도 동치이다.
\end{lemma}

\begin{proof}
(1)과 (2)의 동치를 증명하자. 함의 (1) $\Rightarrow$ (2)는 즉시 따른다.
다음 두 그림이 보조정리와 같은 그림이라고 하자.
$$
\xymatrix{
X \times_Y V \ar[d] \ar[r] & V \ar[d] \\
X \ar[r]^f & Y
}
\quad\quad
\xymatrix{
X \times_Y V' \ar[d] \ar[r] & V' \ar[d] \\
X \ar[r]^f & Y
}
$$
$V \to Y$가 전사이고 $X \times_Y V \to V$가 성질 $\mathcal{P}$를
갖는다고 가정하자. (2)가 (1)을 함의함을 보이려면
$X \times_Y V' \to V'$가 $\mathcal{P}$를 가짐을 증명해야 한다.
이를 위해 다음 그림을 생각하자.
$$
\xymatrix{
X \times_Y V \ar[d] &
(X \times_Y V) \times_X (X \times_Y V') \ar[l] \ar[d] \ar[r] &
X \times_Y V' \ar[d] \\
V &
V \times_Y V' \ar[l] \ar[r] &
V'
}
$$
% P09-ERR-0729: frozen source invokes locality on the source although the lemma assumes locality on the target.
$\mathcal{P}$가 원천에 관해 에탈 국소적이라는 가정에 의해
$\mathcal{P}$는 에탈 밑변환으로 보존된다. Descent, Lemma
\ref{descent-lemma-pullback-property-local-target}를 보라. 따라서 왼쪽 수직
화살이 $\mathcal{P}$를 가지면 가운데 수직 화살도 그렇다.
% P09-ERR-0730: frozen source switches from the defined V,V' to undefined U,U' here.
$U \times_X U' \to U'$는 전사 에탈이므로 $U'$의 에탈 덮개를 정의한다.
% P09-ERR-0731: frozen source concludes for the left vertical arrow where the proof goal is the right one.
따라서 $\mathcal{P}$가 에탈 위상에서 표적에 관해 국소적이라는 가정에 의해
왼쪽 수직 화살이 $\mathcal{P}$를 갖는다.

\medskip\noindent
$X$와 $Y$가 표현 가능하면 에탈 덮개로 $\text{id}_Y : Y \to Y$를 택할
수 있으므로 보조정리의 마지막 명제가 참이다.
\end{proof}

\noindent
``코쥘 정칙(각각 $H_1$-정칙, 준정칙) 몰입이라는 것''은 표적에 관해 fpqc
국소적인 스킴 사상의 성질이다. Descent, Lemma
\ref{descent-lemma-descending-property-regular-immersion}을 보라.
따라서 다음 정의는 잘 정의된다.

\begin{definition}
\label{spaces-more-morphisms-definition-regular-immersion}
$S$를 스킴이라 하자. $i : X \to Y$를 $S$ 위 대수공간의 사상이라 하자.
\begin{enumerate}
\item $i$를 {\it 코쥘 정칙 몰입}이라 한다는 것은 $i$가 표현 가능하고 Lemma
\ref{spaces-more-morphisms-lemma-representable-etale-local-target}의 동치 조건들이
$\mathcal{P}(f) =$ ``$f$는 코쥘 정칙 몰입이다''에 대해 성립한다는 뜻이다.
\item $i$를 {\it $H_1$-정칙 몰입}이라 한다는 것은 $i$가 표현 가능하고 Lemma
\ref{spaces-more-morphisms-lemma-representable-etale-local-target}의 동치 조건들이
$\mathcal{P}(f) =$ ``$f$는 $H_1$-정칙 몰입이다''에 대해 성립한다는 뜻이다.
\item $i$를 {\it 준정칙 몰입}이라 한다는 것은 $i$가 표현 가능하고 Lemma
\ref{spaces-more-morphisms-lemma-representable-etale-local-target}의 동치 조건들이
$\mathcal{P}(f) =$ ``$f$는 준정칙 몰입이다''에 대해 성립한다는 뜻이다.
\end{enumerate}
\end{definition}

\begin{lemma}
\label{spaces-more-morphisms-lemma-regular-quasi-regular-immersion}
$S$를 스킴이라 하자. $i : Z \to X$를 $S$ 위 대수공간의 몰입이라 하자.
다음 함의들이 성립한다.
$i$가 코쥘 정칙 $\Rightarrow$
$i$가 $H_1$-정칙 $\Rightarrow$
$i$가 준정칙이다.
\end{lemma}

\begin{proof}
정의를 사용하면 이 보조정리는 Divisors, Lemma
\ref{divisors-lemma-regular-quasi-regular-immersion}으로 즉시 환원된다.
\end{proof}

\begin{lemma}
\label{spaces-more-morphisms-lemma-regular-immersion-noetherian}
$S$를 스킴이라 하자. $i : Z \to X$를 $S$ 위 대수공간의 몰입이라 하자.
$X$가 국소 뇌터라고 가정하자. 그러면
$i$가 코쥘 정칙 $\Leftrightarrow$
$i$가 $H_1$-정칙 $\Leftrightarrow$
$i$가 준정칙이다.
\end{lemma}

\begin{proof}
Definition \ref{spaces-more-morphisms-definition-regular-immersion}과 Properties of Spaces,
Section \ref{spaces-properties-section-types-properties}의 국소 뇌터 대수공간
정의를 사용하면, 이는 스킴의 경우인 Divisors, Lemma
\ref{divisors-lemma-regular-immersion-noetherian}으로 즉시 옮겨진다.
\end{proof}

\begin{lemma}
\label{spaces-more-morphisms-lemma-flat-base-change-regular-immersion}
\begin{slogan}
정칙 몰입은 평탄 밑변환으로 보존된다.
\end{slogan}
$S$를 스킴이라 하자. $i : Z \to X$를 코쥘 정칙,
$H_1$-정칙 또는 준정칙 몰입이라 하며, 이 몰입은 $S$ 위 대수공간의
몰입이라 하자. $X' \to X$를 $S$ 위 대수공간의
평탄 사상이라 하자. 그러면 밑변환 $i' : Z \times_X X' \to X'$은 각각
코쥘 정칙, $H_1$-정칙 또는 준정칙 몰입이다.
\end{lemma}

\begin{proof}
Definition \ref{spaces-more-morphisms-definition-regular-immersion}과 Morphisms of Spaces,
Section \ref{spaces-morphisms-section-flat}의 대수공간의 평탄 사상 정의를
사용하면, 이 보조정리는 스킴의 경우로 환원된다. Divisors, Lemma
\ref{divisors-lemma-flat-base-change-regular-immersion}을 보라.
\end{proof}

\begin{lemma}
\label{spaces-more-morphisms-lemma-quasi-regular-immersion}
$S$를 스킴이라 하자. $i : Z \to X$를 $S$ 위 대수공간의 몰입이라 하자.
그러면 $i$가 준정칙 몰입일 필요충분조건은 다음 조건들이 성립하는 것이다.
\begin{enumerate}
\item $i$는 국소 유한 표시이다.
\item 여법선 층 $\mathcal{C}_{Z/X}$는 유한 국소 자유이다.
\item 사상 (\ref{spaces-more-morphisms-equation-conormal-algebra-quotient})은 동형이다.
\end{enumerate}
\end{lemma}

\begin{proof}
에탈 국소화를 통해 스킴의 경우에서 따른다(Divisors, Lemma
\ref{divisors-lemma-quasi-regular-immersion}). Definition
\ref{spaces-more-morphisms-definition-regular-immersion}과 Lemma
\ref{spaces-more-morphisms-lemma-etale-conormal-algebra}를 사용하라.
\end{proof}

\begin{lemma}
\label{spaces-more-morphisms-lemma-transitivity-conormal-quasi-regular}
$S$를 스킴이라 하자. $Z \to Y \to X$를 $S$ 위 대수공간의 몰입들이라 하자.
$Z \to Y$가 $H_1$-정칙이라고 가정하자. 그러면 Lemma
\ref{spaces-more-morphisms-lemma-transitivity-conormal}의 표준 열
$$
0 \to i^*\mathcal{C}_{Y/X} \to
\mathcal{C}_{Z/X} \to
\mathcal{C}_{Z/Y} \to 0
$$
은 완전하고 (에탈) 국소적으로 분할된다.
\end{lemma}

\begin{proof}
$\mathcal{C}_{Z/Y}$가 유한 국소 자유이므로(Lemma
\ref{spaces-more-morphisms-lemma-quasi-regular-immersion}과 Lemma
\ref{spaces-more-morphisms-lemma-regular-quasi-regular-immersion}을 보라), 이 열이 완전임을
증명하면 충분하다. 이 열은 일반적으로 이미 오른쪽 완전이므로 첫 사상이
단사임을 보이면 충분하다. $X$ 위에서 에탈 국소화한 뒤에는 스킴의 경우로
환원된다. Divisors, Lemma
\ref{divisors-lemma-transitivity-conormal-quasi-regular}를 보라.
\end{proof}

\noindent
준정칙 몰입의 합성은 준정칙이 아닐 수도 있다. Algebra, Remark
\ref{algebra-remark-join-quasi-regular-sequences}를 보라. 다른 종류의 정칙
몰입은 합성으로 보존된다.

\begin{lemma}
\label{spaces-more-morphisms-lemma-composition-regular-immersion}
$S$를 스킴이라 하자. $i : Z \to Y$와 $j : Y \to X$를 $S$ 위
대수공간의 몰입들이라 하자.
\begin{enumerate}
\item $i$와 $j$가 코쥘 정칙 몰입이면 $j \circ i$도 그렇다.
\item $i$와 $j$가 $H_1$-정칙 몰입이면 $j \circ i$도 그렇다.
\item $i$가 $H_1$-정칙 몰입이고 $j$가 준정칙 몰입이면
$j \circ i$는 준정칙 몰입이다.
\end{enumerate}
\end{lemma}

\begin{proof}
스킴의 경우에서 즉시 따른다. Divisors, Lemma
\ref{divisors-lemma-composition-regular-immersion}을 보라.
\end{proof}

\begin{lemma}
\label{spaces-more-morphisms-lemma-permanence-regular-immersion}
$S$를 스킴이라 하자. $i : Z \to Y$와 $j : Y \to X$를 $S$ 위
대수공간의 몰입들이라 하자. $j$가 국소 유한 표시이고, Lemma
\ref{spaces-more-morphisms-lemma-transitivity-conormal}의 열
$$
0 \to i^*\mathcal{C}_{Y/X} \to
\mathcal{C}_{Z/X} \to
\mathcal{C}_{Z/Y} \to 0
$$
이 완전하고 국소적으로 분할된다고 가정하자.
\begin{enumerate}
\item $j \circ i$가 준정칙 몰입이면 $i$도 그렇다.
\item $j \circ i$가 $H_1$-정칙 몰입이면 $i$도 그렇다.
\item $j$와 $j \circ i$가 모두 코쥘 정칙 몰입이면 $i$도 그렇다.
\end{enumerate}
\end{lemma}

\begin{proof}
스킴의 경우에서 즉시 따른다. Divisors, Lemma
\ref{divisors-lemma-permanence-regular-immersion}을 보라.
\end{proof}

\begin{lemma}
\label{spaces-more-morphisms-lemma-extra-permanence-regular-immersion-noetherian}
$S$를 스킴이라 하자. $i : Z \to Y$와 $j : Y \to X$를 $S$ 위
대수공간의 몰입들이라 하자. $X$가 국소 뇌터라고 가정하자.
다음 조건들은 동치이다.
\begin{enumerate}
\item $i$와 $j$는 코쥘 정칙 몰입이다.
\item $i$와 $j \circ i$는 코쥘 정칙 몰입이다.
\item $j \circ i$는 코쥘 정칙 몰입이고 여법선 열
$$
0 \to i^*\mathcal{C}_{Y/X} \to
\mathcal{C}_{Z/X} \to
\mathcal{C}_{Z/Y} \to 0
$$
은 완전하고 국소적으로 분할된다.
\end{enumerate}
\end{lemma}

\begin{proof}
스킴의 경우에서 즉시 따른다. Divisors, Lemma
\ref{divisors-lemma-extra-permanence-regular-immersion-noetherian}을 보라.
\end{proof}














\section{상대 유사연접성}
\label{spaces-more-morphisms-section-relative-pseudo-coherence}

\noindent
이 절은 More on Morphisms, Section
\ref{more-morphisms-section-relative-pseudo-coherence}에 대응한다. 다만
대수공간에 대해 이 내용을 다룰 때에는 코호몰로지 층들이 준연접인 유도 범주의
대상만을 사용하기로 한다. 여기에는 두 이유가 있다. (1) 서술이 크게 단순해지고,
(2) 현재로서는 더 일반적인 개념을 사용할 일이 없다.

\begin{remark}
\label{spaces-more-morphisms-remark-match-relative-pseudo-coherence}
$S$를 스킴이라 하자. $f : X \to Y$를 $S$ 위 표현 가능한 대수공간의 국소
유한형 사상이라 하자. 스킴의 사상 $f_0 : X_0 \to Y_0$를 $f$를 표현하는
것으로 두자(어색하지만 임시적인 표기이다). 그러면 $f_0$는
국소 유한형이다. $E$가 $D_\QCoh(\mathcal{O}_X)$의 대상이면 $E$는
유일한 대상 $E_0$를 당겨온 것이고, 이 대상은
$D_\QCoh(\mathcal{O}_{X_0})$에 속한다.
Derived Categories of Spaces, Lemma
\ref{spaces-perfect-lemma-derived-quasi-coherent-small-etale-site}를 보라.
이 상황에서 ``$E$는 $m$-유사연접이며 상대 밑이 $Y$이다''라는 말은
``$E_0$는 $m$-유사연접이며 상대 밑이 $Y_0$이다''라는 뜻으로 사용한다. 이는
More on Morphisms, Section
\ref{more-morphisms-section-relative-pseudo-coherence}에서 정의되었다.
\end{remark}

\begin{lemma}
\label{spaces-more-morphisms-lemma-qcoh-relative-pseudo-coherence-characterize}
$S$를 스킴이라 하자. $f : X \to Y$를 $S$ 위 대수공간의 국소 유한형
사상이라 하자. $m \in \mathbf{Z}$라 하자.
$E \in D_\QCoh(\mathcal{O}_X)$라 하자. Remark
\ref{spaces-more-morphisms-remark-match-relative-pseudo-coherence}에서 설명한 표기 아래 다음은
동치이다.
\begin{enumerate}
\item 모든 가환 그림
$$
\xymatrix{
U \ar[d] \ar[r] & V \ar[d] \\
X \ar[r] & Y
}
$$
에서 $U$, $V$가 스킴이고 수직 화살들이 에탈이면, 복합체 $E|_U$는
$m$-유사연접이며 상대 밑은 $V$이다.
\item (1)과 같은 어떤 가환 그림에서 $U \to X$가 전사이면, 복합체
$E|_U$는 $m$-유사연접이며 상대 밑은 $V$이다.
\item (1)과 같은 모든 가환 그림에서 $U$와 $V$가 아핀이면, 복합체
$R\Gamma(U, E)$는 $\mathcal{O}_X(U)$-가군들의 복합체이고
$m$-유사연접이며 상대 밑은 $\mathcal{O}_Y(V)$이다.
\end{enumerate}
\end{lemma}

\begin{proof}
(1)은 More on Morphisms, Lemma
\ref{more-morphisms-lemma-qcoh-relative-pseudo-coherence-characterize}에 의해
(3)을 함의한다.

\medskip\noindent
(3)을 가정하자. (1)과 같은 임의의 가환 그림을 택하되 $U \to X$가
전사이게 하자. 아핀 열린 덮개 $V = \bigcup V_j$와 아핀 열린 덮개들
$(U \to V)^{-1}(V_j) = \bigcup U_{ij}$를 택하자. (3)과 More on Morphisms,
Lemma \ref{more-morphisms-lemma-qcoh-relative-pseudo-coherence-characterize}에 의해
$E|_U$는 $m$-유사연접이며 상대 밑은 $V$이다. 따라서 (3)은 (2)를 함의한다.

\medskip\noindent
(2)를 가정하자. 다음 가환 그림을 택하자.
$$
\xymatrix{
U \ar[d] \ar[r] & V \ar[d] \\
X \ar[r] & Y
}
$$
여기서 $U$, $V$는 스킴이고 수직 화살들은 에탈이며, 사상 $U \to X$는
전사이고 $E|_U$는 $m$-유사연접이며 상대 밑은 $V$이다. 다음으로 둘째 가환 그림
$$
\xymatrix{
U' \ar[d] \ar[r] & V' \ar[d] \\
X \ar[r] & Y
}
$$
이 주어졌다고 하자. 수직 화살들은 에탈이고 $U', V'$는 스킴이라 하자.
$E|_{U'}$가 $m$-유사연접이며 상대 밑이 $V'$임을 보이고자 한다.
사상 $U'' = U \times_X U' \to U'$는 전사 에탈이고, $U'' \to V'$는
$V'' = V' \times_Y V$를 통해 인수분해되며 이 공간은 $V'$ 위에서 에탈이다.
따라서 $E|_{U''}$이 $m$-유사연접이며 상대 밑이 $V''$임을 보이면 충분하다. More on
Morphisms, Lemmas
\ref{more-morphisms-lemma-relative-pseudo-coherent-descends-fppf}와
\ref{more-morphisms-lemma-relative-pseudo-coherent-post-compose}를 보라.
둘째 보조정리를 한 번 더 사용하면 $E|_{U''}$이 $m$-유사연접이며 상대
밑이 $V$임을 보이면 충분하다. 이는 More on Morphisms, Lemma
\ref{more-morphisms-lemma-pull-relative-pseudo-coherent}와, 스킴의 에탈 사상이
유사연접이라는 사실에서 따른다. More on Morphisms, Lemma
\ref{more-morphisms-lemma-flat-finite-presentation-pseudo-coherent}를 보라.
\end{proof}

\begin{definition}
\label{spaces-more-morphisms-definition-relative-pseudo-coherence}
$S$를 스킴이라 하자. $f : X \to Y$를 $S$ 위 대수공간의 국소 유한형
사상이라 하자. $E$를 $D_\QCoh(\mathcal{O}_X)$의 대상으로 하자.
$\mathcal{F}$를 준연접 $\mathcal{O}_X$-가군이라 하자.
$m \in \mathbf{Z}$를 고정하자.
\begin{enumerate}
\item $E$가 {\it $m$-유사연접이며 상대 밑이 $Y$}라고 한다는 것은 Lemma
\ref{spaces-more-morphisms-lemma-qcoh-relative-pseudo-coherence-characterize}의 동치 조건들이
성립한다는 뜻이다.
\item $E$가 {\it 상대 밑 $Y$에 대해 유사연접}이라고 한다는 것은,
$E$가 $m$-유사연접이고 상대 밑이 $Y$인 것이 모든
$m \in \mathbf{Z}$에 대해 성립한다는 뜻이다.
\item $\mathcal{F}$가 {\it $m$-유사연접이며 상대 밑이 $Y$}라고 한다는 것은,
$\mathcal{F}$를 $D_\QCoh(\mathcal{O}_X)$의 대상으로 보았을 때
$m$-유사연접이고 상대 밑이 $Y$라는 뜻이다.
\item $\mathcal{F}$가 {\it 상대 밑 $Y$에 대해 유사연접}이라고 한다는 것은,
$\mathcal{F}$를 $D_\QCoh(\mathcal{O}_X)$의 대상으로 보았을 때 상대 밑
$Y$에 대해 유사연접이라는 뜻이다.
\end{enumerate}
\end{definition}

\noindent
밑에 대한 유사연접 복합체의 성질 대부분은 스킴의 경우에 대응하는 성질에서
즉시 따른다. 필요한 관련 보조정리를 여기에 추가하겠다.

\begin{lemma}
\label{spaces-more-morphisms-lemma-relative-pseudo-coherent-is-moot}
$S$를 스킴이라 하자. $f : X \to Y$를 $S$ 위 대수공간의 사상이라 하자.
$E$를 $D_\QCoh(\mathcal{O}_X)$의 대상으로 하자. $f$가 평탄이고 국소 유한
표시이면 다음은 동치이다.
\begin{enumerate}
\item $E$는 $Y$에 대해 유사연접이다.
\item $E$는 $X$ 위에서 유사연접이다.
\end{enumerate}
\end{lemma}

\begin{proof}
에탈 국소화와 정의에 의해 $X$와 $Y$가 스킴이라고 가정할 수 있다.
스킴의 경우에는 More on Morphisms, Lemma
\ref{more-morphisms-lemma-check-relative-pseudo-coherence-on-charts}에서 따른다.
\end{proof}


















\section{유사연접 사상}
\label{spaces-more-morphisms-section-pseudo-coherent}

\noindent
이 절은 스킴의 사상에 대해 More on Morphisms, Section
\ref{more-morphisms-section-pseudo-coherent}에 대응한다. 독자는 먼저 그 절에서
스킴의 유사연접 사상을 읽어 보기를 권한다.

\medskip\noindent
스킴 사상의 ``유사연접''이라는 성질은 원천과 표적에 관해 에탈 국소적이다.
이를 보려면 More on Morphisms, Lemmas
\ref{more-morphisms-lemma-descending-property-pseudo-coherent}와
\ref{more-morphisms-lemma-pseudo-coherent-fppf-local-source}, 그리고 Descent,
Lemma \ref{descent-lemma-etale-local-source-target}를 사용하라. Morphisms of
Spaces, Lemma \ref{spaces-morphisms-lemma-local-source-target}에 의해 다음과 같이
대수공간의 유사연접 사상을 정의할 수 있다. 문제의 대수공간들이 표현 가능할
때에는 More on Morphisms, Section
\ref{more-morphisms-section-pseudo-coherent}에서 이미 정의한 개념과 일치한다.

\begin{definition}
\label{spaces-more-morphisms-definition-pseudo-coherent}
$S$를 스킴이라 하자. $f : X \to Y$를 $S$ 위 대수공간의 사상이라 하자.
\begin{enumerate}
\item $f$가 {\it 유사연접}이라고 한다는 것은 Morphisms of Spaces, Lemma
\ref{spaces-morphisms-lemma-local-source-target}의 동치 조건들이
$\mathcal{P} =$ ``유사연접''에 대해 성립한다는 뜻이다.
\item $x \in |X|$라 하자. $f$가 {\it $x$에서 유사연접}이라고 한다는 것은
열린 근방 $X' \subset X$가 $x$에 대해 존재하여
$f|_{X'} : X' \to Y$가 유사연접이라는 뜻이다.
\end{enumerate}
\end{definition}

\noindent
유사연접 사상의 밑변환은 일반적으로 유사연접이 아닐 수 있음에 주의하라.

\begin{lemma}
\label{spaces-more-morphisms-lemma-flat-base-change-pseudo-coherent}
유사연접 사상의 평탄 밑변환은 유사연접이다.
\end{lemma}

\begin{proof}
생략한다. 힌트: 이 보조정리의 스킴 판본인 More on Morphisms, Lemma
\ref{more-morphisms-lemma-flat-base-change-pseudo-coherent}를 사용하라.
\end{proof}

\begin{lemma}
\label{spaces-more-morphisms-lemma-composition-pseudo-coherent}
유사연접 사상들의 합성은 유사연접이다.
\end{lemma}

\begin{proof}
생략한다. 힌트: 이 보조정리의 스킴 판본인 More on Morphisms, Lemma
\ref{more-morphisms-lemma-composition-pseudo-coherent}를 사용하라.
\end{proof}

\begin{lemma}
\label{spaces-more-morphisms-lemma-pseudo-coherent-finite-presentation}
유사연접 사상은 국소 유한 표시이다.
\end{lemma}

\begin{proof}
정의에서 즉시 따른다.
\end{proof}

\begin{lemma}
\label{spaces-more-morphisms-lemma-flat-finite-presentation-pseudo-coherent}
국소 유한 표시인 평탄 사상은 유사연접이다.
\end{lemma}

\begin{proof}
생략한다. 힌트: 이 보조정리의 스킴 판본인 More on Morphisms, Lemma
\ref{more-morphisms-lemma-flat-finite-presentation-pseudo-coherent}를 사용하라.
\end{proof}

\begin{lemma}
\label{spaces-more-morphisms-lemma-permanence-pseudo-coherent}
대수공간 사상 $f : X \to Y$가 어떤 밑 대수공간 $B$ 위에서 유사연접이면,
$f$는 유사연접이다.
\end{lemma}

\begin{proof}
생략한다. 힌트: 이 보조정리의 스킴 판본인 More on Morphisms, Lemma
\ref{more-morphisms-lemma-permanence-pseudo-coherent}를 사용하라.
\end{proof}

\begin{lemma}
\label{spaces-more-morphisms-lemma-Noetherian-pseudo-coherent}
$S$를 스킴이라 하자. $f : X \to Y$를 $S$ 위 대수공간의 사상이라 하자.
$Y$가 국소 뇌터이면, $f$가 유사연접일 필요충분조건은 $f$가 국소 유한형인
것이다.
\end{lemma}

\begin{proof}
생략한다. 힌트: 이 보조정리의 스킴 판본인 More on Morphisms, Lemma
\ref{more-morphisms-lemma-Noetherian-pseudo-coherent}를 사용하라.
\end{proof}








\section{완전 사상}
\label{spaces-more-morphisms-section-perfect}

\noindent
이 절은 스킴의 사상에 대해 More on Morphisms, Section
\ref{more-morphisms-section-perfect}에 대응한다. 독자는 먼저 그 절에서 스킴의
완전 사상을 읽어 보기를 권한다.

\medskip\noindent
스킴 사상의 ``완전''이라는 성질은 원천과 표적에 관해 에탈 국소적이다.
이를 보려면 More on Morphisms, Lemmas
\ref{more-morphisms-lemma-descending-property-perfect}와
\ref{more-morphisms-lemma-perfect-fppf-local-source}, 그리고 Descent, Lemma
\ref{descent-lemma-etale-local-source-target}를 사용하라. Morphisms of Spaces,
Lemma \ref{spaces-morphisms-lemma-local-source-target}에 의해 다음과 같이
대수공간의 완전 사상을 정의할 수 있다. 문제의 대수공간들이 표현 가능할
때에는 More on Morphisms, Section \ref{more-morphisms-section-perfect}에서
이미 정의한 개념과 일치한다.

\begin{definition}
\label{spaces-more-morphisms-definition-perfect}
$S$를 스킴이라 하자. $f : X \to Y$를 $S$ 위 대수공간의 사상이라 하자.
\begin{enumerate}
\item $f$가 {\it 완전}이라고 한다는 것은 Morphisms of Spaces, Lemma
\ref{spaces-morphisms-lemma-local-source-target}의 동치 조건들이
$\mathcal{P} =$ ``완전''에 대해 성립한다는 뜻이다.
\item $x \in |X|$라 하자. $f$가 {\it $x$에서 완전}이라고 한다는 것은
열린 근방 $X' \subset X$가 $x$에 대해 존재하여
$f|_{X'} : X' \to Y$가 완전이라는 뜻이다.
\end{enumerate}
\end{definition}

\noindent
완전 사상은 유사연접이며, 따라서 국소 유한 표시임에 유의하라. 완전 사상의
밑변환은 일반적으로 완전이 아닐 수 있음에 주의하라.

\begin{lemma}
\label{spaces-more-morphisms-lemma-flat-base-change-perfect}
완전 사상의 평탄 밑변환은 완전이다.
\end{lemma}

\begin{proof}
생략한다. 힌트: 이 보조정리의 스킴 판본인 More on Morphisms, Lemma
\ref{more-morphisms-lemma-flat-base-change-perfect}를 사용하라.
\end{proof}

\begin{lemma}
\label{spaces-more-morphisms-lemma-composition-perfect}
완전 사상들의 합성은 완전이다.
\end{lemma}

\begin{proof}
생략한다. 힌트: 이 보조정리의 스킴 판본인 More on Morphisms, Lemma
\ref{more-morphisms-lemma-composition-perfect}를 사용하라.
\end{proof}

\begin{lemma}
\label{spaces-more-morphisms-lemma-flat-finite-presentation-perfect}
$S$를 스킴이라 하자. $f : X \to Y$를 $S$ 위 대수공간의 사상이라 하자.
다음은 동치이다.
\begin{enumerate}
\item $f$는 평탄이고 완전이다.
\item $f$는 평탄이고 국소 유한 표시이다.
\end{enumerate}
\end{lemma}

\begin{proof}
생략한다. 힌트: 이 보조정리의 스킴 판본인 More on Morphisms, Lemma
\ref{more-morphisms-lemma-flat-finite-presentation-perfect}를 사용하라.
\end{proof}

\begin{lemma}
\label{spaces-more-morphisms-lemma-perfect-proper-perfect-direct-image}
$S$를 스킴이라 하자. $Y$를 $S$ 위의 뇌터 대수공간이라 하자.
$f : X \to Y$를 대수공간의 완전 고유 사상이라 하자.
$E \in D(\mathcal{O}_X)$가 완전이라 하자. 그러면 $Rf_*E$는
$D(\mathcal{O}_Y)$의 완전 대상이다.
\end{lemma}

\begin{proof}
Derived Categories of Spaces, Lemma
\ref{spaces-perfect-lemma-perfect-direct-image}를 적용할 수 있다고 주장한다.
조건 (1)과 (2)는 즉시 따른다. 조건 (3)은 $X$ 위에서 국소적이다. 따라서
$X$와 $Y$가 아핀이고 $E$가 $\mathcal{O}_X$-가군의 엄밀히 완전한 복합체로
표현된다고 가정할 수 있다. 그러므로 $\mathcal{O}_X$가 에탈 사이트 위의
$f^{-1}\mathcal{O}_Y$-가군 층으로서 유한 Tor 차원을 가짐을 보이면 충분하다.
Derived Categories of Spaces, Lemma
\ref{spaces-perfect-lemma-tor-dimension-rel}에 의해 자리츠키 사이트에서 이를
확인하면 충분하다. 유한형 스킴 사상에 대해서는 이것이 완전성과 동치이다.
More on Morphisms, Lemma \ref{more-morphisms-lemma-check-perfect-stalks}를 보라.
\end{proof}











\section{국소 완전 교차 사상}
\label{spaces-more-morphisms-section-lci}

\noindent
이 절은 스킴의 사상에 대해 More on Morphisms, Section
\ref{more-morphisms-section-lci}에 대응한다. 독자는 먼저 그 절에서 스킴의
국소 완전 교차 사상을 읽어 보기를 권한다.

\medskip\noindent
스킴 사상의 ``국소 완전 교차 사상이라는 것''은 원천과 표적에 관해 에탈
국소적인 성질이다. 이를 보려면 More on Morphisms, Lemmas
\ref{more-morphisms-lemma-descending-property-lci}와
\ref{more-morphisms-lemma-lci-syntomic-local-source}, 그리고 Descent, Lemma
\ref{descent-lemma-etale-local-source-target}를 사용하라. Morphisms of Spaces,
Lemma \ref{spaces-morphisms-lemma-local-source-target}에 의해 다음과 같이
대수공간의 국소 완전 교차 사상을 정의할 수 있다. 문제의 대수공간들이 표현
가능할 때에는 More on Morphisms, Section \ref{more-morphisms-section-lci}에서
이미 정의한 개념과 일치한다.

\begin{definition}
\label{spaces-more-morphisms-definition-lci}
$S$를 스킴이라 하자. $f : X \to Y$를 $S$ 위 대수공간의 사상이라 하자.
\begin{enumerate}
\item $f$가 {\it 코쥘 사상}, 또는 $f$가 {\it 국소 완전 교차 사상}이라고
한다는 것은 Morphisms of Spaces, Lemma
\ref{spaces-morphisms-lemma-local-source-target}의 동치 조건들이
$\mathcal{P}(f) =$ ``$f$는 국소 완전 교차 사상이다''에 대해 성립한다는 뜻이다.
\item $x \in |X|$라 하자. $f$가 {\it $x$에서 코쥘}이라고 한다는 것은
열린 근방 $X' \subset X$가 $x$에 대해 존재하여
$f|_{X'} : X' \to Y$가 국소 완전 교차 사상이라는 뜻이다.
\end{enumerate}
\end{definition}

\noindent
어떤 의미에서 국소 완전 교차 사상의 정의적 성질은 다음 보조정리의 결과이다.

\begin{lemma}
\label{spaces-more-morphisms-lemma-lci}
$S$를 스킴이라 하자. $f : X \to Y$를 $S$ 위 대수공간의 국소 완전 교차
사상이라 하자. $P$를 $Y$ 위에서 매끄러운 대수공간이라 하자.
$U \to X$를 대수공간의 에탈 사상으로 하고, $i : U \to P$를 $Y$ 위
대수공간의 몰입이라 하자. 그림은 다음과 같다.
$$
\xymatrix{
X \ar[rd] & U \ar[l] \ar[d] \ar[r]_i & P \ar[ld] \\
& Y
}
$$
그러면 $i$는 대수공간의 코쥘 정칙 몰입이다.
\end{lemma}

\begin{proof}
스킴 $V$와 전사 에탈 사상 $V \to Y$를 택하자. 스킴 $W$와 전사 에탈 사상
$W \to P \times_Y V$를 택하자. $U' = U \times_P W$로 놓으면, 이는 $U$ 위의
에탈 스킴이다. Definition \ref{spaces-more-morphisms-definition-regular-immersion}에 의해
$U' \to W$가 스킴의 코쥘 정칙 몰입임을 보여야 한다. 위의 Definition
\ref{spaces-more-morphisms-definition-lci}에 의해 스킴의 사상 $U' \to V$는 국소 완전 교차
사상이다. 따라서 More on Morphisms, Lemma
\ref{more-morphisms-lemma-lci}에서 결과가 따른다.
\end{proof}

\noindent
모든 코쥘 사상이 공통으로 갖는 몇 가지 성질을 여기에 모아 두는 것이 좋겠다.

\begin{lemma}
\label{spaces-more-morphisms-lemma-lci-properties}
$S$를 스킴이라 하자. $f : X \to Y$를 $S$ 위 대수공간의 국소 완전 교차
사상이라 하자. 그러면
\begin{enumerate}
\item $f$는 국소 유한 표시이다.
\item $f$는 유사연접이다.
\item $f$는 완전이다.
\end{enumerate}
\end{lemma}

\begin{proof}
생략한다. 힌트: 이 보조정리의 스킴 판본인 More on Morphisms, Lemma
\ref{more-morphisms-lemma-lci-properties}를 사용하라.
\end{proof}

\noindent
코쥘 사상의 밑변환은 일반적으로 코쥘 사상이 아닐 수 있음에 주의하라.

\begin{lemma}
\label{spaces-more-morphisms-lemma-flat-base-change-lci}
국소 완전 교차 사상의 평탄 밑변환은 국소 완전 교차 사상이다.
\end{lemma}

\begin{proof}
생략한다. 힌트: 이 보조정리의 스킴 판본인 More on Morphisms, Lemma
\ref{more-morphisms-lemma-flat-base-change-lci}를 사용하라.
\end{proof}

\begin{lemma}
\label{spaces-more-morphisms-lemma-composition-lci}
국소 완전 교차 사상들의 합성은 국소 완전 교차 사상이다.
\end{lemma}

\begin{proof}
생략한다. 힌트: 이 보조정리의 스킴 판본인 More on Morphisms, Lemma
\ref{more-morphisms-lemma-composition-lci}를 사용하라.
\end{proof}

\begin{lemma}
\label{spaces-more-morphisms-lemma-flat-lci}
\begin{slogan}
신토믹은 평탄 더하기 lci와 같다(대수공간의 경우).
\end{slogan}
$S$를 스킴이라 하자. $f : X \to Y$를 $S$ 위 대수공간의 사상이라 하자.
다음은 동치이다.
\begin{enumerate}
\item $f$는 평탄이고 국소 완전 교차 사상이다.
\item $f$는 신토믹이다.
\end{enumerate}
\end{lemma}

\begin{proof}
생략한다. 힌트: 이 보조정리의 스킴 판본인 More on Morphisms, Lemma
\ref{more-morphisms-lemma-flat-lci}를 사용하라.
\end{proof}

\begin{lemma}
\label{spaces-more-morphisms-lemma-regular-immersion-lci}
$S$를 스킴이라 하자. $S$ 위 대수공간의 코쥘 정칙 몰입은 국소 완전 교차
사상이다.
\end{lemma}

\begin{proof}
$i : X \to Y$를 $S$ 위 대수공간의 코쥘 정칙 몰입이라 하자. 정의에 의해
전사 에탈 사상 $V \to Y$가 존재하는데, $V$는 스킴이고 $X \times_Y V$도
스킴이며 밑변환 $X \times_Y V \to V$는 스킴의 코쥘 정칙 몰입이다.
More on Morphisms, Lemma
\ref{more-morphisms-lemma-regular-immersion-lci}에 의해
$X \times_Y V \to V$는 국소 완전 교차 사상이다. Definition
\ref{spaces-more-morphisms-definition-lci}에 의해 $i$가 대수공간의 국소 완전 교차 사상이라고
결론 내린다.
\end{proof}

\begin{lemma}
\label{spaces-more-morphisms-lemma-lci-permanence}
$S$를 스킴이라 하자. 다음 그림을
$$
\xymatrix{
X \ar[rr]_f \ar[rd] & & Y \ar[ld] \\
& Z
}
$$
$S$ 위 대수공간 사상들의 가환 그림이라 하자. $Y \to Z$가 매끄럽고
$X \to Z$가 국소 완전 교차 사상이라고 가정하자. 그러면
$f : X \to Y$는 국소 완전 교차 사상이다.
\end{lemma}

\begin{proof}
스킴 $W$와 전사 에탈 사상 $W \to Z$를 택하자. 스킴 $V$와 전사 에탈 사상
$V \to W \times_Z Y$를 택하자. 스킴 $U$와 전사 에탈 사상
$U \to V \times_Y X$를 택하자. 그러면 $U \to W$는 스킴의 국소 완전
교차 사상이고 $V \to W$는 스킴의 매끄러운 사상이다. 스킴에 대한 결과
(More on Morphisms, Lemma \ref{more-morphisms-lemma-lci-permanence})에 의해
$U \to V$가 국소 완전 교차 사상이라고 결론 내린다. 정의에 의해 이는
$f$가 국소 완전 교차 사상이라는 뜻이다.
\end{proof}

\begin{lemma}
\label{spaces-more-morphisms-lemma-descending-property-lci}
성질 $\mathcal{P}(f) =$ ``$f$는 국소 완전 교차 사상이다''는 밑에 관해 fpqc
국소적이다.
\end{lemma}

\begin{proof}
$S$를 스킴이라 하자. $f : X \to Y$를 $S$ 위 대수공간의 사상이라 하자.
$\{Y_i \to Y\}$를 fpqc 덮개라 하자(Topologies on Spaces, Definition
\ref{spaces-topologies-definition-fpqc-covering}).
$f_i : X_i \to Y_i$를 $f$의 밑변환이라 하되, 밑변환 사상은
$Y_i \to Y$라 하자.
$f$가 국소 완전 교차 사상이면 Lemma \ref{spaces-more-morphisms-lemma-flat-base-change-lci}에 의해
각 $f_i$도 국소 완전 교차 사상이다.

\medskip\noindent
역으로 각 $f_i$가 국소 완전 교차 사상이라고 가정하자. 덮개를 세분으로
대체할 수 있다. 이는 평탄 밑변환이 국소 완전 교차 사상이라는 성질을
보존하기 때문이다. 따라서 $Y_i$가 각 $i$에 대해 스킴이라고 가정할 수 있다.
Topologies on Spaces, Lemma
\ref{spaces-topologies-lemma-refine-fpqc-schemes}를 보라. 스킴 $V$와 전사
에탈 사상 $V \to Y$를 택하자. 스킴 $U$와 전사 에탈 사상
$U \to V \times_Y X$를 택하자. $U \to V$가 스킴의 국소 완전 교차
사상임을 보여야 한다. Topologies on Spaces, Lemma
\ref{spaces-topologies-lemma-recognize-fpqc-covering}에 의해
$\{Y_i \times_Y V \to V\}$는 스킴의 fpqc 덮개이다. 스킴의 경우
(More on Morphisms, Lemma \ref{more-morphisms-lemma-descending-property-lci})에
의해, 다음 밑변환이 국소 완전 교차 사상임을 증명하면 충분하다.
% P09-ERR-0732: frozen source's displayed base-change arrow ends at V rather than V x_Y Y_i.
$$
U \times_Y Y_i = U \times_V (V \times_Y Y_i) \longrightarrow V
$$
이는 $U \to V$를 $V \times_Y Y_i \to V$에 의해 밑변환한 것이다.
이를 다음 합성으로 쓸 수 있다.
$$
U \times_Y Y_i \longrightarrow
(V \times_Y X) \times_Y Y_i = V \times_Y X_i \longrightarrow
V \times_Y Y_i
$$
첫 화살은 스킴의 에탈 사상이다($U \to V \times_Y X$의 밑변환이기 때문이다).
둘째 화살은 $f_i$의 평탄 밑변환이므로 스킴의 국소 완전 교차 사상이다.
국소 완전 교차 사상이라는 것은 원천에 관해 신토믹 국소적이고, 에탈 사상은
신토믹이므로 결과가 따른다(More on Morphisms, Lemma
\ref{more-morphisms-lemma-lci-syntomic-local-source}와 Morphisms, Lemma
\ref{morphisms-lemma-etale-syntomic}).
\end{proof}

\begin{lemma}
\label{spaces-more-morphisms-lemma-lci-syntomic-local-source}
성질 $\mathcal{P}(f) =$ ``$f$는 국소 완전 교차 사상이다''는 원천에 관해
신토믹 국소적이다.
\end{lemma}

\begin{proof}
Descent on Spaces, Lemma \ref{spaces-descent-lemma-transfer-from-schemes}와
More on Morphisms, Lemma
\ref{more-morphisms-lemma-lci-syntomic-local-source}에서 따른다.
\end{proof}

\begin{lemma}
\label{spaces-more-morphisms-lemma-base-change-lci-fibres}
$S$를 스킴이라 하자. 다음 가환 그림을 생각하자.
$$
\xymatrix{
X \ar[rr]_f \ar[rd]_p & & Y \ar[ld]^q \\
& Z
}
$$
이는 $S$ 위 대수공간의 그림이라 하자. $p$와 $q$가 모두 평탄이고 국소
유한 표시라고 가정하자. 그러면 열린 부분공간 $U(f) \subset X$가 존재하여
$|U(f)| \subset |X|$는 $f$가 코쥘인 점들의 집합이다. 또한 대수공간의 임의의
사상 $Z' \to Z$에 대해, $f' : X' \to Y'$가 $f$의 밑변환이고 그
밑변환 사상이 $Z' \to Z$이면, $U(f')$은 $U(f)$의 역상이며 이 역상은
사영 $X' \to X$ 아래에서 취한다.
\end{lemma}

\begin{proof}
이 보조정리는 More on Morphisms, Lemma
\ref{more-morphisms-lemma-base-change-lci-fibres}에 대응하며, 실제로 그
보조정리에서 이를 유도하겠다. Definition \ref{spaces-more-morphisms-definition-lci}에 의해 집합
$\{x \in |X| : f \text{가 코쥘인 점은 }x\}$는 $|X|$에서 열려 있다.
따라서 Properties of Spaces, Lemma
\ref{spaces-properties-lemma-open-subspaces}에 의해 이는 열린 부분공간
$U(f)$에 대응하며, 이 공간은 $X$의 부분공간이다. 그러므로 마지막 명제만
증명하면 된다.

\medskip\noindent
스킴 $W$와 전사 에탈 사상 $W \to Z$를 택하자. 스킴 $V$와 전사 에탈 사상
$V \to W \times_Z Y$를 택하자. 스킴 $U$와 전사 에탈 사상
$U \to V \times_Y X$를 택하자. 마지막으로 스킴 $W'$과 전사 에탈 사상
$W' \to W \times_Z Z'$를 택하자. $V' = W' \times_W V$ 및
$U' = W' \times_W U$로 놓으면 전사 에탈 사상 $V' \to Y'$와
$U' \to X'$을 얻는다. 대수공간의 에탈 사상이 결합 위상공간의 열린 사상을
유도한다는 사실을 더 언급하지 않고 사용하겠다(Properties of Spaces, Lemma
\ref{spaces-properties-lemma-etale-open}를 보라).
정의에 의해 $U(f)$는 $|X|$에서 집합 $T$의 상임에 유의하자. 이 집합은
$U$의 점들 가운데 스킴 사상 $U \to V$가 코쥘인 점들로 이루어진다.
마찬가지로 $U(f')$은 $|X'|$에서 집합 $T'$의 상이다. 이 집합은 $U'$의
점들 가운데 스킴 사상 $U' \to V'$가 코쥘인 점들로 이루어진다. 이제
구성에 의해 그림
$$
\xymatrix{
U' \ar[r] \ar[d] & U \ar[d] \\
V' \ar[r] & V
}
$$
은 스킴의 범주에서 카르테시안이다. 따라서 앞서 말한 More on Morphisms,
Lemma \ref{more-morphisms-lemma-base-change-lci-fibres}에 의해 $T'$은 $T$의
역상이다. $|U'| \to |X'|$가 전사이므로 보조정리가 따른다.
\end{proof}

\begin{lemma}
\label{spaces-more-morphisms-lemma-unramified-lci}
$S$를 스킴이라 하자. $f : X \to Y$를 $S$ 위 대수공간의 국소 완전 교차
사상이라 하자. 그러면 $f$가 비분기일 필요충분조건은 $f$가 형식적으로
비분기인 것이다. 이 경우 여법선 층 $\mathcal{C}_{X/Y}$는 $X$ 위에서 유한
국소 자유이다.
\end{lemma}

\begin{proof}
에탈 국소화를 통해 스킴 사상에 대한 대응 결과에서 따른다. More on Morphisms,
Lemma \ref{more-morphisms-lemma-unramified-lci}와 Lemma
\ref{spaces-more-morphisms-lemma-universal-thickening-localize}를 보라. 이 보조정리의 상황에서 사상
$V \to U$는 정의에 의해 비분기이고 국소 완전 교차 사상임에 유의하라.
\end{proof}

\begin{lemma}
\label{spaces-more-morphisms-lemma-transitivity-conormal-lci}
$S$를 스킴이라 하자. $Z \to Y \to X$를 $S$ 위 대수공간의 형식적으로
비분기인 사상들이라 하자. $Z \to Y$가 국소 완전 교차 사상이라고 가정하자.
Lemma \ref{spaces-more-morphisms-lemma-transitivity-conormal}의 완전열
$$
0 \to i^*\mathcal{C}_{Y/X} \to
\mathcal{C}_{Z/X} \to
\mathcal{C}_{Z/Y} \to 0
$$
은 짧은 완전열이다.
\end{lemma}

\begin{proof}
스킴 $U$와 전사 에탈 사상 $U \to X$를 택하자. 스킴 $V$와 전사 에탈 사상
$V \to U \times_X Y$를 택하자. 스킴 $W$와 전사 에탈 사상
$W \to V \times_Y Z$를 택하자. Lemma
\ref{spaces-more-morphisms-lemma-universal-thickening-localize}에 의해 사상들 $W \to V$와
$V \to U$는 형식적으로 비분기이다. 또한 열
$i^*\mathcal{C}_{Y/X} \to \mathcal{C}_{Z/X} \to \mathcal{C}_{Z/Y} \to 0$을
제한하면 대응 열
$i^*\mathcal{C}_{V/U} \to \mathcal{C}_{W/U} \to \mathcal{C}_{W/V} \to 0$을
얻으며, 이는 $W \to V \to U$에 대한 열이다. 따라서 스킴에 대한 결과
(More on Morphisms, Lemma
\ref{more-morphisms-lemma-transitivity-conormal-lci})에서 결론이 따른다.
실제로 정의에 의해 사상 $W \to V$는 국소 완전 교차 사상이다.
\end{proof}








\section{사상이 언제 동형인가?}
\label{spaces-more-morphisms-section-when-isomorphism}

\noindent
더 일반적으로 ``사상이 언제 성질 $\mathcal{P}$를 갖는가?''라고 물을 수 있다.
더 정확한 질문은 다음과 같다. 대수공간의 가환 그림
$$
\xymatrix{
X \ar[rr]_f \ar[rd]_p & & Y \ar[ld]^q \\
& Z
}
$$
이 주어졌다고 하자. 다음 두 성질을 갖는 대수공간의 단사 사상 $W \to Z$가
존재하는가?
\begin{enumerate}
\item 밑변환 $f_W : X_W \to Y_W$는 성질 $\mathcal{P}$를 갖는다.
\item 대수공간의 임의의 사상 $Z' \to Z$가 $W$를 통해 인수분해될
필요충분조건은 밑변환 $f_{Z'} : X_{Z'} \to Y_{Z'}$가 성질
$\mathcal{P}$를 갖는 것이다.
\end{enumerate}
많은 경우 $W \to Z$가 존재하면 이는 몰입, 열린 몰입 또는 닫힌 몰입이다.

\medskip\noindent
이 질문의 답은 사상들 $f$, $p$, $q$의 보조 성질에 따라 달라질 수 있다.
한 예는 $\mathcal{P}(f) =$ ``$f$는 평탄이다''이다. 스킴 사상의 경우
$Y = S$일 때 이를 ``More on Flatness'' 장에서 자세히 논의했으며, More on
Flatness, Section \ref{flat-section-flattening-functors}에서 시작한다.

\begin{lemma}
\label{spaces-more-morphisms-lemma-where-unramified}
대수공간의 가환 그림
$$
\xymatrix{
X \ar[rr]_f \ar[rd]_p & & Y \ar[ld]^q \\
& Z
}
$$
을 생각하자. $p$가 국소 유한형이고 닫힌 사상이라고 가정하자. 그러면 열린
부분공간 $W \subset Z$가 존재하여, 사상 $Z' \to Z$가 $W$를 통해
인수분해될 필요충분조건은 밑변환 $f_{Z'} : X_{Z'} \to Y_{Z'}$가 비분기인
것이다.
\end{lemma}

\begin{proof}
Morphisms of Spaces, Lemma \ref{spaces-morphisms-lemma-where-unramified}에 의해
열린 부분공간 $U(f) \subset X$가 존재하며, 이는 $f$가 비분기인 점들의
집합이다. 또한 $U(f)$의 형성은 임의의 밑변환과 교환한다. $W \subset Z$를
다음 바탕 점 집합을 갖는 열린 부분공간이라 하자(Properties of Spaces,
Lemma \ref{spaces-properties-lemma-open-subspaces}를 보라).
$$
|W| = |Z| \setminus |p|\left(|X| \setminus |U(f)|\right)
$$
즉, $z \in |Z|$가 $W$의 점일 필요충분조건은 $f$가 $X$의 모든 점 가운데
$z$ 위에 놓인 점에서 비분기인 것이다. $p$가 닫혀 있다고 가정했으므로 이는 열린
집합임에 유의하자. $U(f)$의 형성이 임의의 밑변환과 교환하므로 Properties
of Spaces, Lemma \ref{spaces-properties-lemma-factor-through-open-subspace}를
사용하면 $W$가 원하는 보편 성질을 가짐을 즉시 알 수 있다.
\end{proof}









\begin{lemma}
\label{spaces-more-morphisms-lemma-where-unramified-universally-injective}
대수공간의 가환 그림
$$
\xymatrix{
X \ar[rr]_f \ar[rd]_p & & Y \ar[ld]^q \\
& Z
}
$$
을 생각하자. 다음을 가정하자.
\begin{enumerate}
\item $p$는 국소 유한형이다.
\item $p$는 닫힌 사상이다.
\item $p_2 : X \times_Y X \to Z$는 닫힌 사상이다.
\end{enumerate}
그러면 열린 부분공간 $W \subset Z$가 존재하여, 사상 $Z' \to Z$가 $W$를
통해 인수분해될 필요충분조건은 밑변환
$f_{Z'} : X_{Z'} \to Y_{Z'}$가 비분기이고 보편 단사인 것이다.
\end{lemma}

\begin{proof}
Lemma \ref{spaces-more-morphisms-lemma-where-unramified}에서 얻은 열린 부분공간으로 $Z$를 바꾸면
$f$가 이미 비분기라고 가정할 수 있다. 이 과정은 가정 (2)나 (3)을 훼손하지
않는다. Morphisms of Spaces, Lemma
\ref{spaces-morphisms-lemma-diagonal-unramified-morphism}에 의해
$\Delta_{X/Y} : X \to X \times_Y X$는 열린 몰입이다. 이는 임의의 밑변환
뒤에도 성립한다. 따라서 Morphisms of Spaces, Lemma
\ref{spaces-morphisms-lemma-universally-injective}에 의해 $f_{Z'}$가 보편
단사일 필요충분조건은 대각 사상의 밑변환
$X_{Z'} \to (X \times_Y X)_{Z'}$가 동형인 것이다. $W \subset Z$를 다음
바탕 점 집합을 갖는 열린 부분공간이라 하자(Properties of Spaces, Lemma
\ref{spaces-properties-lemma-open-subspaces}를 보라).
$$
|W| = |Z| \setminus
|p_2|\left(|X \times_Y X| \setminus \Im(|\Delta_{X/Y}|)\right)
$$
즉, $z \in |Z|$가 $W$의 점일 필요충분조건은
$|X \times_Y X| \to |Z|$의 $z$ 위 올이
$|X| \to |X \times_Y X|$의 상에 포함되는 것이다. 그러면 위의 논의에서
$p^{-1}(W) \to q^{-1}(W)$로 주어진 $f$의 제한이 비분기이고 보편 단사임이
분명하다.

\medskip\noindent
거꾸로 $f_{Z'}$가 비분기이고 보편 단사라고 가정하자. $Z' \to Z$가 $W$를
통해 인수분해됨을 보이려면 $|Z'| \to |Z|$의 상이 $|W|$에 포함됨을 보이면
충분하다. Properties of Spaces, Lemma
\ref{spaces-properties-lemma-factor-through-open-subspace}를 보라. 따라서
$Z'$가 체의 스펙트럼인 경우를 증명하면 충분하다. $z \in |Z|$를
$|Z'| \to |Z|$의 상이라 하자. 위의 논의에 의해
$$
|X_{Z'}| \longrightarrow |(X \times_Y X)_{Z'}|
$$
는 전사이다. Properties of Spaces, Lemma
\ref{spaces-properties-lemma-points-cartesian}에 의해 가환 그림
$$
\xymatrix{
|X_{Z'}| \ar[d] \ar[r] &
|(X \times_Y X)_{Z'}| \ar[d] \\
|p|^{-1}(\{z\}) \ar[r] &
|p_2|^{-1}(\{z\})
}
$$
에서 수직 화살표들은 전사이다. 따라서 원하는 대로 $z \in |W|$이다.
\end{proof}

\begin{lemma}
\label{spaces-more-morphisms-lemma-where-closed-immersion}
대수공간의 가환 그림
$$
\xymatrix{
X \ar[rr]_f \ar[rd]_p & & Y \ar[ld]^q \\
& Z
}
$$
을 생각하자. 다음을 가정하자.
\begin{enumerate}
\item $p$는 국소 유한형이다.
\item $p$는 보편 닫힌 사상이다.
\item $q : Y \to Z$는 분리 사상이다.
\end{enumerate}
그러면 열린 부분공간 $W \subset Z$가 존재하여, 사상 $Z' \to Z$가 $W$를
통해 인수분해될 필요충분조건은 밑변환
$f_{Z'} : X_{Z'} \to Y_{Z'}$가 닫힌 몰입인 것이다.
\end{lemma}

\begin{proof}
닫힌 몰입을 보편 닫힌, 비분기, 보편 단사인 사상으로 특징짓는 Lemma
\ref{spaces-more-morphisms-lemma-characterize-closed-immersion}를 사용하겠다. 먼저 $p$가 보편
닫힌 사상이고 $q$가 분리 사상이므로 $f$는 보편 닫힌 사상이다. Morphisms
of Spaces, Lemma
\ref{spaces-morphisms-lemma-universally-closed-permanence}를 보라. 따라서
$f$의 모든 밑변환도 보편 닫힌 사상이다. Morphisms of Spaces, Lemma
\ref{spaces-morphisms-lemma-base-change-universally-closed}를 보라. 그러므로
보조정리의 증명을 끝내려면 Lemma
\ref{spaces-more-morphisms-lemma-where-unramified-universally-injective}의 가정들이 성립함을 보이면
충분하다. 사영 $\text{pr}_0 : X \times_Y X \to X$는 $f$의 밑변환이므로
보편 닫힌 사상이다. Morphisms of Spaces, Lemma
\ref{spaces-morphisms-lemma-base-change-universally-closed}를 보라. 따라서
$X \times_Y X \to Z$는 보편 닫힌 사상들의 합성이므로 보편 닫힌 사상이다
(Morphisms of Spaces, Lemma
\ref{spaces-morphisms-lemma-composition-universally-closed}를 보라). 이것으로
증명이 끝난다.
\end{proof}

\begin{lemma}
\label{spaces-more-morphisms-lemma-where-flat}
대수공간의 가환 그림
$$
\xymatrix{
X \ar[rr]_f \ar[rd]_p & & Y \ar[ld]^q \\
& Z
}
$$
을 생각하자. 다음을 가정하자.
\begin{enumerate}
\item $p$는 국소 유한 표시이다.
\item $p$는 평탄하다.
\item $p$는 닫힌 사상이다.
\item $q$는 국소 유한형이다.
\end{enumerate}
그러면 열린 부분공간 $W \subset Z$가 존재하여, 사상 $Z' \to Z$가 $W$를
통해 인수분해될 필요충분조건은 밑변환
$f_{Z'} : X_{Z'} \to Y_{Z'}$가 평탄한 것이다.
\end{lemma}

\begin{proof}
Lemma \ref{spaces-more-morphisms-lemma-base-change-flatness-fibres}에 의해 집합
$$
A = \{x \in |X| : X\text{ flat at }x \text{ over }Y\}.
$$
는 $|X|$에서 열려 있고 그 형성은 임의의 밑변환과 교환한다. $W \subset Z$를
다음 바탕 점 집합을 갖는 열린 부분공간이라 하자(Properties of Spaces,
Lemma \ref{spaces-properties-lemma-open-subspaces}를 보라).
$$
|W| = |Z| \setminus |p|\left(|X| \setminus A\right)
$$
즉, $z \in |Z|$가 $W$의 점일 필요충분조건은 $|X| \to |Z|$의 $z$ 위
전체 올이 $A$에 포함되는 것이다. $p$가 닫힌 사상이므로 이는 열린 집합이다.
$A$의 형성이 임의의 밑변환과 교환하므로 $W$가 원하는 성질을 갖는다.
\end{proof}

\begin{lemma}
\label{spaces-more-morphisms-lemma-where-surjective-flat}
대수공간의 가환 그림
$$
\xymatrix{
X \ar[rr]_f \ar[rd]_p & & Y \ar[ld]^q \\
& Z
}
$$
을 생각하자. 다음을 가정하자.
\begin{enumerate}
\item $p$는 국소 유한 표시이다.
\item $p$는 평탄하다.
\item $p$는 닫힌 사상이다.
\item $q$는 국소 유한형이다.
\item $q$는 닫힌 사상이다.
\end{enumerate}
그러면 열린 부분공간 $W \subset Z$가 존재하여, 사상 $Z' \to Z$가 $W$를
통해 인수분해될 필요충분조건은 밑변환
$f_{Z'} : X_{Z'} \to Y_{Z'}$가 전사이고 평탄한 것이다.
\end{lemma}

\begin{proof}
Lemma \ref{spaces-more-morphisms-lemma-where-flat}에 의해 $f$가 평탄하다고 가정할 수 있다.
Morphisms of Spaces, Lemma
\ref{spaces-morphisms-lemma-finite-presentation-permanence}에 의해 $f$는 국소
유한 표시임에 유의하자. 따라서 $f$는 열린 사상이다. Morphisms of Spaces,
Lemma \ref{spaces-morphisms-lemma-fppf-open}을 보라. $W \subset Z$를 다음
바탕 점 집합을 갖는 열린 부분공간이라 하자(Properties of Spaces, Lemma
\ref{spaces-properties-lemma-open-subspaces}를 보라).
$$
|W| = |Z| \setminus |q|\left(|Y| \setminus |f|(|X|)\right).
$$
달리 말하면 $z \in |Z|$에 대하여 $z \in |W|$일 필요충분조건은
$|Y| \to |Z|$의 $z$ 위 전체 올이 $|X| \to |Y|$의 상에 포함되는 것이다.
$q$가 닫힌 사상이므로 이 집합은 $|Z|$에서 열린다. 사상
$X_W \to Y_W$는 구성에 의해 전사이다. 끝으로 $X_{Z'} \to Y_{Z'}$가
전사라고 가정하자. $Z' \to Z$가 $W$를 통해 인수분해됨을 보이려면
$|Z'| \to |Z|$의 상이 $|W|$에 포함됨을 보이면 충분하다. Properties of
Spaces, Lemma \ref{spaces-properties-lemma-factor-through-open-subspace}를
보라. 따라서 $Z'$가 체의 스펙트럼인 경우를 증명하면 충분하다.
$z \in |Z|$를 $|Z'| \to |Z|$의 상이라 하자. Properties of Spaces,
Lemma \ref{spaces-properties-lemma-points-cartesian}에 의해 가환 그림
$$
\xymatrix{
|X_{Z'}| \ar[d] \ar[r] &
|Y_{Z'}| \ar[d] \\
|p|^{-1}(\{z\}) \ar[r] &
|q|^{-1}(\{z\})
}
$$
에서 수직 화살표들은 전사이다. 따라서 원하는 대로 $z \in |W|$이다.
\end{proof}

\begin{lemma}
\label{spaces-more-morphisms-lemma-where-isomorphism}
대수공간의 가환 그림
$$
\xymatrix{
X \ar[rr]_f \ar[rd]_p & & Y \ar[ld]^q \\
& Z
}
$$
을 생각하자. 다음을 가정하자.
\begin{enumerate}
\item $p$는 국소 유한 표시이다.
\item $p$는 평탄하다.
\item $p$는 보편 닫힌 사상이다.
\item $q$는 국소 유한형이다.
\item $q$는 닫힌 사상이다.
\item $q$는 분리 사상이다.
\end{enumerate}
그러면 열린 부분공간 $W \subset Z$가 존재하여, 사상 $Z' \to Z$가 $W$를
통해 인수분해될 필요충분조건은 밑변환
$f_{Z'} : X_{Z'} \to Y_{Z'}$가 동형인 것이다.
\end{lemma}

\begin{proof}
Lemma \ref{spaces-more-morphisms-lemma-where-surjective-flat}에 의해 열린 부분공간
$W_1 \subset Z$가 존재하여, $f_{Z'}$가 전사이고 평탄할 필요충분조건은
$Z' \to Z$가 $W_1$을 통해 인수분해되는 것이다. Lemma
\ref{spaces-more-morphisms-lemma-where-closed-immersion}에 의해 열린 부분공간 $W_2 \subset Z$가
존재하여, $f_{Z'}$가 닫힌 몰입일 필요충분조건은 $Z' \to Z$가 $W_2$를
통해 인수분해되는 것이다. $W = W_1 \cap W_2$가 원하는 부분공간이라고
주장한다. 분명히 $f_{Z'}$가 동형이면 $Z' \to Z$는 $W$를 통해
인수분해된다. 따라서 $f_W$가 동형임을 보이면 충분하다. 구성에 의해
$f_W$는 전사이고 평탄한 닫힌 몰입이다. 특히 $f_W$는 표현 가능하다.
스킴의 전사 평탄 닫힌 몰입은 동형이므로(Morphisms, Lemma
\ref{morphisms-lemma-characterize-flat-closed-immersions}를 보라) 결론이
따른다. (실제로 $f_W$는 국소 유한 표시이고 따라서 열린 사상이므로, 원한다면
이 보조정리를 사용하지 않아도 된다.)
\end{proof}

\begin{lemma}
\label{spaces-more-morphisms-lemma-where-lci}
대수공간의 가환 그림
$$
\xymatrix{
X \ar[rr]_f \ar[rd]_p & & Y \ar[ld]^q \\
& Z
}
$$
을 생각하자. 다음을 가정하자.
\begin{enumerate}
\item $p$는 평탄하고 국소 유한 표시이다.
\item $p$는 닫힌 사상이다.
\item $q$는 평탄하고 국소 유한 표시이다.
\end{enumerate}
그러면 열린 부분공간 $W \subset Z$가 존재하여, 사상 $Z' \to Z$가 $W$를
통해 인수분해될 필요충분조건은 밑변환
$f_{Z'} : X_{Z'} \to Y_{Z'}$가 국소 완전 교차 사상인 것이다.
\end{lemma}

\begin{proof}
Lemma \ref{spaces-more-morphisms-lemma-base-change-lci-fibres}에 의해 열린 부분공간
$U(f) \subset X$가 존재하며, 이는 $f$가 코쥘인 점들의 집합이다. 또한
$U(f)$의 형성은 임의의 밑변환과 교환한다. $W \subset Z$를 다음 바탕 점
집합을 갖는 열린 부분공간이라 하자(Properties of Spaces, Lemma
\ref{spaces-properties-lemma-open-subspaces}를 보라).
$$
|W| = |Z| \setminus |p|\left(|X| \setminus |U(f)|\right)
$$
즉, $z \in |Z|$가 $W$의 점일 필요충분조건은 $f$가 $X$의 모든 점 가운데
$z$ 위에 놓인 점에서 코쥘인 것이다. $p$가 닫힌 사상이라고 가정했으므로
이는 열린 집합이다. $U(f)$의 형성이 임의의 밑변환과 교환하므로 Properties
of Spaces, Lemma \ref{spaces-properties-lemma-factor-through-open-subspace}를
사용하면 $W$가 원하는 보편 성질을 가짐을 즉시 알 수 있다.
\end{proof}

\section{미분 층과 여법선 층의 완전열}
\label{spaces-more-morphisms-section-exact}

\noindent
이 절에서는 여법선 층과 미분 층의 완전열에 관한 몇 가지 결과를 모은다.
어떤 의미에서 이들은 모두 대수공간의 합성 가능한 사상들에 결부된 여접
복합체의 삼각형을 실현한 것이다.

\medskip\noindent
아래 열들의 각 사상은 Lemma \ref{spaces-more-morphisms-lemma-functoriality-differentials} 또는
Lemma \ref{spaces-more-morphisms-lemma-universal-thickening-functorial}에서 구성한 사상이다.
$S$를 스킴이라 하자. $g : Z \to Y$와 $f : Y \to X$를 $S$ 위 대수공간의
사상이라 하자.
\begin{enumerate}
\item 표준적인 완전열
$$
g^*\Omega_{Y/X} \to \Omega_{Z/X} \to \Omega_{Z/Y} \to 0,
$$
이 있다. Lemma \ref{spaces-more-morphisms-lemma-triangle-differentials}를 보라.
$g : Z \to Y$가 형식적으로 매끄러우면 이 열은 짧은 완전열이다. Lemma
\ref{spaces-more-morphisms-lemma-triangle-differentials-formally-smooth}를 보라.
\item $g$가 형식적으로 비분기이면 표준적인 완전열
$$
\mathcal{C}_{Z/Y} \to g^*\Omega_{Y/X} \to \Omega_{Z/X} \to 0,
$$
이 있다. Lemma \ref{spaces-more-morphisms-lemma-universally-unramified-differentials-sequence}를
보라. $f \circ g : Z \to X$가 형식적으로 매끄러우면 이 열은 짧은
완전열이다. Lemma
\ref{spaces-more-morphisms-lemma-differentials-formally-unramified-formally-smooth}를 보라.
\item $g$와 $f \circ g$가 형식적으로 비분기이면 표준적인 완전열
$$
\mathcal{C}_{Z/X} \to \mathcal{C}_{Z/Y} \to g^*\Omega_{Y/X} \to 0,
$$
이 있다. Lemma \ref{spaces-more-morphisms-lemma-two-unramified-morphisms}를 보라.
$f : Y \to X$가 형식적으로 매끄러우면 이 열은 짧은 완전열이다. Lemma
\ref{spaces-more-morphisms-lemma-two-unramified-morphisms-formally-smooth}를 보라.
\item $g$와 $f$가 형식적으로 비분기이면 표준적인 완전열
$$
g^*\mathcal{C}_{Y/X} \to \mathcal{C}_{Z/X} \to \mathcal{C}_{Z/Y} \to 0.
$$
이 있다. Lemma \ref{spaces-more-morphisms-lemma-transitivity-conormal-unramified}를 보라.
$g : Z \to Y$가 국소 완전 교차 사상이면 이 열은 짧은 완전열이다. Lemma
\ref{spaces-more-morphisms-lemma-transitivity-conormal-lci}를 보라.
\end{enumerate}










\section{유사연접 복합체의 특징짓기, II}
\label{spaces-more-morphisms-section-characterize-pseudo-coherent}

\noindent
이 절에서는 유사연접 복합체를 코호몰로지로 특징짓는 방법을 논의한다.
대수공간 위 유사연접 복합체에 관한 앞선 내용은 Derived Categories of
Spaces, Section \ref{spaces-perfect-section-spell-out} 및 Derived Categories
of Spaces, Section \ref{spaces-perfect-section-pseudo-coherent-hocolim}에서
찾을 수 있다. 이 절의 스킴에 대한 대응 결과는 More on Morphisms, Section
\ref{more-morphisms-section-characterize-pseudo-coherent}이다. 기본 도구는
차우 보조정리의 유도 버전을 사용하여 사영공간의 경우로 환원하는 것이다.
Lemma \ref{spaces-more-morphisms-lemma-derived-chow}를 보라.

\begin{lemma}
\label{spaces-more-morphisms-lemma-case-of-tor-independence}
$S$를 스킴이라 하자. 대수공간의 가환 그림
$$
\xymatrix{
Z' \ar[d] \ar[r] & Y' \ar[d] \\
X' \ar[r] & B'
}
$$
을 생각하자. 이 그림은 $S$ 위에 있다. $B \to B'$를 사상이라 하자. $X$와
$Y$로 각각 $X'$와 $Y'$의
$B$로의 밑변환을 나타내자. $Y' \to B'$와 $Z' \to X'$가 평탄하다고
가정하자. 그러면 $X \times_B Y$와 $Z'$는 $X' \times_{B'} Y'$ 위에서
Tor-독립이다.
\end{lemma}

\begin{proof}
Derived Categories of Spaces, Lemma
\ref{spaces-perfect-lemma-tor-independent}에 의해 Tor-독립성은
$X \times_B Y$와 $Z'$ 위에서 에탈 국소적으로 확인할 수 있다.
이로써\footnote{논증을 더 자세히 적으면 다음과 같다. 전사 에탈 사상
$W' \to B'$를 택하되, $W'$는 스킴이라 하자. 전사 에탈 사상
$W \to B \times_{B'} W'$를 택하되, $W$는 스킴이라 하자. 전사 에탈 사상
$U' \to X' \times_{B'} W'$를 택하되, $U'$은 스킴이라 하자. 전사 에탈 사상
$V' \to Y' \times_{B'} W'$를 택하되, $V'$은 스킴이라 하자. 사상
$U' \times_{W'} V' \to X' \times_{B'} Y'$가 전사 에탈임에 유의하자.
전사 에탈 사상
$T' \to Z' \times_{X' \times_{B'} Y'} U' \times_{W'} V'$를 택하되,
$T'$은 스킴이라 하자.
$U$와 $V$로 각각 $U'$과 $V'$의 $W$로의 밑변환을 나타내자. 그러면
보조정리의 주장은 대수공간으로서 $X \times_B Y$와 $Z'$가
$X' \times_{B'} Y'$ 위에서 Tor-독립일 필요충분조건이 스킴으로서
$U \times_W V$와 $T'$가 $U' \times_{W'} V'$ 위에서 Tor-독립이라는 것이다.
따라서 꼭짓점이 $T', U', V', W'$인 정사각형과 $W \to W'$에 의한 밑변환에
대해 보조정리를 증명하면 충분하다. $Y' \to B'$와 $Z' \to X'$의 평탄성은
$V' \to W'$와 $T' \to U'$의 평탄성을 함의한다.}
보조정리는 스킴의 경우로 환원되며, 그 결과는 More on Morphisms, Lemma
\ref{more-morphisms-lemma-case-of-tor-independence}이다.
\end{proof}

\begin{lemma}[유도 차우 보조정리]
\label{spaces-more-morphisms-lemma-derived-chow}
$A$를 환이라 하자. $X$를 $A$ 위에서 유한 표시이고 분리인 대수공간이라
하자. $x \in |X|$라 하자. 그러면 다음 조건들을 만족하는 $n \geq 0$,
닫힌 부분공간 $Z \subset X \times_A \mathbf{P}^n_A$, 점 $z \in |Z|$,
열린 부분 $V \subset \mathbf{P}^n_A$, 그리고 대상 $E$가 존재하며, 이 대상은
$D(\mathcal{O}_{X \times_A \mathbf{P}^n_A})$에 속한다.
\begin{enumerate}
\item $Z \to X \times_A \mathbf{P}^n_A$는 유한 표시이다.
\item $c : Z \to \mathbf{P}^n_A$는 $V$ 위에서 닫힌 몰입이다.
$W = c^{-1}(V)$로 놓는다.
\item $b : Z \to X$의 $W$로의 제한은 에탈이고, $z \in W$이며,
$b(z) = x$이다.
\item $E|_{X \times_A V} \cong
(b, c)_*\mathcal{O}_Z|_{X \times_A V}$이다.
\item $E$는 유사연접이고 그 지지는 $Z$에 포함된다.
\end{enumerate}
\end{lemma}

\begin{proof}
유한형인 $\mathbf{Z}$-부분대수 $A' \subset A$와 대수공간 $X'$을 찾을 수
있는데, 이는 $A'$ 위에서 유한 표시이고 분리이며 $A$로 밑변환하면 $X$가
된다.
Limits of Spaces, Lemmas
\ref{spaces-limits-lemma-descend-finite-presentation}와
\ref{spaces-limits-lemma-descend-separated-morphism}를 보라.
$x' \in |X'|$를 $x$의 상이라 하자. $(X'/A', x')$에 대하여 보조정리를
증명할 수 있다면 $(X/A, x)$에 대한 보조정리가 따른다. 실제로
$n', Z', z', V', E'$가 $(X'/A', x')$에 대한 해를 준다면 다음과 같이 놓을
수 있다. $n = n'$로 놓고, $Z \subset X \times \mathbf{P}^n$를 $Z'$의
역상으로 놓고,
% P09-ERR-0753
$z \in Z$를 $x$로 가는 유일한 점으로 놓고, $V \subset \mathbf{P}^n_A$를
$V'$의 역상으로 놓으며, $E$를 $E'$의 유도 역상으로 놓는다. Cohomology on
Sites, Lemma \ref{sites-cohomology-lemma-pseudo-coherent-pullback}에 의해
$E$가 유사연접임에 유의하자. 이제 (5)만 확인하면 된다.
% P09-ERR-0754
이를 위해 $W = c^{-1}(V)$와 $W' = (c')^{-1}(V')$로 놓고 카르테시안
정사각형
$$
\xymatrix{
W \ar[d]_{(b, c)} \ar[r] & W' \ar[d]^{(b', c')} \\
X \times_A V \ar[r] & X' \times_{A'} V'
}
$$
을 생각하자. Lemma \ref{spaces-more-morphisms-lemma-case-of-tor-independence}에 의해
$X \times_A V$와 $W'$는 $X' \times_{A'} V'$ 위에서 Tor-독립이다. 따라서
$(b', c')_*\mathcal{O}_{W'}$의 $X \times_A V$로의 유도 역상은 Derived
Categories of Spaces, Lemma
\ref{spaces-perfect-lemma-compare-base-change}에 의해
$(b, c)_*\mathcal{O}_W$이다. 여기서는
$R(b', c')_*\mathcal{O}_{Z'} = (b', c')_*\mathcal{O}_{Z'}$라는 사실도
사용한다. 이는 $(b', c')$가 닫힌 몰입이기 때문이며,
$(b, c)_*\mathcal{O}_Z$에 대해서도 마찬가지이다.
% P09-ERR-0755
$E'|_{U' \times_{A'} V'} = (b', c')_*\mathcal{O}_{W'}$이므로
$E|_{U \times_A V} = (b, c)_*\mathcal{O}_W$를 얻고 (5)가 성립한다.
% P09-ERR-0754
이로써 다음 문단에서 기술하는 상황으로 환원된다.

\medskip\noindent
$A$가 $\mathbf{Z}$ 위 유한형이라고 가정하자. 에탈 사상 $U \to X$를
택하되, $U$는 아핀 스킴이고 점 $u \in U$는 $x$로 간다고 하자. 그러면 $U$는
$A$ 위 유한형이다. 닫힌 몰입 $U \to \mathbf{A}^n_A$를 택하고,
% P09-ERR-0756
그 합성으로 얻는 몰입을 $j : U \to \mathbf{P}^n_A$로 나타내자. 여기서
합성에는 열린 몰입 $\mathbf{A}^n_A \to \mathbf{P}^n_A$를 사용한다. 다음
사상의 스킴론적 폐포를 $Z$라 하자.
$$
(\text{id}_U, j) : U \longrightarrow X \times_A \mathbf{P}^n_A
$$
$z \in Z$를 $u$의 상이라 하자. $Y \subset \mathbf{P}^n_A$를 $j$의
스킴론적 폐포라 하자. 그러면 $Z \subset X \times_A Y$는 다음 사상의
스킴론적 폐포임이 분명하다.
$(\text{id}_U, j) : U \to X \times_A Y$.
$X$가 분리이므로 사상 $X \times_A Y \to Y$도 분리이다. 따라서 Morphisms
of Spaces, Lemma
\ref{spaces-morphisms-lemma-scheme-theoretic-image-of-partial-section}에 의해
$Z \to Y$는 열린 부분스킴 $j(U) \subset Y$ 위에서 동형이다.
$V \subset \mathbf{P}^n_A$를 $V \cap Y = j(U)$인 열린 부분이라 하자.
그러면 (2)가 성립하고 $W = (\text{id}_U, j)(U)$이며, 따라서 (3)도
성립한다. $A$가 뇌터 환이므로 (1)이 성립한다.

\medskip\noindent
$A$가 뇌터 환이므로 $X$와 $X \times_A \mathbf{P}^n_A$는 뇌터
대수공간이다. 따라서 이 경우 $E = (b, c)_*\mathcal{O}_Z$로 택할 수 있다.
(4)는 분명하고, (5)에 대해서는 Derived Categories of Spaces, Lemma
\ref{spaces-perfect-lemma-identify-pseudo-coherent-noetherian}를 보라.
이것으로 증명이 끝난다.
\end{proof}

\begin{lemma}
\label{spaces-more-morphisms-lemma-compute-Fourier-Mukai-for-derived-chow}
$X/A$, $x \in |X|$, 그리고 $n, Z, z, V, E$가 Lemma
\ref{spaces-more-morphisms-lemma-derived-chow}에서와 같다고 하자. 임의의
$K \in D_\QCoh(\mathcal{O}_X)$에 대하여
$$
Rq_*(Lp^*K \otimes^\mathbf{L} E)|_V = R(W \to V)_*K|_W
$$
이다. 여기서 $p : X \times_A \mathbf{P}^n_A \to X$와
$q : X \times_A \mathbf{P}^n_A \to \mathbf{P}^n_A$는 사영이고,
사상 $W \to V$는 유한 표시 닫힌 몰입 $c|_W : W \to V$이다.
\end{lemma}

\begin{proof}
$W = c^{-1}(V)$이고 $c$가 $V$ 위에서 닫힌 몰입이므로 $c|_W$는 닫힌
몰입이다. $W$와 $V$가 $A$ 위 유한 표시이므로 이 사상도 유한 표시이다.
Morphisms of Spaces, Lemma
\ref{spaces-morphisms-lemma-finite-presentation-permanence}를 보라. 먼저
$$
Rq_*(Lp^*K \otimes^\mathbf{L} E)|_V =
Rq'_*\left((Lp^*K \otimes^\mathbf{L} E)|_{X \times_A V}\right)
$$
이다. 전체 직접상의 형성은 국소화와 교환하므로 여기서
$q' : X \times_A V \to V$는 사영이다.
% P09-ERR-0757
주어진 닫힌 몰입을 $i = (b, c)|_W : W \to X \times_A V$로 나타내자.
그러면 성질 (5)에 의해
% P09-ERR-1457
$$
Rq'_*\left((Lp^*K \otimes^\mathbf{L} E)|_{X \times_A V}\right) =
Rq'_*(Lp^*K|_{X \times_A V} \otimes^\mathbf{L} i_*\mathcal{O}_W)
$$
이다. $i$가 닫힌 몰입이므로
$i_*\mathcal{O}_W = Ri_*\mathcal{O}_W$이다. Derived Categories of Spaces,
Lemma \ref{spaces-perfect-lemma-cohomology-base-change}를 사용하면 이를
다음과 같이 다시 쓸 수 있다.
$$
Rq'_* Ri_* Li^* Lp^*K|_{X \times_A V} =
R(q' \circ i)_* Lb^*K|_W =
R(W \to V)_* K|_W
$$
이것이 원하는 식이다. ($W$로 제한하는 것과 $W \to X$를 따라 유도 역상을
취하는 것은 같은데, 이는 $W$가 $X$ 위 에탈이기 때문이다.)
\end{proof}

\begin{lemma}
\label{spaces-more-morphisms-lemma-characterize-pseudo-coherent}
$A$를 환이라 하자. $X$를 $A$ 위에서 유한 표시이고 분리인 대수공간이라
하자. $K \in D_\QCoh(\mathcal{O}_X)$라 하자.
$R\Gamma(X, E \otimes^\mathbf{L} K)$가 $D(A)$에서 유사연접이고 이것이 모든
유사연접 대상 $E$에 대해 성립하며 그 대상이 $D(\mathcal{O}_X)$에 속한다고
하자. 그러면 $K$는
$A$에 대하여 유사연접이다(Definition
\ref{spaces-more-morphisms-definition-relative-pseudo-coherence}).
\end{lemma}

\begin{proof}
$K \in D_\QCoh(\mathcal{O}_X)$이고
$R\Gamma(X, E \otimes^\mathbf{L} K)$가 $D(A)$에서 유사연접이며 이것이 모든
유사연접 대상 $E$에 대해 성립하고 그 대상이 $D(\mathcal{O}_X)$에 속한다고
가정하자.
$x \in |X|$라 하자. $K$가 $A$에 대하여 유사연접임을 $x$의 에탈 근방에서
유사연접임을 보이겠다. 그러면 상대 유사연접성의 정의에 의해 보조정리가
따른다.

\medskip\noindent
Lemma \ref{spaces-more-morphisms-lemma-derived-chow}에서와 같이 $n, Z, z, V, E$를 택하자.
사영들을 $p : X \times \mathbf{P}^n \to X$ 및
$q : X \times \mathbf{P}^n \to \mathbf{P}^n_A$로 나타내자. 그러면 임의의
$i \in \mathbf{Z}$에 대하여
\begin{align*}
& R\Gamma(\mathbf{P}^n_A,
Rq_*(Lp^*K \otimes^\mathbf{L} E)
\otimes^\mathbf{L}
\mathcal{O}_{\mathbf{P}^n_A}(i)) \\
& =
R\Gamma(X \times \mathbf{P}^n,
Lp^*K \otimes^\mathbf{L} E \otimes^\mathbf{L}
Lq^*\mathcal{O}_{\mathbf{P}^n_A}(i)) \\
& =
R\Gamma(X,
K \otimes^\mathbf{L} Rp_*(E \otimes^\mathbf{L}
Lq^*\mathcal{O}_{\mathbf{P}^n_A}(i)))
\end{align*}
이다. Derived Categories of Spaces, Lemma
\ref{spaces-perfect-lemma-cohomology-base-change}를 보라. Derived Categories
of Spaces, Lemma
\ref{spaces-perfect-lemma-flat-proper-pseudo-coherent-direct-image-general}에
의해 복합체
$Rp_*(E \otimes^\mathbf{L} Lq^*\mathcal{O}_{\mathbf{P}^n_A}(i))$는 $X$ 위에서
유사연접이다. 따라서 가정에 의해 위 표시식은 $D(A)$의 유사연접 대상이다.
Derived Categories of Schemes, Lemma
\ref{perfect-lemma-pseudo-coherent-on-projective-space}에 의해
$Rq_*(Lp^*K \otimes^\mathbf{L} E)$가 $\mathbf{P}^n_A$ 위에서 유사연접이라는
결론을 얻는다. Lemma \ref{spaces-more-morphisms-lemma-compute-Fourier-Mukai-for-derived-chow}에 의해
% P09-ERR-0758
$$
Rq_*(Lp^*K \otimes^\mathbf{L} E)|_{X \times_A V} =
R(W \to V)_*K|_W
$$
이다. $W \to V$가 $\mathbf{P}^n_A$의 열린 부분스킴으로 들어가는 닫힌
몰입이므로, 이는 $K|_W$가 $A$에 대하여 유사연접임을 뜻한다. 예를 들어
More on Morphisms, Lemma
\ref{more-morphisms-lemma-check-relative-pseudo-coherence-on-charts}를 보라.
\end{proof}

\begin{lemma}
\label{spaces-more-morphisms-lemma-characterize-pseudo-coh-improved}
$A$를 환이라 하자. $X$를 $A$ 위에서 유한 표시이고 분리인 대수공간이라
하자. $K \in D_\QCoh(\mathcal{O}_X)$라 하자.
$R \Gamma (X, E \otimes ^{\mathbf{L}} K)$가 $D(A)$에서 유사연접이고 이것이
모든 완전 대상 $E \in D(\mathcal{O}_X)$에 대해 성립한다고 하자. 그러면 $K$는
$A$에 대하여 유사연접이다.
\end{lemma}

\begin{proof}
Lemma \ref{spaces-more-morphisms-lemma-characterize-pseudo-coherent}에 비추어 보면,
$R \Gamma (X, E \otimes ^{\mathbf{L}} K)$가 $D(A)$에서 유사연접이고 이것이
모든 유사연접 대상 $E \in D(\mathcal{O}_X)$에 대해 성립함을 보이면
충분하다. Derived Categories of Spaces, Proposition
\ref{spaces-perfect-proposition-detecting-bounded-above}에 의해
$K \in D^-_\QCoh (\mathcal{O}_X)$이다. 이제 Derived Categories of Spaces,
Lemma \ref{spaces-perfect-lemma-perfect-enough}에 의해 결과가 따른다.
\end{proof}










\section{상대적으로 완전한 대상}
\label{spaces-more-morphisms-section-relatively-perfect}

\noindent
이 절에서는 \cite{lieblich-complexes}의 개념을 도입한다. 이 개념은 스킴의
사상에 대해서 Derived Categories of Schemes, Section
\ref{perfect-section-relatively-perfect}에서 논의했다.

\begin{definition}
\label{spaces-more-morphisms-definition-relatively-perfect}
$S$를 스킴이라 하자. $f : X \to Y$를 $S$ 위 대수공간의 평탄하고 국소
유한 표시인 사상이라 하자. 대상 $E$가 $D(\mathcal{O}_X)$에 속한다고 하자.
이를 {\it $Y$에 대하여 완전하다} 또는 {\it $Y$-완전하다}고 하는 것은 다음
두 조건이 성립할 때이다. $E$는 유사연접이고(Cohomology on Sites,
Definition \ref{sites-cohomology-definition-pseudo-coherent}), $E$는
$D(f^{-1}\mathcal{O}_Y)$의 대상으로서 국소적으로 유한 Tor 차원을 갖는다
(Cohomology on Sites, Definition
\ref{sites-cohomology-definition-tor-amplitude}).
\end{definition}

\noindent
이에 관한 논의는 Derived Categories of Schemes, Remark
\ref{perfect-remark-discuss-rel-perfect}를 보라. 여기서는 Lemma
\ref{spaces-more-morphisms-lemma-relative-pseudo-coherent-is-moot}에 의해 $E$가 유사연접이라는
것이 $E$가 $Y$에 대하여 유사연접이라는 것과 같다는 점만 언급한다. 또한
$E$의 유사연접성은 $E \in D_\QCoh(\mathcal{O}_X)$를 함의한다. Derived
Categories of Spaces, Lemma \ref{spaces-perfect-lemma-pseudo-coherent}를 보라.

\begin{example}
\label{spaces-more-morphisms-example-relatively-perfect-field}
$k$를 체라 하자. $X$를 $k$ 위 유한 표시인 대수공간이라 하자(특히 $X$는
준콤팩트이다). 그러면 대상 $E$가 $D(\mathcal{O}_X)$에 속할 때 $k$-완전일
필요충분조건은 정의에 의해 유계이고 유사연접인 것이며, 이는 다시
$D^b_{\textit{Coh}}(X)$에 속하는 것과 동치이다(Derived Categories of
Spaces, Lemma
\ref{spaces-perfect-lemma-identify-pseudo-coherent-noetherian}). 따라서
상대적으로 완전하다는 것은 ``올 위에서 완전하다''는 뜻이 {\bf 아니다}.
\end{example}

\noindent
이에 대응하는 대수적 개념은 More on Algebra, Section
\ref{more-algebra-section-relatively-perfect}에서 연구한다. 대수공간에 대한
개념과 대수적 개념은 다음과 같이 연결된다.

\begin{lemma}
\label{spaces-more-morphisms-lemma-affine-locally-rel-perfect}
$S$를 스킴이라 하자. $f : X \to Y$를 $S$ 위 대수공간의 평탄하고 국소
유한 표시인 사상이라 하자. $E \in D_\QCoh(\mathcal{O}_X)$라 하자. 다음은
동치이다.
\begin{enumerate}
\item $E$는 $Y$-완전하다.
\item 다음 가환 그림
$$
\xymatrix{
U \ar[d] \ar[r]_g & V \ar[d] \\
X \ar[r]^f & Y
}
$$
에서 $U$, $V$가 스킴이고 수직 화살표들이 에탈일 때마다, 복합체 $E|_U$는
Derived Categories of Schemes, Definition
\ref{perfect-definition-relatively-perfect}의 의미에서 $V$-완전하다.
\item (2)와 같은 어떤 가환 그림에서 $U \to X$가 전사이고 복합체 $E|_U$가
Derived Categories of Schemes, Definition
\ref{perfect-definition-relatively-perfect}의 의미에서 $V$-완전하다.
\item (2)와 같은 모든 가환 그림에서 $U$와 $V$가 아핀이면 복합체
$R\Gamma(U, E)$는 $\mathcal{O}_Y(V)$-완전하다.
\end{enumerate}
\end{lemma}

\begin{proof}
보조정리의 (2), (3), (4)를 해석하기 위해 다음을 관찰하자.
$E|_U$가 $D_\QCoh(\mathcal{O}_U)$의 대상일 때 이에 대응하는 $E_0$는
$D_\QCoh(\mathcal{O}_{U_0})$의 대상이고, 여기서 $U_0$는 $U$의 바탕
스킴이다. Derived
Categories of Spaces, Lemma
\ref{spaces-perfect-lemma-derived-quasi-coherent-small-etale-site}를 보라.
또한 이 경우 $E_0$가 유사연접일 필요충분조건은 $E|_U$가 유사연접인
것이다. Derived Categories of Spaces, Lemma
\ref{spaces-perfect-lemma-descend-pseudo-coherent}를 보라. 그리고
$E|_U$가 $f^{-1}\mathcal{O}_Y|_U = g^{-1}\mathcal{O}_V$ 위에서 국소적으로
유한 Tor 차원을 가질 필요충분조건은 $E_0$가
$g_0^{-1}\mathcal{O}_{V_0}$ 위에서 국소적으로 유한 Tor 차원을 갖는
것이다. Derived Categories of Spaces, Lemma
\ref{spaces-perfect-lemma-tor-dimension-rel}을 보라. 여기서
$g_0 : U_0 \to V_0$는 $g : U \to V$를 나타내는 스킴 사상이다(표기는
Derived Categories of Spaces, Remark
\ref{spaces-perfect-remark-match-total-direct-images}를 따른다). 끝으로
``유사연접이다''라는 성질은 에탈 국소적이고, 물론 ``국소적으로 유한 Tor
차원을 갖는다''라는 성질도 에탈 국소적이다. 따라서 $Y$-완전성은 에탈
국소적으로 확인하면 충분하며, 위 논의에 의해 (1)은 (2)를 함의하고 (3)은
(1)을 함의한다. Derived Categories of Schemes, Lemma
\ref{perfect-lemma-affine-locally-rel-perfect}에 의해 (4)가 (2) 및 (3)과
동치이므로 증명이 끝난다.
\end{proof}

\begin{lemma}
\label{spaces-more-morphisms-lemma-triangulated}
$S$를 스킴이라 하자. $f : X \to Y$를 $S$ 위 대수공간의 평탄하고 국소
유한 표시인 사상이라 하자. $D(\mathcal{O}_X)$에서 $Y$-완전 대상들로 이루어진
충만한 부분범주는 포화된\footnote{Derived Categories, Definition
\ref{derived-definition-saturated}.} 삼각 부분범주이다.
\end{lemma}

\begin{proof}
이는 Cohomology on Sites, Lemmas
\ref{sites-cohomology-lemma-cone-pseudo-coherent},
\ref{sites-cohomology-lemma-summands-pseudo-coherent},
\ref{sites-cohomology-lemma-cone-tor-amplitude}, 그리고
\ref{sites-cohomology-lemma-summands-tor-amplitude}에서 따른다.
\end{proof}

\begin{lemma}
\label{spaces-more-morphisms-lemma-perfect-relatively-perfect}
$S$를 스킴이라 하자. $f : X \to Y$를 $S$ 위 대수공간의 평탄하고 국소
유한 표시인 사상이라 하자. $D(\mathcal{O}_X)$의 완전 대상은 $Y$-완전하다.
$K, M \in D(\mathcal{O}_X)$이면, $K \otimes_{\mathcal{O}_X}^\mathbf{L} M$이
$Y$-완전하기 위한 충분조건은 $K$가 완전하고 $M$이 $Y$-완전한 것이다.
\end{lemma}

\begin{proof}
Lemma \ref{spaces-more-morphisms-lemma-affine-locally-rel-perfect}를 사용하여 스킴의 경우로 환원한
뒤 Derived Categories of Schemes, Lemma
\ref{perfect-lemma-perfect-relatively-perfect}를 적용한다.
\end{proof}

\begin{lemma}
\label{spaces-more-morphisms-lemma-base-change-relatively-perfect}
$S$를 스킴이라 하자. $f : X \to Y$를 $S$ 위 대수공간의 평탄하고 국소
유한 표시인 사상이라 하자. $g : Y' \to Y$를 $S$ 위 대수공간의 사상이라
하자. $X' = Y' \times_Y X$로 놓고 사영을 $g' : X' \to X$로 나타내자.
$K \in D(\mathcal{O}_X)$가 $Y$-완전이면 $L(g')^*K$는 $Y'$-완전이다.
\end{lemma}

\begin{proof}
Lemma \ref{spaces-more-morphisms-lemma-affine-locally-rel-perfect}를 사용하여 스킴의 경우로 환원한
뒤 Derived Categories of Schemes, Lemma
\ref{perfect-lemma-base-change-relatively-perfect}를 적용한다.
\end{proof}

\begin{situation}
\label{spaces-more-morphisms-situation-relative-descent}
$S$를 스킴이라 하자. $Y = \lim_{i \in I} Y_i$를 $S$ 위 대수공간의 유향계의
극한이라 하며, 전이 사상 $g_{i'i} : Y_{i'} \to Y_i$는 아핀이라고 하자.
$Y_i$가 준콤팩트이고 준분리라고 가정하며 이는 모든 $i \in I$에 대해
성립한다고 하자. 사영을
$g_i : Y \to Y_i$로 나타낸다. 원소 $0 \in I$와 유한 표시 평탄 사상
$X_0 \to Y_0$를 고정한다. $X_i = Y_i \times_{Y_0} X_0$ 및
$X = Y \times_{Y_0} X_0$로 놓고, 전이 사상 $f_{i'i} : X_{i'} \to X_i$와
사영 $f_i : X \to X_i$를 나타낸다.
\end{situation}

\begin{lemma}
\label{spaces-more-morphisms-lemma-relative-descend-homomorphisms}
Situation \ref{spaces-more-morphisms-situation-relative-descent}의 상황이라 하자.
$K_0$와 $L_0$를 $D(\mathcal{O}_{X_0})$의 대상이라 하자.
$K_i = Lf_{i0}^*K_0$ 및 $L_i = Lf_{i0}^*L_0$로 놓되 이는 $i \geq 0$에
대하여 정의하고, $K = Lf_0^*K_0$ 및 $L = Lf_0^*L_0$로 놓자. 그러면 사상
$$
\colim_{i \geq 0} \Hom_{D(\mathcal{O}_{X_i})}(K_i, L_i)
\longrightarrow
\Hom_{D(\mathcal{O}_X)}(K, L)
$$
은 다음 조건에서 동형이다. 즉, $K_0$가 유사연접이고
$L_0 \in D_\QCoh(\mathcal{O}_{X_0})$가
$D((X_0 \to Y_0)^{-1}\mathcal{O}_{Y_0})$의 대상으로서 (국소적으로) 유한
Tor 차원을 갖는다.
\end{lemma}

\begin{proof}
각 준콤팩트 준분리 대상 $U_0$를 생각하되 이는
$(X_0)_{spaces, \etale}$의 대상이라고 하자. 다음 조건을 $P$라 하자.
$$
\colim_{i \geq 0} \Hom_{D(\mathcal{O}_{U_i})}(K_i|_{U_i}, L_i|_{U_i})
\longrightarrow
\Hom_{D(\mathcal{O}_U)}(K|_U, L|_U)
$$
이 동형이며, 여기서 $U = X \times_{X_0} U_0$이고
$U_i = X_i \times_{X_0} U_0$이다. $P$가 각 $U_0$에 대하여 성립함을
증명하겠다.

\medskip\noindent
$(U_0 \subset W_0, V_0 \to W_0)$가 $(X_0)_{spaces, \etale}$의 기본
구별 정사각형이고 $P$가 $U_0, V_0, U_0 \times_{W_0} V_0$에 대하여
성립한다고 하자. 그러면 유도 범주의 Hom에 대한 마이어--피토리스에 의해
$P$는 $W_0$에 대해서도 성립한다. Derived Categories of Spaces, Lemma
\ref{spaces-perfect-lemma-mayer-vietoris-hom}을 보라.

\medskip\noindent
먼저 $U_0 = W_0 \times_{Y_0} X_0$인 경우를 생각하자. 여기서 $W_0$는
$(Y_0)_{spaces, \etale}$의 준콤팩트 준분리 대상이다. Derived Categories of Spaces,
Lemma \ref{spaces-perfect-lemma-induction-principle}의 귀납 원리를 이
$W_0$들과 앞 문단에 적용하면 $P$를 $U_0 = W_0 \times_{Y_0} X_0$인 경우에
증명하면 충분함을 알 수 있으며, 여기서 $W_0$는 아핀이다.
달리 말하면 $Y_0$가 아핀인 경우로 환원했다. 다음으로 $X_0$의 모든
준콤팩트 준분리 열린 부분에 귀납 원리를 다시 적용하면 $X_0$도 아핀인
경우로 환원된다.

\medskip\noindent
$X_0$와 $Y_0$가 아핀이면 Derived Categories of Schemes, Lemma
\ref{perfect-lemma-relative-descend-homomorphisms}에서 증명한 스킴의
경우가 된다. 독자는 Derived Categories of Spaces, Lemmas
\ref{spaces-perfect-lemma-pseudo-coherent},
\ref{spaces-perfect-lemma-derived-quasi-coherent-small-etale-site},
\ref{spaces-perfect-lemma-descend-pseudo-coherent}, 그리고
\ref{spaces-perfect-lemma-tor-dimension-rel}을 사용하여 명제를 스킴과 스킴
위 가군의 유도 범주만을 포함하는 명제로 옮길 수 있다.
\end{proof}

\begin{lemma}
\label{spaces-more-morphisms-lemma-descend-relatively-perfect}
Situation \ref{spaces-more-morphisms-situation-relative-descent}에서 $Y$-완전 대상들로 이루어진
$D(\mathcal{O}_X)$의 범주는 $Y_i$-완전 대상들로 이루어진
$D(\mathcal{O}_{X_i})$의 범주들의 여극한이다.
\end{lemma}

\begin{proof}
각 준콤팩트 준분리 대상 $U_0$를 생각하되 이는
$(X_0)_{spaces, \etale}$의 대상이라고 하자. 다음 함자가 동치라는 조건을
$P$라 하자.
$$
\colim_{i \geq 0} D_{Y_i\text{-perfect}}(\mathcal{O}_{U_i})
\longrightarrow
D_{Y\text{-perfect}}(\mathcal{O}_U)
$$
여기서 $U = X \times_{X_0} U_0$이고 $U_i = X_i \times_{X_0} U_0$이다.
Lemma \ref{spaces-more-morphisms-lemma-relative-descend-homomorphisms}에 의해 이 함자가 이미
충실충만임을 알고 있다. 따라서 본질적 전사성만 증명하면 충분하다.

\medskip\noindent
$(U_0 \subset W_0, V_0 \to W_0)$가 $(X_0)_{spaces, \etale}$의 기본
구별 정사각형이고 $P$가 $U_0, V_0, U_0 \times_{W_0} V_0$에 대하여
성립한다고 하자. $P$가 $W_0$에 대해서도 성립한다고 주장한다. 표기
$U_i = X_i \times_{X_0} U_0$, $U = X \times_{X_0} U_0$를 사용하고,
$V_0$ 및 $W_0$에 대해서도 비슷한 표기를 사용하겠다. 기호 $f_i$를 다음
모든 사상에 남용하여 사용하겠다. 즉, $U \to U_i$, $V \to V_i$,
$U \times_W V \to U_i \times_{W_i} V_i$, 그리고 $W \to W_i$이다. $E$가
$Y$-완전 대상이고 $D(\mathcal{O}_W)$에 속한다고 하자. 목표는 $E$가 이 함자의
본질적 상에 속함을 보이는 것이다. 가정에 의해 $i \geq 0$, $Y_i$-완전
대상 $E_{U, i}$를 $U_i$ 위에서, 또 $Y_i$-완전 대상 $E_{V, i}$를 $V_i$
위에서 찾을 수 있고, 아울러
동형 $Lf_i^*E_{U, i} \to E|_U$ 및 $Lf_i^*E_{V, i} \to E|_V$를 찾을 수 있다.
다음 사상들을 이 동형들에 수반된 사상이라 하자.
$$
a : E_{U, i} \to (Rf_{i, *}E)|_{U_i}
\quad\text{and}\quad
b : E_{V, i} \to (Rf_{i, *}E)|_{V_i}
$$
즉, 각각 $Lf_i^*E_{U, i} \to E|_U$와
$Lf_i^*E_{V, i} \to E|_V$에 수반된다. 충실충만성에 의해 $i$를 늘린 뒤,
다음 동형을 찾을 수 있다.
$c : E_{U, i}|_{U_i \times_{W_i} V_i} \to
E_{V, i}|_{U_i \times_{W_i} V_i}$.
이 동형의 역상은 다음 식별이다.
$$
Lf_i^*E_{U, i}|_{U \times_W V} \to E|_{U \times_W V} \to
Lf_i^*E_{V, i}|_{U \times_W V}.
$$
Derived Categories of Spaces, Lemma \ref{spaces-perfect-lemma-glue}를
적용하면 대상 $E_i$를 $W_i$ 위에서 얻고 사상 $d : E_i \to Rf_{i, *}E$를
얻는데, 이는 사상 $a$와 $b$로 각각 $U_i$와 $V_i$ 위에서 제한된다. 그러면
$E_i$가 $Y_i$-완전임이 분명하다(상대적 완전성은 에탈 국소적 성질이다).
또한 $d$는 동형 $Lf_i^*E_i \to E$에 수반된다.

\medskip\noindent
Lemma \ref{spaces-more-morphisms-lemma-relative-descend-homomorphisms}의 증명과 정확히 같은 논증에서
귀납 원리(Derived Categories of Spaces, Lemma
\ref{spaces-perfect-lemma-induction-principle})를 사용하면 $X_0$와 $Y_0$가
모두 아핀인 경우로 환원된다. 먼저 $(Y_0)_{spaces, \etale}$의 준콤팩트
준분리 대상들을 사용하여 $Y_0$가 아핀인 경우로 환원하고, 이어서
$(X_0)_{spaces, \etale}$의 준콤팩트 준분리 대상들을 사용하여 $X_0$가
아핀인 경우로 환원한다. 아핀인 경우에는 Derived Categories of Schemes,
Lemma \ref{perfect-lemma-descend-relatively-perfect}인 스킴의 경우에서 결과가
따른다. 스킴의 경우로 옮기는 작업은 Lemma
\ref{spaces-more-morphisms-lemma-affine-locally-rel-perfect}가 해 준다.
\end{proof}

\begin{lemma}
\label{spaces-more-morphisms-lemma-derived-pushforward-rel-perfect}
$S$를 스킴이라 하자. $f : X \to Y$를 $S$ 위 대수공간의 평탄하고 고유하며
유한 표시인 사상이라 하자. $E \in D(\mathcal{O}_X)$가 $Y$-완전이라고 하자.
그러면 $Rf_*E$는 $D(\mathcal{O}_Y)$의 완전 대상이고, 그 형성은 임의의
밑변환과 교환한다.
\end{lemma}

\begin{proof}
밑변환에 관한 명제는 Derived Categories of Spaces, Lemma
\ref{spaces-perfect-lemma-base-change-tensor}이다(여기서
$\mathcal{G}^\bullet$은 $\mathcal{O}_X$와 같고 그 차수는 $0$이다). 따라서
$Rf_*E$가 완전 대상임을 보이면 충분하다. 극한 논증으로 $Y$가 뇌터 아핀인
경우로 환원하겠다.

\medskip\noindent
문제는 $Y$ 위 에탈 국소적이므로 $Y$가 아핀이라고 가정할 수 있다.
$Y = \Spec(R)$라 쓰자. $R = \colim R_i$를 뇌터 환 $R_i$들의 여과 여극한으로
쓰자. Limits of Spaces, Lemma
\ref{spaces-limits-lemma-descend-finite-presentation}에 의해 어떤 $i$와
$X_i$가 존재하며, 이는 $R_i$ 위 유한 표시인 대수공간이고 이를 $R$로 밑변환하면
$X$가 된다. Limits of Spaces, Lemmas
\ref{spaces-limits-lemma-eventually-proper}와
\ref{spaces-limits-lemma-descend-flat}에 의해 $X_i$가 $R_i$ 위에서 고유하고
평탄하다고 가정할 수 있다. Lemma \ref{spaces-more-morphisms-lemma-descend-relatively-perfect}에
의해 $R_i$-완전 대상 $E_i$가 존재한다고 가정할 수 있는데, 이는
$D(\mathcal{O}_{X_i})$의 대상이고 $X$로의 역상이 $E$이다. Derived
Categories of Spaces, Lemma \ref{spaces-perfect-lemma-perfect-direct-image}를
$X_i \to \Spec(R_i)$와 $E_i$에 적용하고 이미 증명한 밑변환 성질을 사용하면
결과를 얻는다.
\end{proof}

\begin{lemma}
\label{spaces-more-morphisms-lemma-compute-ext-rel-perfect}
$S$를 스킴이라 하자. $f : X \to Y$를 $S$ 위 대수공간의 사상이라 하자.
$E, K \in D(\mathcal{O}_X)$라 하자. 다음을 가정하자.
\begin{enumerate}
\item $Y$는 준콤팩트이고 준분리이다.
\item $f$는 고유하고 평탄하며 유한 표시이다.
\item $E$는 $Y$-완전하다.
\item $K$는 유사연접이다.
\end{enumerate}
그러면 다음을 만족하는 유사연접 대상 $L \in D(\mathcal{O}_Y)$이 존재한다.
$$
Rf_*R\SheafHom(K, E) = R\SheafHom(L, \mathcal{O}_Y)
$$
임의의 밑변환 뒤에도 같은 명제가 성립한다. 즉, 다음 그림이 주어지면
$$
\vcenter{
\xymatrix{
X' \ar[r]_{g'} \ar[d]_{f'} &
X \ar[d]^f \\
Y' \ar[r]^g &
Y
}
}
\quad\quad
\begin{matrix}
\text{카르테시안일 때 다음 식이 성립한다. } \\
Rf'_*R\SheafHom(L(g')^*K, L(g')^*E) \\
= R\SheafHom(Lg^*L, \mathcal{O}_{Y'})
\end{matrix}
$$
이다.
\end{lemma}

\begin{proof}
$Y$가 준콤팩트이고 준분리이므로 $X$도 그러하다. Derived Categories of
Spaces, Lemma \ref{spaces-perfect-lemma-pseudo-coherent-hocolim}에 의해
$K = \text{hocolim} K_n$으로 쓸 수 있으며, 여기서 $K_n$은 완전하고
$K_n \to K$는 절단 $\tau_{\geq -n}$ 위에서 동형을 유도한다.
$K_n^\vee$를 쌍대 완전 복합체라 하자(Cohomology on Sites, Lemma
\ref{sites-cohomology-lemma-dual-perfect-complex}). 완전 대상들의 역계
$$
\ldots \to K_3^\vee \to K_2^\vee \to K_1^\vee
$$
를 얻는다. Lemma \ref{spaces-more-morphisms-lemma-perfect-relatively-perfect}에 의해
$K_n^\vee \otimes_{\mathcal{O}_X} E$는 $Y$-완전하다. 따라서 Lemma
\ref{spaces-more-morphisms-lemma-derived-pushforward-rel-perfect}를
$K_n^\vee \otimes_{\mathcal{O}_X} E$에 적용하여 완전 복합체들의 역계
$$
\ldots \to M_3 \to M_2 \to M_1
$$
를 얻는다. 이들은 $Y$ 위에 있고, 여기서
$$
M_n = Rf_*(K_n^\vee \otimes_{\mathcal{O}_X}^\mathbf{L} E) =
Rf_*R\SheafHom(K_n, E)
$$
이다. 또한 이 복합체들의 형성은 임의의 밑변환과 교환한다. 즉,
$Lg^*M_n =
Rf'_*((L(g')^*K_n)^\vee \otimes_{\mathcal{O}_{X'}}^\mathbf{L} L(g')^*E) =
Rf'_*R\SheafHom(L(g')^*K_n, L(g')^*E)$이다.

\medskip\noindent
$K_n \to K$가 $\tau_{\geq -n}$ 위에서 동형을 유도하므로
$K_n \to K_{n + 1}$도 $\tau_{\geq -n}$ 위에서 동형을 유도한다. 따라서
$K_{n + 1}^\vee \to K_n^\vee$는 $\tau_{\leq n}$ 위에서 동형을 유도한다.
실제로 $K_n^\vee = R\SheafHom(K_n, \mathcal{O}_X)$이다. $E$가 Tor 진폭
$[a, b]$를 갖는다고 가정하되, 이는 $f^{-1}\mathcal{O}_Y$-가군의 복합체로서의
진폭이다. 그러면 임의의 밑변환 뒤에도 마찬가지이다. Derived Categories
of Spaces, Lemma
\ref{spaces-perfect-lemma-tor-independence-and-tor-amplitude}를 보라. 사상
$$
K_{n + 1}^\vee \otimes_{\mathcal{O}_X} E \to
K_n^\vee \otimes_{\mathcal{O}_X} E
$$
이 $\tau_{\leq n + a}$ 위에서 동형을 유도하며, 임의의 밑변환 뒤에도
마찬가지임을 얻는다. 우유도 함자 $Rf_*$를 적용하면 사상
$M_{n + 1} \to M_n$이 $\tau_{\leq n + a}$ 위에서 동형을 유도하고,
임의의 밑변환 뒤에도 마찬가지이다. 구별 삼각형을 택하자.
$$
M_{n + 1} \to M_n \to C_n \to M_{n + 1}[1]
$$
$y \in |Y|$를 택하자. 기본 에탈 근방 $(U, u) \to (Y, y)$를 택할 수 있다.
Decent Spaces, Lemma
\ref{decent-spaces-lemma-decent-space-elementary-etale-neighbourhood}를 보라.
$Y'$을 $u$에서의 잉여체의 스펙트럼으로 놓자. 역상을 취하면
$C_n|_U \otimes_{\mathcal{O}_U}^\mathbf{L} \kappa(u)$의 영이 아닌 코호몰로지는
차수 $\geq n + a$에만 존재한다. More on Algebra, Lemma
\ref{more-algebra-lemma-lift-perfect-from-residue-field}에 의해 완전 복합체
$C_n|_U$는 Tor 진폭 $[n + a, m_n]$를 가지며, 여기서 $m_n$은 어떤 정수이고
필요하면 $U$를 줄일 수 있다. 따라서 $C_n$은 Tor 진폭 $[n + a, m_n]$를
가지며, 여기서 $m_n$은 어떤 정수이다($Y$가 준콤팩트이기 때문이다). 특히 쌍대
완전 복합체 $C_n^\vee$는 Tor 진폭 $[-m_n, -n - a]$를 갖는다.

\medskip\noindent
$L_n = M_n^\vee$를 쌍대 완전 복합체라 하자. 앞 문단의 결론은 사상
$L_n \to L_{n + 1}$이 $\tau_{\geq -n - a}$ 위에서 동형을 유도한다는
것이다. 따라서 $L = \text{hocolim} L_n$은 유사연접이다. Derived Categories
of Spaces, Lemma \ref{spaces-perfect-lemma-pseudo-coherent-hocolim}를 보라.
또한
$$
R\SheafHom(K, E) = R\SheafHom(\text{hocolim} K_n, E) =
R\lim R\SheafHom(K_n, E) = R\lim K_n^\vee \otimes_{\mathcal{O}_X} E
$$
이고(Cohomology on Sites, Lemma
\ref{sites-cohomology-lemma-colim-and-lim-of-duals}), $R\lim$은 $Rf_*$와
교환하므로
$$
Rf_*R\SheafHom(K, E) = R\lim M_n = R\lim R\SheafHom(L_n, \mathcal{O}_Y) =
R\SheafHom(L, \mathcal{O}_Y)
$$
이다. 이로써 $Y$ 위의 식을 증명했다. $M_n$의 구성이 밑변환과 호환되므로
% P09-ERR-0766
임의의 밑변환 뒤에도 이 식이 계속 성립한다.
\end{proof}

\begin{remark}
\label{spaces-more-morphisms-remark-compare-L}
독자는 Lemma \ref{spaces-more-morphisms-lemma-compute-ext-rel-perfect}와 Derived Categories of
Spaces, Lemma \ref{spaces-perfect-lemma-compute-ext} 사이의 유사성을 눈치챘을
것이다. 실제로 Lemma \ref{spaces-more-morphisms-lemma-compute-ext-rel-perfect}의 유사연접 복합체
$L$은 다음과 같이 특징지을 수 있다. 이는 $Y$ 위의 유일한 유사연접
복합체로서 함자적인 동형
$$
\Ext^i_{\mathcal{O}_Y}(L, \mathcal{F}) \longrightarrow
\Ext^i_{\mathcal{O}_X}(K,
E \otimes_{\mathcal{O}_X}^\mathbf{L} Lf^*\mathcal{F})
$$
을 가지며, 이 동형들은 $\mathcal{F}$가 $\QCoh(\mathcal{O}_Y)$를 달릴 때
경계 사상들과 호환된다. 이 결과가 필요해지면 여기에 정확한 명제를 세우고
자세한 증명을 제시하겠다.
\end{remark}

\begin{lemma}
\label{spaces-more-morphisms-lemma-bounded-on-fibres}
$S$를 스킴이라 하자. $X$를 $S$ 위 대수공간이라 하고, 구조 사상
$f : X \to S$가 평탄하고 국소 유한 표시라고 하자. $E$를
$D(\mathcal{O}_X)$의 유사연접 대상이라 하자. 다음은 동치이다.
\begin{enumerate}
\item $E$는 $S$-완전하다.
\item $E$는 국소적으로 아래로 유계이고, 모든 점 $s \in S$에 대하여
대상 $L(X_s \to X)^*E$가 $D(\mathcal{O}_{X_s})$에 속하며 국소적으로
아래로 유계이다.
\end{enumerate}
\end{lemma}

\begin{proof}
모든 것이 국소적이므로 $X$와 $S$가 아핀인 경우로 즉시 환원된다. Lemma
\ref{spaces-more-morphisms-lemma-affine-locally-rel-perfect}를 보라. 이 경우는 Derived Categories
of Schemes, Lemma \ref{perfect-lemma-bounded-on-fibres}에서 다룬다.
\end{proof}

\begin{lemma}
\label{spaces-more-morphisms-lemma-characterize-relatively-perfect}
$A$를 환이라 하자. $X$를 $A$ 위에서 유한 표시이고 평탄하며 분리인
대수공간이라 하자. $K \in D_\QCoh(\mathcal{O}_X)$라 하자.
$R \Gamma (X, E \otimes^\mathbf{L} K)$가 $D(A)$에서 완전이고 이것이 모든
완전 대상 $E \in D(\mathcal{O}_X)$에 대하여 성립하면 $K$는
$\Spec(A)$-완전하다.
\end{lemma}

\begin{proof}
Lemma \ref{spaces-more-morphisms-lemma-characterize-pseudo-coh-improved}에 의해 $K$는 $A$에 대하여
유사연접이다. Lemma \ref{spaces-more-morphisms-lemma-relative-pseudo-coherent-is-moot}에 의해
$K$는 $D(\mathcal{O}_X)$에서 유사연접이다. Derived Categories of Spaces,
Proposition \ref{spaces-perfect-proposition-detecting-bounded-below}에 의해
$K$는 $D^-(\mathcal{O}_X)$에 속한다. $\mathfrak{p}$를 $A$의 소 아이디얼이라
하고,
% P09-ERR-0767
$i : Y \to X$로 $\mathfrak{p}$ 위의 스킴론적 올의 포함을 나타내자. 즉,
% P09-ERR-0768
$Y$는 $\kappa(\mathfrak p)$ 위의 스킴이다. Lemma
\ref{spaces-more-morphisms-lemma-bounded-on-fibres}에 의해 $Li^*(K)$가 아래로 유계임을 보이면
증명이 끝난다. $G \in D_{perf} (\mathcal{O}_X)$를
$D_\QCoh (\mathcal{O}_X)$를 생성하는 완전 복합체라 하자. Derived
Categories of Spaces, Theorem
\ref{spaces-perfect-theorem-bondal-van-den-Bergh}를 보라. 다음이 성립한다.
\begin{align*}
R\Hom _{\mathcal{O}_Y}(Li^*(G), Li^*(K))
& =
R\Gamma(Y, Li^*(G ^\vee \otimes ^\mathbf{L} K)) \\
& =
R\Gamma(X, G^\vee \otimes ^{\mathbf{L}} K)
\otimes^\mathbf{L}_A \kappa(\mathfrak{p})
\end{align*}
첫 번째 등식에서는 $Li^*$가 완전 대상과 쌍대를 보존한다는 사실 및
Cohomology on Sites, Lemma
\ref{sites-cohomology-lemma-dual-perfect-complex}를 사용한다. 몇 가지 세부사항은
생략한다. 두 번째 등식은 $X$가 $A$ 위에서 평탄하므로 Derived Categories
of Spaces, Lemma \ref{spaces-perfect-lemma-compare-base-change}에서 따른다.
가정에 의해 이는 $D(\kappa(\mathfrak{p}))$의 완전 대상이다. 대상
$Li^*(G) \in D_{perf}(\mathcal{O}_Y)$는 $D_\QCoh(\mathcal{O}_Y)$를 생성한다.
Derived Categories of Spaces, Remark
\ref{spaces-perfect-remark-pullback-generator}를 보라. 따라서 Derived
Categories of Spaces, Proposition
\ref{spaces-perfect-proposition-detecting-bounded-below}에 의해 이제
$Li^*(K)$가 아래로 유계이고, 결론이 따른다.
\end{proof}










\section{세제곱 정리}
\label{spaces-more-morphisms-section-theorem-cube}

\noindent
이 절은 More on Morphisms, Section
\ref{more-morphisms-section-theorem-cube}에 대응한다. 다음 보조정리는 그
가정 아래에서 피카르 함자의 대각이 국소 닫힌 몰입들로 표현 가능함을 말한다.

\begin{lemma}
\label{spaces-more-morphisms-lemma-diagonal-picard-flat-proper}
$S$를 스킴이라 하자. $f : X \to Y$를 $S$ 위 대수공간의 평탄하고 고유하며
유한 표시인 사상이라 하자. $\mathcal{E}$를 유한 국소 자유
$\mathcal{O}_X$-가군이라 하자. 사상 $g : Y' \to Y$에 대하여 밑변환 그림
$$
\xymatrix{
X' \ar[d]_{f'} \ar[r]_{g'} & X \ar[d]^f \\
Y' \ar[r]^g & Y
}
$$
을 생각하자. $\mathcal{O}_{Y'} \to f'_*\mathcal{O}_{X'}$가 모든
$g : Y' \to Y$에 대하여 동형이라고 가정하자. 그러면
유한 표시인 몰입 $j : Z \to Y$가 존재하여, 사상 $g : Y' \to Y$가 $Z$를
통해 인수분해될 필요충분조건은
유한 국소 자유 $\mathcal{O}_{Y'}$-가군 $\mathcal{N}$이 존재하여
% P09-ERR-0769
$(f')^*\mathcal{N} \cong (g')^*\mathcal{L}$을 만족하는 것이다.
\end{lemma}

\begin{proof}
$y : \Spec(k) \to Y$를 체 값 점이라 하자. 그러면 $X_y$, 즉 $f$의
$y$에서의 올은
$H^0(X_y, \mathcal{O}_{X_y}) = k$라는 가정에 의해 연결이다. 따라서
$\mathcal{E}$의 계수는 올들 위에서 상수이다. $f$가 열리고(Morphisms of
Spaces, Lemma \ref{spaces-morphisms-lemma-fppf-open}) 닫혀 있으므로
$Y = \coprod Y_r$인 분해가 존재한다. 이는 $Y$를 열린 동시에 닫힌
부분공간들로 분해하며, $\mathcal{E}$는 상수 계수 $r$을 갖고 이는 $Y_r$의 역상 위에서
갖는다. 따라서 $\mathcal{E}$가 상수 계수 $r$을 갖는다고 가정할 수 있다.
$\mathcal{E}^\vee = \SheafHom(\mathcal{E}, \mathcal{O}_X)$로 나타내는 대상은
쌍대 계수 $r$ 가군이다.

\medskip\noindent
코호몰로지와 밑변환, 더 정확히는 Derived Categories of Spaces, Lemma
\ref{spaces-perfect-lemma-flat-proper-perfect-direct-image-general}에 의해
$E = Rf_*\mathcal{E}$는 $Y$의 유도 범주의 완전 대상이고, 그 형성은 임의의
밑변환과 교환한다. $E' = Rf_*\mathcal{E}^\vee$에 대해서도 마찬가지이다.
차수 $< 0$에서는 코호몰로지가 전혀 없으므로 $E$와 $E'$은 (국소적으로)
Tor 진폭 $[0, b]$를 가지며, 여기서 $b$는 어떤 수이다. 임의의 $g : Y' \to Y$에
대하여
$f'_*((g')^*\mathcal{E}) = H^0(Lg^*E)$ 및
$f'_*((g')^*\mathcal{E}^\vee) = H^0(Lg^*E')$임에 유의하자. $j : Z \to Y$와
$j' : Z' \to Y$를 Derived Categories of Spaces, Lemma
\ref{spaces-perfect-lemma-locally-closed-where-H0-locally-free}에서 $E$와
$E'$에 대하여 $a = 0$ 및 $r = r$로 구성한 국소 닫힌 몰입이라 하자. 이들은
$H^0(Lj^*E)$와 $H^0((j')^*E')$가 역상과 호환되는 계수 $r$의 국소 자유
가군이라는 성질로 특징지어진다.

\medskip\noindent
$g : Y' \to Y$를 사상이라 하자. 보조정리와 같은 $\mathcal{N}$이 존재하면,
사영 공식(Cohomology on Sites, Lemma
\ref{sites-cohomology-lemma-projection-formula})을 사용하여 가군들
$$
f'_*((g')^*\mathcal{E}) \cong
f'_*((f')^*\mathcal{N}) \cong
\mathcal{N} \otimes_{\mathcal{O}_{Y'}} f'_*\mathcal{O}_{X'} \cong
\mathcal{N}\quad\text{and similarly }\quad
f'_*((g')^*\mathcal{E}^\vee) \cong \mathcal{N}^\vee
$$
이 계수 $r$의 국소 자유이고, 더 나아가 계수 $r$의 국소 자유로 남는데 이는
임의의 밑변환 $Y'' \to Y'$ 뒤에도 성립한다. 따라서 이 경우
$g : Y' \to Y$는 $j$와 $j'$을 통해 인수분해된다. 그러므로 $Y$를
$Z \times_Y Z'$으로 바꾸어 $f_*\mathcal{E}$와 $f_*\mathcal{E}^\vee$가
국소 자유 $\mathcal{O}_Y$-가군이고 계수 $r$이며, 그 형성이 임의의 밑변환과
교환한다고 가정할 수 있다.

\medskip\noindent
이 상황에서 $g : Y' \to Y$가 사상이고 보조정리와 같은 $\mathcal{N}$이
존재하면 다음 사상, 즉 차수 $0$의 컵곱은 완전 쌍대이다.
$$
f'_*((g')^*\mathcal{E})
\otimes_{\mathcal{O}_{Y'}}
f'_*((g')^*\mathcal{E}^\vee)
\longrightarrow \mathcal{O}_{Y'}
$$
거꾸로 이 컵곱 사상이 완전 쌍대라면 $Y'$ 위 국소적으로
% P09-ERR-0769
$\sigma_1, \ldots, \sigma_r$가 $f'_*((g')^*\mathcal{L})$의 단면들의 기저이고,
$\tau_1, \ldots, \tau_r$가 $f'_*((g')^*\mathcal{E}^\vee)$의 단면들의
기저이며, $\sigma_i \tau_j = \delta_{ij}$를 만족한다. $\sigma_i$를
$(g')^*\mathcal{L}$의 단면으로 생각하되 이는 $X'$ 위에 있고, $\tau_j$를
$(g')^*\mathcal{L}^\vee$의 단면으로 생각하되 이 또한 $X'$ 위에 있다고 하자.
그러면 다음 사상이 동형임을 얻는다.
$$
\sigma_1, \ldots, \sigma_r :
\mathcal{O}_{X'}^{\oplus r}
\longrightarrow
(g')^*\mathcal{E}
$$
그 역은 다음 사상이다.
$$
\tau_1, \ldots, \tau_r :
(g')^*\mathcal{E}
\longrightarrow
\mathcal{O}_{X'}^{\oplus r}
$$
달리 말하면 $(f')^*f'_*(g')^*\mathcal{E} \cong (g')^*\mathcal{E}$이다.
한편 컵곱이 비퇴화라는 조건은 준콤팩트 열린 부분스킴, 즉 적절한 행렬식이
영이 아닌 자취를 골라낸다. 이것으로 증명이 끝난다.
\end{proof}








\section{복합체의 유한성 성질의 하강}
\label{spaces-more-morphisms-section-descent-finiteness}

\noindent
이 절은 More on Morphisms, Section
\ref{more-morphisms-section-descent-finiteness} 및 Derived Categories of
Schemes, Section \ref{perfect-section-descent-finiteness}에 대응한다.

\begin{lemma}
\label{spaces-more-morphisms-lemma-pseudo-coherent-descends-fpqc}
$S$를 스킴이라 하자. $\{f_i : X_i \to X\}$를 $S$ 위 대수공간의 fpqc
피복이라 하자. $E \in D_\QCoh(\mathcal{O}_X)$라 하고
$m \in \mathbf{Z}$라 하자. 그러면 $E$가 $m$-유사연접일 필요충분조건은
각 $Lf_i^*E$가 $m$-유사연접인 것이다.
\end{lemma}

\begin{proof}
역상은 언제나 $m$-유사연접성을 보존한다. Cohomology on Sites, Lemma
\ref{sites-cohomology-lemma-pseudo-coherent-pullback}을 보라. 따라서
$Lf_i^*E$가 $m$-유사연접이라고 가정하고 $E$가 $m$-유사연접임을 보이면
충분하다. 먼저 각 $X_i$가 스킴이라고 가정할 수 있으며, 이는 모든 $i$에
대하여 성립한다.
Topologies on Spaces, Lemma
\ref{spaces-topologies-lemma-refine-fpqc-schemes}를 보라. 다음으로 전사 에탈
사상 $U \to X$를 택하되 $U$는 스킴이라 하자.
% P09-ERR-0772
그러면 $U_i = U \times_X X_i$는 스킴이고, 스킴의 fpqc 피복
$\{U_i \to U\}$를 얻는다. Topologies on Spaces, Lemma
\ref{spaces-topologies-lemma-recognize-fpqc-covering}을 보라. 스킴의 경우에
의해 $\{U_i \to U\}_{i \in I}$에 대한 결과는 참이다. Derived Categories
of Schemes, Lemma \ref{perfect-lemma-pseudo-coherent-descends-fpqc}를 보라.
한편 제한 $E|_U$는 $D_\QCoh(\mathcal{O}_U)$의 대상에서 온다. 여기서는
$U$의 자리츠키 위상과 ``통상적인'' 구조층을 사용하여 이를 정의한다.
Derived Categories of Spaces, Lemma
\ref{spaces-perfect-lemma-derived-quasi-coherent-small-etale-site}를 보라.
에탈 유사연접성과 자리츠키 유사연접성은 Derived Categories of Spaces,
Lemma \ref{spaces-perfect-lemma-descend-pseudo-coherent}에 의해 일치하므로
보조정리가 따른다.
\end{proof}

\begin{lemma}
\label{spaces-more-morphisms-lemma-tor-amplitude-descends-fppf}
$S$를 스킴이라 하자. $\{g_i : Y_i \to Y\}$를 $S$ 위 대수공간의 fpqc
피복이라 하자. $f : X \to Y$를 대수공간의 사상이라 하고,
$X_i = Y_i \times_Y X$로 놓으며 사영들을 $f_i : X_i \to Y_i$ 및
$g'_i : X_i \to X$로 나타내자. $E \in D_\QCoh(\mathcal{O}_X)$라 하고
$a, b \in \mathbf{Z}$라 하자. 다음은 동치이다.
\begin{enumerate}
\item $E$는 Tor 진폭 $[a, b]$를 가지며, 이는
$D(f^{-1}\mathcal{O}_Y)$의 대상으로서의 진폭이다.
\item $L(g'_i)^*E$는 Tor 진폭 $[a, b]$를 가지며, 이는
$D(f_i^{-1}\mathcal{O}_{Y_i})$의 대상으로서의 진폭이고 모든 $i$에 대하여
성립한다.
% P09-ERR-0773
\end{enumerate}
``Tor 진폭 $[a, b]$를 갖는다''를 ``국소적으로 유한 Tor 차원을 갖는다''로
바꾸어도 참이다.
\end{lemma}

\begin{proof}
Derived Categories of Spaces, Lemma
\ref{spaces-perfect-lemma-tor-independence-and-tor-amplitude}에 의해 역상은
``Tor 진폭 $[a, b]$를 갖는다''라는 성질을 보존한다.
% P09-ERR-0774
$Y_i$와 $X$는 $Y$ 위에서 Tor-독립인데, 이는 $Y_i \to Y$가 평탄하기
때문이다. (2)를 가정하고 (1)을 증명하자. Tor 차원은 줄기에서 계산할 수
있다. Cohomology on Sites, Lemma
\ref{sites-cohomology-lemma-tor-amplitude-stalk} 및 Properties of Spaces,
Theorem \ref{spaces-properties-theorem-exactness-stalks}를 보라.
$\overline{x}$를 $X$의 기하점이라 하자. 어떤 $i$와 기하점
$\overline{x}_i$를 택하되 이는 $X_i$에 있고, 그 상 $\overline{x}$는 $X$에
있다고 하자.
그러면
% P09-ERR-0775
$$
(L(g_i')^*E)_{\overline{x}_i} =
E_{\overline{x}}
\otimes_{\mathcal{O}_{X, \overline{x}}}^\mathbf{L}
\mathcal{O}_{X_, \overline{x}_i}
$$
$\overline{y}_i$를 $Y_i$에, $\overline{y}$를 $Y$에 있는 점이라 하며,
이들은 각각 $\overline{x}_i$와 $\overline{x}$의 상이라 하자. $X$와 $Y_i$는 $Y$ 위에서
Tor-독립이므로 More on Algebra, Lemma
\ref{more-algebra-lemma-base-change-comparison}을 적용하면 표시식의 우변이
$$
E_{\overline{x}}
\otimes_{\mathcal{O}_{Y, \overline{y}}}^\mathbf{L}
\mathcal{O}_{Y_i, \overline{y}_i}
$$
와 $D(\mathcal{O}_{Y_i, \overline{y}_i})$에서 같음을 알 수 있다. 이 대상의
Tor 진폭이 $[a, b]$에 들어간다고 가정했으므로
% P09-ERR-0776
$E_{\overline{x}}$의 $D(\mathcal{O}_{Y, \overline{y}})$에서의 Tor 진폭도
$[a, b]$에 들어간다. More on Algebra, Lemma
\ref{more-algebra-lemma-flat-descent-tor-amplitude}를 보라. 따라서 (1)이
따른다.

\medskip\noindent
기초적인 위상수학을 사용하면 ``국소적으로 유한 Tor 차원을 갖는다''라는
경우도 따른다.
\end{proof}

\noindent
다음 보조정리들은 사실 이 절에 속하지 않는다.

\begin{lemma}
\label{spaces-more-morphisms-lemma-thickening-pseudo-coherent}
$S$를 스킴이라 하자. $i : X \to X'$를 대수공간의 유한 차수 두껍히기라
하자. $K' \in D(\mathcal{O}_{X'})$가 $K = Li^*K'$이 유사연접이라는 성질을
갖는 대상이라 하자. 그러면 $K'$은 유사연접이다.
\end{lemma}

\begin{proof}
먼저 $K'$의 코호몰로지 층들이 준연접임을 증명한다. 독자는 이 부분을
건너뛰어도 좋다. 이를 위해 Section \ref{spaces-more-morphisms-section-thickenings}에 의해 일차
두껍히기의 경우로 환원할 수 있다. $\mathcal{I} \subset \mathcal{O}_{X'}$를
$X$를 잘라내는 준연접 아이디얼 층이라 하자. 짧은 완전열
$$
0 \to \mathcal{I} \to \mathcal{O}_{X'} \to i_*\mathcal{O}_X \to 0
$$
에 $K'$을 텐서하면 구별 삼각형
$$
K' \otimes_{\mathcal{O}_{X'}}^\mathbf{L} \mathcal{I}
\to K' \to
K' \otimes_{\mathcal{O}_{X'}}^\mathbf{L} i_*\mathcal{O}_X
\to
(K' \otimes_{\mathcal{O}_{X'}}^\mathbf{L} \mathcal{I})[1]
$$
을 얻는다. $i_* = Ri_*$이고 일차 두껍히기이므로 $\mathcal{I}$를 준연접
$\mathcal{O}_X$-가군으로 볼 수 있다. 따라서 위 삼각형을
$$
i_*(K \otimes_{\mathcal{O}_X}^\mathbf{L} \mathcal{I})
\to K' \to
i_*K \to
i_*(K \otimes_{\mathcal{O}_X}^\mathbf{L} \mathcal{I})[1]
$$
로 다시 쓸 수 있다. 각 항의 식별에는 Cohomology of Spaces, Lemma
\ref{spaces-cohomology-lemma-projection-formula-finite}를 사용하라. $K$가
$D_\QCoh(\mathcal{O}_X)$에 속하므로 $K'$은
$D_\QCoh(\mathcal{O}_{X'})$에 속한다. 여기서는 Derived Categories of
Spaces, Lemmas \ref{spaces-perfect-lemma-pseudo-coherent},
\ref{spaces-perfect-lemma-quasi-coherence-tensor-product}, 그리고
\ref{spaces-perfect-lemma-quasi-coherence-direct-image}를 사용한다.

\medskip\noindent
$K'$이 $D_\QCoh(\mathcal{O}_{X'})$에 속한다고 가정하자. 문제는 $X'$ 위
에탈 국소적이므로 $X'$이 아핀이라고 가정할 수 있다. 이 경우 결과는 스킴의
경우(More on Morphisms, Lemma
\ref{more-morphisms-lemma-thickening-pseudo-coherent})에서 따른다. 스킴의
언어로 옮기는 데에는 Derived Categories of Spaces, Lemmas
\ref{spaces-perfect-lemma-derived-quasi-coherent-small-etale-site}와
\ref{spaces-perfect-lemma-descend-pseudo-coherent}, 그리고 Remark
\ref{spaces-perfect-remark-match-total-direct-images}를 사용한다.
\end{proof}

\begin{lemma}
\label{spaces-more-morphisms-lemma-thickening-relatively-perfect}
$S$를 스킴이라 하자. 대수공간의 카르테시안 그림
$$
\xymatrix{
X \ar[r]_i \ar[d]_f & X' \ar[d]^{f'} \\
Y \ar[r]^j & Y'
}
$$
을 생각하자. 이 그림은 $S$ 위에 있다. $X' \to Y'$이 평탄하고 국소 유한
표시이며, $Y \to Y'$이
유한 차수 두껍히기라고 가정하자. $E' \in D(\mathcal{O}_{X'})$라 하자.
$E = Li^*(E')$가 $Y$-완전이면 $E'$은 $Y'$-완전이다.
\end{lemma}

\begin{proof}
$Y$-완전하다는 것이 $E$에 대하여 뜻하는 바는 $E$가 유사연접이고
$f^{-1}\mathcal{O}_Y$-가군의 복합체로서 국소적으로 유한 Tor 차원을
갖는다는 뜻임을 상기하자(Definition
\ref{spaces-more-morphisms-definition-relatively-perfect}). Lemma
\ref{spaces-more-morphisms-lemma-thickening-pseudo-coherent}에 의해 $E'$은 유사연접이다. 특히
$E'$은 $D_\QCoh(\mathcal{O}_{X'})$에 속한다. Derived Categories of Spaces,
Lemma \ref{spaces-perfect-lemma-pseudo-coherent}를 보라. Lemma
\ref{spaces-more-morphisms-lemma-affine-locally-rel-perfect}에 의해 스킴의 경우로 환원된다.
스킴의 경우는 More on Morphisms, Lemma
\ref{more-morphisms-lemma-thickening-relatively-perfect}이다.
\end{proof}
\begin{lemma}
\label{spaces-more-morphisms-lemma-henselian-relatively-perfect}
쌍 $(R, I)$가 환과 야코브슨 근기에 포함되는 아이디얼 $I$로
이루어져 있다고 하자. $S = \Spec(R)$ 및 $S_0 = \Spec(R/I)$로 놓는다.
$X$를 $R$ 위의 대수공간이라 하고, 그 구조 사상
$f : X \to S$가 고유하고 평탄하며 유한 표시라고 하자.
% P09-ERR-0777
$X_0 = S_0 \times_S X$로 표기한다. $E \in D(\mathcal{O}_X)$가
유사연접이라고 하자. $E_0$로 표기되는 $E$의 $X_0$으로의 유도 제한이
$S_0$-완전이면 $E$는 $S$-완전이다.
\end{lemma}

\begin{proof}
전사 에탈 사상 $U \to X$를 택하되 $U$는 아핀으로 택한다.
폐몰입 $U \to \mathbf{A}^d_S$를 택한다.
$U_0 = S_0 \times_S U$로 놓는다.
복합체 $E_0|_{U_0}$의 Tor 진폭은 $[a, b]$ 안에 있으며, 여기서
$a, b \in \mathbf{Z}$이다.
기하적 점 $\overline{x}$를 $X$에서 택하자.
$E_{\overline{x}}$의 $R$ 위 Tor 진폭이
$[a - d, b]$ 안에 있음을 보이겠다.
Tor 진폭은 다음 결과들에 의해 줄기에서 판별할 수 있으므로 이것으로 증명이 끝난다:
Cohomology on Sites, Lemma \ref{sites-cohomology-lemma-tor-amplitude-stalk}
및 Properties of Spaces, Theorem
\ref{spaces-properties-theorem-exactness-stalks}.

\medskip\noindent
$x \in |X|$를 $\overline{x}$가 정하는 점이라 하자.
$|X| \to |S|$가 닫힌 사상임을 상기하자(이는 고유 사상의 정의에 따른다).
$I$가 야코브슨 근기에 포함되므로 $S$의 공집합이 아닌 모든 닫힌
부분집합은 닫힌 부분스킴 $S_0$의 한 점을 포함한다.
따라서 특수화 $x \leadsto x_0$를 $|X|$ 안에서 찾을 수 있으며
$x_0 \in |X_0|$이다. $u_0 \in U_0$를 택하여 $x_0$로 가게 한다.
Decent Spaces, Lemma \ref{decent-spaces-lemma-generalizations-lift-flat}
(또는 에탈 사상에 직접 적용되는 Decent Spaces, Lemma
\ref{decent-spaces-lemma-specialization})에 의해
특수화 $u \leadsto u_0$를 $U$ 안에서 찾아 $u$가 $x$로 가게 한다.
$\overline{x}$를 기하적 점 $\overline{u}$로 올릴 수 있으며,
이는 $U$의 점으로서 $u$ 위에 놓인다.
그러면 $E_{\overline{x}} = (E|_U)_{\overline{u}}$이다.

\medskip\noindent
$U = \Spec(A)$로 쓰자. 그러면 $A$는 평탄하고 유한 표시인
$R$-대수이며, 어떤 다항식 $R$-대수의 몫이다. 이 다항식 대수의
변수 수는 $d$이다.
% P09-ERR-0778
제한 $E|_U$는 다음 결과들에 의해
Derived Categories of Spaces, Lemmas
\ref{spaces-perfect-lemma-pseudo-coherent},
\ref{spaces-perfect-lemma-derived-quasi-coherent-small-etale-site}, and
\ref{spaces-perfect-lemma-descend-pseudo-coherent}
and
Derived Categories of Schemes, Lemma
\ref{perfect-lemma-affine-compare-bounded} and
\ref{perfect-lemma-pseudo-coherent-affine}
유사연접 대상 $K$에 대응하며, 이 대상은 $D(A)$에 속한다.
$E_0$가 $K \otimes_A^\mathbf{L} A/IA$에 대응함을 관찰하자.
$\mathfrak q \subset \mathfrak q_0 \subset A$를 $u \leadsto u_0$에
대응하는 소아이디얼들이라 하자. 그러면
$$
E_{\overline{x}} =
(E|_U)_{\overline{u}} =
E_u \otimes_{\mathcal{O}_{U, u}}^\mathbf{L} \mathcal{O}_{U, \overline{u}} =
K_{\mathfrak q} \otimes_{A_\mathfrak q}^\mathbf{L} A_{\mathfrak q}^{sh}
$$
이다(몇 가지 세부사항은 생략한다). $A_\mathfrak q \to A_\mathfrak q^{sh}$가
평탄하므로, 이를 $R$-가군으로 보았을 때의 Tor 진폭은
$K_\mathfrak q$를 $R$-가군으로 보았을 때의 Tor 진폭과 같다
(More on Algebra, Lemma \ref{more-algebra-lemma-no-change-tor-amplitude}).
또한 $K_{\mathfrak q}$는 $K_{\mathfrak q_0}$의 국소화이다.
따라서 $K_{\mathfrak q_0}$의 Tor 진폭이 $[a - d, b]$ 안에 있음을,
이를 $R$-가군의 복합체로 보아 증명하면 충분하다.

\medskip\noindent
$I \subset \mathfrak p_0 \subset R$을 $f(x_0)$에 대응하는
소아이디얼이라 하자. 그러면
\begin{align*}
K \otimes_R^\mathbf{L} \kappa(\mathfrak p_0)
& =
(K \otimes_R^\mathbf{L} R/I) \otimes_{R/I}^\mathbf{L}
\kappa(\mathfrak p_0) \\
& =
(K \otimes_A^\mathbf{L} A/IA) \otimes_{R/I}^\mathbf{L} \kappa(\mathfrak p_0)
\end{align*}
이다. 두 번째 등식은 $R \to A$가 평탄하기 때문이다.
$a, b$의 선택에 의해 이 복합체의 코호몰로지는 구간 $[a, b]$의
차수에서만 존재한다. 이제 마지막으로 More on Algebra, Lemma
\ref{more-algebra-lemma-lift-from-fibre-relatively-perfect}를
$R \to A$, $\mathfrak q_0$, $\mathfrak p_0$, $K$에 적용하여 결론을 얻는다.
% P09-ERR-0779
\end{proof}






\section{마디 곡선족}
\label{spaces-more-morphisms-section-families-nodal}

\noindent
이 절은 Algebraic Curves, Section
\ref{curves-section-families-nodal}의 연속이다.
용어의 선택에 대해서도 그 절을 보라.

\medskip\noindent
스킴 사상의 ``상대차원 $1$에서 기껏해야 마디형''이라는 성질은
소스와 타깃에 대해 에탈 국소적이다. 다음을 보라:
Descent, Lemma \ref{descent-lemma-etale-local-source-target}
및
Algebraic Curves, Lemmas \ref{curves-lemma-nodal-family-etale-local-source},
\ref{curves-lemma-nodal-family-fpqc-local-target}, and
\ref{curves-lemma-nodal-family-postcompose-etale}.
또한 이는 기저변환에 대해 안정적이고 타깃 위 fpqc 국소적이다. 다음을 보라:
Algebraic Curves, Lemmas \ref{curves-lemma-base-change-nodal-family} and
\ref{curves-lemma-nodal-family-fpqc-local-target}. 따라서
Morphisms of Spaces, Lemma \ref{spaces-morphisms-lemma-local-source-target}에 의해
대수공간에 대해서도 다음과 같이 상대차원 $1$에서 기껏해야 마디형인
사상의 개념을 정의할 수 있다. 사상이 표현가능한 경우 이 개념은
Morphisms of Spaces, Section \ref{spaces-morphisms-section-representable}에서
이미 정의한 개념과 일치한다.

\begin{definition}
\label{spaces-more-morphisms-definition-nodal-family}
$S$를 스킴이라 하자.
$f : X \to Y$를 $S$ 위 대수공간들의 사상이라 하자.
$f$가 {\it 상대차원 $1$에서 기껏해야 마디형}이라는 것은
Morphisms of Spaces, Lemma \ref{spaces-morphisms-lemma-local-source-target}의
동치 조건들이 $\mathcal{P} =$``상대차원 $1$에서 기껏해야 마디형''에 대해
성립한다는 뜻이다.
\end{definition}

\begin{lemma}
\label{spaces-more-morphisms-lemma-base-change-nodal}
상대차원 $1$에서 기껏해야 마디형이라는 성질은 기저변환에 의해 보존된다.
\end{lemma}

\begin{proof}
다음을 보라:
Morphisms of Spaces, Remark \ref{spaces-morphisms-remark-base-change-P}
및
Algebraic Curves, Lemma \ref{curves-lemma-base-change-nodal-family}.
\end{proof}

\begin{lemma}
\label{spaces-more-morphisms-lemma-nodal-local}
$S$를 스킴이라 하자.
$f : X \to Y$를 $S$ 위 대수공간들의 사상이라 하자.
다음 조건들은 동치이다.
\begin{enumerate}
\item $f$는 상대차원 $1$에서 기껏해야 마디형이다.
\item 모든 스킴 $Z$와 임의의 사상 $Z \to Y$에 대하여 사상
$Z \times_Y X \to Z$는 상대차원 $1$에서 기껏해야 마디형이다.
\item 모든 아핀 스킴 $Z$와 임의의 사상
$Z \to Y$에 대하여 사상 $Z \times_Y X \to Z$는
상대차원 $1$에서 기껏해야 마디형이다.
\item 스킴 $V$와 전사 에탈 사상 $V \to Y$가 존재하여
$V \times_Y X \to V$가 상대차원 $1$에서 기껏해야 마디형이다.
\item 스킴 $U$와 전사 에탈 사상
$\varphi : U \to X$가 존재하여 합성 $f \circ \varphi$가
상대차원 $1$에서 기껏해야 마디형이다.
\item 모든 가환 그림
$$
\xymatrix{
U \ar[d] \ar[r] & V \ar[d] \\
X \ar[r] & Y
}
$$
에 대하여, $U$, $V$가 스킴이고 수직 화살표들이 에탈이면
위쪽 수평 화살표는 상대차원 $1$에서 기껏해야 마디형이다.
\item 가환 그림
$$
\xymatrix{
U \ar[d] \ar[r] & V \ar[d] \\
X \ar[r] & Y
}
$$
이 존재하여 $U$, $V$가 스킴이고 수직 화살표들이 에탈이며
$U \to X$가 전사이고 위쪽 수평 화살표가
상대차원 $1$에서 기껏해야 마디형이다.
\item 자리츠키 덮개들 $Y = \bigcup_{i \in I} Y_i$ 및
$f^{-1}(Y_i) = \bigcup X_{ij}$가 존재하여 각 사상
$X_{ij} \to Y_i$가 상대차원 $1$에서 기껏해야 마디형이다.
\end{enumerate}
\end{lemma}

\begin{proof}
생략한다.
\end{proof}

\noindent
다음 보조정리는 사상이 상대차원 $1$에서 기껏해야 마디형인지 여부를
올에서 판정할 수 있음을 말해 준다.

\begin{lemma}
\label{spaces-more-morphisms-lemma-locus-where-nodal}
$S$를 스킴이라 하자.
$f : X \to Y$를 $S$ 위 대수공간들의 사상이라 하고, 이 사상은 평탄하며
국소 유한 표시라고 하자. 그러면 최대 열린 부분공간 $X' \subset X$가
존재하여 $f|_{X'} : X' \to Y$가 상대차원 $1$에서 기껏해야 마디형이다.
더 나아가 $X'$의 형성은 임의의 기저변환과 가환한다.
\end{lemma}

\begin{proof}
가환 그림
$$
\xymatrix{
U \ar[d] \ar[r]_h & V \ar[d] \\
X \ar[r]^f & Y
}
$$
을 택하되, $U$, $V$는 스킴이고 수직 화살표들은 에탈이며
$U \to X$는 전사라고 하자. 스킴의 경우에 대한 보조정리
(Algebraic Curves, Lemma \ref{curves-lemma-locus-where-nodal})에 의해
최대 열린 부분스킴 $U' \subset U$를 얻되,
$h|_{U'} : U' \to V$가 상대차원 $1$에서 기껏해야 마디형이고
$U'$의 형성은 기저변환과 가환한다.
$X' \subset X$를 그 점들이 열린 부분집합
$\Im(|U'| \to |X|)$에 대응하는 열린 부분공간이라 하자.
Lemma \ref{spaces-more-morphisms-lemma-nodal-local}에 의해 $X' \to Y$가
상대차원 $1$에서 기껏해야 마디형이며 $X'$는 이 성질을 갖는
가장 큰 열린 부분공간임을 알 수 있다(이는 또한 $U'$가 $X'$의
$U$ 안에서의 역상임을 뜻하지만 여기서는 필요하지 않다).
기저변환 뒤에도 같은 사실이 성립하므로 증명이 끝난다.
\end{proof}




\section{해소 성질}
\label{spaces-more-morphisms-section-resolution-property}

\noindent
Derived Categories of Spaces, Section
\ref{spaces-perfect-section-resolution-property}의 논의를 이어 간다.

\begin{situation}
\label{spaces-more-morphisms-situation-descend-resolution-property}
$S$를 스킴이라 하자. $X$를 $S$ 위의 준콤팩트 준분리 대수공간이라 하자.
전사 에탈 사상 $V \to X$를 택하되 $V$는 아핀 스킴으로 택한다(이러한 사상은
Properties of Spaces, Lemma
\ref{spaces-properties-lemma-quasi-compact-affine-cover}에 의해 존재한다).
가환 그림
$$
\xymatrix{
V \ar[rd]_\varphi \ar[rr]_j & & Y \ar[ld]^\pi \\
& X
}
$$
을 택하되 $j$는 열린 몰입이고 $\pi$는 대수공간들의 유한 사상이라고 하자
(이러한 그림은 Lemma
\ref{spaces-more-morphisms-lemma-quasi-finite-separated-pass-through-finite}에 의해 존재한다).
$\mathcal{I} \subset \mathcal{O}_Y$를 $Y$ 위의 유한형 준연접 아이디얼 층으로
택하여 $V(\mathcal{I}) = Y \setminus j(V)$가 되게 한다(이러한 아이디얼 층은
Limits of Spaces, Lemma
\ref{spaces-limits-lemma-quasi-coherent-finite-type-ideals}에 의해 존재한다).
\end{situation}

\begin{lemma}
\label{spaces-more-morphisms-lemma-surjection-in-Noetherian-case}
Situation \ref{spaces-more-morphisms-situation-descend-resolution-property}에서 $X$가 뇌터라고 가정하자.
그러면 임의의 연접 $\mathcal{O}_X$-가군 $\mathcal{F}$에 대하여
$r \geq 0$, 정수 $n_1, \ldots, n_r \geq 0$, 그리고 전사
$$
\bigoplus\nolimits_{i = 1, \ldots, r} \pi_*(\mathcal{I}^{n_i})
\longrightarrow \mathcal{F}
$$
가 존재하며, 이는 $\mathcal{O}_X$-가군의 사상이다.
\end{lemma}

\begin{proof}
$\omega_{Y/X}$를 연접 $\mathcal{O}_Y$-가군으로 표기한다. 이 가군에는 동형
$$
\pi_*\omega_{Y/X} \cong
\SheafHom_{\mathcal{O}_X}(\pi_*\mathcal{O}_Y, \mathcal{O}_X)
$$
이 있으며, 이는 $\pi_*\mathcal{O}_Y$-가군의 동형이다.
다음을 보라: Morphisms of Spaces, Lemma
\ref{spaces-morphisms-lemma-affine-equivalence-modules} 및
Descent on Spaces, Lemma \ref{spaces-descent-lemma-finite-over-finite-module}.
표준 사상 $\mathcal{O}_X \to \pi_*\mathcal{O}_Y$는 표준 사상
$$
\text{Tr}_\pi : \pi_*\omega_{Y/X} \longrightarrow \mathcal{O}_X
$$
를 유도한다. $V$가 뇌터 아핀이므로 단면들
$$
s_1, \ldots, s_r \in
\Gamma(V, \pi^*\mathcal{F} \otimes_{\mathcal{O}_Y} \omega_{Y/X})
$$
을 택하여 연접 가군
$\pi^*\mathcal{F} \otimes_{\mathcal{O}_X} \omega_{Y/X}$를 $V$ 위에서
생성하게 할 수 있다.
Cohomology of Spaces, Lemma \ref{spaces-cohomology-lemma-homs-over-open}에 의해
정수 $n_i \geq 0$를 택하여 $s_i$가 사상
$s_i' : \mathcal{I}^{n_i} \to
\pi^*\mathcal{F} \otimes_{\mathcal{O}_Y} \omega_{Y/X}$로 연장되게 할 수 있다.
$X$로 전방상하여 사상들을 얻는다:
$$
\sigma_i :
\pi_*\mathcal{I}^{n_i}
\xrightarrow{\pi_*s'_i}
\pi_*(\pi^*\mathcal{F} \otimes_{\mathcal{O}_Y} \omega_{Y/X}) =
\mathcal{F} \otimes_{\mathcal{O}_X} \pi_*\omega_{Y/X}
\xrightarrow{\text{Tr}_\pi} \mathcal{F}
$$
여기서 등호는 Cohomology of Spaces, Lemma
\ref{spaces-cohomology-lemma-finite-rings}에 따른다.
증명을 끝내기 위해 이 사상들의 합이 전사임을 보이겠다.

\medskip\noindent
$x \in |X|$를 $X$의 점이라 하자. $v \in |V|$를 $x$로 가는 점이라 하자.
에탈 근방 $(U, u) \to (X, x)$를 택하여
$$
U \times_X Y = W \coprod W'
$$
가 되게 할 수 있다(대수공간들의 분리합). 여기서 $W \to U$는 동형이고,
유일한 점 $w \in W$는 $u$ 위에 놓이며 $v$라는 $V \subset Y$의 점으로 간다.
이를 확인하려면 Lemma
\ref{spaces-more-morphisms-lemma-etale-splits-off-just-one-quasi-finite-part} 및
\'Etale Morphisms, Lemma \ref{etale-lemma-etale-etale-local}을 사용한다.
필요하다면 $U$를 더 줄여서 사영에 의해 $W$가 $V \subset Y$ 안으로
가게 할 수 있다. $\omega_{Y/X}$의 형성은 에탈 국소화와 가환하므로
$$
\pi_*\omega_{Y/X}|_U = (\pi|_W)_*\omega_{W/U} \oplus (\pi|_{W'})_*\omega_{W'/U}
$$
이다. $(\pi|_W)_*\omega_{W/U} = \mathcal{O}_U$이고, 이 동형은 위 분해의
첫 번째 직합인자에 제한한 대각합 사상 $\text{Tr}_\pi|_U$로 주어진다.
$W$가 $V$ 안으로 가므로 $\mathcal{I}^{n_i}|_W = \mathcal{O}_W$이다.
따라서
$$
\pi_*(\mathcal{I}^{n_i})|_U = \mathcal{O}_U \oplus
(W' \to U)_*(\mathcal{I}^{n_i}|_{W'})
$$
이다. 그림을 추적하면(세부사항은 생략한다) $\sigma_i|_U$가 직합인자
$\mathcal{O}_U$ 위에서 사상 $\mathcal{O}_U \to \mathcal{F}$이며, 이는 단면
$$
s_i|_W  \in
\Gamma(W,\pi^*\mathcal{F} \otimes_{\mathcal{O}_Y} \omega_{Y/X})
=
\Gamma(W, \mathcal{F}|_W \otimes_{\mathcal{O}_W} \omega_{W/U})
=
\Gamma(U, \mathcal{F})
$$
에 대응함을 알 수 있다.
단면들 $s_i$가 가군
$\pi^*\mathcal{F} \otimes_{\mathcal{O}_Y} \omega_{Y/X}$를 $V$ 위에서 생성하고
$W$가 $V$ 안으로 가므로 $\bigoplus \sigma_i$의 $U$로의 제한은 전사이다.
$x$가 임의의 점이었으므로 증명이 끝난다.
\end{proof}

\begin{lemma}
\label{spaces-more-morphisms-lemma-resolution-property-in-Noetherian-case}
Situation \ref{spaces-more-morphisms-situation-descend-resolution-property}에서 $X$가 뇌터라고 가정하자.
그러면 $X$가 해소 성질을 갖는 것과 $\pi_*\mathcal{I}$가 유한 국소 자유
$\mathcal{O}_X$-가군의 몫인 것은 동치이다.
\end{lemma}

\begin{proof}
가군 $\pi_*\mathcal{I}$는 Cohomology of Spaces, Lemma
\ref{spaces-cohomology-lemma-finite-pushforward-coherent}에 의해 연접이다.
따라서 $X$가 해소 성질을 가지면 $\pi_*\mathcal{I}$는 유한 국소 자유
$\mathcal{O}_X$-가군의 몫이다. 역으로 전사
$\mathcal{E} \to \pi_*\mathcal{I}$가 주어졌다고 하자. 여기서
$\mathcal{O}_X$-가군 $\mathcal{E}$는 유한 국소 자유이다.
모든 $n \geq 1$에 대하여 전사
$$
\pi_*\mathcal{I} \otimes_{\mathcal{O}_X} \pi_*\mathcal{I}^n
\longrightarrow
\pi_*\mathcal{I}^{n + 1}
$$
가 있음을 관찰하자. 따라서 $\mathcal{E}^{\otimes n}$은
$\pi_*\mathcal{I}^n$으로 전사하며, 이는 모든 $n \geq 1$에 대해 성립한다.
이를 Lemma \ref{spaces-more-morphisms-lemma-surjection-in-Noetherian-case}의 결과와 합치면
$X$가 해소 성질을 가짐을 얻는다.
\end{proof}

\begin{lemma}
\label{spaces-more-morphisms-lemma-resolution-property-one-sheaf}
Situation \ref{spaces-more-morphisms-situation-descend-resolution-property}에서
대수공간 $X$가 해소 성질을 갖는 것과 $\pi_*\mathcal{I}$가 유한 국소 자유
$\mathcal{O}_X$-가군의 몫인 것은 동치이다.
\end{lemma}

\begin{proof}
전방상 $\pi_*\mathcal{G}$를 생각하자. 유한형 준연접
$\mathcal{O}_Y$-가군 $\mathcal{G}$에서 온 이 전방상은 Descent on Spaces, Lemma
\ref{spaces-descent-lemma-finite-over-finite-module}에 의해 유한형 준연접
$\mathcal{O}_X$-가군이다. 특히 $X$가 해소 성질을 가지면
$\pi_*\mathcal{I}$는 Derived Categories of Spaces, Definition
\ref{spaces-perfect-definition-resolution-property}에 의해 유한 국소 자유
$\mathcal{O}_X$-가군의 몫이다.

\medskip\noindent
전사 $\mathcal{E} \to \pi_*\mathcal{I}$가 있다고 가정하자. 여기서
$\mathcal{O}_X$-가군 $\mathcal{E}$는 유한 국소 자유이다.
증명의 나머지에서는 Lemma
\ref{spaces-more-morphisms-lemma-resolution-property-in-Noetherian-case}에서 다룬 뇌터 경우로
환원하여 $X$가 해소 성질을 가짐을 보인다.
독자는 증명의 나머지를 건너뛰어도 좋다.

\medskip\noindent
첫 번째 환원으로 $X$를 $\Spec(\mathbf{Z})$ 위의 대수공간으로 볼 수 있다.
Spaces, Definition \ref{spaces-definition-base-change}을 보라.
(이는 보조정리의 조건이나 결론에 영향을 주지 않는다.)

\medskip\noindent
Limits of Spaces, Lemma
\ref{spaces-limits-lemma-finite-in-finite-and-finite-presentation}에 의해
$Y = \lim Y_i$로 쓸 수 있다. 여기서 $Y_i$는 $X$ 위 유한이고 유한 표시이며
전이 사상들은 닫힌 몰입이다. 닫힌 부분공간
$Z = V(\mathcal{I})$를 $Y$ 안에서 생각하자. $\mathcal{I}$가 유한형이므로 사상
$Z \to Y$는 유한 표시이다. 따라서 어떤 $i$와 유한 표시 사상
$Z_i \to Y_i$를 찾아, 이를 $Y$로 기저변환하면 $Z \to Y$가 되게 할 수 있다.
Limits of Spaces, Lemma
\ref{spaces-limits-lemma-descend-finite-presentation}을 보라.
$i' \geq i$에 대하여 $Z_{i'} = Z_i \times_{Y_i} Y_{i'}$로 표기한다.
$i$를 증가시켜 $Z_i \to Y_i$가 (유한 표시) 닫힌 몰입이라고 가정할 수 있다.
Limits of Spaces, Lemma
\ref{spaces-limits-lemma-descend-closed-immersion}을 보라.
$\mathcal{I}_i \subset \mathcal{O}_{Y_i}$를 $Z_i$의 아이디얼 층으로 표기하고
$\pi_i : Y_i \to X$를 구조 사상으로 표기한다. $i' \geq i$에 대해서도
마찬가지로 표기한다. $Z = \lim_{i' \geq i} Z_{i'}$이므로
$$
\pi_*\mathcal{I} = \colim \pi_{i', *}\mathcal{I}_{i'}
$$
이다. 이 계의 전이 사상들은 모두 전사이다. 이는 사상들
$\pi_{i, *}\mathcal{O}_{Y_i} \to \pi_{i', *}\mathcal{O}_{Y_{i'}}$의 전사성과
$Z_{i'} = Z_i \times_{Y_i} Y_{i'}$라는 사실에서 따른다.
Cohomology of Spaces, Lemma
\ref{spaces-cohomology-lemma-finite-presentation-quasi-compact-colimit}에 의해
어떤 $i' \geq i$에 대하여 사상 $\mathcal{E} \to \pi_*\mathcal{I}$를
사상 $\mathcal{E} \to \pi_{i', *}\mathcal{I}_{i'}$로 올릴 수 있다.
$i'$를 증가시키면 이 사상
$\mathcal{E} \to \pi_{i', *}\mathcal{I}_{i'}$가 전사가 된다고 가정할 수 있다
(그렇지 않으면 전사 전이 사상들을 갖는 여핵들의 여극한이 영이 아니기 때문이다).
이로써 다음 문단에서 다루는 경우로 환원된다.

\medskip\noindent
$X$가 $\mathbf{Z}$ 위의 대수공간이고 $Y \to X$가 유한 표시라고 가정하자.
절대 뇌터 근사에 의해 $X = \lim X_i$를 유향극한으로 쓸 수 있다.
여기서 각 $X_i$는 $\mathbf{Z}$ 위 유한형 준분리 대수공간이고 전이 사상들은
아핀이다. Limits of Spaces, Proposition
\ref{spaces-limits-proposition-approximate}을 보라.
$\pi : Y \to X$가 유한 표시이므로 어떤 $i$와 유한 표시 사상
$\pi_i : Y_i \to X_i$를 찾아, 이를 $X$로 기저변환하면 $\pi$가 되게 할 수 있다.
Limits of Spaces, Lemma
\ref{spaces-limits-lemma-descend-finite-presentation}을 보라.
$i$를 증가시켜 $\pi_i$가 유한이라고 가정할 수 있다. Limits of Spaces, Lemma
\ref{spaces-limits-lemma-descend-finite}을 보라.
다음으로 유한 국소 자유 $\mathcal{O}_{X_i}$-가군 $\mathcal{E}_i$가 존재하여
이를 $X$로 당기면 $\mathcal{E}$가 된다고 가정할 수 있다. Limits of Spaces,
Lemma \ref{spaces-limits-lemma-descend-invertible-modules}을 보라.
또한 사상 $\mathcal{E}_i \to \pi_{i, *}\mathcal{O}_{Y_i}$가 존재하여
이를 $X$로 당기면 합성
$\mathcal{E} \to \pi_*\mathcal{I} \to \pi_*\mathcal{O}_Y$가 된다고 가정할 수 있다.
Limits of Spaces, Lemma
\ref{spaces-limits-lemma-descend-modules-finite-presentation}을 보라.
여핵
$$
\mathcal{E}_i \to \pi_{i, *}\mathcal{O}_{Y_i} \to \mathcal{Q}_i \to 0
$$
은 연접 $\mathcal{O}_{Y_i}$-가군이며, 이를 $X$로 당기면
(유한 표시) 여핵 $\mathcal{Q}$가 된다. 이 여핵은 사상
$\mathcal{E} \to \pi_*\mathcal{O}_Y$의 여핵이다.
다시 말해 $\mathcal{Q} = \pi_*(\mathcal{O}_Y/\mathcal{I})$이다.
사상
$$
\mathcal{E}_i
\otimes_{\mathcal{O}_{X_i}}
\pi_{i, *} \mathcal{O}_{Y_i}
\longrightarrow
\pi_{i, *}\mathcal{O}_{Y_i}
\otimes_{\mathcal{O}_{X_i}}
\pi_{i, *} \mathcal{O}_{Y_i}
\to
\pi_{i, *} \mathcal{O}_{Y_i}
\to
\mathcal{Q}_i
$$
을 생각하자. 여기서 두 번째 화살표는
$\pi_{i, *}\mathcal{O}_{Y_i}$의 대수 구조로 주어진다.
이 사상을 $Y$로 당기면 영이다. 그 이유는 사상
$\mathcal{E} \to \pi_*\mathcal{O}_Y$의 상이 아이디얼
$\pi_*\mathcal{I}$이기 때문이다.
따라서 Limits of Spaces, Lemma
\ref{spaces-limits-lemma-descend-modules-finite-presentation}에 의해
$i$를 증가시켜 위에 표시한 합성이 영이라고 가정할 수 있다.
이는 정확히 사상 $\mathcal{E}_i \to \pi_{i, *}\mathcal{O}_{Y_i}$의 상이
$\pi_{i, *}\mathcal{I}_i$ 꼴이라는 뜻이다. 여기서
$\mathcal{I}_i \subset \mathcal{O}_{Y_i}$는 어떤 연접 아이디얼 층이다.
$\mathcal{E}_i \to \pi_{i, *}\mathcal{O}_{Y_i}$를 당기면
$\mathcal{E} \to \pi_*\mathcal{O}_Y$가 되므로, $\mathcal{I}_i$를 $Y$로
당긴 것은 $\mathcal{I}$를 생성한다.
$V_i \subset Y_i$를 그 여집합이 $V(\mathcal{I}_i) \subset Y_i$인
열린 부분공간으로 표기한다. 위의 설명에 의해 $V$는 $V_i$의 역상이다.
$i$를 증가시켜 $V_i$가 아핀이고
$\pi_i|_{V_i} : V_i \to X_i$가 에탈이라고 가정할 수 있다. 다음을 보라:
Limits of Spaces, Lemmas
\ref{spaces-limits-lemma-limit-is-affine} and
\ref{spaces-limits-lemma-descend-etale}.
이제 Lemma \ref{spaces-more-morphisms-lemma-resolution-property-in-Noetherian-case}를 적용하여
$X_i$가 해소 성질을 가짐을 얻는다. $X \to X_i$가 아핀이므로
Derived Categories of Spaces, Lemma
\ref{spaces-perfect-lemma-resolution-property-goes-up-affine}에 의해
$X$도 해소 성질을 갖는다.
\end{proof}

\begin{lemma}
\label{spaces-more-morphisms-lemma-resolution-property-descends}
$S$를 스킴이라 하자.
$X = \lim X_i$를 $S$ 위 준콤팩트 준분리 대수공간들로 이루어지고
아핀 전이 사상들을 갖는 직계의 극한이라 하자.
그러면 $X$가 해소 성질을 갖는 것과 $X_i$가 어떤 $i$에 대해
해소 성질을 갖는 것은 동치이다.
\end{lemma}

\begin{proof}
$X_i$가 해소 성질을 가지면 Derived Categories of Spaces, Lemma
\ref{spaces-perfect-lemma-resolution-property-goes-up-affine}에 의해
$X$도 해소 성질을 갖는다.
$X$가 해소 성질을 갖는다고 가정하자.
$i \in I$를 택한다. 아핀 스킴 $V_i$와 전사 에탈 사상
$V_i \to X_i$를 택할 수 있다
(Properties of Spaces, Lemma
\ref{spaces-properties-lemma-quasi-compact-affine-cover}).
몰입 $j : V_i \to Y_i$를 택하되 $Y_i$는 $X_i$ 위 유한이고 유한 표시가
되게 할 수 있다
(Lemma \ref{spaces-more-morphisms-lemma-quasi-finite-separated-pass-through-finite-addendum}).
유한형 준연접 아이디얼 $\mathcal{I}_i \subset \mathcal{O}_{Y_i}$를 택하여
$V_i = Y_i \setminus V(\mathcal{I}_i)$가 되게 할 수 있다
(Limits of Spaces, Lemma
\ref{spaces-limits-lemma-quasi-coherent-finite-type-ideals}).
$V \to Y \to X$를 $V_i \to Y_i \to X_i$의 $X$로의 기저변환으로 표기한다.
$\mathcal{I} \subset \mathcal{O}_Y$를 아이디얼 $\mathcal{I}_i$의 당김으로
표기한다. Lemma \ref{spaces-more-morphisms-lemma-resolution-property-one-sheaf}의 쉬운 방향에 의해
유한 국소 자유 $\mathcal{O}_X$-가군 $\mathcal{E}$와 전사
$\mathcal{E} \to \pi_*\mathcal{I}$가 존재한다.
$\pi_i : Y_i \to X_i$가 유한이고 유한 표시이므로
$\pi : Y \to X$도 유한이고 유한 표시이다. 또한 $\mathcal{O}_{X_i}$-가군들
$\pi_{i, *}\mathcal{O}_{Y_i}$와
$\pi_{i, *}(\mathcal{O}_{Y_i}/\mathcal{I}_i)$는 유한 표시이고, 이를 $X$로
당기면 각각 $\pi_*\mathcal{O}_Y$와
$\pi_*(\mathcal{O}_Y/\mathcal{I})$를 얻는다. 따라서 Limits of Spaces, Lemma
\ref{spaces-limits-lemma-descend-modules-finite-presentation}에 의해
$i$를 증가시켜 유한 국소 자유 $\mathcal{O}_{X_i}$-가군
$\mathcal{E}_i$와 사상 $\mathcal{E}_i \to \pi_{i, *}\mathcal{O}_{Y_i}$를
찾을 수 있으며, 이를 $X$로 기저변환하면 합성
$\mathcal{E} \to \pi_*\mathcal{I} \to \pi_*\mathcal{O}_Y$를 회복한다.
유한 표시 $\mathcal{O}_{X_i}$-가군들
$\Coker(\mathcal{E}_i \to \pi_{i, *}\mathcal{O}_{Y_i})$와
$\pi_{i, *}(\mathcal{O}_{Y_i}/\mathcal{I}_i)$를 $X$로 당기면
$\pi_*\mathcal{O}_Y$의 몫으로서 서로 일치한다. 따라서 Limits of Spaces,
Lemma \ref{spaces-limits-lemma-descend-modules-finite-presentation}에 의해
이들이 일치한다고 가정할 수 있다. 다시 말해, 사상
% P09-ERR-KO-C76-13764
$\mathcal{E}_i \to \pi_{i, *}\mathcal{O}_{X_i}$의 상은
$\pi_{i, *}\mathcal{I}_i$와 같다.
그러면 Lemma \ref{spaces-more-morphisms-lemma-resolution-property-one-sheaf}에 의해
$X_i$가 해소 성질을 가짐을 얻는다.
\end{proof}

\begin{lemma}
\label{spaces-more-morphisms-lemma-resolution-property-affine-diagonal}
$S$를 스킴이라 하자.
$X$를 해소 성질을 갖는 준콤팩트 준분리 대수공간이라 하자.
그러면 $X$는 $\mathbf{Z}$ 위 아핀 대각사상을 갖는다
(Properties of Spaces, Definition
\ref{spaces-properties-definition-separated}에서처럼).
\end{lemma}

\begin{proof}
스킴의 경우와 같이 증명할 수도 있지만, 여기서는 스킴의 경우로부터
보조정리를 도출한다. 먼저 $S = \Spec(\mathbf{Z})$라고 가정할 수 있고
그렇게 가정한다. 다음으로 스킴 $Y$와 전사 정수적 사상
$f : Y \to X$를 택한다. Decent Spaces, Lemma
\ref{decent-spaces-lemma-there-is-a-scheme-integral-over}를 보라.
그러면 $f$는 아핀이므로 Derived Categories of Spaces, Lemma
\ref{spaces-perfect-lemma-resolution-property-goes-up-affine}에 의해
$Y$는 해소 성질을 갖는다. 따라서 스킴의 경우에 의해 스킴 $Y$는
아핀 대각사상을 갖는다. Derived Categories of Schemes, Lemma
\ref{perfect-lemma-resolution-property-affine-diagonal}를 보라.
이제 가환 그림
$$
\xymatrix{
Y \ar[d] \ar[rr]_{\Delta_Y} & & Y \times_{\mathbf{Z}} Y \ar[d] \\
X \ar[rr]^{\Delta_X} & & X \times_{\mathbf{Z}} X
}
$$
을 생각하자. 오른쪽 수직 화살표는 정수적이며, 특히 아핀이다.
$W \to X \times_{\mathbf{Z}} X$를 사상이라 하고 $W$는 아핀이라고 하자.
그러면
$$
Y \times_{X \times_{\mathbf{Z}} X} W =
Y \times_{\Delta_Y, Y \times_{\mathbf{Z}} Y}
(Y \times_{\mathbf{Z}} Y) \times_{X \times_{\mathbf{Z}} X} W
$$
가 아핀임을 알 수 있다. 한편 $Y \to X$는 정수적이고 전사이므로
$$
Y \times_{X \times_{\mathbf{Z}} X} W
\longrightarrow
X \times_{X \times_{\mathbf{Z}} X} W
$$
는 $Y \to X$를 $W$로 기저변환한 것으로서 정수적 전사이다.
Limits of Spaces, Proposition \ref{spaces-limits-proposition-affine}에 의해
이 화살표의 타깃은 아핀이다. 따라서 원하는 대로 $\Delta_X$는 아핀이다.
\end{proof}
\section{블로업과 해소 성질}
\label{spaces-more-morphisms-section-resolution-property-by-blowing-up}

\noindent
블로업 뒤에 해소 성질이 성립함을 증명한다.

\begin{lemma}
\label{spaces-more-morphisms-lemma-blow-up-resolution-property}
$S$를 스킴이라 하자. $X$를 $S$ 위의 준콤팩트 준분리 대수공간이라 하자.
$|X|$가 유한 개의 기약 성분을 갖는다고 가정하자.
조밀한 준콤팩트 열린집합 $U \subset X$와 $U$-허용 블로업 $X' \to X$가
존재하여 대수공간 $X'$가 해소 성질을 갖는다.
\end{lemma}

\begin{proof}
Limits of Spaces, Lemma
\ref{spaces-limits-lemma-there-is-a-scheme-finite-over-refined}에 의해
전사이고 유한이며 유한 표시인 사상 $f : Y \to X$가 존재하되 $Y$는
스킴이다. 또한 준콤팩트 조밀 열린집합 $U \subset X$가 존재하여
$f^{-1}(U) \to U$가 유한 에탈이다.
More on Morphisms, Lemma \ref{more-morphisms-lemma-blow-up-ample-family}에 의해
준콤팩트 조밀 열린집합 $V \subset Y$와 $V$-허용 블로업 $Y' \to Y$가
존재하여 $Y'$은 가역 가군들의 앰플 족을 갖는다.
$U$를 줄여서 $f^{-1}(U) \subset V$라고 가정할 수 있다(세부사항은 생략한다).
따라서 $f' : Y' \to X$는 $U$ 위에서 유한 에탈이고, 특히 사상
$(f')^{-1}(U) \to U$는 유한 국소 자유이다.
Lemma \ref{spaces-more-morphisms-lemma-finite-after-blowing-up}에 의해 $U$-허용 블로업
$X' \to X$가 존재한다. $Y''$을 $Y'$의 진변환이라 하면 이는 $X'$ 위
유한 국소 자유이다.
그림으로 쓰면
$$
\xymatrix{
Y'' \ar[d] \ar[r]_g &
Y' \ar[r] &
Y \ar[d] \\
X' \ar[rr] & & X
}
$$
이다. $g : Y'' \to Y'$은 $X' \to X$의 중심의 역상에서의 블로업이므로
(Divisors on Spaces, Lemma \ref{spaces-divisors-lemma-strict-transform}),
$g : Y'' \to Y'$은 사영적이고 어떤 $g$-앰플 가역 가군이 $Y''$ 위에
존재한다. 따라서 More on Morphisms, Lemma
\ref{more-morphisms-lemma-ample-family-ample-relative}에 의해
$Y''$은 가역 가군들의 앰플 족을 갖는다.
따라서 $Y''$은 해소 성질을 갖는다. Derived Categories of Schemes, Lemma
\ref{perfect-lemma-resolution-property-ample-family}를 보라.
Derived Categories of Spaces, Lemma
\ref{spaces-perfect-lemma-resolution-property-goes-down-finite-flat}에 의해
$X'$도 해소 성질을 갖는다.
\end{proof}

\begin{lemma}
\label{spaces-more-morphisms-lemma-blow-up-resolution-property-filtered}
$S$를 스킴이라 하자. $X$를 $S$ 위의 준콤팩트 준분리 대수공간이라 하자.
$t \geq 0$와 닫힌 부분공간들
$$
X \supset Z_0 \supset Z_1 \supset \ldots \supset Z_t = \emptyset
$$
이 존재하여 $Z_i \to X$는 유한 표시이고, $Z_0 \subset X$는 두껍히기이며,
각 $i = 0, \ldots t - 1$에 대하여
% P09-ERR-KO-C76-13872
$(Z_i \setminus Z_{i - 1})$-허용 블로업 $Z'_i \to Z_i$가 존재하여
$Z'_i$는 해소 성질을 갖는다.
\end{lemma}

\begin{proof}
이 문단에서는 절대 뇌터 근사를 사용하여
$\Spec(\mathbf{Z})$ 위 유한 표시인 대수공간의 경우로 환원한다.
$X$를 $\Spec(\mathbf{Z})$ 위의 대수공간으로 볼 수 있다. 다음을 보라:
Spaces, Definition \ref{spaces-definition-base-change} 및
Properties of Spaces, Definition \ref{spaces-properties-definition-separated}.
따라서 Limits of Spaces, Proposition
\ref{spaces-limits-proposition-approximate}을 적용할 수 있다.
아핀 사상 $X \to X_0$를 찾을 수 있으며, 여기서 $X_0$는
$\mathbf{Z}$ 위 유한 표시이다.
$X_0$에 대해 보조정리를 증명할 수 있다면 층화와 블로업들의 중심을
$X$로 당겨 $X$에 대한 결과를 얻을 수 있다. 여기서는 해소 성질이 아핀
사상을 따라 올라간다는 사실(Derived Categories of Spaces, Lemma
\ref{spaces-perfect-lemma-resolution-property-goes-up-affine})과
아핀 사상의 진변환이 아핀이라는 사실을 사용한다. 세부사항은 생략한다.
이로써 다음 문단에서 다루는 경우로 환원된다.

\medskip\noindent
$X$가 $\mathbf{Z}$ 위 유한 표시라고 가정하자.
그러면 $X$는 뇌터이고 $|X|$는 유한 차원이며 유한 개의 기약 성분을 갖는
뇌터 위상공간이다. 따라서 $\dim(|X|)$에 관한 귀납법을 사용할 수 있다.
Lemma \ref{spaces-more-morphisms-lemma-blow-up-resolution-property}에 의해 조밀 열린집합
$U \subset X$와 $U$-허용 블로업 $X' \to X$가 존재하여 $X'$은 해소
성질을 갖는다. $Z_0 = X$로 놓고, $Z_1 \subset X$를
$|Z_1| = |X| \setminus |U|$인 기약 닫힌 부분공간이라 하자.
귀납법에 의해 정수 $t \geq 0$와 닫힌 부분공간들에 의한 여과
$$
Z_1 \supset Z_{1, 0} \supset Z_{1, 1} \supset \ldots
\supset Z_{1, t} = \emptyset
$$
를 얻는다. 여기서 $Z_{1, 0} \to Z_1$은 두껍히기이고,
$(Z_{1, i} \setminus Z_{1, i + 1})$-허용 블로업
$Z'_{1, i} \to Z_{1, i}$들이 존재하여 $Z'_{1, i}$는 해소 성질을 갖는다.
$Z_1$이 기약이므로 $Z_1 = Z_{1, 0}$이다.
따라서 $Z_i = Z_{1, i - 1}$ 및 $Z'_i = Z'_{1, i - 1}$로 놓을 수 있으며,
이는 $i \geq 1$에 대하여 적용된다. 이로써 보조정리가 증명되었다.
\end{proof}
